% !TeX program = pdflatex
% Resumen ejecutivo de la propuesta curricular (formato OPES--CONARE).
% Compilar desde esta carpeta: python scripts/generate_resumen_ejecutivo.py (regenera las tablas de los apartados 10 a 12)
% y luego pdflatex resumen-ejecutivo (dos veces).
\documentclass[11pt,a4paper]{article}

\usepackage[utf8]{inputenc}
\usepackage[T1]{fontenc}
\usepackage{lmodern}
\renewcommand{\familydefault}{\sfdefault}
\usepackage[spanish, es-tabla]{babel}
\usepackage[left=1.9cm,right=1.9cm,top=2.2cm,bottom=2.2cm,headheight=14pt]{geometry}
\usepackage[svgnames, dvipsnames, table]{xcolor}
\usepackage{amsmath, amsfonts, amssymb}
\usepackage{graphicx}
\usepackage{tikz}
\usetikzlibrary{calc, positioning}
\usepackage{sectsty}
\usepackage{longtable}
\usepackage{pdflscape}
\usepackage{array}
\usepackage{float}
\usepackage{fancyhdr}
\usepackage{enumitem}
\usepackage[colorlinks=true, allcolors=Azul-UCR]{hyperref}

%%%%%%%%%%%%%%%%%%%%%%%%%%%%%%%%%%%%%%%%%%%%%%%%%%
%%%%%%%%%% Colores UCR y estilos %%%%%%%%%%%%%%%%%%
%%%%%%%%%%%%%%%%%%%%%%%%%%%%%%%%%%%%%%%%%%%%%%%%%%
\definecolor{Celeste-UCR}{RGB}{0, 192, 243}
\definecolor{Azul-UCR}{RGB}{0, 93, 164}
\definecolor{Verde-UCR}{RGB}{109, 192, 103}
\definecolor{Amarillo-UCR}{RGB}{255, 224, 106}

\sectionfont{\color{Azul-UCR}}
\subsectionfont{\color{Azul-UCR}}
\subsubsectionfont{\color{Azul-UCR}}

\setlength{\parindent}{0pt}
\setlength{\parskip}{5pt}
\setlist{itemsep=2pt, topsep=3pt}
\setcounter{secnumdepth}{1}

\pagestyle{fancy}
\fancyhf{}
\renewcommand{\headrulewidth}{0.4pt}
\fancyhead[L]{\small\color{Azul-UCR} Resumen ejecutivo --- Bachillerato y Licenciatura en Matemática}
\fancyhead[R]{\small\color{Azul-UCR} Universidad de Costa Rica}
\fancyfoot[C]{\small\thepage}

\newcommand{\N}{\mathbb{N}}
\newcommand{\Z}{\mathbb{Z}}
\newcommand{\C}{\mathbb{C}}
\newcommand{\Q}{\mathbb{Q}}
\newcommand{\I}{\mathbb{I}}
\newcommand{\R}{\mathbb{R}}
\newcommand{\op}[1]{\operatorname{#1}}
\newcommand{\bb}[1]{\mathbb{#1}}
\newcommand{\fr}[1]{\mathfrak{#1}}
\newcommand{\legendre}[2]{\left( \dfrac{#1}{#2} \right)}
\newcommand{\mcd}{\operatorname{mcd}}
\newcommand{\MA}[1]{MA-#1}
\newcommand{\rojo}[1]{{#1}}

% Información que la propuesta no contiene y que debe completar la unidad académica.
\newcommand{\pendiente}[1]{%
  \par\smallskip\noindent\fcolorbox{Amarillo-UCR!90!black}{Amarillo-UCR!35}{%
    \parbox{\dimexpr\linewidth-2\fboxsep-2\fboxrule\relax}{\small\textbf{Por completar por la unidad académica:} #1}}\par\smallskip}

\newcommand{\encabezadotabla}[1]{\rowcolor{Azul-UCR}\textcolor{white}{\textbf{#1}}}

% Anexo de currículos: bloque uniforme por persona (una página por persona).
\newsavebox{\cvbox}
\newcommand{\cvsec}[1]{\par\vspace{2pt}{\bfseries\color{Azul-UCR}#1}\par\nopagebreak}
\newenvironment{cvlist}{\begin{itemize}[leftmargin=1.1em, itemsep=0pt, topsep=1pt, parsep=0pt, partopsep=0pt]}{\end{itemize}}
\newenvironment{cvblock}[2]
  {\par\medskip\noindent\begin{lrbox}{\cvbox}\begin{minipage}{\dimexpr\linewidth-2\fboxsep-2\fboxrule\relax}\small
   {\large\bfseries\color{Azul-UCR} #1}\par\nopagebreak{\itshape #2}\par\nopagebreak\vspace{1pt}\hrule\vspace{1pt}}
  {\end{minipage}\end{lrbox}%
   \ifdim\dimexpr\ht\cvbox+\dp\cvbox\relax>0.95\textheight\PackageWarning{cvblock}{Bloque mayor a una pagina}\fi
   \noindent\fcolorbox{Celeste-UCR}{white}{\usebox{\cvbox}}\par\clearpage}

\title{Resumen ejecutivo de la propuesta curricular}
\author{Escuela de Matemática}
\date{\today}

\begin{document}

%%%%%%%%%%%%%%%%%%%%%%%%%%%%%%%%%%%%%%%%%%%%%%%%%%
%%%%%%%%%%%%%%%%%%%% Portada %%%%%%%%%%%%%%%%%%%%%
%%%%%%%%%%%%%%%%%%%%%%%%%%%%%%%%%%%%%%%%%%%%%%%%%%
% Portada OPES--CONARE (dos páginas). Réplica de Documentos/pagina2-resumen.docx
% Medidas tomadas del Word (A4): marco de 88.2% x 90.9% de la página, centrado.
\definecolor{PortadaCeleste}{HTML}{00C0F3}
\definecolor{PortadaAzulTitulo}{HTML}{002060}

\newcommand{\portadaFuente}[2]{\fontfamily{phv}\fontseries{b}\selectfont\fontsize{#1}{#2}\selectfont}

% Numeración como en el Word: portada = 0, segunda página = 1, cuerpo desde 2.
\setcounter{page}{0}

%%%%%%%%%%%%%%%%%%%%%%%%%%%%%%%%%%%%%%%%%%%%%%%%%%
%%%%%%%%%%%%%%%% Página 1 %%%%%%%%%%%%%%%%%%%%%%%
%%%%%%%%%%%%%%%%%%%%%%%%%%%%%%%%%%%%%%%%%%%%%%%%%%
\thispagestyle{empty}
\null
\begin{tikzpicture}[remember picture, overlay]
  % Marco: x de 1.24 a 19.76 cm; y de 1.35 a 28.35 cm (desde el borde superior)
  \def\mL{1.24cm}\def\mR{19.76cm}
  % Banda celeste superior
  \fill[PortadaCeleste]
    ([xshift=\mL, yshift=-1.35cm]current page.north west) rectangle
    ([xshift=\mR, yshift=-5.41cm]current page.north west);
  % Banda celeste central
  \fill[PortadaCeleste]
    ([xshift=\mL, yshift=-13.47cm]current page.north west) rectangle
    ([xshift=\mR, yshift=-23.66cm]current page.north west);
  % Franja celeste inferior (bajo los logos)
  \fill[PortadaCeleste]
    ([xshift=\mL, yshift=-27.76cm]current page.north west) rectangle
    ([xshift=\mR, yshift=-28.35cm]current page.north west);

  % Título (Arial negrita 26 pt, azul 002060, centrado en la banda blanca)
  \node[anchor=center, inner sep=0pt, text width=16.0cm, align=center]
    at ([xshift=10.5cm, yshift=-9.44cm]current page.north west)
    {\color{PortadaAzulTitulo}\portadaFuente{27}{30}%
     RESUMEN EJECUTIVO PARA EL\\
     REDISEÑO DE LA CARRERA DE\\
     BACHILLERATO Y LICENCIATURA\\
     EN MATEMÁTICA CON ÉNFASIS EN\\
     MATEMÁTICA PURA O\\
     MATEMÁTICA APLICADA};

  % Logos institucionales: UCR, Vicerrectoría de Docencia y CEA
  \node[anchor=north west, inner sep=0pt]
    at ([xshift=2.54cm, yshift=-24.75cm]current page.north west)
    {\includegraphics[width=5.04cm, height=1.907cm]{assets/logos/logo-ucr-nombre.png}};
  \node[anchor=north west, inner sep=0pt]
    at ([xshift=9.25cm, yshift=-24.73cm]current page.north west)
    {\includegraphics[width=2.5cm, height=1.889cm]{assets/logos/logo-vicerrectoria-docencia.png}};
  \node[anchor=north west, inner sep=0pt]
    at ([xshift=14.06cm, yshift=-24.10cm]current page.north west)
    {\includegraphics[width=4.403cm, height=2.983cm]{assets/logos/logo-cea.png}};
\end{tikzpicture}
\clearpage

%%%%%%%%%%%%%%%%%%%%%%%%%%%%%%%%%%%%%%%%%%%%%%%%%%
%%%%%%%%%%%%%%%% Página 2 %%%%%%%%%%%%%%%%%%%%%%%
%%%%%%%%%%%%%%%%%%%%%%%%%%%%%%%%%%%%%%%%%%%%%%%%%%
\thispagestyle{empty}
\null
\begin{tikzpicture}[remember picture, overlay]
  % Número de página en el encabezado (derecha)
  \node[anchor=north east, inner sep=0pt]
    at ([xshift=-2.54cm, yshift=-1.32cm]current page.north east)
    {\fontfamily{phv}\selectfont\fontsize{10}{12}\selectfont\thepage};

  \node[anchor=north, inner sep=0pt]
    at ([yshift=-2.65cm]current page.north)
    {\begin{minipage}{15.92cm}
      \centering
      \portadaFuente{12}{16.45}%
      \setlength{\parindent}{0pt}\setlength{\parskip}{9.4pt}%
      Universidad de Costa Rica\par
      \mbox{}\par\mbox{}\par\mbox{}\par
      Sede Rodrigo Facio\par
      \mbox{}\par\mbox{}\par
      Facultad de Ciencias\par
      Escuela de Matemática\par
      \mbox{}\par\mbox{}\par
      BACHILLERATO Y LICENCIATURA EN MATEMÁTICA CON ÉNFASIS EN MATEMÁTICA PURA O MATEMÁTICA APLICADA\par
      \mbox{}\par
      Aprobado por Asamblea de Asamblea de Escuela n.º EMat-526-2026\par
      Acuerdo n° 2\par
      \mbox{}\par
      Trámite Facultad de Ciencias n° de oficio EMat-2180-2026\par
      \mbox{}\par
      Comisión que rediseñó el plan de estudios:\par
      Adrián Barquero Sánchez\par
      Juan Calvo Alpízar\par
      Samaria Montenegro Guzmán\par
      José Alexander Ramírez González\par
      \mbox{}\par\mbox{}\par
      Noviembre 2026\par
    \end{minipage}};
\end{tikzpicture}
\clearpage

Este documento resume la propuesta de rediseño curricular de la carrera \textbf{Bachillerato y Licenciatura en Matemática con énfasis en Matemática Pura o Matemática Aplicada} de la Universidad de Costa Rica, y se organiza según los doce apartados de la plantilla de resumen para el rediseño de carreras de la Oficina de Planificación de la Educación Superior (OPES) del CONARE. Se agregan dos secciones complementarias (otros aspectos y currículos vitae de las personas docentes) y un anexo con el resumen de los currículos. Las secciones sin contenido corresponden a información que aún no se ha incorporado. El detalle completo se encuentra en el documento \emph{Propuesta de modificación curricular} aprobado por la Asamblea de Escuela de Matemática el 28 de octubre de 2026.

%%%%%%%%%%%%%%%%%%%%%%%%%%%%%%%%%%%%%%%%%%%%%%%%%%
\section{Datos generales}

\begingroup
\renewcommand{\arraystretch}{1.35}
\arrayrulecolor{Azul-UCR!80}
\begin{longtable}{|>{\columncolor{Celeste-UCR!15}\raggedright\arraybackslash\bfseries}p{4.3cm}|>{\raggedright\arraybackslash}p{12.5cm}|}
\hline
Universidad & Universidad de Costa Rica (UCR). \\ \hline
Sede y recinto académico al que pertenece & Sede Rodrigo Facio. \\ \hline
Unidad académica responsable & Escuela de Matemática, Facultad de Ciencias. La carrera se imparte en la Sede Rodrigo Facio (Campus Universitario). \\ \hline
Nombre del grado & Bachillerato y Licenciatura en Matemática con énfasis en Matemática Pura o Matemática Aplicada. \\ \hline
Grados y títulos &
\begin{minipage}[t]{12.2cm}
\begin{itemize}[leftmargin=1.2em, itemsep=1pt, topsep=2pt, parsep=0pt]
  \item Bachillerato en Matemática con énfasis en Matemática Pura.
  \item Bachillerato en Matemática con énfasis en Matemática Aplicada.
  \item Licenciatura en Matemática con énfasis en Matemática Pura.
  \item Licenciatura en Matemática con énfasis en Matemática Aplicada.
\end{itemize}
\end{minipage} \\ \hline
Duración &
\begin{minipage}[t]{12.2cm}
\begin{itemize}[leftmargin=1.2em, itemsep=1pt, topsep=2pt, parsep=0pt]
  \item \textbf{Número de ciclos:} 8 ciclos para el bachillerato (4 años) y 2 ciclos adicionales para la licenciatura (1 año más); en total, 10 ciclos.
  \item \textbf{Semanas por ciclo:} 16 semanas lectivas.
  \item \textbf{Ciclos por año:} 2 (I y II ciclo lectivo).
  \item \textbf{Periodicidad de la oferta académica:} anual
  \item \textbf{Número de promociones:} indefinidas.
\end{itemize}
\end{minipage} \\ \hline
Modalidad & Presencial.\\ \hline
Convenio o cartas de entendimiento  & No se requiere \\ \hline
Datos de contacto de la unidad académica responsable y la instancia curricular a las que se remitirá el acuerdo del CONARE & \begin{minipage}[t]{12.2cm} \begin{itemize}[leftmargin=1.2em, itemsep=1pt, topsep=2pt, parsep=0pt] 
\item Dr. Pedro Méndez Hernández, director, Escuela de Matemática, \href{mailto:direccion.em@ucr.ac.cr}{direccion.em@ucr.ac.cr}
\item Dr. José Alexander Ramírez González, coordinador comisión de rediseño, \href{mailto:alexander.ramirez_g@ucr.ac.cr}{alexander.ramirez\_g@ucr.ac.cr}
\item Centro de Evaluación Académica, Universidad de Costa Rica, \href{mailto:cea@ucr.ac.cr}{cea@ucr.ac.cr}, teléfono 2511-1220 
\end{itemize}
\end{minipage}\\ \hline
\end{longtable}
\endgroup

\subsection*{Resumen de modificaciones}
El rediseño del plan de estudios incorpora ajustes orientados a fortalecer la coherencia académica y la pertinencia de la formación matemática. Se reorganiza la estructura curricular en un bloque común (ciclos I a III) y dos énfasis, en Matemática Pura y en Matemática Aplicada, a partir del ciclo IV, de modo que cada persona estudiante profundice en la ruta de su interés. Asimismo, se establecen niveles progresivos de profundización de los contenidos: cursos iniciales de exploración y fundamentación del razonamiento matemático (Matemática Exploratoria e Introducción a las Demostraciones), cursos obligatorios en áreas fundamentales, optativas de núcleo y de tópicos, y seminarios de estudios dirigidos en la licenciatura. Se promueve una mayor articulación entre teoría, práctica e investigación mediante las prácticas profesionales (en modalidad académica o profesional), los seminarios de estudios dirigidos y el curso de Comunicación en las Ciencias. Adicionalmente, se actualizan contenidos en áreas estratégicas y se amplía la oferta optativa con cursos de tópicos. 

\subsection*{Presupuesto y finanzas}
Cuenta con 128 docentes, 41 de ellos con grado de doctorado.
. 

La Universidad de Costa Rica, a través de la Escuela de Matemática, indica que dispone de los recursos necesarios para el rediseño y la implementación de la carrera de Bachillerato y Licenciatura en Matemática con énfasis en Matemática Pura o Matemática Aplicada mediante la siguiente planificación:
\begin{itemize}
    \item {\color{Azul-UCR}\textbf{Talento Humano (Recurso Humano)}}
    \begin{itemize}
        \item \textbf{Cuerpo Docente:} La carrera cuenta con un equipo consolidado de aproximadamente 45 académicos con amplia trayectoria en docencia, investigación y acción social.
        \item \textbf{Formación de Alto Nivel:} La mayoría del equipo docente posee formación de posgrado (41 del equipo cuenta con doctorado) en áreas clave o de conocimiento como: el álgebra, la teoría de números, la geometría, el análisis, la probabilidad, las ecuaciones diferenciales, la lógica, la combinatoria, el análisis numérico y otras.
        \item \textbf{Capacidad Instalada:} El personal actual cubre las áreas sustantivas de la carrera, y existen docentes en proceso de formación, etc..
    \end{itemize}
    \item {\color{Azul-UCR}\textbf{Infraestructura y Equipo}}
    \begin{itemize}
        \item \textbf{Edificio Propio:} La carrera dispone de espacio en dos edificios propios (uno de ellos compartido con Física), lo que garantiza condiciones físicas adecuadas.
        \item \textbf{Aulas Equipadas:} Cuenta con xxxx aulas con capacidad para xxxx estudiantes cada una, equipadas con proyectores, computadoras y acceso a plataformas institucionales.
        \item \textbf{Espacios de Uso Múltiple:} Dispone de un miniauditorio para 30 personas y aulas multiuso para xx personas, utilizados para conferencias, cursos especiales y defensas de tesis.
        \item \textbf{Tecnología:} Todas las aulas cuentan con tecnología instalada, licencias, etc…, facilitando la participación en actividades académicas.
    \end{itemize}
    \item {\color{Azul-UCR}\textbf{Recursos Presupuestarios y Financieros}}
    \begin{enumerate}
        \item \textbf{Sostenibilidad y Eficiencia:} La propuesta de reestructuración no requiere ampliaciones extraordinarias de presupuesto. Al contrario, el nuevo diseño optimiza la carga docente, logrando una reducción de xxx horas semanales (equivalente a xx tiempos completos) en comparación con el plan vigente.
        \item \textbf{Uso Racional:} Esta optimización genera un margen presupuestario que permite fortalecer la investigación y la innovación metodológica sin implicar una mayor erogación financiera para la Unidad Académica.
        \item \textbf{Apoyo Administrativo:} Se cuenta con personal de apoyo en áreas de asuntos estudiantiles, recepción, servicios de limpieza y soporte tecnológico (área RID), además de acceso al Sistema de Bibliotecas (SIBDI).
    \end{enumerate}
\end{itemize}


%%%%%%%%%%%%%%%%%%%%%%%%%%%%%%%%%%%%%%%%%%%%%%%%%%
\section{Justificación}

La propuesta responde a la necesidad de contar con personas profesionales en matemática con excelencia académica, ética y sentido crítico, capaces de aplicar, investigar y comunicar su conocimiento de forma colaborativa. Atiende las demandas de investigación y docencia universitaria, de modelación matemática para otras disciplinas, de análisis de datos y de consultoría cuantitativa en el sector productivo, y contribuye al proyecto humanista de la Universidad mediante la generación de capital humano avanzado, la innovación en sectores estratégicos y la cultura científica del país.

La matemática ocupa un lugar central en el desarrollo científico, tecnológico y social contemporáneo. La digitalización de la información, el avance de la inteligencia artificial, el manejo de datos a gran escala y la modelación de fenómenos complejos han ampliado el objeto de estudio de la disciplina y los ámbitos en que la persona matemática puede incidir. La UNESCO proclamó en 2019 el 14 de marzo como Día Internacional de las Matemáticas y subraya su papel en el logro de los Objetivos de Desarrollo Sostenible y en la formación de capacidades para enfrentar desafíos como el cambio climático, la salud pública y la innovación tecnológica.

La demanda laboral respalda esta pertinencia. En Costa Rica, la empleabilidad de las carreras STEM fue de 95,8\,\% frente a 92,3\,\% en las no STEM, y carreras afines como Estadística o Ciencias Actuariales presentan niveles de desempleo muy bajos. A nivel internacional, la Oficina de Trabajo Estadístico de Estados Unidos proyectó un aumento del 28\,\% en los empleos de ocupaciones matemáticas entre 2016 y 2026, impulsado por el \emph{big data} y la analítica de datos.

Las personas profesionales en matemática encuentran oportunidades en instituciones de educación superior y de investigación, en organismos públicos y privados que requieren modelación y análisis cuantitativo, y en sectores como finanzas, banca, economía, tecnología y servicios basados en datos. 

Sobre el contexto universitario en que se desarrolla la carrera actual se tiene:

\begin{itemize}
  \item \textbf{Trayectoria de la carrera.} La Escuela de Matemática ofrece el Bachillerato y la Licenciatura desde 1974 (con antecedentes desde 1940), y desarrolla investigación y docencia de posgrado mediante la Maestría en Matemáticas, el CIMM, el CIMPA y la \emph{Revista de Matemática: Teoría y Aplicaciones}. 
  \item \textbf{Actualización del plan.} El último rediseño integral del plan data de hace más de tres décadas. La propuesta se apoya en el perfil de egreso y los marcos referenciales elaborados entre 2018 y 2020 con apoyo del Centro de Evaluación Académica (CEA), mediante entrevistas, grupos focales, análisis estructurales de la actividad y consultas a la comunidad académica; el perfil fue aprobado en la Asamblea de Escuela en 2020.
  \item \textbf{Situación de la población estudiantil.} Entre 2014 y 2026 se graduaron 72 personas del bachillerato, con una duración real media de 5,3 años frente a los 4 años del plan; solo el 22,2\,\% concluyó en 4 años o menos. De las 543 solicitudes aceptadas entre 2012 y 2018, apenas el 9,9\,\% se había graduado en 2026. La propuesta busca atender el rezago mediante cursos preliminares en el primer ciclo y una introducción más gradual de la abstracción.
  \item \textbf{Admisión.} La oferta de cupos y la demanda han disminuido desde 2002; desde 2017 el cupo de primer ingreso es de 32 plazas, y la ocupación se recuperó a partir de 2024.
\end{itemize}

La creación de la carrera de Educación Matemática (2017) y la separación de las funciones actuariales hacia Ciencias Actuariales redefinen el espacio ocupacional de la persona matemática, orientándola hacia la investigación, la docencia universitaria, la modelación, el análisis de datos y el trabajo interdisciplinario.

La actualización de la formación de profesionales en matemática es fundamental para mantener el rol preponderante que la disciplina tiene en la formación de profesionales de otras carreras, la investigación en matemática y otras disciplinas así como en los programas de acción social donde son necesarios conocimientos de matemáticas.

La propuesta también busca favorecer el éxito académico de la población estudiantil, procurando un enfoque pedagógico más moderno y acorde con el perfil de la población de entrada. 



%%%%%%%%%%%%%%%%%%%%%%%%%%%%%%%%%%%%%%%%%%%%%%%%%%
\section{Objeto de estudio de la profesión/carrera}

La matemática se entiende como el estudio de conceptos y entes abstractos, sus propiedades y las relaciones entre ellos. Estas entidades responden, en su mayoría, a versiones idealizadas de la realidad, aunque el paso del objeto real observado al ente abstracto que lo representa puede ser intrincado y tomar siglos; también ocurre que entidades creadas solo en la abstracción terminen respondiendo a otras realidades. 

La matemática se comprende desde diez áreas de conocimiento, que en el marco referencial se denominan \emph{identidades matemáticas}: Álgebra, Análisis, Geometría, Modelación Matemática, Análisis Numérico, Lógica Matemática, Teoría de Números, Topología, Probabilidad y Ecuaciones Diferenciales. Esta clasificación no pretende abordar la disciplina de forma separada: cada vez más las situaciones se analizan de manera integrada, con el aporte de distintas áreas.

En lo teórico-metodológico, la matemática es una ciencia formal y su conocimiento es colectivo, no personal: se construye mediante el intercambio de ideas y el trabajo colaborativo en congresos, simposios, coloquios y otras reuniones científicas. Se produce mediante el pensamiento y la reflexión sobre conjeturas, el estudio de casos particulares y ejemplos, y el seguimiento de la intuición. La teoría y la práctica están estrechamente ligadas: la práctica es la propia investigación, que se nutre de la de otras personas, y, en el sector productivo, la aplicación de modelos existentes puede generar nuevos problemas que luego se analizan matemáticamente. 

Respecto de su finalidad, la matemática se ha desarrollado históricamente como respuesta a necesidades del quehacer humano. Hoy se genera cada vez más conocimiento matemático en el entorno socioeconómico, el entorno productivo (optimización, clasificación, \emph{big data}, análisis numérico), las otras ciencias básicas (Física, Química, Biología) y el estudio de fenómenos sociodemográficos. Es desarrollada por personas matemáticas dedicadas a la investigación y, cada vez más, por profesionales de disciplinas afines en instituciones financieras y empresas del sector industrial. 

Finalmente, la matemática mantiene una relación estrecha con otras disciplinas: se nutre de aportes de la biología, la demografía, la microbiología, la física, la computación, la filosofía y la economía, y a su vez sus herramientas y teorías sirven a la ecología, la dinámica social, la ingeniería, la química, las finanzas y la arquitectura, entre otras. Su impacto es actual: por ejemplo, la modelación apoya políticas de salud pública, el análisis numérico resuelve los modelos de la física, la geología y la meteorología, la lógica contribuye a la criptografía y a la complejidad algorítmica, y las ecuaciones diferenciales se aplican a la imagen médica, la epidemiología y la mecánica cuántica.

%%%%%%%%%%%%%%%%%%%%%%%%%%%%%%%%%%%%%%%%%%%%%%%%%%
\section{Objetivos de la carrera}

\textbf{Objetivo general.} Formar profesionales en matemática con excelencia académica, ética y sentido crítico, capaces de aplicar, investigar y comunicar su conocimiento de forma colaborativa, para contribuir al desarrollo productivo y social del país desde una perspectiva de bien común.

\textbf{Objetivos específicos.}
\begin{enumerate}
  \item Formar una persona profesional que impacte de manera positiva el entorno productivo, público y privado, mediante la aplicación del conocimiento matemático.
  \item Formar una persona profesional capaz de transmitir sus ideas de manera clara y respetuosa, que interactúe con profesionales de otras disciplinas y muestre gusto por compartir su conocimiento.
  \item Formar una persona profesional con un alto compromiso con la investigación, tanto teórica como aplicada.
  \item Formar una persona profesional con convicciones éticas y pensamiento crítico y reflexivo sobre su práctica profesional.
  \item Formar una persona profesional comprometida con el bien común y la justicia, que utilice su conocimiento para construir una sociedad más justa y solidaria.
  \item Incentivar una actitud positiva hacia el trabajo colaborativo con otros especialistas, en la búsqueda de soluciones que mejoren la calidad de vida y el entorno social.
  \item Formar una persona profesional dispuesta a aprender siempre, a buscar nuevas formas de hacer las cosas y a mantenerse actualizada en su área de especialidad.
\end{enumerate}

{\footnotesize Los objetivos corresponden a la Declaración de propósitos (propósito general y propósitos específicos) de la propuesta curricular.\par}

%%%%%%%%%%%%%%%%%%%%%%%%%%%%%%%%%%%%%%%%%%%%%%%%%%
\section{Perfil académico-profesional}

En la presente sección se presenta el perfil de egreso.

%\newpage 

\subsection*{Bachillerato en matemática, énfasis en matemática pura}


\begin{table}[H]
\centering
\renewcommand{\arraystretch}{1.3}
\begin{tabular}{|c|p{9cm}|}
\hline
\cellcolor[HTML]{00C0F3}\textbf{Código(s)} & \cellcolor[HTML]{00C0F3}\textbf{Enunciado de la dimensión declarativa} \\
\hline
CDP01 & Conoce y comprende los fundamentos matemáticos —incluyendo la lógica, la teoría de números y la geometría euclídea— y los principios formales que sustentan el razonamiento y la argumentación rigurosa. \\
\hline
CDP02 & Conoce el uso de software especializado, los fundamentos de la programación y la computación científica, así como las técnicas básicas de análisis numérico para formular, explorar y resolver problemas matemáticos mediante enfoques computacionales. \\
\hline
CDP03 & Conoce los conceptos, métodos y aplicaciones básicos del cálculo y del álgebra lineal como fundamentos de la formación matemática. \\
\hline
CDP04 & Conoce los conceptos y técnicas de las ecuaciones diferenciales, la teoría de la medida, la probabilidad y los procesos estocásticos, así como del análisis complejo y funcional. \\
\hline
CDP05 & Conoce los fundamentos y resultados básicos de la geometría y la topología. \\
\hline
CDP06 & Conoce las estructuras y resultados fundamentales del álgebra abstracta. 
\\
\hline
CDP07 & Conoce las estructuras y resultados fundamentales de algún área de especialización en matemática pura más allá de las antes expuestas. 
\\
\hline
\end{tabular}
\caption{Conocimientos declarativos o saber conocer.}
\end{table}

%\newpage

\begin{table}[H]
\centering
\renewcommand{\arraystretch}{1.3}
\begin{tabular}{|c|p{9cm}|}
\hline
\cellcolor[HTML]{00C0F3}\textbf{Código(s)}  & \cellcolor[HTML]{00C0F3}\textbf{Enunciado de la dimensión procedimental} \\
\hline
CPP01 & Aplica pensamiento operacional, abstracto y formal para analizar, resolver y sintetizar problemas matemáticos. \\
\hline
CPP02 & Formula conjeturas, resuelve problemas, interpreta resultados y se comunica de manera coherente en lenguaje matemático formal. \\
\hline
CPP03 & Utiliza tecnologías y software científico para explorar, validar y resolver problemas matemáticos.\\
\hline
CPP04 & Participa activamente en discusiones científicas y divulga el conocimiento matemático de forma accesible a diferentes audiencias.\\
\hline
CPP05 & Formula y demuestra rigurosamente resultados matemáticos, generalizando estructuras y construyendo argumentos lógicos válidos dentro de un sistema axiomático.\\
\hline
\end{tabular}
\caption{Conocimientos procedimentales o saber hacer.}
\end{table}

\begin{table}[H]
\centering
\renewcommand{\arraystretch}{1.3}
\begin{tabular}{|c|p{9cm}|}
\hline
\cellcolor[HTML]{00C0F3}\textbf{Código(s)} & \cellcolor[HTML]{00C0F3}\textbf{Enunciado de la dimensión actitudinal} \\
\hline
CAP01 & Manifiesta respeto por la diversidad humana, matemática, científica y cultural, fomentando un ambiente inclusivo y de valoración mutua. \\
\hline
CAP02 & Se comunica con claridad, colabora en equipos académicos y profesionales, y mantiene una disposición al autoaprendizaje y al liderazgo.\\
\hline
CAP03 & Demuestra conciencia social y participación ciudadana, con una visión sociocultural e integral de la sociedad en la que la ciencia y la matemática ocupan un lugar relevante.\\
\hline
CAP04 & Ejerce pensamiento crítico, actúa con integridad ética y evalúa con rigor las prácticas científicas y académicas.\\
\hline
CAP05 & Utiliza su conocimiento matemático para resolver problemas académicos y profesionales, adaptándolo a distintos contextos y necesidades.\\
\hline
CAP06 & Mantiene un interés permanente por los desarrollos de la investigación científica y matemática, valorando la actualización como parte de su formación profesional.\\
\hline
CAP07 & Muestra paciencia y perseverancia ante la abstracción sostenida, valorando el rigor y la elegancia argumentativa como fines en sí mismos, incluso cuando la solución de un problema no es inmediata.\\
\hline
\end{tabular}
\caption{Conocimientos actitudinales o saber ser.}
\end{table}


\subsection*{Bachillerato en matemática, énfasis en matemática aplicada}



\begin{table}[H]
\centering
\renewcommand{\arraystretch}{1.3}
\begin{tabular}{|c|p{9cm}|}
\hline
\cellcolor[HTML]{00C0F3}\textbf{Código(s)} & \cellcolor[HTML]{00C0F3}\textbf{Enunciado de la dimensión declarativa} \\
\hline
CDA01 & Conoce y comprende los fundamentos matemáticos —incluyendo la lógica, la teoría de números y la geometría euclídea— y los principios formales que sustentan el razonamiento y la argumentación rigurosa. \\
\hline
CDA02 & Conoce el uso de software especializado, los fundamentos de la programación y la computación científica, así como las técnicas básicas de análisis numérico para formular, explorar y resolver problemas matemáticos mediante enfoques computacionales. \\
\hline
CDA03 & Conoce los conceptos, métodos y aplicaciones básicos del cálculo y del álgebra lineal como fundamentos de la formación matemática. \\
\hline
CDA04 & Conoce los conceptos y técnicas de las ecuaciones diferenciales, la teoría de la medida, la probabilidad y los procesos estocásticos, así como del análisis complejo y funcional. \\
\hline
CDA05 & Conoce los principios de la matemática aplicada, por ejemplo, por medio de la modelación, la estadística o el análisis de datos, para resolver problemas y tomar decisiones en ámbitos como la industria, la banca y otros sectores.
\\
\hline
CDA06 & Conoce las principales tecnologías de aplicación de la matemática, por ejemplo, la programación y computación científica para aplicar conceptos teóricos a problemas reales.
\\
\hline
CDA07 & Conoce las estructuras y resultados fundamentales de algún área de especialización en matemática aplicada más allá de las antes expuestas. 
\\
\hline
\end{tabular}
\caption{Conocimientos declarativos o saber conocer.}
\end{table}

%\newpage

\begin{table}[H]
\centering
\renewcommand{\arraystretch}{1.3}
\begin{tabular}{|c|p{9cm}|}
\hline
\cellcolor[HTML]{00C0F3}\textbf{Código(s)}  & \cellcolor[HTML]{00C0F3}\textbf{Enunciado de la dimensión procedimental} \\
\hline
CPA01 & Aplica pensamiento operacional, abstracto y formal para analizar, resolver y sintetizar problemas matemáticos. \\
\hline
CPA02 & Formula conjeturas, resuelve problemas, interpreta resultados y se comunica de manera coherente en lenguaje matemático formal. \\
\hline
CPA03 & Utiliza tecnologías y software científico para explorar, validar y resolver problemas matemáticos.\\
\hline
CPA04 & Participa activamente en discusiones científicas y divulga el conocimiento matemático de forma accesible a diferentes audiencias.\\
\hline
CPA05 & Formula modelos matemáticos a partir de problemas del entorno, los resuelve mediante herramientas analíticas o computacionales, e interpreta y valida los resultados en el contexto original.\\
\hline
\end{tabular}
\caption{Conocimientos procedimentales o saber hacer.}
\end{table}


%\newpage

\begin{table}[H]
\centering
\renewcommand{\arraystretch}{1.3}
\begin{tabular}{|c|p{9cm}|}
\hline
\cellcolor[HTML]{00C0F3}\textbf{Código(s)} & \cellcolor[HTML]{00C0F3}\textbf{Enunciado de la dimensión actitudinal} \\
\hline
CAA01 & Manifiesta respeto por la diversidad humana, matemática, científica y cultural, fomentando un ambiente inclusivo y de valoración mutua. \\
\hline
CAA02 & Se comunica con claridad, colabora en equipos académicos y profesionales, y mantiene una disposición al autoaprendizaje y al liderazgo.\\
\hline
CAA03 & Demuestra conciencia social y participación ciudadana, con una visión sociocultural e integral de la sociedad en la que la ciencia y la matemática ocupan un lugar relevante.\\
\hline
CAA04 & Ejerce pensamiento crítico, actúa con integridad ética y evalúa con rigor las prácticas científicas y académicas.\\
\hline
CAA05 & Utiliza su conocimiento matemático para resolver problemas académicos y profesionales, adaptándolo a distintos contextos y necesidades.\\
\hline
CAA06 & Mantiene un interés permanente por los desarrollos de la investigación científica y matemática, valorando la actualización como parte de su formación profesional.\\
\hline
CAA07 & Asume con responsabilidad ética las implicaciones prácticas de los modelos y datos que utiliza, y muestra apertura al diálogo interdisciplinario y tolerancia a la incertidumbre propia de los fenómenos reales.\\
\hline
\end{tabular}
\caption{Conocimientos actitudinales o saber ser.}
\end{table}


\subsection*{Licenciatura en matemática, énfasis en matemática pura}


En adición a los conocimientos enunciados para el bachillerato, con la licenciatura se tendrían los siguientes.

\begin{table}[H]
\centering
\renewcommand{\arraystretch}{1.3}
\begin{tabular}{|c|p{9cm}|}
\hline
\cellcolor[HTML]{00C0F3}\textbf{Código(s)} & \cellcolor[HTML]{00C0F3}\textbf{Enunciado de la dimensión declarativa} \\
\hline
CDP08 & Conoce estructuras y resultados de nivel intermedio para algunas áreas medulares de especialización en matemática pura. 
\\
\hline
\end{tabular}
\caption{Conocimientos declarativos o saber conocer.}
\end{table}

%\newpage

\begin{table}[H]
\centering
\renewcommand{\arraystretch}{1.3}
\begin{tabular}{|c|p{9cm}|}
\hline
\cellcolor[HTML]{00C0F3}\textbf{Código(s)}  & \cellcolor[HTML]{00C0F3}\textbf{Enunciado de la dimensión procedimental} \\
\hline
CP05 & Aplica pensamiento operacional, abstracto y formal, con fuentes transversales de varias áreas, para formular, analizar y resolver problemas matemáticos. \\
\hline
\end{tabular}
\caption{Conocimientos procedimentales o saber hacer.}
\end{table}


\begin{table}[H]
\centering
\renewcommand{\arraystretch}{1.3}
\begin{tabular}{|c|p{9cm}|}
\hline
\cellcolor[HTML]{00C0F3}\textbf{Código(s)} & \cellcolor[HTML]{00C0F3}\textbf{Enunciado de la dimensión actitudinal} \\
\hline
CA06 & Muestra iniciativa y autonomía para formular preguntas, explorar problemas matemáticos y profundizar en ellos más allá de los procedimientos conocidos.\\
\hline
\end{tabular}
\caption{Conocimientos actitudinales o saber ser.}
\end{table}

\subsection*{Licenciatura en matemática, énfasis en matemática aplicada}

En adición a los conocimientos enunciados para el bachillerato, con la licenciatura se tendrían los siguientes.

\begin{table}[H]
\centering
\renewcommand{\arraystretch}{1.3}
\begin{tabular}{|c|p{9cm}|}
\hline
\cellcolor[HTML]{00C0F3}\textbf{Código(s)} & \cellcolor[HTML]{00C0F3}\textbf{Enunciado de la dimensión declarativa} \\
\hline
CDA08 & Conoce estructuras y resultados de nivel intermedio para algunas áreas medulares de especialización en matemática aplicada. 
\\
\hline
\end{tabular}
\caption{Conocimientos declarativos o saber conocer.}
\end{table}

%\newpage

\begin{table}[H]
\centering
\renewcommand{\arraystretch}{1.3}
\begin{tabular}{|c|p{9cm}|}
\hline
\cellcolor[HTML]{00C0F3}\textbf{Código(s)}  & \cellcolor[HTML]{00C0F3}\textbf{Enunciado de la dimensión procedimental} \\
\hline
CP05 & Aplica pensamiento operacional, abstracto y formal, con fuentes transversales de varias áreas, para formular, analizar y resolver problemas matemáticos. \\
\hline
\end{tabular}
\caption{Conocimientos procedimentales o saber hacer.}
\end{table}


\begin{table}[H]
\centering
\renewcommand{\arraystretch}{1.3}
\begin{tabular}{|c|p{9cm}|}
\hline
\cellcolor[HTML]{00C0F3}\textbf{Código(s)} & \cellcolor[HTML]{00C0F3}\textbf{Enunciado de la dimensión actitudinal} \\
\hline
CA06 & Muestra iniciativa y autonomía para formular preguntas, explorar problemas matemáticos y profundizar en ellos más allá de los procedimientos conocidos.\\
\hline
\end{tabular}
\caption{Conocimientos actitudinales o saber ser.}
\end{table}

%%%%%%%%%%%%%%%%%%%%%%%%%%%%%%%%%%%%%%%%%%%%%%%%%%
\section{Marco de cualificaciones}

El Marco de Cualificaciones para la Educación Superior Centroamericana (MCESCA), publicado por el Consejo Superior Universitario Centroamericano (CSUCA) en 2018, establece los resultados de aprendizaje esperados para los niveles de técnico superior, bachillerato (CINE 6A), licenciatura (CINE 6B), maestría y doctorado. Los resultados se organizan en cinco descriptores: saberes disciplinarios y profesionales; aplicación de conocimientos, resolución de problemas e innovación; autonomía con responsabilidad personal, laboral y social; comunicación; e interacción profesional, cultural y social.

Se presenta una tabla para el Bachillerato y otra para la Licenciatura. En ellas, los dos énfasis se muestran juntos y se distinguen por la letra del código (P: Matemática Pura; A: Matemática Aplicada). La columna «Resultado MCESCA» transcribe los resultados de las Tablas 6 y 7 del MCESCA. La columna «Documento plan de estudios» transcribe, sin modificar palabra alguna, los enunciados del perfil de egreso (sección 5) que responden al resultado, con su código; cuando el perfil no tiene un enunciado explícito, se indica en qué parte del plan se atiende (cursos, propósitos o requisitos de graduación).

\begingroup
\newcommand{\qe}[2]{\par\textbf{#1:} ``#2''}
\newcommand{\cur}[1]{\par\textit{Cursos:} #1}
\newcommand{\nt}[1]{\par\textit{Nota:} #1}
\newcommand{\mrow}[3]{#1 & #2 & #3 \\ \hline}
\newcommand{\mcesca}[2]{%
  \begingroup
  \footnotesize
  \setlength{\tabcolsep}{4pt}%
  \setlength{\parskip}{2pt}%
  \arrayrulecolor{Azul-UCR!80}%
  \renewcommand{\arraystretch}{1.2}%
  \rowcolors{2}{white}{Celeste-UCR!12}%
  \begin{longtable}{|>{\raggedright\arraybackslash}p{2.2cm}|>{\raggedright\arraybackslash}p{5.1cm}|>{\raggedright\arraybackslash}p{8.4cm}|}
  \caption{#1}\label{#2}\\
  \hline
  \rowcolor{Azul-UCR}
  \textcolor{white}{\textbf{Descriptor}} & \textcolor{white}{\textbf{Resultado MCESCA}} & \textcolor{white}{\textbf{Documento plan de estudios}} \\ \hline
  \endfirsthead
  \hline
  \rowcolor{Azul-UCR}
  \textcolor{white}{\textbf{Descriptor}} & \textcolor{white}{\textbf{Resultado MCESCA}} & \textcolor{white}{\textbf{Documento plan de estudios}} \\ \hline
  \endhead
  \hline
  \endfoot
}
\newcommand{\finmcesca}{\end{longtable}\endgroup}

\subsection*{Bachillerato en Matemática (CINE 6A), ambos énfasis}

\mcesca{Análisis del plan de estudios respecto del MCESCA: nivel Bachillerato}{res:mcesca-bach}
\mrow{Saberes disciplinarios y profesionales}
{Demuestra la comprensión y explica los principios y conceptos que sustentan los procedimientos, procesos, sistemas y metodologías de la profesión.}
{\qe{CDP01 / CDA01}{Conoce y comprende los fundamentos matemáticos —incluyendo la lógica, la teoría de números y la geometría euclídea— y los principios formales que sustentan el razonamiento y la argumentación rigurosa.}
 \qe{CDP03 / CDA03}{Conoce los conceptos, métodos y aplicaciones básicos del cálculo y del álgebra lineal como fundamentos de la formación matemática.}
 \nt{se complementan con CDP04 a CDP07 (énfasis pura) y CDA04, CDA05 y CDA07 (énfasis aplicada), que cubren ecuaciones diferenciales, medida, probabilidad, análisis complejo y funcional, geometría, topología, álgebra abstracta y la especialización de cada énfasis.}
 \cur{MA-0252 Introducción a las Demostraciones; MA-0251 Cálculo en una Variable; MA-0351 Cálculo en Varias Variables; MA-0361 Álgebra Lineal I; MA-0461 Álgebra Lineal II.}}
\mrow{Saberes disciplinarios y profesionales}
{Demuestra conocimientos de cultura humanística y social que le brindan una visión amplia de su profesión y del entorno.}
{ \qe{CAP03 / CAA03}{Demuestra conciencia social y participación ciudadana, con una visión sociocultural e integral de la sociedad en la que la ciencia y la matemática ocupan un lugar relevante.}
 \cur{EG-I y EG-II Curso Integrado de Humanidades I y II; EG- Curso de Arte; RP- Repertorio; SR-I y SR-II Seminario de Realidad Nacional I y II; optativas de otra disciplina de Filosofía (por ejemplo, F2002 Filosofía de la Matemática).}}
\mrow{Saberes disciplinarios y profesionales}
{Muestra conocimiento sobre el riesgo de desastres en el campo de su especialidad y propone soluciones para reducirlo.}
{\nt{el perfil no tiene un enunciado explícito para este resultado. Se identifican solo aportes indirectos: la modelación matemática del énfasis aplicada y las optativas de otra disciplina relacionadas con fenómenos naturales. Se señala como aspecto por fortalecer.}
 \cur{MA-0647 Modelación Matemática (énfasis aplicada); FS-4101 Introducción a la Meteorología; G0111 Fundamentos de Geología; G0214 Geología General (optativas de otra disciplina).}}
\mrow{Saberes disciplinarios y profesionales}
{Domina las tecnologías y herramientas propias de su profesión para la innovación de procedimientos, procesos, sistemas y metodologías para la mejora de resultados de su área.}
{\qe{CDP02 / CDA02}{Conoce el uso de software especializado, los fundamentos de la programación y la computación científica, así como las técnicas básicas de análisis numérico para formular, explorar y resolver problemas matemáticos mediante enfoques computacionales.}
 \qe{CPP03 / CPA03}{Utiliza tecnologías y software científico para explorar, validar y resolver problemas matemáticos.}
 \qe{CDA06}{Conoce las principales tecnologías de aplicación de la matemática, por ejemplo, la programación y computación científica para aplicar conceptos teóricos a problemas reales.}
 \cur{MA-0341 Matemática Computacional; MA-0501 Análisis Numérico I; CA-0204 Herramientas de Ciencia de Datos I (énfasis aplicada).}}
\mrow{Aplicación de conocimientos, resolución de problemas e innovación}
{Demuestra pensamiento crítico, actitud de indagación y rigor analítico en el planteamiento y la resolución de problemas.}
{\qe{CPP01 / CPA01}{Aplica pensamiento operacional, abstracto y formal para analizar, resolver y sintetizar problemas matemáticos.}
 \qe{CAP04 / CAA04}{Ejerce pensamiento crítico, actúa con integridad ética y evalúa con rigor las prácticas científicas y académicas.}
 \qe{CPP05}{Formula y demuestra rigurosamente resultados matemáticos, generalizando estructuras y construyendo argumentos lógicos válidos dentro de un sistema axiomático.}
 \cur{MA-0152 Matemática Exploratoria; MA-0252 Introducción a las Demostraciones; MA-0451 y MA-0555 Principios de Análisis I y II; MA-0716 Análisis Real I.}}
\mrow{Aplicación de conocimientos, resolución de problemas e innovación}
{Aplica conocimiento teórico-práctico, en diferentes contextos para identificar y resolver problemas de su ámbito laboral y colabora en la resolución de problemas complejos.}
{\qe{CAP05 / CAA05}{Utiliza su conocimiento matemático para resolver problemas académicos y profesionales, adaptándolo a distintos contextos y necesidades.}
 \qe{CPP02 / CPA02}{Formula conjeturas, resuelve problemas, interpreta resultados y se comunica de manera coherente en lenguaje matemático formal.}
 \qe{CPA05}{Formula modelos matemáticos a partir de problemas del entorno, los resuelve mediante herramientas analíticas o computacionales, e interpreta y valida los resultados en el contexto original.}
 \cur{MA-0647 Modelación Matemática (énfasis aplicada); MA-0881 Práctica Profesional en Matemática Pura; MA-0883 Práctica Profesional en Matemática Aplicada.}}
\mrow{Aplicación de conocimientos, resolución de problemas e innovación}
{Propone e implementa nuevos procedimientos aplicables a la solución de problemas y mejora de su campo profesional.}
{\qe{CPP02 / CPA02}{Formula conjeturas, resuelve problemas, interpreta resultados y se comunica de manera coherente en lenguaje matemático formal.}
 \qe{CPA05}{Formula modelos matemáticos a partir de problemas del entorno, los resuelve mediante herramientas analíticas o computacionales, e interpreta y valida los resultados en el contexto original.}
 \cur{MA-0647 Modelación Matemática, que incluye un proyecto final integrador (énfasis aplicada); MA-0881 y MA-0883 Práctica Profesional.}}
\mrow{Aplicación de conocimientos, resolución de problemas e innovación}
{Toma decisiones profesionales con base en datos en información pertinente, válida y confiable.}
{\qe{CDA05}{Conoce los principios de la matemática aplicada, por ejemplo, por medio de la modelación, la estadística o el análisis de datos, para resolver problemas y tomar decisiones en ámbitos como la industria, la banca y otros sectores.}
 \qe{CAA07}{Asume con responsabilidad ética las implicaciones prácticas de los modelos y datos que utiliza, y muestra apertura al diálogo interdisciplinario y tolerancia a la incertidumbre propia de los fenómenos reales.}
 \cur{CA-0721 Probabilidad; CA-0303 Estadística Actuarial I; CA-0411 Análisis de Datos I (énfasis aplicada).}
 \nt{en el énfasis pura se atiende de forma más limitada, mediante CPP02 / CPA02 (interpreta resultados) y optativas como MA-0841 Probabilidad, MA-0833 Tópicos de Estadística y XS-3150 Diseño de Experimentos.}}
\mrow{Aplicación de conocimientos, resolución de problemas e innovación}
{Demuestra destreza y habilidad en el uso de herramientas metodológicas, tecnológicas, equipo especializado y en la lectura de datos para el ejercicio de su profesión en diversos contextos.}
{\qe{CPP03 / CPA03}{Utiliza tecnologías y software científico para explorar, validar y resolver problemas matemáticos.}
 \cur{MA-0341 Matemática Computacional; MA-0501 Análisis Numérico I; MA-0647 Modelación Matemática, con validación e interpretación con datos (énfasis aplicada); CA-0204 Herramientas de Ciencia de Datos I y CA-0411 Análisis de Datos I (énfasis aplicada).}}
\mrow{Autonomía con responsabilidad personal, laboral y social}
{Demuestra responsabilidad en la determinación de resultados personales y laborales de acuerdo a su función tomando como referencias las normativas legales y éticas de su campo profesional.}
{\qe{CAP04 / CAA04}{Ejerce pensamiento crítico, actúa con integridad ética y evalúa con rigor las prácticas científicas y académicas.}
 \qe{CAA07}{Asume con responsabilidad ética las implicaciones prácticas de los modelos y datos que utiliza, y muestra apertura al diálogo interdisciplinario y tolerancia a la incertidumbre propia de los fenómenos reales.}
 \cur{MA-0880 Comunicación en las Ciencias, que incluye comunicación científica y ética profesional; MA-0881 y MA-0883 Práctica Profesional, que busca desarrollar autonomía y responsabilidad.}}
\mrow{Autonomía con responsabilidad personal, laboral y social}
{Identifica sus necesidades de actualización, capacitación y formación, durante su proceso formativo y en el ejercicio profesional, y busca los medios para cubrirlas por medios formales e informales, nacionales e internacionales, presenciales y en línea.}
{\qe{CAP06 / CAA06}{Mantiene un interés permanente por los desarrollos de la investigación científica y matemática, valorando la actualización como parte de su formación profesional.}
 \qe{CAP02 / CAA02}{Se comunica con claridad, colabora en equipos académicos y profesionales, y mantiene una disposición al autoaprendizaje y al liderazgo.}
 \cur{MA-0880 Comunicación en las Ciencias (búsqueda bibliográfica, gestión de fuentes y planificación de carrera).}
 \nt{como requisito de graduación se asiste a diez actividades académicas (coloquios, seminarios, conferencias, entre otras).}}
\mrow{Autonomía con responsabilidad personal, laboral y social}
{Evalúa su desempeño profesional con base en referentes de calidad, nacionales e internacionales, y la incidencia de sus decisiones en aspectos económicos, sociales, ambientales y de reducción de riesgos.}
{\qe{CAP04 / CAA04}{Ejerce pensamiento crítico, actúa con integridad ética y evalúa con rigor las prácticas científicas y académicas.}
 \qe{CAP03 / CAA03}{Demuestra conciencia social y participación ciudadana, con una visión sociocultural e integral de la sociedad en la que la ciencia y la matemática ocupan un lugar relevante.}
 \cur{MA-0647 Modelación Matemática, con validación, incertidumbre e interpretación (énfasis aplicada).}
 \nt{el perfil no menciona referentes de calidad nacionales e internacionales; el resultado se atiende de forma parcial.}}
\mrow{Autonomía con responsabilidad personal, laboral y social}
{Emprende proyectos profesionales para la generación de negocios y para el beneficio social con criterios de pertinencia, calidad, innovación y sustentabilidad.}
{\nt{el perfil no tiene un enunciado específico para este resultado. Se relaciona con el objetivo específico de formar una persona profesional que impacte de manera positiva el entorno productivo, público y privado, y se atiende de forma indirecta en las prácticas profesionales y, como opción, en cursos de otra disciplina.}
 \cur{MA-0881 y MA-0883 Práctica Profesional (modalidad académica o profesional); EC-1100 Introducción a la Economía (optativa de otra disciplina).}}
\mrow{Comunicación}
{Comunica, en varios lenguajes y formatos, y de forma pertinente, propositiva y clara, información de su campo profesional a los diferentes actores, equipos de trabajo y destinatarios de las acciones profesionales.}
{\qe{CPP04 / CPA04}{Participa activamente en discusiones científicas y divulga el conocimiento matemático de forma accesible a diferentes audiencias.}
 \qe{CAP02 / CAA02}{Se comunica con claridad, colabora en equipos académicos y profesionales, y mantiene una disposición al autoaprendizaje y al liderazgo.}
 \cur{MA-0880 Comunicación en las Ciencias.}}
\mrow{Comunicación}
{Se comunica correctamente en su lengua oficial y utiliza una lengua extranjera con el dominio requerido para el ejercicio de su profesión.}
{
\qe{CPP04 / CPA04}{Participa activamente en discusiones científicas y divulga el conocimiento matemático de forma accesible a diferentes audiencias.}
 \cur{LM-3039 Inglés para Estadística I; LM-3040 Inglés para Estadística II; MA-0880 Comunicación en las Ciencias.}}
\mrow{Comunicación}
{Utiliza tecnologías digitales para el manejo e interpretación de datos e información de forma apropiada a su nivel y su profesión.}
{\qe{CPP03 / CPA03}{Utiliza tecnologías y software científico para explorar, validar y resolver problemas matemáticos.}
 \cur{MA-0341 Matemática Computacional; MA-0501 Análisis Numérico I; CA-0204 Herramientas de Ciencia de Datos I (énfasis aplicada).}}
\mrow{Interacción profesional, cultural y social}
{Demuestra habilidades colaborativas y cooperativas en el campo profesional, cultural y social.}
{\qe{CAP02 / CAA02}{Se comunica con claridad, colabora en equipos académicos y profesionales, y mantiene una disposición al autoaprendizaje y al liderazgo.}
 \qe{CAP01 / CAA01}{Manifiesta respeto por la diversidad humana, matemática, científica y cultural, fomentando un ambiente inclusivo y de valoración mutua.}
 \cur{SR-I y SR-II Seminario de Realidad Nacional I y II.}
 \nt{como requisito de graduación se cumple el Trabajo Comunal Universitario (TCU).}}
\mrow{Interacción profesional, cultural y social}
{Lidera y colabora proactivamente en equipos de trabajo y en comunidades profesionales para el logro de objetivos y mejoramiento de la calidad de vida.}
{\qe{CAP02 / CAA02}{Se comunica con claridad, colabora en equipos académicos y profesionales, y mantiene una disposición al autoaprendizaje y al liderazgo.}
 \nt{se relaciona con el objetivo específico de incentivar una actitud positiva hacia el trabajo colaborativo con otros especialistas, en la búsqueda de soluciones que mejoren la calidad de vida y el entorno social.}
 \cur{MA-0881 y MA-0883 Práctica Profesional.}}
\mrow{Interacción profesional, cultural y social}
{Muestra respeto hacia la diversidad en todas sus manifestaciones y contribuye al bien común.}
{\qe{CAP01 / CAA01}{Manifiesta respeto por la diversidad humana, matemática, científica y cultural, fomentando un ambiente inclusivo y de valoración mutua.}
 \qe{CAP03 / CAA03}{Demuestra conciencia social y participación ciudadana, con una visión sociocultural e integral de la sociedad en la que la ciencia y la matemática ocupan un lugar relevante.}
 \nt{se relaciona con el objetivo específico de formar una persona profesional comprometida con el bien común y la justicia, y con el TCU.}}
\mrow{Interacción profesional, cultural y social}
{Participa en redes de colaboración que fortalezcan su campo profesional.}
{\qe{CPP04 / CPA04}{Participa activamente en discusiones científicas y divulga el conocimiento matemático de forma accesible a diferentes audiencias.}
 \qe{CAP06 / CAA06}{Mantiene un interés permanente por los desarrollos de la investigación científica y matemática, valorando la actualización como parte de su formación profesional.}
 \nt{la asistencia a diez actividades académicas (coloquios, seminarios y conferencias) como requisito de graduación favorece el contacto con la comunidad matemática.}}
\finmcesca

\textbf{Saberes disciplinarios y profesionales.} El perfil responde de manera directa al primer resultado mediante los enunciados declarativos CDP01 a CDP07 y CDA01 a CDA07, que abarcan los fundamentos matemáticos, el cálculo y el álgebra lineal, el análisis, la geometría, la topología y el álgebra abstracta (énfasis pura), o los principios y tecnologías de la matemática aplicada (énfasis aplicada). El dominio de herramientas tecnológicas se apoya en CDP02 / CDA02 y CPP03 / CPA03, y se desarrolla en MA-0341 y MA-0501. La cultura humanística y social se atiende en los cursos de Estudios Generales y en los seminarios de realidad nacional, junto con CAP03 / CAA03, aunque el perfil no declara un enunciado propio. El riesgo de desastres no se aborda de forma explícita; solo se identifican aportes indirectos, por lo que se señala como aspecto por fortalecer.

\textbf{Aplicación de conocimientos, resolución de problemas e innovación.} El pensamiento crítico y el rigor analítico se reflejan en CPP01 / CPA01, CAP04 / CAA04 y, en el énfasis pura, en CPP05, y se trabajan desde el bloque común (MA-0152, MA-0252) y en la secuencia de análisis. La resolución de problemas en distintos contextos se plasma en CAP05 / CAA05 y CPP02 / CPA02; en el énfasis aplicada se refuerza con CPA05, CDA05 y los cursos de modelación y de ciencia de datos, y en ambos énfasis con la Práctica Profesional. La toma de decisiones con base en datos se atiende con mayor fuerza en el énfasis aplicada; en el énfasis pura se aborda de forma más limitada.

\textbf{Autonomía con responsabilidad personal, laboral y social.} La responsabilidad y la ética profesional se expresan en CAP04 / CAA04 y CAA07, y se trabajan en MA-0880 y en la Práctica Profesional. La actualización permanente se refleja en CAP06 / CAA06 y CAP02 / CAA02, y se promueve con la asistencia a actividades académicas. La evaluación del desempeño se atiende de manera parcial, y el emprendimiento no tiene un enunciado propio en el perfil, aunque se relaciona con el objetivo de incidir positivamente en el entorno productivo y se aborda de forma indirecta en las prácticas profesionales.

\textbf{Comunicación.} La comunicación en distintos lenguajes y formatos se refleja en CPP04 / CPA04 y CAP02 / CAA02, y se desarrolla en MA-0880. En cuanto a la lengua extranjera, el perfil no declara un enunciado específico, pero el plan incluye LM-3039 y LM-3040 para ambos énfasis. El uso de tecnologías digitales para el manejo de datos se sustenta en CPP03 / CPA03 y en los cursos de matemática computacional, análisis numérico y ciencia de datos.

\textbf{Interacción profesional, cultural y social.} Los enunciados CAP01 / CAA01, CAP02 / CAA02, CAP03 / CAA03 y CPP04 / CPA04 respaldan el trabajo colaborativo, el liderazgo, el respeto a la diversidad, el compromiso social y la participación en discusiones científicas. Esto se complementa con los objetivos específicos de la carrera sobre trabajo colaborativo y bien común, con los seminarios de realidad nacional, con el TCU y con la asistencia a actividades académicas.

\subsection*{Licenciatura en Matemática (CINE 6B), ambos énfasis}

La Licenciatura mantiene los enunciados del perfil del bachillerato y agrega solo tres rasgos propios: CDP08 / CDA08, CP05 y CA06. Por ello, en la Tabla~\ref{res:mcesca-lic} se indica, para cada resultado del nivel de licenciatura, qué se conserva del bachillerato y qué agrega este grado.

\mcesca{Análisis del plan de estudios respecto del MCESCA: nivel Licenciatura}{res:mcesca-lic}
\mrow{Saberes disciplinarios y profesionales}
{Comprende en forma crítica el cuerpo conceptual, metodológico, procedimental y normativo, que le permite el ejercicio de su profesión en el contexto nacional e internacional.}
{\textit{Agrega la licenciatura:}\par
 \qe{CDP08}{Conoce estructuras y resultados de nivel intermedio para algunas áreas medulares de especialización en matemática pura.}
 \qe{CDA08}{Conoce estructuras y resultados de nivel intermedio para algunas áreas medulares de especialización en matemática aplicada.}
 \qe{CP05}{Aplica pensamiento operacional, abstracto y formal, con fuentes transversales de varias áreas, para formular, analizar y resolver problemas matemáticos.}
 \cur{MA-0982 y MA-1081 Seminario de Estudios Dirigidos en Matemática Pura I y II; MA-0984 y MA-1082 Seminario de Estudios Dirigidos en Matemática Aplicada I y II; OPT-001 u OPT-002 y OPT-003 (optativas).}}
\mrow{Saberes disciplinarios y profesionales}
{Demuestra conocimientos como cultura humanística y derechos fundamentales, ambiente, entre otros que le brindan una visión amplia de su profesión, de las relaciones de esta con otros saberes y de su entorno.}
{Igual que en el bachillerato (CAP03 / CAA03 y cursos de Estudios Generales y de Realidad Nacional). La licenciatura no agrega elementos propios, y el perfil no tiene un enunciado específico sobre derechos fundamentales ni ambiente.}
\mrow{Saberes disciplinarios y profesionales}
{Muestra conocimiento sobre el riesgo de desastres en el campo de su especialidad y propone soluciones para reducirlo.}
{Igual que en el bachillerato: el perfil no tiene un enunciado explícito y se señala como aspecto por fortalecer. Las optativas generales (OPT-003) permiten profundizar, por ejemplo, con MA-0794 Tópicos de Modelación.}
\mrow{Saberes disciplinarios y profesionales}
{Identifica oportunidades y riesgos para la innovación y adaptación de conocimientos y tecnologías para resolver problemas.}
{\textit{Agrega la licenciatura:}\par
 \qe{CP05}{Aplica pensamiento operacional, abstracto y formal, con fuentes transversales de varias áreas, para formular, analizar y resolver problemas matemáticos.}
 \qe{CA06}{Muestra iniciativa y autonomía para formular preguntas, explorar problemas matemáticos y profundizar en ellos más allá de los procedimientos conocidos.}}
\mrow{Aplicación de conocimientos, resolución de problemas e innovación}
{Demuestra pensamiento crítico, actitud investigativa y rigor analítico en el planteamiento y la resolución de problemas complejos.}
{Se conservan CPP01 / CPA01, CAP04 / CAA04 y CPP05 del bachillerato.\par
 \textit{Agrega la licenciatura:}\par
 \qe{CA06}{Muestra iniciativa y autonomía para formular preguntas, explorar problemas matemáticos y profundizar en ellos más allá de los procedimientos conocidos.}
 \cur{Seminarios de Estudios Dirigidos (MA-0982, MA-1081, MA-0984, MA-1082), cuyo objetivo es desarrollar de manera autónoma actividades asociadas a la investigación.}}
\mrow{Aplicación de conocimientos, resolución de problemas e innovación}
{Aplica los conocimientos de su disciplina en la elaboración, fundamentación y defensa de argumentos para prevenir y resolver problemas complejos en su campo profesional, identificando y aplicando innovaciones.}
{Se conservan CAP05 / CAA05 y CPP02 / CPA02 del bachillerato.\par
 \textit{Agrega la licenciatura:}\par
 \qe{CP05}{Aplica pensamiento operacional, abstracto y formal, con fuentes transversales de varias áreas, para formular, analizar y resolver problemas matemáticos.}
 \nt{el Trabajo Final de Graduación (TFG) y los seminarios de estudios dirigidos.}}
\mrow{Aplicación de conocimientos, resolución de problemas e innovación}
{Propone e implementa nuevos procedimientos y metodologías aplicables a la solución de problemas complejos y mejora de su campo profesional.}
{Se conservan CPP02 / CPA02 y CPA05 del bachillerato.\par
 \textit{Agrega la licenciatura:}\par
 \qe{CA06}{Muestra iniciativa y autonomía para formular preguntas, explorar problemas matemáticos y profundizar en ellos más allá de los procedimientos conocidos.}
 \cur{Seminarios de Estudios Dirigidos de cada énfasis.}}
\mrow{Aplicación de conocimientos, resolución de problemas e innovación}
{Toma decisiones profesionales con base en fundamentos teóricos, datos e información pertinente, válida y confiable.}
{Se conservan CDA05 y CAA07 del bachillerato (énfasis aplicada).\par
 \textit{Agrega la licenciatura:} CDP08 / CDA08, que aportan fundamentos teóricos de nivel intermedio (enunciados transcritos en la primera fila de esta tabla).}
\mrow{Aplicación de conocimientos, resolución de problemas e innovación}
{Demuestra destreza y habilidad en la selección, uso y adaptación de herramientas metodológicas, tecnológicas, equipos especializados y en la lectura e interpretación de datos, pertinentes al contexto de su ejercicio profesional.}
{Se conservan CDP02 / CDA02 y CPP03 / CPA03 del bachillerato.\par
 \textit{Agrega la licenciatura:} CA06 (iniciativa y autonomía, enunciado transcrito en la fila anterior de saberes) y los Seminarios de Estudios Dirigidos de cada énfasis.}
\mrow{Autonomía con responsabilidad personal, laboral y social}
{Demuestra responsabilidad y autonomía profesional en la determinación de resultados personales y laborales de acuerdo a su función tomando como referencias las normativas legales y éticas de su campo profesional.}
{Se conservan CAP04 / CAA04 y CAA07 del bachillerato.\par
 \textit{Agrega la licenciatura:}\par
 \qe{CA06}{Muestra iniciativa y autonomía para formular preguntas, explorar problemas matemáticos y profundizar en ellos más allá de los procedimientos conocidos.}
 \cur{Seminarios de Estudios Dirigidos de cada énfasis.}}
\mrow{Autonomía con responsabilidad personal, laboral y social}
{Identifica sus necesidades de actualización, capacitación y formación, durante su proceso formativo y en el ejercicio profesional, y busca los medios para cubrirlas por medios formales e informales, nacionales e internacionales, presenciales y en línea.}
{Igual que en el bachillerato (CAP06 / CAA06 y CAP02 / CAA02). La licenciatura agrega CA06 y la asistencia a cinco actividades académicas distintas de las empleadas para el bachillerato.}
\mrow{Autonomía con responsabilidad personal, laboral y social}
{Evalúa su desempeño profesional con base en referentes de calidad, nacionales e internacionales, y la incidencia de sus decisiones en los aspectos humanos, sociales, ambientales y de reducción de riesgos.}
{Igual que en el bachillerato (CAP04 / CAA04 y CAP03 / CAA03), con la misma salvedad sobre los referentes de calidad. La licenciatura no agrega elementos propios.}
\mrow{Autonomía con responsabilidad personal, laboral y social}
{Emprende proyectos profesionales para la generación de negocios y para el beneficio social con criterios de pertinencia, calidad, innovación y sustentabilidad.}
{Igual que en el bachillerato: el perfil no tiene un enunciado específico y se señala como aspecto por fortalecer. La licenciatura no agrega elementos propios.}
\mrow{Comunicación}
{Comunica a diversos públicos, información de su campo profesional, en varios lenguajes y formatos de manera asertiva, clara, rigurosa y precisa, con el uso apropiado de recursos tecnológicos.}
{Se conservan CPP04 / CPA04 y CAP02 / CAA02 del bachillerato, y el curso MA-0880 Comunicación en las Ciencias. La licenciatura agrega, como requisito de graduación, el Trabajo Final de Graduación (TFG).}
\mrow{Comunicación}
{Se comunica correctamente en su lengua oficial y utiliza una lengua extranjera con el dominio requerido para el ejercicio de su profesión.}
{Igual que en el bachillerato (LM-3039 Inglés para Estadística I; LM-3040 Inglés para Estadística II; MA-0880 Comunicación en las Ciencias). El perfil no tiene un enunciado específico y la licenciatura no agrega elementos propios.}
\mrow{Comunicación}
{Utiliza tecnologías digitales para el manejo e interpretación de datos e información de forma apropiada a su nivel y su profesión.}
{Igual que en el bachillerato (CPP03 / CPA03 y los cursos de matemática computacional, análisis numérico y ciencia de datos).}
\mrow{Interacción profesional, cultural y social}
{Demuestra habilidades colaborativas y cooperativas en el campo profesional, cultural y social.}
{Igual que en el bachillerato (CAP01 / CAA01 y CAP02 / CAA02). La licenciatura no agrega elementos propios.}
\mrow{Interacción profesional, cultural y social}
{Lidera y colabora proactivamente en equipos de trabajo y en comunidades profesionales para el logro de objetivos y mejoramiento de la calidad de vida.}
{Igual que en el bachillerato (CAP02 / CAA02 y el objetivo específico sobre trabajo colaborativo). La licenciatura no agrega elementos propios.}
\mrow{Interacción profesional, cultural y social}
{Muestra respeto hacia la diversidad en todas sus manifestaciones y contribuye al bien común.}
{Igual que en el bachillerato (CAP01 / CAA01 y CAP03 / CAA03). La licenciatura no agrega elementos propios.}
\mrow{Interacción profesional, cultural y social}
{Participa en redes de colaboración que fortalezcan su campo profesional.}
{Igual que en el bachillerato (CPP04 / CPA04 y CAP06 / CAA06). La licenciatura agrega la asistencia a cinco actividades académicas adicionales y los Seminarios de Estudios Dirigidos.}
\finmcesca

\textbf{Saberes disciplinarios y profesionales.} La licenciatura profundiza el primer resultado con CDP08 / CDA08, que incorporan estructuras y resultados de nivel intermedio de áreas medulares de cada énfasis, y con CP05, que añade el trabajo con fuentes transversales de varias áreas. Los seminarios de estudios dirigidos y las optativas de núcleo y generales concretan esa profundización. La cultura humanística y el riesgo de desastres siguen siendo los del bachillerato.

\textbf{Aplicación de conocimientos, resolución de problemas e innovación.} El grado agrega CA06 (iniciativa y autonomía para explorar problemas más allá de los procedimientos conocidos) y CP05, que respaldan el tratamiento de problemas complejos y la propuesta de nuevos procedimientos. Estos rasgos se desarrollan en los seminarios de estudios dirigidos, orientados a la investigación autónoma, y en el TFG.

\textbf{Autonomía con responsabilidad personal, laboral y social.} CA06 y los seminarios de estudios dirigidos refuerzan la autonomía profesional. Las demás exigencias del nivel (actualización, evaluación del desempeño y emprendimiento) se atienden como en el bachillerato, con la asistencia adicional a cinco actividades académicas.

\textbf{Comunicación.} La licenciatura mantiene los enunciados y cursos del bachillerato y agrega la comunicación de resultados mediante el TFG. La lengua extranjera sigue atendiéndose con los cursos de inglés del bloque común.

\textbf{Interacción profesional, cultural y social.} Este descriptor se atiende con los mismos enunciados del bachillerato; la licenciatura suma la participación en cinco actividades académicas adicionales, que fortalecen el contacto con redes de colaboración.

\textbf{Síntesis.} El perfil de egreso no tiene enunciados específicos para tres resultados del MCESCA: riesgo de desastres, emprendimiento y lengua extranjera. Este último se atiende mediante los cursos de inglés del plan; los otros dos se identifican como aspectos por fortalecer en futuras revisiones del plan.
\endgroup

%%%%%%%%%%%%%%%%%%%%%%%%%%%%%%%%%%%%%%%%%%%%%%%%%%
\section{Campo de inserción profesional}

Los principales campos de inserción identificados para las personas graduadas son:
\begin{itemize}
  \item \textbf{Investigación académica} (teórica y aplicada), en universidades, centros de investigación y redes científicas nacionales e internacionales.
  \item \textbf{Docencia universitaria}, en programas de pregrado y posgrado en matemática y disciplinas afines.
  \item \textbf{Modelación matemática}, en ingeniería, física, biología, epidemiología, ciencias ambientales y otras áreas.
  \item \textbf{Análisis de datos y ciencia de datos}, en empresas, instituciones financieras, organismos públicos y proyectos de investigación interdisciplinaria.
  \item \textbf{Sector productivo e industrial}, en consultoría cuantitativa, optimización, criptografía, simulación y apoyo a la innovación tecnológica.
  \item \textbf{Trabajo interdisciplinario}, como puente entre equipos de distintas disciplinas en proyectos científicos, tecnológicos o de política pública.
\end{itemize}

%%%%%%%%%%%%%%%%%%%%%%%%%%%%%%%%%%%%%%%%%%%%%%%%%%
\section{Requisitos de ingreso}

El ingreso a la carrera se rige por los requisitos y procedimientos establecidos por la Universidad de Costa Rica para la admisión a carreras de grado, mediante el proceso institucional que incluye la aplicación de la Prueba de Aptitud Académica (PAA).

Además, las personas aspirantes deben obtener una calificación mínima en la Prueba de Habilidades Cuantitativas (PHC), de acuerdo con lo establecido anualmente en la resolución correspondiente de la Vicerrectoría de Docencia.

Habilidades y características deseables
\begin{itemize}
    \item Habilidad para el razonamiento numérico.
    \item Interés por la matemática.
    \item Capacidades de razonamiento abstracto.
    \item Disposición para el estudio y la investigación científica.
    \item Habilidad para resolver problemas.
    \item Inclinación por el conocimiento.
    \item Dominio de elementos básicos de la matemática.
\end{itemize}

La carrera no contempla cursos nivelatorios.

%%%%%%%%%%%%%%%%%%%%%%%%%%%%%%%%%%%%%%%%%%%%%%%%%%
\section{Requisitos de graduación}

\subsection*{Bachillerato en Matemática (ambos énfasis)}
\begin{itemize}
  \item Aprobar el plan de estudios de los ciclos I a VIII: 130 créditos en el énfasis en Matemática Pura o 131 créditos en el énfasis en Matemática Aplicada. Esto incluye la Práctica Profesional del último ciclo (modalidad académica o profesional).
  \item Asistir, a lo largo de la carrera, a diez (10) actividades académicas válidas tales como charlas del Coloquio del Departamento de Matemática, seminarios, conferencias, ciclos especiales u otras actividades que disponga y facilite la Dirección del Departamento de Matemática.
  \item Cumplir los requisitos que establece la normativa universitaria, tales como el Trabajo Comunal Universitario (TCU).
\end{itemize}

\subsection*{Licenciatura en Matemática (ambos énfasis)}
\begin{itemize}
  \item Aprobar los 30 créditos del plan de licenciatura (ciclos IX y X): dos seminarios de estudios dirigidos del énfasis y cursos optativos.
  \item Asistir a cinco (5) actividades académicas válidas, distintas de las empleadas para el bachillerato, tales como charlas del Coloquio del Departamento de Matemática, seminarios, conferencias, ciclos especiales u otras actividades que disponga y facilite la Dirección del Departamento de Matemática.
  \item Cumplir los requisitos de la normativa universitaria para la obtención del grado de licenciatura, tales como el Trabajo Final de Graduación (TFG).
\end{itemize}

%%%%%%%%%%%%%%%%%%%%%%%%%%%%%%%%%%%%%%%%%%%%%%%%%%
\section{Listado de cursos}

Las tablas siguientes presentan el nombre, los créditos y el ciclo de cada curso que conforma la estructura curricular. Los cursos con sigla MA y CA son de la Escuela de Matemática; sin embargo los primeros corresponden al Departamento de Matemática y los segundos, al Departamento de Ciencias Actuariales, ambos adscritos a dicha Escuela. En el presente plan, el Departamento de Ciencias Actuariales ofrece cursos al Departamento de Matemática. Por su parte, los cursos optativos de otras disciplinas (OPT-) son ofrecidos por otras unidades académicas de la Universidad de Costa Rica.

\setcounter{table}{0}
% Generado por scripts/generate_resumen_ejecutivo.py. No editar a mano.

\subsection*{Cursos de la estructura curricular}

\begingroup
\small
\arrayrulecolor{Azul-UCR!80}%
\renewcommand{\arraystretch}{1.2}%
\rowcolors{2}{white}{Celeste-UCR!12}%
\begin{longtable}{|c|l|>{\raggedright\arraybackslash}p{8.6cm}|c|}
\caption{Bloque común (ciclos I a III)}\label{res:malla-comun}\\
\hline
\rowcolor{Azul-UCR}
\textcolor{white}{\textbf{Ciclo}} & \textcolor{white}{\textbf{Sigla}} & \textcolor{white}{\textbf{Nombre del curso}} & \textcolor{white}{\textbf{Créditos}} \\ \hline
\endfirsthead
\hline
\rowcolor{Azul-UCR}
\textcolor{white}{\textbf{Ciclo}} & \textcolor{white}{\textbf{Sigla}} & \textcolor{white}{\textbf{Nombre del curso}} & \textcolor{white}{\textbf{Créditos}} \\ \hline
\endhead
\hline
\endfoot
1 & EF- & Actividad Deportiva & 0 \\ \hline
1 & EG-I & Curso Integrado de Humanidades I & 6 \\ \hline
1 & MA-0151 & Fundamentos de Álgebra, Trigonometría y Geometría Analítica & 3 \\ \hline
1 & MA-0152 & Matemática Exploratoria & 3 \\ \hline
1 & RP- & Repertorio & 3 \\ \hline
\rowcolor{Verde-UCR!28}\multicolumn{4}{|l|}{\textit{Créditos del ciclo 1: 15}} \\ \hline
2 & EG- & Curso de Arte & 2 \\ \hline
2 & EG-II & Curso Integrado de Humanidades II & 6 \\ \hline
2 & LM-3039 & Inglés para Estadística I & 3 \\ \hline
2 & MA-0251 & Cálculo en una Variable & 4 \\ \hline
2 & MA-0252 & Introducción a las Demostraciones & 3 \\ \hline
\rowcolor{Verde-UCR!28}\multicolumn{4}{|l|}{\textit{Créditos del ciclo 2: 18}} \\ \hline
3 & SR-I & Seminario de Realidad Nacional I & 2 \\ \hline
3 & MA-0361 & Álgebra Lineal I & 4 \\ \hline
3 & LM-3040 & Inglés para Estadística II & 3 \\ \hline
3 & MA-0351 & Cálculo en Varias Variables & 4 \\ \hline
3 & MA-0341 & Matemática Computacional & 4 \\ \hline
\rowcolor{Verde-UCR!28}\multicolumn{4}{|l|}{\textit{Créditos del ciclo 3: 17}} \\ \hline
\rowcolor{Amarillo-UCR!45}\multicolumn{4}{|r|}{\textbf{Total créditos bloque común: 50}} \\ \hline
\end{longtable}
\endgroup

\begingroup
\small
\arrayrulecolor{Azul-UCR!80}%
\renewcommand{\arraystretch}{1.2}%
\rowcolors{2}{white}{Celeste-UCR!12}%
\begin{longtable}{|c|l|>{\raggedright\arraybackslash}p{8.6cm}|c|}
\caption{Énfasis en Matemática pura (ciclos IV a VIII)}\label{res:malla-pura}\\
\hline
\rowcolor{Azul-UCR}
\textcolor{white}{\textbf{Ciclo}} & \textcolor{white}{\textbf{Sigla}} & \textcolor{white}{\textbf{Nombre del curso}} & \textcolor{white}{\textbf{Créditos}} \\ \hline
\endfirsthead
\hline
\rowcolor{Azul-UCR}
\textcolor{white}{\textbf{Ciclo}} & \textcolor{white}{\textbf{Sigla}} & \textcolor{white}{\textbf{Nombre del curso}} & \textcolor{white}{\textbf{Créditos}} \\ \hline
\endhead
\hline
\endfoot
4 & SR-II & Seminario de Realidad Nacional II & 2 \\ \hline
4 & MA-0461 & Álgebra Lineal II & 4 \\ \hline
4 & MA-0451 & Principios de Análisis I & 4 \\ \hline
4 & MA-0471 & Introducción a la Geometría Diferencial & 3 \\ \hline
4 & MA-0496 & Teoría de Números & 3 \\ \hline
\rowcolor{Verde-UCR!28}\multicolumn{4}{|l|}{\textit{Créditos del ciclo 4: 16}} \\ \hline
5 & MA-0571 & Álgebra Abstracta I & 4 \\ \hline
5 & MA-0555 & Principios de Análisis II & 4 \\ \hline
5 & MA-0556 & Ecuaciones Diferenciales & 4 \\ \hline
5 & OPT- & Optativa de Otra Disciplina & 3 \\ \hline
\rowcolor{Verde-UCR!28}\multicolumn{4}{|l|}{\textit{Créditos del ciclo 5: 15}} \\ \hline
6 & MA-0661 & Álgebra Abstracta II & 5 \\ \hline
6 & MA-0635 & Topología General & 4 \\ \hline
6 & MA-0501 & Análisis Numérico I & 4 \\ \hline
6 & OPT- & Optativa de Otra Disciplina & 3 \\ \hline
\rowcolor{Verde-UCR!28}\multicolumn{4}{|l|}{\textit{Créditos del ciclo 6: 16}} \\ \hline
7 & MA-0716 & Análisis Real I & 5 \\ \hline
7 & MA-0735 & Análisis Complejo & 4 \\ \hline
7 & MA-0880 & Comunicación en las Ciencias & 2 \\ \hline
7 & OPT-001 & Optativa Núcleo en Matemática Pura & 4 \\ \hline
\rowcolor{Verde-UCR!28}\multicolumn{4}{|l|}{\textit{Créditos del ciclo 7: 15}} \\ \hline
8 & MA-0881 & Práctica Profesional en Matemática Pura & 6 \\ \hline
8 & OPT-001 & Optativa Núcleo en Matemática Pura & 8 \\ \hline
8 & OPT-003 & Optativa General en Matemática & 4 \\ \hline
\rowcolor{Verde-UCR!28}\multicolumn{4}{|l|}{\textit{Créditos del ciclo 8: 18}} \\ \hline
\rowcolor{Amarillo-UCR!45}\multicolumn{4}{|r|}{\textbf{Total créditos bachillerato pura: 130}} \\ \hline
\end{longtable}
\endgroup

\begingroup
\small
\arrayrulecolor{Azul-UCR!80}%
\renewcommand{\arraystretch}{1.2}%
\rowcolors{2}{white}{Celeste-UCR!12}%
\begin{longtable}{|c|l|>{\raggedright\arraybackslash}p{8.6cm}|c|}
\caption{Énfasis en Matemática aplicada (ciclos IV a VIII)}\label{res:malla-aplicada}\\
\hline
\rowcolor{Azul-UCR}
\textcolor{white}{\textbf{Ciclo}} & \textcolor{white}{\textbf{Sigla}} & \textcolor{white}{\textbf{Nombre del curso}} & \textcolor{white}{\textbf{Créditos}} \\ \hline
\endfirsthead
\hline
\rowcolor{Azul-UCR}
\textcolor{white}{\textbf{Ciclo}} & \textcolor{white}{\textbf{Sigla}} & \textcolor{white}{\textbf{Nombre del curso}} & \textcolor{white}{\textbf{Créditos}} \\ \hline
\endhead
\hline
\endfoot
4 & MA-0461 & Álgebra Lineal II & 4 \\ \hline
4 & MA-0451 & Principios de Análisis I & 4 \\ \hline
4 & CA-0721 & Probabilidad & 4 \\ \hline
4 & CA-0204 & Herramientas de Ciencia de Datos I & 4 \\ \hline
\rowcolor{Verde-UCR!28}\multicolumn{4}{|l|}{\textit{Créditos del ciclo 4: 16}} \\ \hline
5 & SR-II & Seminario de Realidad Nacional II & 2 \\ \hline
5 & MA-0555 & Principios de Análisis II & 4 \\ \hline
5 & MA-0556 & Ecuaciones Diferenciales & 4 \\ \hline
5 & CA-0303 & Estadística Actuarial I & 4 \\ \hline
5 & OPT- & Optativa de Otra Disciplina & 3 \\ \hline
\rowcolor{Verde-UCR!28}\multicolumn{4}{|l|}{\textit{Créditos del ciclo 5: 17}} \\ \hline
6 & MA-0501 & Análisis Numérico I & 4 \\ \hline
6 & MA-0635 & Topología General & 4 \\ \hline
6 & MA-0647 & Modelación Matemática & 4 \\ \hline
6 & OPT- & Optativa de Otra Disciplina & 3 \\ \hline
\rowcolor{Verde-UCR!28}\multicolumn{4}{|l|}{\textit{Créditos del ciclo 6: 15}} \\ \hline
7 & MA-0716 & Análisis Real I & 5 \\ \hline
7 & MA-0735 & Análisis Complejo & 4 \\ \hline
7 & CA-0411 & Análisis de Datos I & 4 \\ \hline
7 & MA-0880 & Comunicación en las Ciencias & 2 \\ \hline
\rowcolor{Verde-UCR!28}\multicolumn{4}{|l|}{\textit{Créditos del ciclo 7: 15}} \\ \hline
8 & MA-0883 & Práctica Profesional en Matemática Aplicada & 6 \\ \hline
8 & OPT-002 & Optativa Núcleo en Matemática Aplicada & 8 \\ \hline
8 & OPT-003 & Optativa General en Matemática & 4 \\ \hline
\rowcolor{Verde-UCR!28}\multicolumn{4}{|l|}{\textit{Créditos del ciclo 8: 18}} \\ \hline
\rowcolor{Amarillo-UCR!45}\multicolumn{4}{|r|}{\textbf{Total créditos bachillerato aplicada: 131}} \\ \hline
\end{longtable}
\endgroup

\begingroup
\small
\arrayrulecolor{Azul-UCR!80}%
\renewcommand{\arraystretch}{1.2}%
\rowcolors{2}{white}{Celeste-UCR!12}%
\begin{longtable}{|c|l|>{\raggedright\arraybackslash}p{8.6cm}|c|}
\caption{Énfasis en Matemática pura (licenciatura, ciclos IX y X)}\label{res:lic-pura}\\
\hline
\rowcolor{Azul-UCR}
\textcolor{white}{\textbf{Ciclo}} & \textcolor{white}{\textbf{Sigla}} & \textcolor{white}{\textbf{Nombre del curso}} & \textcolor{white}{\textbf{Créditos}} \\ \hline
\endfirsthead
\hline
\rowcolor{Azul-UCR}
\textcolor{white}{\textbf{Ciclo}} & \textcolor{white}{\textbf{Sigla}} & \textcolor{white}{\textbf{Nombre del curso}} & \textcolor{white}{\textbf{Créditos}} \\ \hline
\endhead
\hline
\endfoot
9 & MA-0982 & Seminario de Estudios Dirigidos en Matemática Pura I & 7 \\ \hline
9 & OPT-001 & Optativa Núcleo en Matemática Pura & 4 \\ \hline
9 & OPT-003 & Optativa General en Matemática & 4 \\ \hline
\rowcolor{Verde-UCR!28}\multicolumn{4}{|l|}{\textit{Créditos del ciclo 9: 15}} \\ \hline
10 & MA-1081 & Seminario de Estudios Dirigidos en Matemática Pura II & 7 \\ \hline
10 & OPT-003 & Optativa General en Matemática & 8 \\ \hline
\rowcolor{Verde-UCR!28}\multicolumn{4}{|l|}{\textit{Créditos del ciclo 10: 15}} \\ \hline
\rowcolor{Amarillo-UCR!45}\multicolumn{4}{|r|}{\textbf{Total créditos licenciatura pura: 30}} \\ \hline
\end{longtable}
\endgroup

\begingroup
\small
\arrayrulecolor{Azul-UCR!80}%
\renewcommand{\arraystretch}{1.2}%
\rowcolors{2}{white}{Celeste-UCR!12}%
\begin{longtable}{|c|l|>{\raggedright\arraybackslash}p{8.6cm}|c|}
\caption{Énfasis en Matemática aplicada (licenciatura, ciclos IX y X)}\label{res:lic-aplicada}\\
\hline
\rowcolor{Azul-UCR}
\textcolor{white}{\textbf{Ciclo}} & \textcolor{white}{\textbf{Sigla}} & \textcolor{white}{\textbf{Nombre del curso}} & \textcolor{white}{\textbf{Créditos}} \\ \hline
\endfirsthead
\hline
\rowcolor{Azul-UCR}
\textcolor{white}{\textbf{Ciclo}} & \textcolor{white}{\textbf{Sigla}} & \textcolor{white}{\textbf{Nombre del curso}} & \textcolor{white}{\textbf{Créditos}} \\ \hline
\endhead
\hline
\endfoot
9 & MA-0984 & Seminario de Estudios Dirigidos en Matemática Aplicada I & 7 \\ \hline
9 & OPT-002 & Optativa Núcleo en Matemática Aplicada & 4 \\ \hline
9 & OPT-003 & Optativa General en Matemática & 4 \\ \hline
\rowcolor{Verde-UCR!28}\multicolumn{4}{|l|}{\textit{Créditos del ciclo 9: 15}} \\ \hline
10 & MA-1082 & Seminario de Estudios Dirigidos en Matemática Aplicada II & 7 \\ \hline
10 & OPT-003 & Optativa General en Matemática & 8 \\ \hline
\rowcolor{Verde-UCR!28}\multicolumn{4}{|l|}{\textit{Créditos del ciclo 10: 15}} \\ \hline
\rowcolor{Amarillo-UCR!45}\multicolumn{4}{|r|}{\textbf{Total créditos licenciatura aplicada: 30}} \\ \hline
\end{longtable}
\endgroup





\subsection*{Cursos optativos de Matemática}

\begingroup
\small
\arrayrulecolor{Azul-UCR!80}%
\renewcommand{\arraystretch}{1.2}%
\rowcolors{2}{white}{Celeste-UCR!12}%
\begin{longtable}{|l|>{\raggedright\arraybackslash}p{9.4cm}|c|}
\caption{Optativas de núcleo en Matemática Pura (OPT-001)}\label{res:opt001}\\
\hline
\rowcolor{Azul-UCR}
\textcolor{white}{\textbf{Sigla}} & \textcolor{white}{\textbf{Nombre del curso}} & \textcolor{white}{\textbf{Créditos}} \\ \hline
\endfirsthead
\hline
\rowcolor{Azul-UCR}
\textcolor{white}{\textbf{Sigla}} & \textcolor{white}{\textbf{Nombre del curso}} & \textcolor{white}{\textbf{Créditos}} \\ \hline
\endhead
\hline
\endfoot
MA-0526 & Combinatoria & 4 \\ \hline
MA-0528 & Sistemas Dinámicos & 4 \\ \hline
MA-0838 & Geometría Diferencial & 4 \\ \hline
MA-0837 & Geometría Algebraica & 4 \\ \hline
MA-0734 & Optimización & 4 \\ \hline
MA-0756 & Ecuaciones Diferenciales Parciales & 5 \\ \hline
MA-0781 & Lógica & 5 \\ \hline
MA-0783 & Integración & 4 \\ \hline
MA-0788 & Análisis Real II & 4 \\ \hline
MA-0835 & Topología Algebraica & 4 \\ \hline
MA-0834 & Análisis Funcional & 4 \\ \hline
MA-0863 & Teoría Analítica de Números & 4 \\ \hline
MA-0862 & Análisis Armónico & 4 \\ \hline
MA-0861 & Teoría de Modelos & 4 \\ \hline
MA-0841 & Probabilidad & 4 \\ \hline
MA-0849 & Teoría de Módulos & 5 \\ \hline
MA-0858 & Teoría de Conjuntos & 4 \\ \hline
MA-0876 & Teoría Algebraica de Números & 4 \\ \hline
MA-0889 & Álgebra Conmutativa & 4 \\ \hline
MA-0920 & Ecuaciones en Derivadas Parciales Numéricas & 4 \\ \hline
\end{longtable}
\endgroup

\begingroup
\small
\arrayrulecolor{Azul-UCR!80}%
\renewcommand{\arraystretch}{1.2}%
\rowcolors{2}{white}{Celeste-UCR!12}%
\begin{longtable}{|l|>{\raggedright\arraybackslash}p{9.4cm}|c|}
\caption{Optativas de núcleo en Matemática Aplicada (OPT-002)}\label{res:opt002}\\
\hline
\rowcolor{Azul-UCR}
\textcolor{white}{\textbf{Sigla}} & \textcolor{white}{\textbf{Nombre del curso}} & \textcolor{white}{\textbf{Créditos}} \\ \hline
\endfirsthead
\hline
\rowcolor{Azul-UCR}
\textcolor{white}{\textbf{Sigla}} & \textcolor{white}{\textbf{Nombre del curso}} & \textcolor{white}{\textbf{Créditos}} \\ \hline
\endhead
\hline
\endfoot
CA-0305 & Herramientas de Ciencia de Datos II & 4 \\ \hline
CA-0307 & Estadística Actuarial II & 4 \\ \hline
CA-0414 & Modelos Lineales y de Sobrevivencia & 4 \\ \hline
CA-0512 & Análisis de Datos II & 4 \\ \hline
CA-0602 & Introducción a la Optimización & 4 \\ \hline
CA-0415 & Series de Tiempo & 4 \\ \hline
MA-0477 & Ecuaciones Diferenciales Parciales Aplicadas & 4 \\ \hline
\end{longtable}
\endgroup

\begingroup
\small
\arrayrulecolor{Azul-UCR!80}%
\renewcommand{\arraystretch}{1.2}%
\rowcolors{2}{white}{Celeste-UCR!12}%
\begin{longtable}{|l|>{\raggedright\arraybackslash}p{9.4cm}|c|}
\caption{Optativas generales en Matemática (OPT-003)}\label{res:opt003}\\
\hline
\rowcolor{Azul-UCR}
\textcolor{white}{\textbf{Sigla}} & \textcolor{white}{\textbf{Nombre del curso}} & \textcolor{white}{\textbf{Créditos}} \\ \hline
\endfirsthead
\hline
\rowcolor{Azul-UCR}
\textcolor{white}{\textbf{Sigla}} & \textcolor{white}{\textbf{Nombre del curso}} & \textcolor{white}{\textbf{Créditos}} \\ \hline
\endhead
\hline
\endfoot
CA-0305 & Herramientas de Ciencia de Datos II & 4 \\ \hline
% CA-0303 & Estadística Actuarial I & 4 \\ \hline
CA-0307 & Estadística Actuarial II & 4 \\ \hline
CA-0414 & Modelos Lineales y de Sobrevivencia & 4 \\ \hline
CA-0512 & Análisis de Datos II & 4 \\ \hline
CA-0602 & Introducción a la Optimización & 4 \\ \hline
CA-0415 & Series de Tiempo & 4 \\ \hline
MA-0477 & Ecuaciones Diferenciales Parciales Aplicadas & 4 \\ \hline
MA-0526 & Combinatoria & 4 \\ \hline
MA-0528 & Sistemas Dinámicos & 4 \\ \hline
MA-0838 & Geometría Diferencial & 4 \\ \hline
MA-0837 & Geometría Algebraica & 4 \\ \hline
MA-0734 & Optimización & 4 \\ \hline
MA-0756 & Ecuaciones Diferenciales Parciales & 5 \\ \hline
MA-0759 & Tópicos de Matemática Financiera & 4 \\ \hline
MA-0781 & Lógica & 5 \\ \hline
MA-0783 & Integración & 4 \\ \hline
MA-0788 & Análisis Real II & 4 \\ \hline
MA-0790 & Tópicos de Análisis & 4 \\ \hline
MA-0791 & Tópicos de Probabilidad & 4 \\ \hline
MA-0792 & Tópicos de Ecuaciones Diferenciales & 4 \\ \hline
MA-0793 & Tópicos de Geometría & 4 \\ \hline
MA-0794 & Tópicos de Modelación & 4 \\ \hline
MA-0795 & Tópicos de Computación Científica & 4 \\ \hline
MA-0796 & Tópicos de Álgebra & 4 \\ \hline
MA-0797 & Tópicos de Matemática Discreta & 4 \\ \hline
MA-0798 & Tópicos de Análisis de Datos & 4 \\ \hline
MA-0836 & Tópicos de Aprendizaje Automático & 4 \\ \hline
MA-0835 & Topología Algebraica & 4 \\ \hline
MA-0834 & Análisis Funcional & 4 \\ \hline
MA-0863 & Teoría Analítica de Números & 4 \\ \hline
MA-0862 & Análisis Armónico & 4 \\ \hline
MA-0861 & Teoría de Modelos & 4 \\ \hline
MA-0830 & Tópicos de Teoría de Números & 4 \\ \hline
MA-0831 & Tópicos de Lógica & 4 \\ \hline
MA-0832 & Tópicos de Topología & 4 \\ \hline
MA-0833 & Tópicos de Estadística & 4 \\ \hline
MA-0841 & Probabilidad & 4 \\ \hline
MA-0849 & Teoría de Módulos & 5 \\ \hline
MA-0858 & Teoría de Conjuntos & 4 \\ \hline
MA-0876 & Teoría Algebraica de Números & 4 \\ \hline
MA-0889 & Álgebra Conmutativa & 4 \\ \hline
MA-0920 & Ecuaciones en Derivadas Parciales Numéricas & 4 \\ \hline
MA-0928 & Procesos Estocásticos & 4 \\ \hline
\end{longtable}
\endgroup



\subsection*{Cursos de servicio que se ofrecen para el plan de estudios}

\begingroup
\small
\arrayrulecolor{Azul-UCR!80}%
\renewcommand{\arraystretch}{1.2}%
\rowcolors{2}{white}{Celeste-UCR!12}%
\begin{longtable}{|l|>{\raggedright\arraybackslash}p{6cm}|c|>{\raggedright\arraybackslash}p{4.2cm}|}
\caption{Cursos de servicio}\label{res:opt-otra}\\
\hline
\rowcolor{Azul-UCR}
\textcolor{white}{\textbf{Sigla}} & \textcolor{white}{\textbf{Nombre del curso}} & \textcolor{white}{\textbf{Créditos}} & \textcolor{white}{\textbf{Unidad académica}} \\ \hline
\endfirsthead
\hline
\rowcolor{Azul-UCR}
\textcolor{white}{\textbf{Sigla}} & \textcolor{white}{\textbf{Nombre del curso}} & \textcolor{white}{\textbf{Créditos}} & \textcolor{white}{\textbf{Unidad académica}} \\ \hline
\endhead
\hline
\endfoot
EG-I & Curso Integrado de Humanidades I  & 6 & Estudios generales \\ \hline
EG-II & Curso Integrado de Humanidades II & 6 & Estudios generales \\ \hline
SR-I & Seminario de Realidad Nacional I & 2 & Estudios generales \\ \hline
SR-II & Seminario de Realidad Nacional II & 2 & Estudios generales \\ \hline
EG- & Curso de Arte & 2 & Estudios generales \\ \hline
EF- & Actividad Deportiva & 0 & Escuela de Educación Física y Deportes \\ \hline
RP- & Repertorio & 3 & Estudios generales \\ \hline
B0106 & Biología General & 3 & Biología \\ \hline
B0107 & Laboratorio de Biología General & 1 & Biología \\ \hline
CA-0204 & Herramientas de Ciencia de Datos I & 4 & Escuela de Matemática \\ \hline
CA-0303 & Estadística Actuarial I & 4 & Escuela de Matemática \\ \hline
CA-0305 & Herramientas de Ciencia de Datos II & 4 & Escuela de Matemática \\ \hline
CA-0307 & Estadística Actuarial II & 4 & Escuela de Matemática \\ \hline
CA-0411 & Análisis de Datos I & 4 & Escuela de Matemática \\ \hline
CA-0414 & Modelos Lineales y de Sobrevivencia & 4 & Escuela de Matemática \\ \hline
CA-0415 & Series de Tiempo & 4 & Escuela de Matemática \\ \hline
CA-0512 & Análisis de Datos II & 4 & Escuela de Matemática \\ \hline
CA-0602 & Introducción a la Optimización & 4 & Escuela de Matemática \\ \hline
CA-0721 & Probabilidad & 4 & Escuela de Matemática \\ \hline
CI-0111 & Estructuras Discretas & 4 & Ciencias de la Computación e Informática \\ \hline
CI-0112 & Programación I & 4 & Ciencias de la Computación e Informática \\ \hline
CI-0113 & Programación II & 4 & Ciencias de la Computación e Informática \\ \hline
CI-0116 & Análisis de Algoritmos y Estructuras de Datos & 4 & Ciencias de la Computación e Informática \\ \hline
CI-0202 & Principios de Informática & 4 & Ciencias de la Computación e Informática \\ \hline
EC-1100 & Introducción a la Economía & 4 & Economía \\ \hline
EC-2100 & Teoría Microeconómica I & 4 & Economía \\ \hline
EC-3200 & Teoría Macroeconómica I & 4 & Economía \\ \hline
F0047 & Filosofía de la Física & 3 & Filosofía \\ \hline
F2001 & Seminario: Filosofía de la Física & 3 & Filosofía \\ \hline
F2002 & Filosofía de la Matemática & 3 & Filosofía \\ \hline
F2012 & Seminario de Epistemología & 3 & Filosofía \\ \hline
F2014 & Teoría del Método & 3 & Filosofía \\ \hline
F2019 & Seminario: Teorías y Revoluciones Científicas & 3 & Filosofía \\ \hline
F2022 & Filosofía de la Mente & 3 & Filosofía \\ \hline
F2334 & Filosofía de la Ciencia & 3 & Filosofía \\ \hline
F2454 & Introducción a la Epistemología & 3 & Filosofía \\ \hline
F2604 & Historia de la Ciencia & 3 & Filosofía \\ \hline
F4754 & Filosofía del Lenguaje & 3 & Filosofía \\ \hline
F5504 & Introducción a la Ética & 3 & Filosofía \\ \hline
FS-0210 & Física General I & 3 & Física \\ \hline
FS-0211 & Laboratorio de Física General I & 1 & Física \\ \hline
FS-0230 & Física I & 3 & Física \\ \hline
FS-0310 & Física General II & 3 & Física \\ \hline
FS-0311 & Laboratorio de Física General II & 1 & Física \\ \hline
FS-0330 & Física II & 3 & Física \\ \hline
FS-0331 & Laboratorio de Física Experimental I & 2 & Física \\ \hline
FS-0410 & Física General III & 3 & Física \\ \hline
FS-0411 & Laboratorio de Física General III & 1 & Física \\ \hline
FS-4101 & Introducción a la Meteorología & 3 & Física \\ \hline
G0111 & Fundamentos de Geología & 4 & Centroamericana de Geología \\ \hline
G0214 & Geología General & 4 & Centroamericana de Geología \\ \hline
QU-0100 & Química General I & 3 & Química \\ \hline
QU-0101 & Laboratorio de Química General I & 1 & Química \\ \hline
QU-0102 & Química General II & 3 & Química \\ \hline
QU-0103 & Laboratorio de Química General II & 1 & Química \\ \hline
QU-0200 & Química Analítica Cuantitativa I & 3 & Química \\ \hline
QU-0201 & Laboratorio de Química Analítica Cuantitativa I & 2 & Química \\ \hline
QU-0210 & Fundamentos de Química Orgánica & 4 & Química \\ \hline
QU-0211 & Laboratorio de Fundamentos de Química Orgánica & 1 & Química \\ \hline
XS-3150 & Diseño de Experimentos & 4 & Estadística \\ \hline
\end{longtable}
\endgroup


%%%%%%%%%%%%%%%%%%%%%%%%%%%%%%%%%%%%%%%%%%%%%%%%%%
\section{Descripción de cursos}

Para cada curso propio de la carrera se presentan sus créditos, su descripción, su objetivo general y una síntesis de la temática. Los programas completos (objetivos específicos, contenidos, metodología, evaluación y bibliografía) se incluyen en la propuesta curricular de la página 152 a la 390. 

% Generado por scripts/generate_resumen_ejecutivo.py. No editar a mano.
% Lista: por curso, créditos, descripción, objetivo general y temática resumida.

\begingroup
\small
\begin{itemize}[leftmargin=1.3em, itemsep=10pt, topsep=4pt, label={\color{Azul-UCR}$\blacksquare$}]
\item {\color{Azul-UCR}\textbf{MA-0151 Fundamentos de Álgebra, Trigonometría y Geometría Analítica}}\par\nopagebreak\smallskip
\begin{itemize}[leftmargin=1.1em, itemsep=3pt, topsep=2pt, parsep=2pt, label={\textendash}]
\item \textbf{Créditos:} 3.
\item \textbf{Descripción:} Curso del primer ciclo que uniformiza los conocimientos y habilidades esenciales del estudiantado que inicia la carrera, dadas las grandes diferencias en la formación secundaria. Desarrolla las habilidades operacionales y el pensamiento intuitivo y geométrico a partir de conceptos básicos del álgebra, la trigonometría y la geometría vectorial y analítica, priorizando el trabajo con ejemplos concretos. Se complementa con MA-0152 Matemática Exploratoria como base para estructuras de mayor complejidad en cursos posteriores.
\item \textbf{Objetivo general:} Aplicar las nociones, intuiciones y conceptos básicos del álgebra, trigonometría y geometría analítica para la resolución de problemas.
\item \textbf{Temática resumida:} Números reales y sus propiedades básicas; Trigonometría; Números complejos; Geometría vectorial y analítica básica; Proyectos de descubrimiento.
\end{itemize}

\item {\color{Azul-UCR}\textbf{MA-0152 Matemática Exploratoria}}\par\nopagebreak\smallskip
\begin{itemize}[leftmargin=1.1em, itemsep=3pt, topsep=2pt, parsep=2pt, label={\textendash}]
\item \textbf{Créditos:} 3.
\item \textbf{Descripción:} Curso de primer ciclo concebido como puerta de entrada al quehacer matemático: a partir de problemas concretos de teoría de números, combinatoria y geometría, el estudiantado explora casos, detecta patrones, formula conjeturas y las justifica o las refuta con un contraejemplo, apoyándose en la representación visual y geométrica. Trabaja de manera intuitiva el razonamiento inductivo y deductivo e incorpora conjuntos, funciones, grafos, propiedades de los enteros y los cuantificadores, de modo que la persona estudiante gane confianza y herramientas para enfrentar problemas nuevos y aprecie la elegancia de las matemáticas.
\item \textbf{Objetivo general:} Iniciar al estudiantado en el quehacer matemático mediante la resolución de problemas de teoría de números, combinatoria y geometría, a través de la exploración, la conjetura y la justificación de resultados.
\item \textbf{Temática resumida:} Introducción a los modos de trabajo y desarrollo de la matemática; Estrategias generales de resolución de problemas; Resolución de problemas concretos; Introducción a la formalización de cuantificadores.
\end{itemize}

\item {\color{Azul-UCR}\textbf{MA-0251 Cálculo en una Variable}}\par\nopagebreak\smallskip
\begin{itemize}[leftmargin=1.1em, itemsep=3pt, topsep=2pt, parsep=2pt, label={\textendash}]
\item \textbf{Créditos:} 4.
\item \textbf{Descripción:} Primer curso de Cálculo, que introduce de manera concreta y práctica el cálculo diferencial e integral en una variable, con énfasis en las habilidades operacionales y el pensamiento intuitivo y geométrico. Tiene como eje unificador el concepto de límite, a partir del cual se definen la continuidad, la derivada y la integral definida, y desarrolla técnicas para calcular límites, derivadas e integrales. Se complementa con MA-0152 Matemática Exploratoria.
\item \textbf{Objetivo general:} Aplicar las nociones, intuiciones y conceptos básicos del cálculo diferencial e integral en una variable para la resolución de problemas.
\item \textbf{Temática resumida:} Introducción a las funciones; Límites y continuidad; Derivación; Integración de Riemann.
\end{itemize}

\item {\color{Azul-UCR}\textbf{MA-0252 Introducción a las Demostraciones}}\par\nopagebreak\smallskip
\begin{itemize}[leftmargin=1.1em, itemsep=3pt, topsep=2pt, parsep=2pt, label={\textendash}]
\item \textbf{Créditos:} 3.
\item \textbf{Descripción:} Curso del segundo ciclo que continúa la transición, iniciada en MA-0152 Matemática Exploratoria, entre la matemática operacional y la matemática formal. Hace énfasis en las bases de la lógica matemática para redactar demostraciones correctamente, y presenta fundamentos de conjuntos, relaciones y funciones, teoría de números elemental e inducción matemática, apoyándose en las nociones de razonamiento inductivo y deductivo ya trabajadas.
\item \textbf{Objetivo general:} Desarrollar en el estudiantado las habilidades de comprensión y aplicación de estrategias y métodos de demostración en matemáticas.
\item \textbf{Temática resumida:} Lógica básica y estrategias de demostración; Teoría de conjuntos; Relaciones y funciones; Inducción matemática y recurrencia; Teoría de números.
\end{itemize}

\item {\color{Azul-UCR}\textbf{MA-0341 Matemática Computacional}}\par\nopagebreak\smallskip
\begin{itemize}[leftmargin=1.1em, itemsep=3pt, topsep=2pt, parsep=2pt, label={\textendash}]
\item \textbf{Créditos:} 4.
\item \textbf{Descripción:} Curso de tercer ciclo orientado a desarrollar habilidades algorítmicas mediante programación funcional en un lenguaje de alto nivel. Retoma contenidos de MA-0252, MA-0251 y MA-0361 para aplicar algoritmos básicos a la resolución de problemas, analizar su desempeño y comparar resultados teóricos con sus aproximaciones, fomentando de manera transversal las habilidades investigativas. Es requisito de MA-0501 Análisis Numérico I.
\item \textbf{Objetivo general:} Implementar el pensamiento algorítmico en la resolución de problemas matemáticos, mediante la ejecución de programas en lenguajes de programación de alto nivel.
\item \textbf{Temática resumida:} Introducción a lenguajes de programación; Introducción a algoritmos; Visualizaciones; Algoritmos básicos en teoría de números; Algoritmos básicos en cálculo diferencial e integral; Algoritmos básicos en álgebra lineal y aplicaciones.
\end{itemize}

\item {\color{Azul-UCR}\textbf{MA-0351 Cálculo en Varias Variables}}\par\nopagebreak\smallskip
\begin{itemize}[leftmargin=1.1em, itemsep=3pt, topsep=2pt, parsep=2pt, label={\textendash}]
\item \textbf{Créditos:} 4.
\item \textbf{Descripción:} Segundo curso de Cálculo, que presenta de manera concreta y práctica los conceptos fundamentales con énfasis en el pensamiento intuitivo y geométrico. Inicia con sucesiones y series y continúa con el cálculo diferencial e integral en varias variables, priorizando los ejemplos concretos. Se complementa con MA-0361 Álgebra Lineal I en el estudio de los espacios euclídeos y sus transformaciones.
\item \textbf{Objetivo general:} Aplicar las nociones, intuiciones y conceptos básicos del cálculo diferencial e integral en varias variables para la resolución de problemas.
\item \textbf{Temática resumida:} Sucesiones y series; Derivación de funciones de varias variables; Integración en varias variables operacional.
\end{itemize}

\item {\color{Azul-UCR}\textbf{MA-0361 Álgebra Lineal I}}\par\nopagebreak\smallskip
\begin{itemize}[leftmargin=1.1em, itemsep=3pt, topsep=2pt, parsep=2pt, label={\textendash}]
\item \textbf{Créditos:} 4.
\item \textbf{Descripción:} Primero de dos cursos de álgebra lineal, centrado en las habilidades operacionales y el pensamiento intuitivo y geométrico, que introduce además un mayor nivel de abstracción mediante la demostración de resultados elementales. Parte de conceptos concretos como sistemas de ecuaciones lineales, vectores, matrices y determinantes, y luego los fundamenta con espacios vectoriales y transformaciones lineales. Brinda herramientas básicas para MA-0351 Cálculo en Varias Variables y el resto de la carrera.
\item \textbf{Objetivo general:} Aplicar las nociones, intuiciones y conceptos básicos del álgebra lineal para la resolución de problemas.
\item \textbf{Temática resumida:} Sistemas de ecuaciones lineales; Vectores y matrices; Determinantes; Espacios vectoriales; Transformaciones lineales.
\end{itemize}

\item {\color{Azul-UCR}\textbf{MA-0451 Principios de Análisis I}}\par\nopagebreak\smallskip
\begin{itemize}[leftmargin=1.1em, itemsep=3pt, topsep=2pt, parsep=2pt, label={\textendash}]
\item \textbf{Créditos:} 4.
\item \textbf{Descripción:} Primero de dos cursos que introducen los principios del análisis matemático, aquí en la recta real, con énfasis en el formalismo, la abstracción y las técnicas de razonamiento y demostración propias del área. Parte de la completitud de los números reales y sus implicaciones, apoyándose en las intuiciones de los cursos de Cálculo y en las bases de lógica y demostración de MA-0152 y MA-0252.
\item \textbf{Objetivo general:} Aplicar los conceptos básicos de completitud, convergencia y diferenciación para la demostración de propiedades del análisis en una variable real.
\item \textbf{Temática resumida:} Completitud y propiedades de los números reales; Sucesiones y series numéricas; Límites y continuidad; Diferenciación; Sucesiones y series de funciones; Construcción de las funciones trascendentes elementales.
\end{itemize}

\item {\color{Azul-UCR}\textbf{MA-0461 Álgebra Lineal II}}\par\nopagebreak\smallskip
\begin{itemize}[leftmargin=1.1em, itemsep=3pt, topsep=2pt, parsep=2pt, label={\textendash}]
\item \textbf{Créditos:} 4.
\item \textbf{Descripción:} Segundo curso de álgebra lineal, de tercer ciclo, que parte de espacios vectoriales de dimensión finita y transformaciones lineales para agregar la estructura de producto escalar, lo que permite el análisis geométrico y el uso de bases ortogonales. Estudia luego la acción de las transformaciones lineales sobre subespacios, con la teoría espectral de matrices simétricas y la forma normal de Jordan como resultados centrales, de uso ubicuo en matemática y otras ciencias.
\item \textbf{Objetivo general:} Aplicar las nociones, intuiciones y conceptos del álgebra lineal para la resolución de problemas.
\item \textbf{Temática resumida:} Polinomios; Dualidad; Valores y vectores propios; Espacios con producto interno; Triangularización, diagonalización y formas canónicas; Formas bilineales y multilineales.
\end{itemize}

\item {\color{Azul-UCR}\textbf{MA-0471 Introducción a la Geometría Diferencial}}\par\nopagebreak\smallskip
\begin{itemize}[leftmargin=1.1em, itemsep=3pt, topsep=2pt, parsep=2pt, label={\textendash}]
\item \textbf{Créditos:} 3.
\item \textbf{Descripción:} Estudio riguroso e intuitivo, con herramientas analíticas, de la geometría de curvas y superficies en el espacio: parametrización, curvatura y torsión de curvas, y superficies regulares en $\mathbb{R}^3$ con sus propiedades locales y globales. A partir de la primera y segunda formas fundamentales y de la aplicación de Gauss analiza la curvatura, y con las isometrías, el Teorema Egregium, el transporte paralelo y las geodésicas introduce la geometría intrínseca. Equilibra formalismo e intuición y sirve de base concreta para el curso de variedades diferenciales.
\item \textbf{Objetivo general:} Analizar de forma rigurosa e intuitiva la geometría diferencial de curvas y superficies, mediante cálculos y demostraciones sobre sus propiedades locales y globales, tanto extrínsecas como intrínsecas.
\item \textbf{Temática resumida:} Curvas; Superficies; Aplicación de Gauss; Geometría intrínseca de superficies; Geometría proyectiva.
\end{itemize}

\item {\color{Azul-UCR}\textbf{MA-0496 Teoría de Números}}\par\nopagebreak\smallskip
\begin{itemize}[leftmargin=1.1em, itemsep=3pt, topsep=2pt, parsep=2pt, label={\textendash}]
\item \textbf{Créditos:} 3.
\item \textbf{Descripción:} Curso de teoría de números elemental para el segundo año de la carrera. Recorre la divisibilidad, los números primos y las congruencias, las funciones aritméticas, las raíces primitivas, las congruencias cuadráticas y la Ley de Reciprocidad Cuadrática de Gauss, y culmina con la aritmética de los enteros en cuerpos cuadráticos y ecuaciones diofánticas clásicas. Sirve de puente hacia la secuencia de álgebra abstracta, pues sus temas ejemplifican anillos, módulos, grupos y cuerpos.
\item \textbf{Objetivo general:} Aplicar los conceptos básicos de la teoría de números elemental para la resolución de problemas.
\item \textbf{Temática resumida:} Repaso de divisibilidad en los enteros; Números primos; Congruencias; Funciones aritméticas; Raíces primitivas; Congruencias cuadráticas y la Ley de Reciprocidad cuadrática; Enteros en cuerpos de números cuadráticos.
\end{itemize}

\item {\color{Azul-UCR}\textbf{MA-0501 Análisis Numérico I}}\par\nopagebreak\smallskip
\begin{itemize}[leftmargin=1.1em, itemsep=3pt, topsep=2pt, parsep=2pt, label={\textendash}]
\item \textbf{Créditos:} 4.
\item \textbf{Descripción:} Curso de segundo ciclo del tercer año, dirigido a Matemática y Ciencias Actuariales, en el que se aprende a resolver problemas matemáticos con la teoría y las herramientas del análisis numérico, sobre la base de MA-0341. Combina temas formales y demostraciones con resolución computacional de interpolación, integración y derivación numéricas, álgebra lineal numérica (factorizaciones y sistemas lineales) y ecuaciones diferenciales ordinarias, y es base de optativos posteriores de análisis numérico.
\item \textbf{Objetivo general:} Aplicar la teoría de análisis numérico y sus herramientas básicas en la resolución de problemas y modelos matemáticos, mediante el estudio teórico de algoritmos y la implementación de algoritmos en un lenguaje de alto nivel.
\item \textbf{Temática resumida:} Error numérico; Aproximación de funciones; Ecuaciones no lineales; Integración; Álgebra lineal numérica; Ecuaciones diferenciales ordinarias; Tema opcional.
\end{itemize}

\item {\color{Azul-UCR}\textbf{MA-0556 Ecuaciones Diferenciales}}\par\nopagebreak\smallskip
\begin{itemize}[leftmargin=1.1em, itemsep=3pt, topsep=2pt, parsep=2pt, label={\textendash}]
\item \textbf{Créditos:} 4.
\item \textbf{Descripción:} Curso de quinto ciclo que proporciona los fundamentos teóricos de las ecuaciones diferenciales ordinarias, abordadas primero de manera intuitiva y luego con rigor. Su eje es el teorema de existencia, unicidad y continuidad respecto a parámetros, seguido de estabilidad como introducción a los sistemas dinámicos y de los problemas de valor de frontera (teoría de Sturm-Liouville), que preparan para las ecuaciones en derivadas parciales. Requiere dominio de cálculo, análisis real y álgebra lineal.
\item \textbf{Objetivo general:} Establecer criterios para determinar la existencia, unicidad y estabilidad de soluciones de ecuaciones diferenciales ordinarias.
\item \textbf{Temática resumida:} Introducción y ecuaciones diferenciales ordinarias de primer orden; Sistemas de ecuaciones diferenciales lineales; La transformada de Laplace; Teoría de existencia y unicidad; Continuación de soluciones; Dependencia de valores iniciales y parámetros; Tema a elegir por la persona docente.
\end{itemize}

\item {\color{Azul-UCR}\textbf{MA-0555 Principios de Análisis II}}\par\nopagebreak\smallskip
\begin{itemize}[leftmargin=1.1em, itemsep=3pt, topsep=2pt, parsep=2pt, label={\textendash}]
\item \textbf{Créditos:} 4.
\item \textbf{Descripción:} Segundo de dos cursos que introducen los principios del análisis matemático, aquí en varias variables reales, con énfasis en el formalismo y las técnicas de demostración. Formaliza y extiende el cálculo diferencial e integral a varias variables, apoyándose en las intuiciones de MA-0251 y MA-0351 y en la lógica y estrategias de demostración de MA-0152 y MA-0252.
\item \textbf{Objetivo general:} Aplicar los conceptos básicos de convergencia, diferenciación e integración para la demostración de propiedades del análisis en una y varias variables reales.
\item \textbf{Temática resumida:} Integración de Riemann en $\mathbb{R}$; Topología métrica de $\mathbb{R}^n$; Funciones de varias variables; Diferenciación en varias variables; Cálculo y análisis vectorial.
\end{itemize}

\item {\color{Azul-UCR}\textbf{MA-0635 Topología General}}\par\nopagebreak\smallskip
\begin{itemize}[leftmargin=1.1em, itemsep=3pt, topsep=2pt, parsep=2pt, label={\textendash}]
\item \textbf{Créditos:} 4.
\item \textbf{Descripción:} Curso de sexto ciclo con las bases de la topología general desde una perspectiva rigurosa. Parte de los espacios métricos para generalizar continuidad, conjuntos abiertos y cerrados, compacidad y conexidad, y profundiza en axiomas de separación, topologías producto, subespacios, homotopía y grupo fundamental. Ofrece un lenguaje unificador para el análisis, la geometría, el álgebra y áreas emergentes como la topología de datos.
\item \textbf{Objetivo general:} Comprender los conceptos fundamentales de los espacios métricos y de los espacios topológicos, así como las funciones continuas entre ellos, y desarrollar herramientas, tanto conjuntistas como algebraicas, para analizar y distinguir dichos espacios.
\item \textbf{Temática resumida:} Repaso de teoría de conjuntos; Espacios métricos; Espacios topológicos; Axiomas de separación; Homotopía y el grupo fundamental.
\end{itemize}

\item {\color{Azul-UCR}\textbf{MA-0571 Álgebra Abstracta I}}\par\nopagebreak\smallskip
\begin{itemize}[leftmargin=1.1em, itemsep=3pt, topsep=2pt, parsep=2pt, label={\textendash}]
\item \textbf{Créditos:} 4.
\item \textbf{Descripción:} Curso introductorio de álgebra abstracta que, a diferencia del enfoque tradicional, inicia con la teoría de anillos aprovechando la familiaridad previa con ejemplos concretos de anillos y módulos. Estudia estructuras, homomorfismos y propiedades fundamentales de los anillos, y en su parte final se desplaza hacia la teoría de módulos como generalización de los espacios vectoriales.
\item \textbf{Objetivo general:} Introducir a la persona estudiante a los conceptos y métodos propios del álgebra abstracta a través del estudio de la teoría de anillos y la teoría de módulos.
\item \textbf{Temática resumida:} Definiciones, ejemplos y propiedades básicas de anillos; Morfismos, ideales y cocientes de anillos; Teoría de divisibilidad y factorización en anillos conmutativos; Anillos de polinomios y factorización; Definiciones, ejemplos y propiedades básicas de módulos; Grupos.
\end{itemize}

\item {\color{Azul-UCR}\textbf{MA-0661 Álgebra Abstracta II}}\par\nopagebreak\smallskip
\begin{itemize}[leftmargin=1.1em, itemsep=3pt, topsep=2pt, parsep=2pt, label={\textendash}]
\item \textbf{Créditos:} 5.
\item \textbf{Descripción:} Segundo curso de álgebra abstracta, centrado en grupos, cuerpos y la teoría de Galois. Estudia la estructura interna de los grupos mediante acciones de grupo y los teoremas de Sylow, introduce la teoría de cuerpos y extensiones (grado, elementos algebraicos y trascendentes) y aborda automorfismos, extensiones normales y separables, hasta el Teorema Fundamental de la Teoría de Galois y los cuerpos finitos, privilegiando la comprensión estructural y la abstracción.
\item \textbf{Objetivo general:} Analizar y aplicar los conceptos fundamentales de la teoría de grupos, cuerpos y extensiones de cuerpos, para comprender la estructura de las extensiones algebraicas y su relación con los grupos de automorfismos, mediante el uso de resultados clásicos de álgebra.
\item \textbf{Temática resumida:} Teoremas de estructura de grupos; Cuerpos; Automorfismos; Extensiones normales; Extensiones separables e inseparables; Teorema Fundamental de la Teoría de Galois; Cuerpos finitos.
\end{itemize}

\item {\color{Azul-UCR}\textbf{MA-0647 Modelación Matemática}}\par\nopagebreak\smallskip
\begin{itemize}[leftmargin=1.1em, itemsep=3pt, topsep=2pt, parsep=2pt, label={\textendash}]
\item \textbf{Créditos:} 4.
\item \textbf{Descripción:} Curso de sexto ciclo del énfasis en matemática aplicada que integra los conocimientos de ecuaciones diferenciales, álgebra lineal, probabilidad y computación científica en un enfoque sistemático de modelación. El eje es un problema real que se traduce a un modelo tratable, se analiza, se interpreta en su contexto y se valida con datos cuando es posible, con énfasis en modelos continuos (ecuaciones diferenciales y sistemas dinámicos), complementados con modelos discretos, simulación numérica y validación estadística elemental. La segunda mitad admite ejes temáticos según la persona docente y un proyecto final integrador.
\item \textbf{Objetivo general:} Modelar problemas determinísticos a través de ecuaciones diferenciales y sistemas dinámicos para situaciones de ciencias, ingeniería o contextos interdisciplinarios, evaluando la validez del modelo y comunicando resultados con claridad.
\item \textbf{Temática resumida:} El ciclo de modelación matemática; Modelación continua con ecuaciones diferenciales y sistemas dinámicos; Modelación discreta y simulación computacional; Validación, incertidumbre e interpretación con datos; Temas complementarios a elección de la persona docente; Proyecto final integrador (transversal, semanas 11--15).
\end{itemize}

\item {\color{Azul-UCR}\textbf{MA-0716 Análisis Real I}}\par\nopagebreak\smallskip
\begin{itemize}[leftmargin=1.1em, itemsep=3pt, topsep=2pt, parsep=2pt, label={\textendash}]
\item \textbf{Créditos:} 5.
\item \textbf{Descripción:} Curso del tercer año con las bases del análisis real: la teoría de la medida y la integración de Lebesgue en $\mathbb{R}^n$, precedida por una introducción a la variación acotada y la integral de Riemann-Stieltjes. Marca la transición del análisis en espacios euclídeos a la teoría de espacios abstractos, con mayor rigor, y permite construir de forma concreta espacios normados y de Hilbert. Requiere topología básica de $\mathbb{R}^n$ y cálculo en varias variables.
\item \textbf{Objetivo general:} Aplicar los conceptos y resultados básicos del análisis real para la resolución de problemas.
\item \textbf{Temática resumida:} La integral de Riemann-Stieltjes; Medida de Lebesgue en $\R^n$; Funciones Lebesgue medibles; La integral de Lebesgue; Espacios $L^p$.
\end{itemize}

\item {\color{Azul-UCR}\textbf{MA-0735 Análisis Complejo}}\par\nopagebreak\smallskip
\begin{itemize}[leftmargin=1.1em, itemsep=3pt, topsep=2pt, parsep=2pt, label={\textendash}]
\item \textbf{Créditos:} 4.
\item \textbf{Descripción:} Curso del tercer año con las bases del análisis complejo, el estudio de las funciones de una variable compleja. Muestra las propiedades que la existencia de derivadas complejas impone a estas funciones, imposibles en el análisis real, y sus aplicaciones en la matemática, la física y la ingeniería, incluida su presencia en demostraciones como las del Teorema Fundamental del Álgebra, el Teorema de los Números Primos y el Último Teorema de Fermat.
\item \textbf{Objetivo general:} Aplicar los conceptos y resultados básicos del análisis complejo para la resolución de problemas teóricos y prácticos.
\item \textbf{Temática resumida:} Números complejos y funciones en el plano complejo; Diferenciación en el plano complejo; Integración compleja, el Teorema de Cauchy y sus aplicaciones; Singularidades, funciones meromorfas y el logaritmo complejo; Funciones enteras; Tópicos adicionales (cubrir al menos uno a criterio de la persona docente).
\end{itemize}

\item {\color{Azul-UCR}\textbf{MA-0880 Comunicación en las Ciencias}}\par\nopagebreak\smallskip
\begin{itemize}[leftmargin=1.1em, itemsep=3pt, topsep=2pt, parsep=2pt, label={\textendash}]
\item \textbf{Créditos:} 2.
\item \textbf{Descripción:} Curso práctico y participativo para los últimos ciclos del bachillerato que desarrolla la escritura y la expresión oral para comunicar conocimiento matemático en contextos profesionales y académicos (docencia, propuestas de investigación, artículos e informes, solicitudes de financiamiento y presentaciones). La persona estudiante prepara y expone textos y presentaciones y recibe retroalimentación; se trabaja redacción técnica, composición con \LaTeX{}, búsqueda bibliográfica y planificación de carrera, con un componente transversal de ética profesional y conducta responsable en la investigación.
\item \textbf{Objetivo general:} Que la persona estudiante sea capaz de comunicarse en forma oral y escrita para transmitir de forma efectiva ideas, resultados y argumentos matemáticos con claridad, rigor y pertinencia en contextos académicos y profesionales.
\item \textbf{Temática resumida:} Comunicación científica y ética profesional; Escritura matemática y científica; Presentación oral y apoyos visuales; Búsqueda bibliográfica y gestión de fuentes; Comunicación en docencia y divulgación; Planificación de carrera y comunicación profesional; Presentaciones integradoras del estudiantado.
\end{itemize}

\item {\color{Azul-UCR}\textbf{MA-0881 Práctica Profesional en Matemática Pura}}\par\nopagebreak\smallskip
\begin{itemize}[leftmargin=1.1em, itemsep=3pt, topsep=2pt, parsep=2pt, label={\textendash}]
\item \textbf{Créditos:} 6.
\item \textbf{Descripción:} Curso del énfasis en matemática pura que ofrece un espacio de formación práctica en el que la persona estudiante aplica conocimientos de las áreas fundamentales de la matemática (álgebra, análisis, geometría, topología, teoría de números u otras afines) en un entorno real de trabajo, bajo supervisión académica del Departamento de Matemática. Admite dos modalidades, formalizadas mediante un plan de trabajo aprobado: la académica, de investigación bajo la dirección de una persona docente o investigadora de la Universidad de Costa Rica o de otra institución, y la profesional, en un ente externo al ámbito académico, donde se requiera rigor formal, modelación teórica o argumentación matemática. En ambas la persona estudiante cuenta con una persona supervisora en el sitio de la práctica y con una persona docente coordinadora, responsable del aval del plan y de la evaluación final.
\item \textbf{Objetivo general:} Que la persona estudiante integre y aplique conocimientos de matemática pura en un entorno académico o profesional, desarrollando autonomía, responsabilidad y capacidad de comunicar resultados.
\item \textbf{Temática resumida:} Los contenidos no se fijan de manera permanente en el plan de estudios: dependen de la modalidad y del plan de trabajo aprobado para cada persona estudiante.
\end{itemize}

\item {\color{Azul-UCR}\textbf{MA-0883 Práctica Profesional en Matemática Aplicada}}\par\nopagebreak\smallskip
\begin{itemize}[leftmargin=1.1em, itemsep=3pt, topsep=2pt, parsep=2pt, label={\textendash}]
\item \textbf{Créditos:} 6.
\item \textbf{Descripción:} Curso del énfasis en matemática aplicada que ofrece un espacio de formación práctica en el que la persona estudiante aplica conocimientos de modelación, análisis numérico, probabilidad, estadística o computación científica en un entorno real de trabajo, bajo supervisión académica del Departamento de Matemática. Admite dos modalidades, formalizadas mediante un plan de trabajo aprobado: la académica, de investigación bajo la dirección de una persona docente o investigadora de la Universidad de Costa Rica o de otra institución, y la profesional, en un ente externo al ámbito académico, con necesidades cuantitativas o de modelación. En ambas la persona estudiante cuenta con una persona supervisora en el sitio de la práctica y con una persona docente coordinadora, responsable del aval del plan y de la evaluación final.
\item \textbf{Objetivo general:} Que la persona estudiante integre y aplique conocimientos de matemática aplicada en un entorno académico o profesional, desarrollando autonomía, responsabilidad y capacidad de comunicar resultados.
\item \textbf{Temática resumida:} Los contenidos no se fijan de manera permanente en el plan de estudios: dependen de la modalidad y del plan de trabajo aprobado para cada persona estudiante.
\end{itemize}

\item {\color{Azul-UCR}\textbf{MA-0982 Seminario de Estudios Dirigidos en Matemática Pura I}}\par\nopagebreak\smallskip
\begin{itemize}[leftmargin=1.1em, itemsep=3pt, topsep=2pt, parsep=2pt, label={\textendash}]
\item \textbf{Créditos:} 7.
\item \textbf{Descripción:} Curso de licenciatura con énfasis en matemática pura que ofrece un espacio de estudios dirigidos en el que la persona estudiante desarrolla, con cierta autonomía y bajo la supervisión de una persona docente del Departamento de Matemáticas, actividades propias de la práctica investigativa: revisión documental, construcción de ejemplos y contraejemplos, demostración de teoremas, programación, experimentación numérica o simbólica y formulación de conjeturas. Los objetivos se acuerdan al inicio del ciclo en un plan de trabajo, el tema debe situarse en áreas de la matemática pura y el curso incluye reuniones periódicas de supervisión, productos intermedios y un informe escrito con presentación oral al cierre, sin exigirse resultados inéditos.
\item \textbf{Objetivo general:} Que la persona estudiante desarrolle, de manera autónoma, actividades asociadas a la investigación en matemática pura.
\item \textbf{Temática resumida:} Los contenidos no se fijan de manera permanente en el plan de estudios: dependen del tema y de los objetivos acordados para cada persona estudiante.
\end{itemize}

\item {\color{Azul-UCR}\textbf{MA-1081 Seminario de Estudios Dirigidos en Matemática Pura II}}\par\nopagebreak\smallskip
\begin{itemize}[leftmargin=1.1em, itemsep=3pt, topsep=2pt, parsep=2pt, label={\textendash}]
\item \textbf{Créditos:} 7.
\item \textbf{Descripción:} Curso de licenciatura con énfasis en matemática pura que ofrece un espacio de estudios dirigidos en el que la persona estudiante desarrolla, con cierta autonomía y bajo la supervisión de una persona docente del Departamento de Matemáticas, actividades propias de la práctica investigativa: revisión documental, construcción de ejemplos y contraejemplos, demostración de teoremas, programación, experimentación numérica o simbólica y formulación de conjeturas. Continúa la formación iniciada en MA-0982, y el tema puede profundizar el del ciclo anterior o abordarse de manera independiente. Los objetivos se acuerdan al inicio del ciclo en un plan de trabajo, el tema debe situarse en áreas de la matemática pura y el curso incluye reuniones periódicas de supervisión, productos intermedios y un informe escrito con presentación oral al cierre, sin exigirse resultados inéditos.
\item \textbf{Objetivo general:} Que la persona estudiante desarrolle, de manera autónoma, actividades asociadas a la investigación en matemática pura.
\item \textbf{Temática resumida:} Los contenidos no se fijan de manera permanente en el plan de estudios: dependen del tema y de los objetivos acordados para cada persona estudiante.
\end{itemize}

\item {\color{Azul-UCR}\textbf{MA-0984 Seminario de Estudios Dirigidos en Matemática Aplicada I}}\par\nopagebreak\smallskip
\begin{itemize}[leftmargin=1.1em, itemsep=3pt, topsep=2pt, parsep=2pt, label={\textendash}]
\item \textbf{Créditos:} 7.
\item \textbf{Descripción:} Curso de licenciatura con énfasis en matemática aplicada que ofrece un espacio de estudios dirigidos en el que la persona estudiante desarrolla, con cierta autonomía y bajo la supervisión de una persona docente del Departamento de Matemáticas, actividades propias de la práctica investigativa: revisión documental, construcción de ejemplos y contraejemplos, demostración de teoremas, programación, experimentación numérica o simbólica y formulación de conjeturas. Los objetivos se acuerdan al inicio del ciclo en un plan de trabajo, el tema debe situarse en áreas de la matemática aplicada y el curso incluye reuniones periódicas de supervisión, productos intermedios y un informe escrito con presentación oral al cierre, sin exigirse resultados inéditos.
\item \textbf{Objetivo general:} Que la persona estudiante desarrolle, de manera autónoma, actividades asociadas a la investigación en matemática aplicada.
\item \textbf{Temática resumida:} Los contenidos no se fijan de manera permanente en el plan de estudios: dependen del tema y de los objetivos acordados para cada persona estudiante.
\end{itemize}

\item {\color{Azul-UCR}\textbf{MA-1082 Seminario de Estudios Dirigidos en Matemática Aplicada II}}\par\nopagebreak\smallskip
\begin{itemize}[leftmargin=1.1em, itemsep=3pt, topsep=2pt, parsep=2pt, label={\textendash}]
\item \textbf{Créditos:} 7.
\item \textbf{Descripción:} Curso de licenciatura con énfasis en matemática aplicada que ofrece un espacio de estudios dirigidos en el que la persona estudiante desarrolla, con cierta autonomía y bajo la supervisión de una persona docente del Departamento de Matemáticas, actividades propias de la práctica investigativa: revisión documental, construcción de ejemplos y contraejemplos, demostración de teoremas, programación, experimentación numérica o simbólica y formulación de conjeturas. Continúa la formación iniciada en MA-0984, y el tema puede profundizar el del ciclo anterior o abordarse de manera independiente. Los objetivos se acuerdan al inicio del ciclo en un plan de trabajo, el tema debe situarse en áreas de la matemática aplicada y el curso incluye reuniones periódicas de supervisión, productos intermedios y un informe escrito con presentación oral al cierre, sin exigirse resultados inéditos.
\item \textbf{Objetivo general:} Que la persona estudiante desarrolle, de manera autónoma, actividades asociadas a la investigación en matemática aplicada.
\item \textbf{Temática resumida:} Los contenidos no se fijan de manera permanente en el plan de estudios: dependen del tema y de los objetivos acordados para cada persona estudiante.
\end{itemize}

\item {\color{Azul-UCR}\textbf{MA-0526 Combinatoria}}\par\nopagebreak\smallskip
\begin{itemize}[leftmargin=1.1em, itemsep=3pt, topsep=2pt, parsep=2pt, label={\textendash}]
\item \textbf{Créditos:} 4.
\item \textbf{Descripción:} Curso sobre estructuras finitas o discretas, mejor descrito por los problemas que aborda, en el que las técnicas para resolver un problema importan más que el resultado mismo. Presenta herramientas básicas de conteo y nociones de matemática discreta (coeficientes binomiales, grafos, funciones generatrices) e introduce problemas de combinatoria extremal.
\item \textbf{Objetivo general:} Desarrollar en la persona estudiante competencias para abordar problemas combinatorios mediante técnicas de conteo, teoría de grafos, métodos probabilísticos y combinatoria extremal.
\item \textbf{Temática resumida:} Herramientas de conteo; Introducción a la teoría de grafos; Método probabilístico; Funciones generatrices; Combinatoria extremal.
\end{itemize}

\item {\color{Azul-UCR}\textbf{MA-0528 Sistemas Dinámicos}}\par\nopagebreak\smallskip
\begin{itemize}[leftmargin=1.1em, itemsep=3pt, topsep=2pt, parsep=2pt, label={\textendash}]
\item \textbf{Créditos:} 4.
\item \textbf{Descripción:} Introducción rigurosa a los sistemas dinámicos y la teoría ergódica, con énfasis en sistemas unidimensionales. Estudia sistemas en tiempo discreto y continuo, puntos fijos y periódicos, conjuntos límite, recurrencia, conjugaciones, difeomorfismos del intervalo y codificación simbólica; en la segunda parte desarrolla medida invariante, recurrencia de Poincaré, promedios de Birkhoff y ergodicidad, aplicados a mapas afines, expansores, el mapa de Gauss y la familia cuadrática.
\item \textbf{Objetivo general:} Que la persona estudiante adquiera las herramientas fundamentales para analizar sistemas dinámicos desde los puntos de vista topológico, diferenciable y ergódico, reconocer estructuras recurrentes, construir modelos simbólicos, estudiar medidas invariantes e interpretar el comportamiento estadístico de las órbitas mediante promedios temporales.
\item \textbf{Temática resumida:} Sistemas dinámicos en tiempo discreto y continuo; Conjuntos límite, recurrencia y conjugaciones; Difeomorfismos del intervalo y linealización local; Codificación simbólica; Medidas invariantes y teoría ergódica; Mapas por ramas, mapa de Gauss y distorsión acotada; La familia cuadrática.
\end{itemize}

\item {\color{Azul-UCR}\textbf{MA-0838 Geometría Diferencial}}\par\nopagebreak\smallskip
\begin{itemize}[leftmargin=1.1em, itemsep=3pt, topsep=2pt, parsep=2pt, label={\textendash}]
\item \textbf{Créditos:} 4.
\item \textbf{Descripción:} Curso optativo de cuarto año de carácter teórico que introduce la geometría de las variedades diferenciales, el marco natural para generalizar curvas y superficies a dimensión arbitraria y estudiarlas de manera intrínseca. Desarrolla espacios tangentes, campos vectoriales, fibrados, formas diferenciales, métricas riemannianas, conexiones, curvatura e integración, con resultados como el teorema del valor regular, el lema de Poincaré y el teorema de Stokes. Se apoya en MA-0471 y MA-0556 y tiene aplicaciones en geometría riemanniana y simpléctica, topología, sistemas dinámicos y física.
\item \textbf{Objetivo general:} Comprender y aplicar las herramientas fundamentales del cálculo diferencial e integral y de la geometría riemanniana en variedades, demostrando resultados centrales.
\item \textbf{Temática resumida:} Variedades diferenciales y topología diferencial; Formas diferenciales; Integración en variedades; Geometría riemanniana.
\end{itemize}

\item {\color{Azul-UCR}\textbf{MA-0837 Geometría Algebraica}}\par\nopagebreak\smallskip
\begin{itemize}[leftmargin=1.1em, itemsep=3pt, topsep=2pt, parsep=2pt, label={\textendash}]
\item \textbf{Créditos:} 4.
\item \textbf{Descripción:} Curso de enfoque clásico, basado en el libro de Fulton, que introduce la geometría algebraica a través de conjuntos algebraicos, variedades y, sobre todo, curvas algebraicas. Parte de ideales y el Teorema de los Ceros de Hilbert, estudia curvas planas, multiplicidades, intersecciones, geometría proyectiva, el teorema de Bézout, morfismos, equivalencia birracional y resolución de singularidades, y alcanza la teoría de divisores, el género y el Teorema de Riemann--Roch; finaliza con cuerpos finitos, la función zeta y las conjeturas de Weil. Equilibra teoría y ejemplos concretos, con apoyo de álgebra computacional.
\item \textbf{Objetivo general:} Introducir a la persona estudiante a los conceptos, métodos y resultados fundamentales de la Geometría Algebraica clásica, con énfasis en variedades y curvas algebraicas, sus singularidades, propiedades proyectivas y birracionales, y sus invariantes geométricos y aritméticos.
\item \textbf{Temática resumida:} Conjuntos algebraicos afines y el Teorema de los Ceros de Hilbert; Variedades afines y propiedades locales; Propiedades locales de curvas planas; Geometría proyectiva; Curvas proyectivas planas e intersecciones; Morfismos y geometría birracional; Resolución de singularidades de curvas; Divisores, género y Teorema de Riemann--Roch; Curvas sobre cuerpos finitos, funciones zeta y conjeturas de Weil.
\end{itemize}

\item {\color{Azul-UCR}\textbf{MA-0734 Optimización}}\par\nopagebreak\smallskip
\begin{itemize}[leftmargin=1.1em, itemsep=3pt, topsep=2pt, parsep=2pt, label={\textendash}]
\item \textbf{Créditos:} 4.
\item \textbf{Descripción:} Primer curso de optimización, con perspectiva de matemática pura, que estudia la existencia y la caracterización de valores extremos sobre un conjunto factible. Articula, desde el análisis en varias variables, el álgebra lineal y la geometría convexa, los problemas sin restricciones, la programación lineal y los problemas con restricciones de igualdad y desigualdad, con las condiciones de optimalidad, los teoremas de separación y la dualidad como eje; los algoritmos solo ilustran la teoría. Es base para análisis convexo, optimización no lineal e investigación de operaciones.
\item \textbf{Objetivo general:} Desarrollar los fundamentos matemáticos de la optimización para establecer la existencia y caracterizar las soluciones de problemas con y sin restricciones.
\item \textbf{Temática resumida:} Fundamentos y optimización sin restricciones; Programación lineal; Optimización con restricciones de igualdad y desigualdad; Proyectos de descubrimiento.
\end{itemize}

\item {\color{Azul-UCR}\textbf{MA-0756 Ecuaciones Diferenciales Parciales}}\par\nopagebreak\smallskip
\begin{itemize}[leftmargin=1.1em, itemsep=3pt, topsep=2pt, parsep=2pt, label={\textendash}]
\item \textbf{Créditos:} 5.
\item \textbf{Descripción:} Curso optativo con una introducción rigurosa a las ecuaciones diferenciales parciales desde el análisis funcional y los espacios de Sobolev. Parte de las ecuaciones clásicas (Laplace, calor y onda) y avanza hacia formulaciones débiles, existencia de soluciones y evolución lineal, siguiendo la estructura de un curso introductorio de posgrado como el de Evans, con énfasis en demostraciones, estimaciones de energía, el teorema de Lax--Milgram y aproximaciones de Galerkin. Supone dominio de medida e integración, espacios $L^p$ y cálculo en varias variables.
\item \textbf{Objetivo general:} Desarrollar competencias teóricas para formular, analizar y resolver problemas de ecuaciones diferenciales parciales mediante métodos clásicos y formulaciones débiles.
\item \textbf{Temática resumida:} Ecuaciones clásicas fundamentales; Espacios de Sobolev; Ecuaciones elípticas de segundo orden; Ecuaciones de evolución lineales.
\end{itemize}

\item {\color{Azul-UCR}\textbf{MA-0781 Lógica}}\par\nopagebreak\smallskip
\begin{itemize}[leftmargin=1.1em, itemsep=3pt, topsep=2pt, parsep=2pt, label={\textendash}]
\item \textbf{Créditos:} 5.
\item \textbf{Descripción:} Curso sobre los fundamentos de la lógica matemática, que estudia los lenguajes y sistemas formales con que se expresan los razonamientos matemáticos y las estructuras que los interpretan. Introduce teoría de conjuntos (ordinales y cardinales), sintaxis y semántica del cálculo de predicados de primer orden, los teoremas de completitud y compacidad, y nociones básicas de teoría de modelos hasta los teoremas de Löwenheim--Skolem y el criterio de Tarski--Vaught.
\item \textbf{Objetivo general:} Desarrollar los conocimientos y herramientas fundamentales de la lógica matemática para analizar y formalizar razonamientos matemáticos y demostrar los principales resultados del cálculo de predicados de primer orden.
\item \textbf{Temática resumida:} Ordinales y cardinales; Lógica de primer orden; Teorema de completitud y teorema de compacidad; Nociones básicas de teorías y estructuras.
\end{itemize}

\item {\color{Azul-UCR}\textbf{MA-0783 Integración}}\par\nopagebreak\smallskip
\begin{itemize}[leftmargin=1.1em, itemsep=3pt, topsep=2pt, parsep=2pt, label={\textendash}]
\item \textbf{Créditos:} 4.
\item \textbf{Descripción:} Curso optativo que profundiza en la medida y la integración abstractas, generalizando la teoría de Lebesgue a espacios de medida más generales. Ofrece un marco abstracto riguroso y lo aplica a desigualdades isoperimétricas, teoremas ergódicos, medidas de Hausdorff y de Haar y transporte óptimo, con énfasis en la resolución de problemas y la conexión con la investigación contemporánea. Supone la teoría de Lebesgue en $\mathbb{R}^n$, espacios $L^p$ y los resultados de MA-0716 y MA-0788.
\item \textbf{Objetivo general:} Desarrollar competencias en teoría abstracta de la medida e integración y en sus aplicaciones clásicas dentro del análisis matemático.
\item \textbf{Temática resumida:} Repaso de integración de Lebesgue; Curvas rectificables e isoperimetría; Teoría abstracta de la medida e integración; Teoremas ergódicos; Medidas de Hausdorff; Medidas de Haar; Transporte óptimo.
\end{itemize}

\item {\color{Azul-UCR}\textbf{MA-0788 Análisis Real II}}\par\nopagebreak\smallskip
\begin{itemize}[leftmargin=1.1em, itemsep=3pt, topsep=2pt, parsep=2pt, label={\textendash}]
\item \textbf{Créditos:} 4.
\item \textbf{Descripción:} Curso del cuarto año que continúa las bases del análisis real, con énfasis en el rigor en el tratamiento y la comunicación de argumentos. Cubre la teoría de diferenciación de Lebesgue, la continuidad absoluta y el Teorema Fundamental del Cálculo, elementos de espacios de Banach y de Hilbert, y una introducción a las series y la transformada de Fourier. Requiere álgebra lineal, topología general y medida e integración de Lebesgue.
\item \textbf{Objetivo general:} Aplicar los conceptos y resultados intermedios del Análisis Real para la resolución de problemas.
\item \textbf{Temática resumida:} Diferenciación; Continuidad absoluta y el Teorema Fundamental del Cálculo; Espacios de Banach; Espacios de Hilbert; Introducción al Análisis de Fourier.
\end{itemize}

\item {\color{Azul-UCR}\textbf{MA-0835 Topología Algebraica}}\par\nopagebreak\smallskip
\begin{itemize}[leftmargin=1.1em, itemsep=3pt, topsep=2pt, parsep=2pt, label={\textendash}]
\item \textbf{Créditos:} 4.
\item \textbf{Descripción:} Curso optativo de enfoque teórico-demostrativo sobre los invariantes algebraicos (grupos de homotopía, homología y cohomología) asociados a espacios topológicos. Desarrolla la homotopía y el grupo fundamental, las homologías cúbica, simplicial y singular y, al final, nociones de cohomología de haces y homotopía superior, con ejemplos en grafos, superficies, espacios proyectivos y complejos simpliciales. Busca que la persona estudiante comprenda la funtorialidad y los isomorfismos clásicos (Hurewicz, Mayer--Vietoris, dualidad).
\item \textbf{Objetivo general:} Que la persona estudiante domine los conceptos y técnicas elementales de la topología algebraica y pueda aplicarlos al análisis de espacios y aplicaciones continuas.
\item \textbf{Temática resumida:} Homotopía (noción básica); Grupo fundamental; Homología; Cohomología de haces (Sheaf Cohomology); Homotopía superior.
\end{itemize}

\item {\color{Azul-UCR}\textbf{MA-0834 Análisis Funcional}}\par\nopagebreak\smallskip
\begin{itemize}[leftmargin=1.1em, itemsep=3pt, topsep=2pt, parsep=2pt, label={\textendash}]
\item \textbf{Créditos:} 4.
\item \textbf{Descripción:} Curso optativo de cuarto año, de énfasis teórico, que introduce el análisis funcional como estudio de espacios de dimensión infinita, principalmente de Banach y de Hilbert, y de los operadores lineales acotados entre ellos. Cubre espacios $L^p$, los teoremas de Hahn--Banach, acotación uniforme, aplicación abierta y gráfica cerrada, dualidad, topologías débiles, el teorema de Alaoglu, operadores compactos y el teorema espectral, y marca el paso del análisis en $\mathbb{R}^n$ al análisis en espacios abstractos.
\item \textbf{Objetivo general:} Que la persona estudiante sea capaz de resolver problemas y desarrollar teoría a partir de los fundamentos del análisis funcional en espacios de Banach y de Hilbert.
\item \textbf{Temática resumida:} Espacios de Banach y espacios $L^p$; Espacios de Hilbert; Operadores lineales acotados; Teoremas fundamentales; Dualidad y topologías débiles; Operadores compactos y teorema espectral.
\end{itemize}

\item {\color{Azul-UCR}\textbf{MA-0863 Teoría Analítica de Números}}\par\nopagebreak\smallskip
\begin{itemize}[leftmargin=1.1em, itemsep=3pt, topsep=2pt, parsep=2pt, label={\textendash}]
\item \textbf{Créditos:} 4.
\item \textbf{Descripción:} Curso que estudia problemas aritméticos con herramientas del análisis real y complejo, con énfasis en el comportamiento promedio de funciones aritméticas y en la distribución de los primos. Recorre funciones aritméticas, convolución de Dirichlet, series de Dirichlet, la transformada de Mellin y la fórmula de Perron, y se centra en la función zeta de Riemann y las funciones $L$ de Dirichlet (productos de Euler, continuación analítica, ecuaciones funcionales y ceros), con el Teorema de los Números Primos, los primos en progresiones aritméticas y los grandes problemas abiertos, como la Hipótesis de Riemann.
\item \textbf{Objetivo general:} Introducir a la persona estudiante a las principales ideas, técnicas y resultados de la Teoría Analítica de Números, enfatizando el uso de herramientas del análisis real y complejo para estudiar funciones aritméticas y la distribución de los números primos.
\item \textbf{Temática resumida:} Funciones aritméticas y convolución de Dirichlet; Funciones sumatorias y órdenes promedio; Series de Dirichlet y herramientas de análisis complejo; La función zeta de Riemann; Ceros de la función zeta y distribución de los números primos; Caracteres y funciones $L$ de Dirichlet; Primos en progresiones aritméticas; Problemas abiertos y desarrollos posteriores.
\end{itemize}

\item {\color{Azul-UCR}\textbf{MA-0862 Análisis Armónico}}\par\nopagebreak\smallskip
\begin{itemize}[leftmargin=1.1em, itemsep=3pt, topsep=2pt, parsep=2pt, label={\textendash}]
\item \textbf{Créditos:} 4.
\item \textbf{Descripción:} Curso optativo de cuarto año, de énfasis teórico, sobre análisis armónico de variable real: la transformada de Fourier, los operadores maximales y las integrales singulares en $\mathbb{R}^n$. Siguiendo en lo esencial a Duoandikoetxea, aborda espacios $L^p$ e interpolación, la transformada de Fourier, funciones maximales, transformadas de Hilbert y de Riesz, integrales singulares de Calderón--Zygmund y la teoría de Littlewood--Paley, y se articula con el análisis funcional y las ecuaciones diferenciales parciales.
\item \textbf{Objetivo general:} Que la persona estudiante sea capaz de resolver problemas y desarrollar teoría a partir de los fundamentos del análisis armónico de variable real en $\R^n$.
\item \textbf{Temática resumida:} Espacios $L^p$; La transformada de Fourier; Funciones maximales; Las transformadas de Hilbert y de Riesz; Integrales singulares de tipo convolución; Teoría de Littlewood--Paley.
\end{itemize}

\item {\color{Azul-UCR}\textbf{MA-0861 Teoría de Modelos}}\par\nopagebreak\smallskip
\begin{itemize}[leftmargin=1.1em, itemsep=3pt, topsep=2pt, parsep=2pt, label={\textendash}]
\item \textbf{Créditos:} 4.
\item \textbf{Descripción:} Curso optativo de lógica matemática que introduce la teoría de modelos, la relación entre estructuras matemáticas y los lenguajes de primer orden que las describen, con el conjunto definible como noción central. Repasa los fundamentos de la lógica de primer orden (sintaxis, semántica, completitud y compacidad) y desarrolla inmersiones elementales, equivalencia elemental, tipos, saturación y eliminación de cuantificadores, ilustrados con cuerpos algebraicamente cerrados, espacios vectoriales, órdenes densos, grupos y grafos, con aplicaciones al álgebra.
\item \textbf{Objetivo general:} Comprender los principios básicos de la teoría de modelos y de sus aplicaciones a las estructuras algebraicas.
\item \textbf{Temática resumida:} Fundamentos de lógica de primer orden; Principios básicos de teoría de modelos; Eliminación de cuantificadores; Aplicaciones al álgebra y la geometría algebraica; Tipos y compacidad.
\end{itemize}

\item {\color{Azul-UCR}\textbf{MA-0841 Probabilidad}}\par\nopagebreak\smallskip
\begin{itemize}[leftmargin=1.1em, itemsep=3pt, topsep=2pt, parsep=2pt, label={\textendash}]
\item \textbf{Créditos:} 4.
\item \textbf{Descripción:} Curso optativo que introduce la teoría general de la probabilidad desde el enfoque axiomático de Kolmogorov, apoyado en $\sigma$-álgebras, medidas de probabilidad e integración respecto a ellas. Supone conocimientos de integración en $\mathbb{R}^n$ (como los de MA-0716), pero no experiencia previa en probabilidad: inicia con una sección de probabilidad elemental y técnicas combinatorias para el cálculo de probabilidades en espacios finitos.
\item \textbf{Objetivo general:} Desarrollar los fundamentos de la teoría de la probabilidad en el marco de la teoría de la medida para modelar incertidumbre, demostrar resultados centrales y aplicarlos en problemas matemáticos e interdisciplinarios.
\item \textbf{Temática resumida:} Probabilidad elemental; Teoría de la medida e integración; Teoría de probabilidad; Independencia; Convergencia de variables aleatorias; Probabilidad condicional.
\end{itemize}

\item {\color{Azul-UCR}\textbf{MA-0849 Teoría de Módulos}}\par\nopagebreak\smallskip
\begin{itemize}[leftmargin=1.1em, itemsep=3pt, topsep=2pt, parsep=2pt, label={\textendash}]
\item \textbf{Créditos:} 5.
\item \textbf{Descripción:} Curso que introduce la teoría de módulos como generalización natural de los espacios vectoriales y los grupos abelianos. Estudia los módulos sobre anillos arbitrarios, con énfasis en los módulos sobre dominios de ideales principales y sus aplicaciones a la clasificación de los grupos abelianos finitamente generados y a la forma normal de Jordan.
\item \textbf{Objetivo general:} Comprender los fundamentos de la teoría de módulos y sus principales aplicaciones algebraicas.
\item \textbf{Temática resumida:} Módulos; Construcciones básicas; Módulos noetherianos y artinianos; Módulos sobre dominios de ideales principales; Módulos proyectivos, inyectivos y planos; Producto tensorial; Introducción al álgebra homológica (Optativo).
\end{itemize}

\item {\color{Azul-UCR}\textbf{MA-0858 Teoría de Conjuntos}}\par\nopagebreak\smallskip
\begin{itemize}[leftmargin=1.1em, itemsep=3pt, topsep=2pt, parsep=2pt, label={\textendash}]
\item \textbf{Créditos:} 4.
\item \textbf{Descripción:} Introducción a la teoría de conjuntos desde un enfoque axiomático, que no requiere un curso previo de lógica formal y usa el lenguaje matemático habitual con uso riguroso de cuantificadores y conectivos. Desarrolla la teoría de Zermelo--Fraenkel con el axioma de elección como fundamento de la matemática moderna, y estudia los números cardinales y ordinales, el teorema de Cantor--Bernstein--Schroeder, la aritmética transfinita y resultados fundamentales sobre el continuo.
\item \textbf{Objetivo general:} Introducir al estudiantado en la teoría axiomática de conjuntos de Zermelo--Fraenkel, desarrollando las herramientas fundamentales como cardinales, ordinales y aritmética transfinita.
\item \textbf{Temática resumida:} Motivación y antecedentes históricos; El sistema axiomático de Zermelo--Fraenkel; Relaciones y funciones; Órdenes; Los números naturales; Números cardinales; Números ordinales; El axioma de elección y la hipótesis del continuo.
\end{itemize}

\item {\color{Azul-UCR}\textbf{MA-0876 Teoría Algebraica de Números}}\par\nopagebreak\smallskip
\begin{itemize}[leftmargin=1.1em, itemsep=3pt, topsep=2pt, parsep=2pt, label={\textendash}]
\item \textbf{Créditos:} 4.
\item \textbf{Descripción:} Curso que introduce la teoría de números algebraicos a partir de la pregunta por la aritmética de los enteros en extensiones finitas de $\mathbb{Q}$. Estudia cuerpos de números, anillos de enteros, factorización de ideales, el grupo de clases (cuya finitud se prueba por métodos geométricos), el teorema de Dirichlet sobre unidades y los primeros aspectos de la función zeta de Dedekind, así como la factorización de primos en extensiones de Galois (ramificación, descomposición, inercia y Frobenius). Combina demostraciones rigurosas con cálculo explícito en \texttt{SageMath}.
\item \textbf{Objetivo general:} Introducir a la persona estudiante a los conceptos, métodos y ejemplos fundamentales de la Teoría de Números Algebraicos, enfatizando la aritmética de los anillos de enteros de cuerpos de números, la factorización de ideales, el grupo de clases, el grupo de unidades, la factorización de primos en extensiones y el uso de herramientas computacionales para explorar ejemplos concretos.
\item \textbf{Temática resumida:} Cuerpos de números y enteros algebraicos; Ideales y factorización; Factorización de ideales primos y ramificación; El grupo de clases y el teorema de Minkowski; Unidades y teorema de Dirichlet; Conteo de ideales y función zeta de Dedekind; Teoría de descomposición e inercia en extensiones de Galois.
\end{itemize}

\item {\color{Azul-UCR}\textbf{MA-0889 Álgebra Conmutativa}}\par\nopagebreak\smallskip
\begin{itemize}[leftmargin=1.1em, itemsep=3pt, topsep=2pt, parsep=2pt, label={\textendash}]
\item \textbf{Créditos:} 4.
\item \textbf{Descripción:} Curso que profundiza en anillos, ideales y módulos a partir de la secuencia de álgebra abstracta, con herramientas útiles en geometría algebraica y teoría de números. Estudia módulos finitamente generados, producto tensorial, anillos noetherianos y artinianos, descomposición primaria, localización, extensiones integrales, normalización de Noether y dimensión de Krull; incluye bases de Gröbner y métodos computacionales, y concluye con anillos locales regulares, dominios de valuación discreta y dominios de Dedekind, combinando teoría, ejemplos, interpretación geométrica y experimentación.
\item \textbf{Objetivo general:} Introducir a la persona estudiante a los conceptos, métodos y resultados fundamentales del Álgebra Conmutativa, enfatizando las propiedades de anillos y módulos noetherianos, la localización, la dependencia integral, la teoría de dimensión, los métodos computacionales y la estructura local de los anillos conmutativos.
\item \textbf{Temática resumida:} Módulos y herramientas algebraicas fundamentales; El diccionario entre álgebra y geometría; Noetherianidad y descomposición primaria; Localización y anillos locales; Dependencia integral y normalización; Dimensión de Krull; Bases de Gröbner y métodos computacionales; Anillos locales regulares y anillos de dimensión uno.
\end{itemize}

\item {\color{Azul-UCR}\textbf{MA-0920 Ecuaciones en Derivadas Parciales Numéricas}}\par\nopagebreak\smallskip
\begin{itemize}[leftmargin=1.1em, itemsep=3pt, topsep=2pt, parsep=2pt, label={\textendash}]
\item \textbf{Créditos:} 4.
\item \textbf{Descripción:} Curso optativo de análisis numérico aplicado a la resolución numérica de ecuaciones diferenciales parciales, con énfasis en la convergencia y la estabilidad de los algoritmos y en su implementación. Estudia los métodos de diferencias finitas y de elemento finito, complementados con álgebra lineal numérica, y acompaña la teoría con programación en un lenguaje científico como \texttt{MATLAB} u \texttt{Octave}, para desarrollar tanto la intuición sobre los métodos como la capacidad de programarlos, evaluarlos y comunicar resultados.
\item \textbf{Objetivo general:} Aplicar la teoría del análisis numérico y sus herramientas básicas para aproximar soluciones numéricas de ecuaciones diferenciales parciales.
\item \textbf{Temática resumida:} Introducción a las ecuaciones diferenciales parciales; Método de diferencias finitas; Método de elemento finito; Resolución numérica de sistemas lineales; Resolución numérica de la ecuación de calor y de onda; Temas adicionales.
\end{itemize}

\item {\color{Azul-UCR}\textbf{MA-0477 Ecuaciones Diferenciales Parciales Aplicadas}}\par\nopagebreak\smallskip
\begin{itemize}[leftmargin=1.1em, itemsep=3pt, topsep=2pt, parsep=2pt, label={\textendash}]
\item \textbf{Créditos:} 4.
\item \textbf{Descripción:} Curso optativo que introduce las ecuaciones diferenciales parciales desde un enfoque clásico y aplicado, centrado en métodos explícitos, principios del máximo, funciones de Green y representaciones integrales, sin desarrollar por completo la teoría abstracta de los espacios de Sobolev. Recorre las ecuaciones de primer orden, los problemas elípticos (Laplace y Poisson), las ecuaciones del calor y de onda, y cierra con una introducción al cálculo de variaciones como puente hacia las soluciones débiles; prepara para modelación, análisis numérico y estadística aplicada.
\item \textbf{Objetivo general:} Aplicar métodos clásicos del análisis de EDP para resolver y analizar ecuaciones lineales que surgen en modelos matemáticos.
\item \textbf{Temática resumida:} EDP de primer orden; Ecuaciones de Laplace y Poisson; Ecuación del calor; Ecuación de onda; Introducción al cálculo de variaciones.
\end{itemize}

\item {\color{Azul-UCR}\textbf{MA-0759 Tópicos de Matemática Financiera}}\par\nopagebreak\smallskip
\begin{itemize}[leftmargin=1.1em, itemsep=3pt, topsep=2pt, parsep=2pt, label={\textendash}]
\item \textbf{Créditos:} 4.
\item \textbf{Descripción:} Curso optativo de formato abierto que permite profundizar en temas de matemática financiera que no se desarrollan de manera sistemática en otros cursos de la carrera. Cada oferta puede orientarse hacia líneas distintas dentro de la matemática financiera, clásicas o contemporáneas, según la especialidad de la persona docente, los avances recientes de la disciplina, las demandas del estudiantado avanzado o las líneas de investigación activas en la Escuela de Matemática. Se espera que la persona estudiante consolide su capacidad de lectura de textos especializados, de demostración o aplicación de resultados y de comunicación matemática con mayor autonomía intelectual.
\item \textbf{Objetivo general:} Profundizar en un tema de la matemática financiera, definido para cada oferta del curso, mediante el estudio guiado de resultados, problemas y, cuando corresponda, aplicaciones.
\item \textbf{Temática resumida:} Los contenidos no se fijan de manera permanente en el plan de estudios: son abiertos por oferta, siempre que pertenezcan al ámbito de Matemática financiera. Al inicio del ciclo lectivo, la persona docente debe comunicar por escrito al estudiantado.
\end{itemize}

\item {\color{Azul-UCR}\textbf{MA-0790 Tópicos de Análisis}}\par\nopagebreak\smallskip
\begin{itemize}[leftmargin=1.1em, itemsep=3pt, topsep=2pt, parsep=2pt, label={\textendash}]
\item \textbf{Créditos:} 4.
\item \textbf{Descripción:} Curso optativo de formato abierto que permite profundizar en temas de análisis matemático que no se desarrollan de manera sistemática en otros cursos de la carrera. Cada oferta puede orientarse hacia líneas distintas dentro del análisis (real, complejo, funcional o numérico), clásicas o contemporáneas, según la especialidad de la persona docente, los avances recientes de la disciplina, las demandas del estudiantado avanzado o las líneas de investigación activas en la Escuela de Matemática. Se espera que la persona estudiante consolide su capacidad de lectura de textos especializados, de demostración o aplicación de resultados y de comunicación matemática con mayor autonomía intelectual.
\item \textbf{Objetivo general:} Profundizar en un tema de análisis matemático, definido para cada oferta del curso, mediante el estudio guiado de resultados, problemas y, cuando corresponda, aplicaciones.
\item \textbf{Temática resumida:} Los contenidos no se fijan de manera permanente en el plan de estudios: son abiertos por oferta, siempre que pertenezcan al ámbito del Análisis matemático.
\end{itemize}

\item {\color{Azul-UCR}\textbf{MA-0791 Tópicos de Probabilidad}}\par\nopagebreak\smallskip
\begin{itemize}[leftmargin=1.1em, itemsep=3pt, topsep=2pt, parsep=2pt, label={\textendash}]
\item \textbf{Créditos:} 4.
\item \textbf{Descripción:} Curso optativo de formato abierto que permite profundizar en temas de probabilidad que no se desarrollan de manera sistemática en otros cursos de la carrera. Cada oferta puede orientarse hacia líneas distintas dentro de la probabilidad, clásicas o contemporáneas, según la especialidad de la persona docente, los avances recientes de la disciplina, las demandas del estudiantado avanzado o las líneas de investigación activas en la Escuela de Matemática. Se espera que la persona estudiante consolide su capacidad de lectura de textos especializados, de demostración o aplicación de resultados y de comunicación matemática con mayor autonomía intelectual.
\item \textbf{Objetivo general:} Profundizar en un tema de la probabilidad, definido para cada oferta del curso, mediante el estudio guiado de resultados, problemas y, cuando corresponda, aplicaciones.
\item \textbf{Temática resumida:} Los contenidos no se fijan de manera permanente en el plan de estudios: son abiertos por oferta, siempre que pertenezcan al ámbito de Probabilidad. Al inicio del ciclo lectivo, la persona docente debe comunicar por escrito al estudiantado.
\end{itemize}

\item {\color{Azul-UCR}\textbf{MA-0792 Tópicos de Ecuaciones Diferenciales}}\par\nopagebreak\smallskip
\begin{itemize}[leftmargin=1.1em, itemsep=3pt, topsep=2pt, parsep=2pt, label={\textendash}]
\item \textbf{Créditos:} 4.
\item \textbf{Descripción:} Curso optativo de formato abierto que permite profundizar en temas de ecuaciones diferenciales que no se desarrollan de manera sistemática en otros cursos de la carrera. Cada oferta puede orientarse hacia líneas distintas dentro de las ecuaciones diferenciales, clásicas o contemporáneas, según la especialidad de la persona docente, los avances recientes de la disciplina, las demandas del estudiantado avanzado o las líneas de investigación activas en la Escuela de Matemática. Se espera que la persona estudiante consolide su capacidad de lectura de textos especializados, de demostración o aplicación de resultados y de comunicación matemática con mayor autonomía intelectual.
\item \textbf{Objetivo general:} Profundizar en un tema de las ecuaciones diferenciales, definido para cada oferta del curso, mediante el estudio guiado de resultados, problemas y, cuando corresponda, aplicaciones.
\item \textbf{Temática resumida:} Los contenidos no se fijan de manera permanente en el plan de estudios: son abiertos por oferta, siempre que pertenezcan al ámbito de Ecuaciones diferenciales. Al inicio del ciclo lectivo, la persona docente debe comunicar por escrito al estudiantado.
\end{itemize}

\item {\color{Azul-UCR}\textbf{MA-0793 Tópicos de Geometría}}\par\nopagebreak\smallskip
\begin{itemize}[leftmargin=1.1em, itemsep=3pt, topsep=2pt, parsep=2pt, label={\textendash}]
\item \textbf{Créditos:} 4.
\item \textbf{Descripción:} Curso optativo de formato abierto que permite profundizar en temas de geometría que no se desarrollan de manera sistemática en otros cursos de la carrera. Cada oferta puede orientarse hacia líneas distintas dentro de la geometría, clásicas o contemporáneas, según la especialidad de la persona docente, los avances recientes de la disciplina, las demandas del estudiantado avanzado o las líneas de investigación activas en la Escuela de Matemática. Se espera que la persona estudiante consolide su capacidad de lectura de textos especializados, de demostración o aplicación de resultados y de comunicación matemática con mayor autonomía intelectual.
\item \textbf{Objetivo general:} Profundizar en un tema de la geometría, definido para cada oferta del curso, mediante el estudio guiado de resultados, problemas y, cuando corresponda, aplicaciones.
\item \textbf{Temática resumida:} Los contenidos no se fijan de manera permanente en el plan de estudios: son abiertos por oferta, siempre que pertenezcan al ámbito de Geometría. Al inicio del ciclo lectivo, la persona docente debe comunicar por escrito al estudiantado.
\end{itemize}

\item {\color{Azul-UCR}\textbf{MA-0794 Tópicos de Modelación}}\par\nopagebreak\smallskip
\begin{itemize}[leftmargin=1.1em, itemsep=3pt, topsep=2pt, parsep=2pt, label={\textendash}]
\item \textbf{Créditos:} 4.
\item \textbf{Descripción:} Curso optativo de formato abierto que permite profundizar en temas de modelación matemática que no se desarrollan de manera sistemática en otros cursos de la carrera. Cada oferta puede orientarse hacia líneas distintas dentro de la modelación matemática, clásicas o contemporáneas, según la especialidad de la persona docente, los avances recientes de la disciplina, las demandas del estudiantado avanzado o las líneas de investigación activas en la Escuela de Matemática. Se espera que la persona estudiante consolide su capacidad de lectura de textos especializados, de demostración o aplicación de resultados y de comunicación matemática con mayor autonomía intelectual.
\item \textbf{Objetivo general:} Profundizar en un tema de la modelación matemática, definido para cada oferta del curso, mediante el estudio guiado de resultados, problemas y, cuando corresponda, aplicaciones.
\item \textbf{Temática resumida:} Los contenidos no se fijan de manera permanente en el plan de estudios: son abiertos por oferta, siempre que pertenezcan al ámbito de Modelación matemática. Al inicio del ciclo lectivo, la persona docente debe comunicar por escrito al estudiantado.
\end{itemize}

\item {\color{Azul-UCR}\textbf{MA-0795 Tópicos de Computación Científica}}\par\nopagebreak\smallskip
\begin{itemize}[leftmargin=1.1em, itemsep=3pt, topsep=2pt, parsep=2pt, label={\textendash}]
\item \textbf{Créditos:} 4.
\item \textbf{Descripción:} Curso optativo de formato abierto que permite profundizar en temas de computación científica que no se desarrollan de manera sistemática en otros cursos de la carrera. Cada oferta puede orientarse hacia líneas distintas dentro de la computación científica, clásicas o contemporáneas, según la especialidad de la persona docente, los avances recientes de la disciplina, las demandas del estudiantado avanzado o las líneas de investigación activas en la Escuela de Matemática. Se espera que la persona estudiante consolide su capacidad de lectura de textos especializados, de demostración o aplicación de resultados y de comunicación matemática con mayor autonomía intelectual.
\item \textbf{Objetivo general:} Profundizar en un tema de la computación científica, definido para cada oferta del curso, mediante el estudio guiado de resultados, problemas y, cuando corresponda, aplicaciones.
\item \textbf{Temática resumida:} Los contenidos no se fijan de manera permanente en el plan de estudios: son abiertos por oferta, siempre que pertenezcan al ámbito de Computación científica. Al inicio del ciclo lectivo, la persona docente debe comunicar por escrito al estudiantado.
\end{itemize}

\item {\color{Azul-UCR}\textbf{MA-0796 Tópicos de Álgebra}}\par\nopagebreak\smallskip
\begin{itemize}[leftmargin=1.1em, itemsep=3pt, topsep=2pt, parsep=2pt, label={\textendash}]
\item \textbf{Créditos:} 4.
\item \textbf{Descripción:} Curso optativo de formato abierto que permite profundizar en temas de álgebra que no se desarrollan de manera sistemática en otros cursos de la carrera. Cada oferta puede orientarse hacia líneas distintas dentro del álgebra, clásicas o contemporáneas, según la especialidad de la persona docente, los avances recientes de la disciplina, las demandas del estudiantado avanzado o las líneas de investigación activas en la Escuela de Matemática. Se espera que la persona estudiante consolide su capacidad de lectura de textos especializados, de demostración o aplicación de resultados y de comunicación matemática con mayor autonomía intelectual.
\item \textbf{Objetivo general:} Profundizar en un tema del álgebra, definido para cada oferta del curso, mediante el estudio guiado de resultados, problemas y, cuando corresponda, aplicaciones.
\item \textbf{Temática resumida:} Los contenidos no se fijan de manera permanente en el plan de estudios: son abiertos por oferta, siempre que pertenezcan al ámbito de Álgebra. Al inicio del ciclo lectivo, la persona docente debe comunicar por escrito al estudiantado.
\end{itemize}

\item {\color{Azul-UCR}\textbf{MA-0797 Tópicos de Matemática Discreta}}\par\nopagebreak\smallskip
\begin{itemize}[leftmargin=1.1em, itemsep=3pt, topsep=2pt, parsep=2pt, label={\textendash}]
\item \textbf{Créditos:} 4.
\item \textbf{Descripción:} Curso optativo de formato abierto que permite profundizar en temas de matemática discreta que no se desarrollan de manera sistemática en otros cursos de la carrera. Cada oferta puede orientarse hacia líneas distintas dentro de la matemática discreta, clásicas o contemporáneas, según la especialidad de la persona docente, los avances recientes de la disciplina, las demandas del estudiantado avanzado o las líneas de investigación activas en la Escuela de Matemática. Se espera que la persona estudiante consolide su capacidad de lectura de textos especializados, de demostración o aplicación de resultados y de comunicación matemática con mayor autonomía intelectual.
\item \textbf{Objetivo general:} Profundizar en un tema de la matemática discreta, definido para cada oferta del curso, mediante el estudio guiado de resultados, problemas y, cuando corresponda, aplicaciones.
\item \textbf{Temática resumida:} Los contenidos no se fijan de manera permanente en el plan de estudios: son abiertos por oferta, siempre que pertenezcan al ámbito de Matemática discreta. Al inicio del ciclo lectivo, la persona docente debe comunicar por escrito al estudiantado.
\end{itemize}

\item {\color{Azul-UCR}\textbf{MA-0798 Tópicos de Análisis de Datos}}\par\nopagebreak\smallskip
\begin{itemize}[leftmargin=1.1em, itemsep=3pt, topsep=2pt, parsep=2pt, label={\textendash}]
\item \textbf{Créditos:} 4.
\item \textbf{Descripción:} Curso optativo de formato abierto que permite profundizar en temas de análisis de datos que no se desarrollan de manera sistemática en otros cursos de la carrera. Cada oferta puede orientarse hacia líneas distintas dentro del análisis de datos, clásicas o contemporáneas, según la especialidad de la persona docente, los avances recientes de la disciplina, las demandas del estudiantado avanzado o las líneas de investigación activas en la Escuela de Matemática. Se espera que la persona estudiante consolide su capacidad de lectura de textos especializados, de demostración o aplicación de resultados y de comunicación matemática con mayor autonomía intelectual.
\item \textbf{Objetivo general:} Profundizar en un tema del análisis de datos, definido para cada oferta del curso, mediante el estudio guiado de resultados, problemas y, cuando corresponda, aplicaciones.
\item \textbf{Temática resumida:} Los contenidos no se fijan de manera permanente en el plan de estudios: son abiertos por oferta, siempre que pertenezcan al ámbito de Análisis de datos. Al inicio del ciclo lectivo, la persona docente debe comunicar por escrito al estudiantado.
\end{itemize}

\item {\color{Azul-UCR}\textbf{MA-0836 Tópicos de Aprendizaje Automático}}\par\nopagebreak\smallskip
\begin{itemize}[leftmargin=1.1em, itemsep=3pt, topsep=2pt, parsep=2pt, label={\textendash}]
\item \textbf{Créditos:} 4.
\item \textbf{Descripción:} Curso optativo de formato abierto que permite profundizar en temas de aprendizaje automático que no se desarrollan de manera sistemática en otros cursos de la carrera. Cada oferta puede orientarse hacia líneas distintas dentro del aprendizaje automático, clásicas o contemporáneas, según la especialidad de la persona docente, los avances recientes de la disciplina, las demandas del estudiantado avanzado o las líneas de investigación activas en la Escuela de Matemática. Se espera que la persona estudiante consolide su capacidad de lectura de textos especializados, de demostración o aplicación de resultados y de comunicación matemática con mayor autonomía intelectual.
\item \textbf{Objetivo general:} Profundizar en un tema del aprendizaje automático, definido para cada oferta del curso, mediante el estudio guiado de resultados, problemas y, cuando corresponda, aplicaciones.
\item \textbf{Temática resumida:} Los contenidos no se fijan de manera permanente en el plan de estudios: son abiertos por oferta, siempre que pertenezcan al ámbito de Aprendizaje automático. Al inicio del ciclo lectivo, la persona docente debe comunicar por escrito al estudiantado.
\end{itemize}

\item {\color{Azul-UCR}\textbf{MA-0830 Tópicos de Teoría de Números}}\par\nopagebreak\smallskip
\begin{itemize}[leftmargin=1.1em, itemsep=3pt, topsep=2pt, parsep=2pt, label={\textendash}]
\item \textbf{Créditos:} 4.
\item \textbf{Descripción:} Curso optativo de formato abierto que permite profundizar en temas de teoría de números que no se desarrollan de manera sistemática en otros cursos de la carrera. Cada oferta puede orientarse hacia líneas distintas dentro de la teoría de números, clásicas o contemporáneas, según la especialidad de la persona docente, los avances recientes de la disciplina, las demandas del estudiantado avanzado o las líneas de investigación activas en la Escuela de Matemática. Se espera que la persona estudiante consolide su capacidad de lectura de textos especializados, de demostración o aplicación de resultados y de comunicación matemática con mayor autonomía intelectual.
\item \textbf{Objetivo general:} Profundizar en un tema de la teoría de números, definido para cada oferta del curso, mediante el estudio guiado de resultados, problemas y, cuando corresponda, aplicaciones.
\item \textbf{Temática resumida:} Los contenidos no se fijan de manera permanente en el plan de estudios: son abiertos por oferta, siempre que pertenezcan al ámbito de Teoría de números. Al inicio del ciclo lectivo, la persona docente debe comunicar por escrito al estudiantado.
\end{itemize}

\item {\color{Azul-UCR}\textbf{MA-0831 Tópicos de Lógica}}\par\nopagebreak\smallskip
\begin{itemize}[leftmargin=1.1em, itemsep=3pt, topsep=2pt, parsep=2pt, label={\textendash}]
\item \textbf{Créditos:} 4.
\item \textbf{Descripción:} Curso optativo de formato abierto que permite profundizar en temas de lógica que no se desarrollan de manera sistemática en otros cursos de la carrera. Cada oferta puede orientarse hacia líneas distintas dentro de la lógica, clásicas o contemporáneas, según la especialidad de la persona docente, los avances recientes de la disciplina, las demandas del estudiantado avanzado o las líneas de investigación activas en la Escuela de Matemática. Se espera que la persona estudiante consolide su capacidad de lectura de textos especializados, de demostración o aplicación de resultados y de comunicación matemática con mayor autonomía intelectual.
\item \textbf{Objetivo general:} Profundizar en un tema de la lógica, definido para cada oferta del curso, mediante el estudio guiado de resultados, problemas y, cuando corresponda, aplicaciones.
\item \textbf{Temática resumida:} Los contenidos no se fijan de manera permanente en el plan de estudios: son abiertos por oferta, siempre que pertenezcan al ámbito de Lógica. Al inicio del ciclo lectivo, la persona docente debe comunicar por escrito al estudiantado.
\end{itemize}

\item {\color{Azul-UCR}\textbf{MA-0832 Tópicos de Topología}}\par\nopagebreak\smallskip
\begin{itemize}[leftmargin=1.1em, itemsep=3pt, topsep=2pt, parsep=2pt, label={\textendash}]
\item \textbf{Créditos:} 4.
\item \textbf{Descripción:} Curso optativo de formato abierto que permite profundizar en temas de topología que no se desarrollan de manera sistemática en otros cursos de la carrera. Cada oferta puede orientarse hacia líneas distintas dentro de la topología, clásicas o contemporáneas, según la especialidad de la persona docente, los avances recientes de la disciplina, las demandas del estudiantado avanzado o las líneas de investigación activas en la Escuela de Matemática. Se espera que la persona estudiante consolide su capacidad de lectura de textos especializados, de demostración o aplicación de resultados y de comunicación matemática con mayor autonomía intelectual.
\item \textbf{Objetivo general:} Profundizar en un tema de la topología, definido para cada oferta del curso, mediante el estudio guiado de resultados, problemas y, cuando corresponda, aplicaciones.
\item \textbf{Temática resumida:} Los contenidos no se fijan de manera permanente en el plan de estudios: son abiertos por oferta, siempre que pertenezcan al ámbito de Topología. Al inicio del ciclo lectivo, la persona docente debe comunicar por escrito al estudiantado.
\end{itemize}

\item {\color{Azul-UCR}\textbf{MA-0833 Tópicos de Estadística}}\par\nopagebreak\smallskip
\begin{itemize}[leftmargin=1.1em, itemsep=3pt, topsep=2pt, parsep=2pt, label={\textendash}]
\item \textbf{Créditos:} 4.
\item \textbf{Descripción:} Curso optativo de formato abierto que permite profundizar en temas de estadística que no se desarrollan de manera sistemática en otros cursos de la carrera. Cada oferta puede orientarse hacia líneas distintas dentro de la estadística, clásicas o contemporáneas, según la especialidad de la persona docente, los avances recientes de la disciplina, las demandas del estudiantado avanzado o las líneas de investigación activas en la Escuela de Matemática. Se espera que la persona estudiante consolide su capacidad de lectura de textos especializados, de demostración o aplicación de resultados y de comunicación matemática con mayor autonomía intelectual.
\item \textbf{Objetivo general:} Profundizar en un tema de la estadística, definido para cada oferta del curso, mediante el estudio guiado de resultados, problemas y, cuando corresponda, aplicaciones.
\item \textbf{Temática resumida:} Los contenidos no se fijan de manera permanente en el plan de estudios: son abiertos por oferta, siempre que pertenezcan al ámbito de Estadística. Al inicio del ciclo lectivo, la persona docente debe comunicar por escrito al estudiantado.
\end{itemize}

\item {\color{Azul-UCR}\textbf{MA-0928 Procesos Estocásticos}}\par\nopagebreak\smallskip
\begin{itemize}[leftmargin=1.1em, itemsep=3pt, topsep=2pt, parsep=2pt, label={\textendash}]
\item \textbf{Créditos:} 4.
\item \textbf{Descripción:} Curso optativo que introduce los procesos estocásticos en tiempo continuo, con énfasis en el movimiento browniano y el cálculo estocástico que de él se deriva. Trata el movimiento browniano, las filtraciones y martingalas, la integración estocástica, la fórmula de Itô, los procesos de Markov y las ecuaciones diferenciales estocásticas. Supone dominio de la probabilidad basada en medida (MA-0841), incluida la convergencia de variables aleatorias, la independencia y la esperanza condicional.
\item \textbf{Objetivo general:} Manejar los conceptos y técnicas básicas de la teoría de procesos estocásticos en tiempo continuo, con especial atención al movimiento browniano y sus aplicaciones.
\item \textbf{Temática resumida:} Procesos gaussianos; Movimiento browniano; Filtraciones, martingalas y semimartingalas continuas; Variación cuadrática e integración estocástica; Fórmula de Itô y consecuencias; Procesos de Markov; Movimiento browniano y ecuaciones en derivadas parciales; Ecuaciones diferenciales estocásticas.
\end{itemize}

\end{itemize}
\endgroup


%%%%%%%%%%%%%%%%%%%%%%%%%%%%%%%%%%%%%%%%%%%%%%%%%%
\section{Tabla de correspondencia del equipo docente}

La Tabla~\ref{res:docentes} presenta, para cada curso, el personal docente del Departamento de Matemática disponible para impartirlo, y la Tabla~\ref{res:docentes-datos} los datos personales y profesionales de dicho personal. Se espera que la gran mayoría de los cursos sean impartidos por docentes con doctorado; se considerará la asignación de docentes con maestría para cursos de índole general o afines a su área de trabajo. Un mismo docente puede figurar en varios cursos.

% Generado por scripts/generate_resumen_ejecutivo.py. No editar a mano.

\begingroup
\small
\arrayrulecolor{Azul-UCR!80}%
\renewcommand{\arraystretch}{1.2}%
\begin{longtable}{|>{\raggedright\arraybackslash}p{6.9cm}|>{\raggedright\arraybackslash}p{9.8cm}|}
\caption{Listado de las actividades de formación (cursos) y personal docente a cargo}\label{res:docentes}\\
\hline
\rowcolor{Azul-UCR}
\textcolor{white}{\textbf{Nombre del curso}} & \textcolor{white}{\textbf{Persona(s) docente(s) a cargo}} \\ \hline
\endfirsthead
\hline
\rowcolor{Azul-UCR}
\textcolor{white}{\textbf{Nombre del curso}} & \textcolor{white}{\textbf{Persona(s) docente(s) a cargo}} \\ \hline
\endhead
\hline
\endfoot
\rowcolor{Celeste-UCR!35}\multicolumn{2}{|l|}{\textbf{Escuela de Matemática, Departamento de Matemática}} \\ \hline
\textbf{MA-0151} Fundamentos de Álgebra, Trigonometría y Geometría Analítica & Villalobos Arias Mario; Guier Acosta Jorge; Sánchez Guevara Jesús; Acuña Larios Jennifer; Lacy Mora Allan \\ \hline
\textbf{MA-0152} Matemática Exploratoria & Calvo Alpízar Juan Gabriel; Zamora Luna Oscar; Mora Sáenz Germán José; Ulloa Esquivel Roberto; Bustamante Medina Ronald \\ \hline
\textbf{MA-0251} Cálculo en una Variable & Acuña Larios Jennifer; Campos Fernández José David; Ramírez González José Alexander; Zamora Luna Oscar; Alfaro Murillo Jorge; Fonseca Mora Christian \\ \hline
\textbf{MA-0252} Introducción a las Demostraciones & Montenegro Guzmán Samaria; Zamora Calero Rafael; Loría Sorio Jennifer de los Ángeles; Ulloa Esquivel Roberto; Mora Sáenz Germán José \\ \hline
\textbf{MA-0341} Matemática Computacional & Calvo Alpízar Juan Gabriel; Rodríguez Rojas Oldemar; Segura Ugalde Esteban; Alfaro Murillo Jorge; Guevara Villalobos Álvaro \\ \hline
\textbf{MA-0351} Cálculo en Varias Variables & Acuña Larios Jennifer; Bolaños Guerrero Raúl Esteban; Jiménez López David; Mena Arias Darío Alberto; Chaves Aguilar Santiago \\ \hline
\textbf{MA-0361} Álgebra Lineal I & Barquero Sánchez Adrián; Bustamante Medina Ronald; Lacy Mora Allan; Rosales Ortega José; Jiménez López David \\ \hline
\textbf{MA-0451} Principios de Análisis I & Acuña Valverde Luis; Cambronero Villalobos Santiago; Campos Fernández José David; Sánchez Chavarría Adriana; Bolaños Guerrero Raúl Esteban \\ \hline
\textbf{MA-0461} Álgebra Lineal II & Guier Acosta Jorge; Hernández Alvarado Alberto; Segura Ugalde Esteban; Campos Salas Daniel; Mora Sáenz Germán José \\ \hline
\textbf{MA-0471} Introducción a la Geometría Diferencial & Rosales Ortega José; Zúñiga Rojas Ronald Alberto; Ugalde Gómez William Javier; Chaves Aguilar Santiago \\ \hline
\textbf{MA-0496} Teoría de Números & Lacy Mora Allan; Coto Mora Edward; Barquero Sánchez Adrián; Ávila Artavia Josué David \\ \hline
\textbf{MA-0556} Ecuaciones Diferenciales & Montalto Cruz Carlos; Loría Sorio Jennifer de los Ángeles; Sánchez Chavarría Adriana; Campos Salas Daniel \\ \hline
\textbf{MA-0555} Principios de Análisis II & Bolaños Guerrero Raúl Esteban; Ugalde Gómez William Javier; Villalobos Alvarado Jorge; Cambronero Villalobos Santiago \\ \hline
\textbf{MA-0635} Topología General & Zamora Calero Rafael; Tall El Hadji Yaya; Villalobos Alvarado Jorge; Rosales Ortega José \\ \hline
\textbf{MA-0571} Álgebra Abstracta I & Barquero Sánchez Adrián; Hernández Alvarado Alberto; Bustamante Medina Ronald; Zúñiga Rojas Ronald Alberto \\ \hline
\textbf{MA-0661} Álgebra Abstracta II & Guier Acosta Jorge; Montenegro Guzmán Samaria; Ávila Artavia Josué David; Coto Mora Edward \\ \hline
\textbf{MA-0647} Modelación Matemática & Montalto Cruz Carlos; Sánchez Peña Fabio; Alfaro Murillo Jorge; Loría Sorio Jennifer de los Ángeles \\ \hline
\textbf{MA-0735} Análisis Complejo & Bolaños Guerrero Raúl Esteban; Méndez Hernández Pedro; Tall El Hadji Yaya; Gutiérrez Pavón Jonathan; Calvo Alpízar Juan Gabriel \\ \hline
\textbf{MA-0716} Análisis Real I & Acuña Valverde Luis; Mena Arias Darío Alberto; Mora Sáenz Germán José; Campos Fernández José David \\ \hline
\textbf{MA-0880} Comunicación en las Ciencias & Mena Arias Darío Alberto; Sánchez Chavarría Adriana; Sánchez Peña Fabio; Alfaro Murillo Jorge \\ \hline
\textbf{MA-0788} Análisis Real II & Cambronero Villalobos Santiago; Ramírez González José Alexander; Fonseca Mora Christian; Méndez Hernández Pedro \\ \hline
\textbf{MA-0876} Teoría Algebraica de Números & Coto Mora Edward; Lacy Mora Allan; Chaves Aguilar Santiago \\ \hline
\textbf{MA-0863} Teoría Analítica de Números & Barquero Sánchez Adrián; Ávila Artavia Josué David; Ulloa Esquivel Roberto \\ \hline
\textbf{MA-0858} Teoría de Conjuntos & Zamora Calero Rafael; Sánchez Guevara Jesús \\ \hline
\textbf{MA-0526} Combinatoria & Zamora Luna Oscar; Tall El Hadji Yaya \\ \hline
\textbf{MA-0528} Sistemas Dinámicos & Sánchez Chavarría Adriana; Tall El Hadji Yaya \\ \hline
\textbf{MA-0783} Integración & Gutiérrez Pavón Jonathan; Mena Arias Darío Alberto; Ugalde Gómez William Javier \\ \hline
\textbf{MA-0838} Geometría Diferencial & Ramírez González José Alexander; Rosales Ortega José; Ugalde Gómez William Javier \\ \hline
\textbf{MA-0837} Geometría Algebraica & Zúñiga Rojas Ronald Alberto; Coto Mora Edward; Ulloa Esquivel Roberto \\ \hline
\textbf{MA-0781} Lógica & Montenegro Guzmán Samaria; Zamora Calero Rafael; Guier Acosta Jorge \\ \hline
\textbf{MA-0734} Optimización & Acuña Larios Jennifer; Guevara Villalobos Álvaro; Segura Ugalde Esteban \\ \hline
\textbf{MA-0755} Ecuaciones Diferenciales Parciales & Méndez Hernández Pedro; Montalto Cruz Carlos; Acuña Valverde Luis; Campos Salas Daniel \\ \hline
\textbf{MA-0835} Topología Algebraica & Zúñiga Rojas Ronald Alberto; Ugalde Gómez William Javier; Ávila Artavia Josué David \\ \hline
\textbf{MA-0834} Análisis Funcional & Cambronero Villalobos Santiago; Jiménez López David; Villalobos Alvarado Jorge \\ \hline
\textbf{MA-0862} Análisis Armónico & Jiménez López David; Mena Arias Darío Alberto; Tall El Hadji Yaya \\ \hline
\textbf{MA-0861} Teoría de Modelos & Bustamante Medina Ronald; Guier Acosta Jorge; Montenegro Guzmán Samaria \\ \hline
\textbf{MA-0841} Probabilidad & Acuña Valverde Luis; Campos Fernández José David; Fonseca Mora Christian \\ \hline
\textbf{MA-0849} Teoría de Módulos & Hernández Alvarado Alberto; Lacy Mora Allan; Sánchez Guevara Jesús \\ \hline
\textbf{MA-0889} Álgebra Conmutativa & Hernández Alvarado Alberto; Ávila Artavia Josué David; Ulloa Esquivel Roberto \\ \hline
\textbf{MA-0920} Ecuaciones en Derivadas Parciales Numéricas & Calvo Alpízar Juan Gabriel; Segura Ugalde Esteban; Jiménez López David; Villalobos Arias Mario \\ \hline
\textbf{MA-0477} Ecuaciones Diferenciales Parciales Aplicadas & Segura Ugalde Esteban; Mora Sáenz Germán José; Villalobos Arias Mario \\ \hline
\textbf{MA-0790} Tópicos de Análisis & Acuña Valverde Luis; Cambronero Villalobos Santiago; Campos Salas Daniel \\ \hline
\textbf{MA-0791} Tópicos de Probabilidad & Acuña Larios Jennifer; Bolaños Guerrero Raúl Esteban; Gutiérrez Pavón Jonathan; Fonseca Mora Christian \\ \hline
\textbf{MA-0792} Tópicos de Ecuaciones Diferenciales & Méndez Hernández Pedro; Sánchez Peña Fabio; Ramírez González José Alexander \\ \hline
\textbf{MA-0793} Tópicos de Geometría & Rosales Ortega José; Chaves Aguilar Santiago; Villalobos Alvarado Jorge \\ \hline
\textbf{MA-0794} Tópicos de Modelación & Montalto Cruz Carlos; Sánchez Peña Fabio; Alfaro Murillo Jorge; Loría Sorio Jennifer de los Ángeles \\ \hline
\textbf{MA-0796} Tópicos de Álgebra & Sánchez Guevara Jesús; Chaves Aguilar Santiago; Hernández Alvarado Alberto \\ \hline
\textbf{MA-0797} Tópicos de Matemática Discreta & Méndez Hernández Pedro; Zamora Luna Oscar; Campos Salas Daniel; Villalobos Arias Mario \\ \hline
\textbf{MA-0830} Tópicos de Teoría de Números & Barquero Sánchez Adrián; Zamora Luna Oscar; Coto Mora Edward \\ \hline
\textbf{MA-0831} Tópicos de Lógica & Bustamante Medina Ronald; Sánchez Guevara Jesús; Montenegro Guzmán Samaria; Zamora Calero Rafael \\ \hline
\textbf{MA-0832} Tópicos de Topología & Zúñiga Rojas Ronald Alberto; Villalobos Alvarado Jorge; Rodríguez Rojas Oldemar \\ \hline
\textbf{MA-0928} Procesos Estocásticos & Campos Fernández José David; Gutiérrez Pavón Jonathan; Fonseca Mora Christian \\ \hline
\textbf{MA-0795} Tópicos de Computación Científica & Calvo Alpízar Juan Gabriel; Rodríguez Rojas Oldemar; Guevara Villalobos Álvaro \\ \hline
\textbf{MA-0798} Tópicos de Análisis de Datos & Rodríguez Rojas Oldemar; Sánchez Chavarría Adriana; Sánchez Peña Fabio \\ \hline
\textbf{MA-0836} Tópicos de Aprendizaje Automático & Guevara Villalobos Álvaro; Rodríguez Rojas Oldemar; Loría Sorio Jennifer de los Ángeles \\ \hline
\textbf{MA-0759} Tópicos de Matemática Financiera & Guevara Villalobos Álvaro; Ramírez González José Alexander; Villalobos Arias Mario \\ \hline
\rowcolor{Celeste-UCR!35}\multicolumn{2}{|l|}{\textbf{Escuela de Matemática, Departamento de Ciencias Actuariales}} \\ \hline
\textbf{CA-0204} Herramientas de Ciencia de Datos I & \textit{Curso ya existente, a cargo del personal docente que actualmente lo imparte.} \\ \hline
\textbf{CA-0721} Probabilidad & \textit{Curso ya existente, a cargo del personal docente que actualmente lo imparte.} \\ \hline
\textbf{CA-0303} Estadística Actuarial I & \textit{Curso ya existente, a cargo del personal docente que actualmente lo imparte.} \\ \hline
\textbf{CA-0304} Herramientas de Ciencia de Datos II & \textit{Curso ya existente, a cargo del personal docente que actualmente lo imparte.} \\ \hline
\textbf{CA-0411a} Análisis de Datos II & \textit{Curso ya existente, a cargo del personal docente que actualmente lo imparte.} \\ \hline
\textbf{CA-0403} Estadística Actuarial II & \textit{Curso ya existente, a cargo del personal docente que actualmente lo imparte.} \\ \hline
\textbf{CA-0512} Análisis de Datos II & \textit{Curso ya existente, a cargo del personal docente que actualmente lo imparte.} \\ \hline
\textbf{CA-0612} Series de Tiempo & \textit{Curso ya existente, a cargo del personal docente que actualmente lo imparte.} \\ \hline
\rowcolor{Celeste-UCR!35}\multicolumn{2}{|l|}{\textbf{Escuela de Matemática, otros cursos ya existentes}} \\ \hline
\textbf{MA-0501} Análisis Numérico I & \textit{Curso ya existente, a cargo del personal docente que actualmente lo imparte.} \\ \hline
\textbf{MA-0602} Introducción a la Optimización & \textit{Curso ya existente, a cargo del personal docente que actualmente lo imparte.} \\ \hline
\end{longtable}
\endgroup


{\footnotesize\noindent\textit{Fuente:} Tabla de correspondencia de profesores por curso de la propuesta de plan de estudios (octubre de 2026).\par}

% Tabla 4 (formato CONARE): datos personales y profesionales del personal docente.
\begin{landscape}
\begingroup
\scriptsize
\setlength{\tabcolsep}{3pt}%
\arrayrulecolor{Azul-UCR!80}%
\renewcommand{\arraystretch}{1.15}%
\begin{longtable}{|>{\raggedright\arraybackslash}p{2.3cm}|>{\raggedright\arraybackslash}p{1.6cm}|>{\raggedright\arraybackslash}p{3.0cm}|>{\raggedright\arraybackslash}p{1.7cm}|>{\raggedright\arraybackslash}p{3.4cm}|>{\raggedright\arraybackslash}p{4.5cm}|>{\raggedright\arraybackslash}p{4.6cm}|>{\raggedright\arraybackslash}p{2.3cm}|}
\caption{Datos personales y profesionales del personal docente de la carrera}\label{res:docentes-datos}\\
\hline
\rowcolor{Azul-UCR}
\multicolumn{3}{|c|}{\textcolor{white}{\textbf{Datos personales}}} & \multicolumn{5}{c|}{\textcolor{white}{\textbf{Datos profesionales}}} \\ \hline
\rowcolor{Azul-UCR}
\textcolor{white}{\textbf{Nombre completo}} & \textcolor{white}{\textbf{Cédula}} & \textcolor{white}{\textbf{Correo electrónico}} & \textcolor{white}{\textbf{Años de experiencia docente universitaria}} & \textcolor{white}{\textbf{Grados y títulos universitarios}} & \textcolor{white}{\textbf{Publicaciones (últimos 5 años)\textsuperscript{a}}} & \textcolor{white}{\textbf{Investigaciones o proyectos (últimos 5 años)\textsuperscript{b}}} & \textcolor{white}{\textbf{Dominio de otro idioma}} \\ \hline
\endfirsthead
\hline
\rowcolor{Azul-UCR}
\multicolumn{3}{|c|}{\textcolor{white}{\textbf{Datos personales}}} & \multicolumn{5}{c|}{\textcolor{white}{\textbf{Datos profesionales}}} \\ \hline
\rowcolor{Azul-UCR}
\textcolor{white}{\textbf{Nombre completo}} & \textcolor{white}{\textbf{Cédula}} & \textcolor{white}{\textbf{Correo electrónico}} & \textcolor{white}{\textbf{Años de experiencia docente universitaria}} & \textcolor{white}{\textbf{Grados y títulos universitarios}} & \textcolor{white}{\textbf{Publicaciones (últimos 5 años)\textsuperscript{a}}} & \textcolor{white}{\textbf{Investigaciones o proyectos (últimos 5 años)\textsuperscript{b}}} & \textcolor{white}{\textbf{Dominio de otro idioma}} \\ \hline
\endhead
\hline
\endfoot
1. Acuña Larios, Jennifer & \textit{Pendiente} & jennifer.\allowbreak{}acunalarios@\allowbreak{}ucr.\allowbreak{}ac.\allowbreak{}cr & 12 & Bachillerato en Matemática;\newline Bachiller en Ciencias Actuariales;\newline Maestría Académica en Matemática Pura & \textit{Ninguna en 2022--2026 según el CV} & \textbullet~Investigación: Estudio actuarial para el Colegio de Enfermeros y Enfermeras de Costa Rica, dos etapas, CIMPA (2022)\newline \textbullet~Investigación: Estudio actuarial para el Archivo Nacional, cálculo de cesantía, CIMPA (2022)\newline \textbullet~Consultoría: Validación de modelos del Banco Nacional de Costa Rica, equipo ACAVAMA, Escuela de Matemática UCR (2021-2023) & Inglés (avanzado); Portugués (avanzado) \\ \hline
2. Acuña Valverde, Luis & \textit{Pendiente} & luis.\allowbreak{}acunavalverde@\allowbreak{}ucr.\allowbreak{}ac.\allowbreak{}cr & 21 & Bachillerato en Matemática;\newline Doctor of Philosophy & \textbullet~Estimates concerning the heat content for the Schrödinger operator related to a subordinate Brownian motion (2025). Journal of Mathematical Analysis and Applications, 543(2)\newline \textbullet~On the Asymptotic Expansion of the Heat Content for Isotropic alpha-Stable Processes Over Convex Bodies (2023). Journal of Geometric Analysis, 33, 25\newline \textbullet~Spectral bounds for the heat flow related to subordinate Brownian motions over bounded regular domains (2022). Preprint (2022) & -- & Inglés (avanzado); Francés (intermedio); Italiano (intermedio) \\ \hline
3. Alfaro Murillo, Jorge & \textit{Pendiente} & jorge.\allowbreak{}alfaro@\allowbreak{}ucr.\allowbreak{}ac.\allowbreak{}cr & 10 & Bach. Matemática;\newline Bach. Economía;\newline Maestría en Estadística;\newline Doctor en Filosofía con esp. en Matemáticas & \textbullet~Cancer evolution in its native host: oncogenic selection and epistasis revealed by human tumor genomes (2026). Journal of Evolutionary Biology (publicación anticipada)\newline \textbullet~Age reshapes selection on melanoma driver mutations (2026). Science Advances (en prensa)\newline \textbullet~Narrative shielding and prebunking: a synergistic strategy against vaccine hesitancy in pediatrics (2026). Expert Review of Vaccines, 25(1):2694390\newline \textit{(6 en total; se listan las 3 principales; el detalle completo está en el CV)} & \textbullet~Investigación: Development of Computational Tools for Mutation Rate Estimation and Selection Inference in Cancer Genomics, investigador principal, proyecto 821-C6-198, Vicerrectoría de Investigación UCR (2026-2028)\newline \textbullet~Consultoría: Coordinador de ED-3480 Servicios de Asesoría y Capacitación en Matemática, CIMPA UCR (2026-actualidad)\newline \textbullet~Consultoría: Estimación conjunta de tasas de mutación somática y coeficientes de selección en cáncer, Townsend Lab, Yale University (2025)\newline \textit{(5 en total; se listan las 3 principales; el detalle completo está en el CV)} & Inglés (fluido); Alemán (avanzado) \\ \hline
4. Ávila Artavia, Josué David & \textit{Pendiente} & josue.\allowbreak{}avila@\allowbreak{}ucr.\allowbreak{}ac.\allowbreak{}cr & 12 & Bachillerato en Matemática;\newline Doctorado académico en Matemática Pura & \textbullet~On the Iwasawa lambda-invariant of the cyclotomic Z2-extension of a family of real quadratic fields in which 2 splits (2026). Preprint (2026), aceptado para publicación en International Journal of Number Theory, arXiv:2605.09111\newline \textbullet~Iwasawa module of the cyclotomic Z2-extension of certain real quadratic fields (2025). The Ramanujan Journal, 67(1), artículo 6 & -- & Inglés (avanzado) \\ \hline
5. Barboza Chinchilla, Luis Alberto & \textit{Pendiente} & luisalberto.\allowbreak{}barboza@\allowbreak{}ucr.\allowbreak{}ac.\allowbreak{}cr & 23 & Bach. Ciencias Actuariales;\newline Maestría en Mat. Aplicada;\newline M.Sc. Statistics Computational Finance;\newline Doctor of Philosophy (Statistics) & \textbullet~Amortized Bayesian inference for spatio-temporal extremes: A copula factor model with autoregression (2026). Spatial Statistics, 76(101052)\newline \textbullet~Hurricane Impact Index: Capturing direct and indirect cyclone effects in Central America (2026). Earth Systems and Environment\newline \textbullet~Bayesian spatio-temporal model with INLA for dengue fever risk prediction in Costa Rica (2023). Environmental and Ecological Statistics\newline \textit{(7 en total; se listan las 3 principales; el detalle completo está en el CV)} & -- & -- \\ \hline
6. Barquero Sánchez, Adrián & \textit{Pendiente} & adrian.\allowbreak{}barquero\_s@\allowbreak{}ucr.\allowbreak{}ac.\allowbreak{}cr & 16 & Bachillerato en Matemática;\newline Ph.D. in Mathematics & \textbullet~The distribution of k-free ideals in ray class groups (2026). Preprint (2026), en revisión, arXiv:2608.13509\newline \textbullet~Counting elliptic curves over Q with bounded naive height (2025). Preprint (2025), en revisión, arXiv:2506.18874\newline \textbullet~On the efficient computation of Fourier coefficients of eta-quotients (2025). Preprint (2025), en revisión, arXiv:2504.01384\newline \textit{(7 en total; se listan las 3 principales; el detalle completo está en el CV)} & \textbullet~Docencia: Instructor de los tutoriales de SageMath, CIMPA School on Combinatorial Structures in Geometry, Graphs and Symmetry, San José (27 de julio a 7 de agosto de 2026)\newline \textbullet~Investigación: Pasantía de investigación y charla en el coloquio, Bucknell University Distinguished Visiting Professor Program (18-26 de abril de 2026)\newline \textbullet~Investigación: Taller de investigación Rethinking Number Theory 6, virtual (16-27 de junio de 2025)\newline \textit{(4 en total; se listan las 3 principales; el detalle completo está en el CV)} & Inglés (avanzado); Alemán (principiante-intermedio); Francés (principiante); Japonés (principiante) \\ \hline
7. Bermúdez Aguilar, Esteban & 1-1121-0325 & esteban.\allowbreak{}bermudez@\allowbreak{}ucr.\allowbreak{}ac.\allowbreak{}cr & 10 & Bachillerato en Ciencias Actuariales;\newline Maestría en Administración de Negocios & \textit{Ninguna en 2022--2026 según el CV} & \textbullet~Consultoría: Asesoría actuarial a varias empresas (2008-actualidad)\newline \textbullet~Docencia: Curso Financiamiento en la seguridad social, CISS-CIESS (2026)\newline \textbullet~Docencia: Curso Lenguaje Actuarial en la Seguridad Social, CISS-CIESS (2024) & -- \\ \hline
8. Bolaños Guerrero, Raúl Esteban & 603660757 & raulesteban.\allowbreak{}bolanos@\allowbreak{}ucr.\allowbreak{}ac.\allowbreak{}cr & 17 & Bach. Ciencias Actuariales;\newline Bach. Matemática;\newline Master of Arts Mathematics;\newline Doctor of Philosophy Mathematics & \textit{Ninguna en 2022--2026 según el CV} & \textbullet~Investigación: C1257 Fluctuaciones gaussianas para la ecuación de calor estocástica con condición inicial aleatoria, participación principal, terminado (11/12/2023-31/12/2023) & Inglés (conversación 100\%, lectura 100\%, escritura 100\%); Portugués (conversación 25\%, lectura 50\%, escritura 25\%) \\ \hline
9. Bustamante Medina, Ronald & \textit{Pendiente} & \textit{Pendiente} & 27 & Bachillerato en Matemática;\newline Doctorado en Lógica Matemática y Fundamentos de la Informática & \textbullet~Groups Definable in Partial Differential Fields with an Automorphism (2026). Israel Journal of Mathematics\newline \textbullet~Zilber's Dichotomy for Differentially Closed Fields with an Automorphism (2024). Preprint & \textbullet~Investigación: Grupos definibles en teorías con operadores (en curso, sin fecha de inicio en el CV)\newline \textbullet~Docencia: Curso de lógica en la Maestría en Bioética, Facultad de Filosofía, Universidad Nacional (mayo-agosto 2022) & -- \\ \hline
10. Calvo Alpízar, Juan Gabriel & \textit{Pendiente} & juan.\allowbreak{}calvo@\allowbreak{}ucr.\allowbreak{}ac.\allowbreak{}cr & 18 & Bach. Matemática;\newline Bach. Enseñanza Mat.;\newline Bach. Ing. Eléctrica;\newline Doctor en Filosofía & \textbullet~Coarse spaces for virtual element methods on irregular 3D subdomain decompositions (2026). Actas de congreso de descomposición de dominio, por aparecer (arXiv)\newline \textbullet~Stochastic dynamics in epidemiological models with interconnected populations (2026). Journal of Biological Systems\newline \textbullet~A Schwarz preconditioner for the smoothed lowest-order virtual element method (2025). Computers and Mathematics with Applications, 197\newline \textit{(17 en total; se listan las 3 principales; el detalle completo está en el CV)} & \textbullet~Investigación (investigador principal): C5054 Iterative substructuring methods for elliptic PDEs (enero 2025-en curso)\newline \textbullet~Investigación (investigador principal): B8290 Machine learning methods for the design of preconditioners (julio 2021-diciembre 2024)\newline \textbullet~Investigación (investigador asociado): C3190 The effect of social determinants on the spread of infectious diseases (abril 2023-en curso)\newline \textit{(5 en total; se listan las 3 principales; el detalle completo está en el CV)} & Inglés (fluido) \\ \hline
11. Cambronero Villalobos, Santiago & 5 0233 0650 & \textit{Pendiente} & 38 & Bach. Matemática;\newline Lic. Matemática;\newline Maestría en Matemática;\newline Doctorado en Matemática & \textbullet~Cylindrical martingale-valued measures, stochastic integration and SPDE's (2024). Stochastic Partial Differential Equations: Analysis and Computations\newline \textbullet~On the suppremum of a family of set functions (2023). Preprint (2023)\newline \textbullet~Itô's Formula for Itô processes defined with respect to a cylindrical-martingale valued measure (2023). Preprint (2023) & -- & -- \\ \hline
12. Campos Fernández, José David & 602830226 & josedavid.\allowbreak{}campos@\allowbreak{}ucr.\allowbreak{}ac.\allowbreak{}cr & 21 & Bach. Enseñanza Mat.;\newline Lic. Enseñanza Mat.;\newline Doctorado en Matemática & \textbullet~The martingale representation theorem for cylindrical martingale valued measures (2026). Preprint (arXiv)\newline \textbullet~Itô's Formula for Itô processes defined with respect to a cylindrical-martingale valued measure (2026). Stochastic Analysis and Applications\newline \textbullet~Markov property and path regularity for the solutions to SPDEs driven by cylindrical-martingale valued measures (2026). Preprint (arXiv)\newline \textit{(5 en total; se listan las 3 principales; el detalle completo está en el CV)} & -- & Inglés (intermedio-avanzado) \\ \hline
13. Campos Salas, Daniel & \textit{Pendiente} & daniel.\allowbreak{}campossalas@\allowbreak{}ucr.\allowbreak{}ac.\allowbreak{}cr & 12 & Bachillerato en Matemática;\newline Master of Science;\newline Doctor of Philosophy & \textit{Ninguna en 2022--2026 según el CV} & \textbullet~Acción social: Proyecto ED-3493 Olimpiadas Universitarias de Matemática, instructor desde 2019 y coordinador general desde 2024 (2019-actualidad)\newline \textbullet~Acción social: Comisión Organizadora de las Olimpiadas Costarricenses de Matemáticas, OLCOMA (desde 2019)\newline \textbullet~Acción social: Cibercoloquio Latinoamericano de Matemáticas, organización (2020-2023)\newline \textit{(5 en total; se listan las 3 principales; el detalle completo está en el CV)} & -- \\ \hline
14. Chaves Aguilar, Santiago & \textit{Pendiente} & anthony.\allowbreak{}chavesaguilar@\allowbreak{}ucr.\allowbreak{}ac.\allowbreak{}cr & 11 & Bachillerato en Matemática;\newline Master in Science;\newline Doctor of Philosophy in Pure Mathematics & \textit{Ninguna en 2022--2026 según el CV} & \textbullet~Investigación: La conjetura de Schmid-Vilonen: un estudio de formas invariantes en G-módulos irreducibles de dimensión infinita, Universidad de Costa Rica (2022-2024)\newline \textbullet~Investigación: Avances en la teoría de factorización de vectores suaves y analíticos de representaciones de grupos de Lie, Proyecto Semilla, Universidad de Costa Rica (2024-2026) & Inglés; Portugués \\ \hline
15. Coto Mora, Edward & \textit{Pendiente} & edward.\allowbreak{}coto@\allowbreak{}ucr.\allowbreak{}ac.\allowbreak{}cr & 10 & Bachillerato en Matemática;\newline Matemáticas y Aplicaciones;\newline Doctorado en Matemáticas & \textbullet~A weak form of Beilinson's conjecture for residues of Hasse-Weil L-functions attached to GSp4 x GL2 (2026). Preprint (2026)\newline \textbullet~Classification of cyclic holomorphs with perfect order subsets (2026). Preprint (2026)\newline \textbullet~169 and Sums of Positive Squares, Part II (2026). Preprint (2026) & -- & Inglés (C1); Francés (C1) \\ \hline
16. Fonseca Mora, Christian & 1-1320-0538 & christianandres.\allowbreak{}fonseca@\allowbreak{}ucr.\allowbreak{}ac.\allowbreak{}cr & 12 & Bach. Matemática;\newline Bach. Economía;\newline Bach. Ciencias Actuariales;\newline Lic. Economía;\newline Doctorado en Matemática & \textbullet~Itô's Formula for Itô processes defined with respect to a cylindrical-martingale valued measure (2026, en prensa). Stochastic Analysis and Applications\newline \textbullet~Itô's formula for Lévy-Itô processes taking values in the dual of nuclear space (2026). Preprint (ArXiv:2608.05514)\newline \textbullet~Markov property and path regularity for the solution to SPDEs driven by a cylindrical martingale-valued measure (2026). Preprint (ArXiv:2606.12381)\newline \textit{(12 en total; se listan las 3 principales; el detalle completo está en el CV)} & -- & -- \\ \hline
17. Guevara Villalobos, Álvaro & \textit{Pendiente} & \textit{Pendiente} & 13 & Bachillerato en Matemática;\newline Maestría en Matemáticas;\newline PhD en Matemáticas & \textit{Ninguna en 2022--2026 según el CV} & -- & Inglés (C2); Alemán (B1); Francés (lectura) \\ \hline
18. Guier Acosta, Jorge & \textit{Pendiente} & \textit{Pendiente} & \textit{Pendiente} & Bachillerato en Matemática;\newline DEA de Lógica Matemática;\newline Doctorat ès Sciences Mathématiques & \textbullet~Projectable reduced f-rings admitting elimination of quantifiers (2026). Séminaires et Congrès, Société Mathématique de France, 35, pp. 113-128\newline \textbullet~Elimination of quantifiers of a theory of real closed rings (2025). Annals of Pure and Applied Logic, 176, pp. 1-30 & -- & Francés (dominio avanzado al oral y al escrito); Inglés (dominio medio al oral y al escrito) \\ \hline
19. Gutiérrez Pavón, Jonathan & 1-1142-0747 & jonathan.\allowbreak{}gutierrez@\allowbreak{}ucr.\allowbreak{}ac.\allowbreak{}cr & 22 & Bach. Enseñanza Mat.;\newline Lic. Enseñanza Mat.;\newline Maestría en Probabilidad;\newline Doctorado en Probabilidad & \textbullet~A formula for the density of local time of the Brox diffusion in a time-window (2025). Random Operators and Stochastic Equations\newline \textbullet~A Schrodinger random operator with semimartingale potential (2023). Random Operators and Stochastic Equations\newline \textbullet~The Killed Brox diffusion (2022). Stochastic Models\newline \textit{(4 en total; se listan las 3 principales; el detalle completo está en el CV)} & -- & -- \\ \hline
20. Hernández Alvarado, Alberto & \textit{Pendiente} & albertojose.\allowbreak{}hernandez@\allowbreak{}ucr.\allowbreak{}ac.\allowbreak{}cr & 10 & Bachillerato en Matemática;\newline Doctorado en Matemática & \textbullet~UMAP projections and the survival of empty space: A geometric approach to high dimensional data (2025). Data Science, Classification and Artificial Intelligence for Modeling Decision Making, Springer\newline \textbullet~Isomorphisms of relative holomorphs and matrix similarity (2024). Bulletin of the Australian Mathematical Society\newline \textbullet~Minimum depth of factorization algebra extensions (2023). Revista Colombiana de Matemáticas 57(1)\newline \textit{(4 en total; se listan las 3 principales; el detalle completo está en el CV)} & \textbullet~Investigación: Advances in factorization theory for smooth and analytic vectors of Lie group representations (2024-2026)\newline \textbullet~Investigación: Indecomposable algebraic modules of a finite dimensional Hopf algebra (2021-2023)\newline \textbullet~Proyecto (coordinador): ACAVAMA, Validación de modelos computacionales para finanzas de la Universidad de Costa Rica, UCR (2021-2024)\newline \textit{(4 en total; se listan las 3 principales; el detalle completo está en el CV)} & -- \\ \hline
21. Jiménez López, David & \textit{Pendiente} & david.\allowbreak{}jimenezlopez@\allowbreak{}ucr.\allowbreak{}ac.\allowbreak{}cr & 22 & Bachillerato en Matemática;\newline Doctorado en Matemáticas & \textit{Ninguna en 2022--2026 según el CV} & -- & Inglés (fluido) \\ \hline
22. Lacy Mora, Allan & \textit{Pendiente} & allan.\allowbreak{}lacy@\allowbreak{}ucr.\allowbreak{}ac.\allowbreak{}cr & 20 & Bachillerato en Matemática;\newline Doctor of Philosophy & \textbullet~Dynatomic Galois groups for a family of quadratic rational maps (2024). International Journal of Number Theory 20(7) & -- & -- \\ \hline
23. Loría Sorio, Jennifer de los Ángeles & \textit{Pendiente} & jenniferdelosange.\allowbreak{}loria@\allowbreak{}ucr.\allowbreak{}ac.\allowbreak{}cr & 4 & Bachillerato en Matemática;\newline Maestría en Matemática;\newline Doctorado en Ciencias & \textbullet~On the estimation of the time-dependent transmission rate in epidemiological models (2025). Inverse Problems, vol. 41, no. 6\newline \textbullet~Time-dependent vaccine efficacy estimation quantified by a mathematical model (2023). Plos One, vol. 18, no. 5\newline \textbullet~A specific prediction model for deep sternal wound infection after coronary artery bypass surgery: Prediction model for DSWI after CABG (2022). Scientific Reports, vol. 12, no. 1 & -- & Portugués (lectura y conversación avanzadas); Inglés (comprensión lectora intermedia y conversación básica) \\ \hline
24. Mena Arias, Darío Alberto & \textit{Pendiente} & dario.\allowbreak{}menaarias@\allowbreak{}ucr.\allowbreak{}ac.\allowbreak{}cr & 16 & Bach. Matemática;\newline Bach. Enseñanza Mat.;\newline Maestría en Matemática;\newline Doctorado en Matemática & \textbullet~Itô's formula for Itô processes defined with respect to a cylindrical-martingale valued measure (2026). Stochastic Analysis and Applications (Stoch. Anal. Appl.)\newline \textbullet~The martingale representation theorem for cylindrical martingale valued measures (2026). Preprint (arXiv:2607.22441)\newline \textbullet~Markov property and path regularity for the solutions to SPDEs driven by cylindrical-martingale valued measures (2026). Preprint (arXiv:2606.12381)\newline \textit{(8 en total; se listan las 3 principales; el detalle completo está en el CV)} & -- & Inglés (dominio del 100 \%) \\ \hline
25. Méndez Hernández, Pedro & \textit{Pendiente} & pedro.\allowbreak{}mendez@\allowbreak{}ucr.\allowbreak{}ac.\allowbreak{}cr & 31 & Bach. Matemática;\newline Lic. Matemática;\newline Ph.D. en Matemática & \textit{Ninguna en 2022--2026 según el CV} & -- & -- \\ \hline
26. Montalto Cruz, Carlos & \textit{Pendiente} & \textit{Pendiente} & 22 & Bachillerato en Matemática;\newline Ph.D en Matemáticas & \textbullet~A SharpView With Enhanced Out-of-Focus Text Legibility for Augmented Reality Systems (2024). IEEE VR & \textbullet~Acción social: Olimpiadas Costarricenses de Matemáticas (OLCOMA), Vicerrectoría de Acción Social, UCR, coordinador (vigente; el CV no indica fechas de vigencia)\newline \textbullet~Investigación: desarrollo de algoritmos de imagen fotoacústica, Laboratorio de Investigación en Biomedicina, UCR (en curso desde setiembre 2017) & Inglés (nivel profesional) \\ \hline
27. Montenegro Guzmán, Samaria & \textit{Pendiente} & samaria.\allowbreak{}montenegroguzman@\allowbreak{}ucr.\allowbreak{}ac.\allowbreak{}cr & 9 & Bach. Matemática;\newline Maestría en Lógica Mat.;\newline Doctorado en Lógica Matemática y Fundamentos de la Informática & \textbullet~Dividing lines among fields (2026). Women in Mathematics in Latin America, Trends in Mathematics, Research Perspectives, Ghent Analysis and PDE Center vol. 14, Birkhäuser/Springer, pp. 111-118\newline \textbullet~Pseudo T-closed fields (2025). Model Theory, 4(1), 1-35\newline \textbullet~Groups definable in partial differential fields with an automorphism (2026). Israel Journal of Mathematics (aceptado para publicación) & \textbullet~Investigación (Fondo UCREA, Vicerrectoría de Investigación, UCR): Generizar para transformar el conocimiento: Construcción de una metodología feminista para el análisis de género de la producción académica en la Universidad de Costa Rica (2022-2024)\newline \textbullet~Investigación (Fondo semilla): Nuevas perspectivas en el estudio de grupos dentro del universo de Shelah, investigadora asociada (2020-2022) & -- \\ \hline
28. Mora Sáenz, Germán José & \textit{Pendiente} & german.\allowbreak{}morasaenz@\allowbreak{}ucr.\allowbreak{}ac.\allowbreak{}cr & \textit{Pendiente} & Bachillerato en Matemática;\newline Master of Science Matemática;\newline Doctor of Philosophy Matemática & \textit{Ninguna en 2022--2026 según el CV} & -- & Inglés \\ \hline
29. Ramírez González, José Alexander & \textit{Pendiente} & alexander.\allowbreak{}ramirez\_g@\allowbreak{}ucr.\allowbreak{}ac.\allowbreak{}cr & 27 & Bach. Matemática;\newline Bach. Ing. Eléctrica;\newline Doctor en Filosofía & \textit{Ninguna en 2022--2026 según el CV} & \textbullet~Consultoría/desarrollo: Quanto Soluciones Analíticas y Cuantitativas, presidente y desarrollador de software para análisis de riesgo (diciembre 2014-presente)\newline \textbullet~Extensión: Inteligencia Artificial y riesgo de crédito: el uso de datos y modelos para la toma de decisiones a lo largo del ciclo de vida, Risk Day, Santo Domingo, República Dominicana (2026)\newline \textbullet~Extensión: Riesgo de liquidez en cooperativas, Congreso sobre Riesgo Integral en Cooperativas, Panamá (2026)\newline \textit{(4 en total; se listan las 3 principales; el detalle completo está en el CV)} & -- \\ \hline
30. Rodríguez Rojas, Oldemar & \textit{Pendiente} & oldemar.\allowbreak{}rodriguez@\allowbreak{}ucr.\allowbreak{}ac.\allowbreak{}cr & \textit{Pendiente} & Bach. Matemática;\newline Lic. Matemática;\newline Lic. Matemáticas de la Decisión;\newline Doctorado en Matemática & \textit{Ninguna en 2022--2026 según el CV} & -- & Inglés (lee, escribe y habla); Francés (lee, escribe y habla) \\ \hline
31. Rosales Ortega, José & \textit{Pendiente} & jose.\allowbreak{}rosales@\allowbreak{}ucr.\allowbreak{}ac.\allowbreak{}cr & 17 & Bachillerato en Matemática;\newline Maestría en Matemática Pura;\newline Doctorado en Matemática Pura & \textbullet~Toeplitz operators with invariant symbols under the action of some subgroups of Sn semidirect Tn on the Fock Space on Cn (aceptado, sin año en el CV). Tales and Trends in Operator Theory, Evolution Equations and Complex Variables: IWOTA 2024, Birkhäuser/Springer\newline \textbullet~Toeplitz operators with symmetric, alternating and anti-symmetric separately radial symbols on the unit ball (2024). Boletín de la Sociedad Matemática Mexicana & \textbullet~Investigación: Operadores Integrales en Espacios de Funciones Analíticas, proyecto CIMPA-UCR C6165, investigador principal (2026-2029)\newline \textbullet~Investigación: Interacciones del análisis y la geometría, proyecto CIMPA-UCR, investigador principal (2023-2024) & -- \\ \hline
32. Sánchez Chavarría, Adriana & \textit{Pendiente} & adriana.\allowbreak{}sanchez\_c@\allowbreak{}ucr.\allowbreak{}ac.\allowbreak{}cr & 9 & Bachillerato en Matemática;\newline Doctorado en Matemática & \textbullet~Analyticity of the Lyapunov exponents of random products of quasi-periodic cocycles (2023). Nonlinearity 36(6)\newline \textbullet~The Continuity Problem of Lyapunov Exponents (2022). International Workshop New Trends in Lyapunov Exponents, Springer International Publishing & \textbullet~Investigación: Proyecto de Educación Continua del CIMPA, investigadora principal (2024-2026)\newline \textbullet~Investigación: Uso de modelos de aprendizaje automático para el pronóstico de caudales extremos (mínimos y máximos), investigadora principal (2023-2025)\newline \textbullet~Investigación: Estudio de sistemas de producción en contexto de industria 4.0 con enfoque al estudio del trabajo apoyado por herramientas de aprendizaje profundo, investigadora asociada (2022-2025)\newline \textit{(5 en total; se listan las 3 principales; el detalle completo está en el CV)} & Inglés (avanzado); Portugués (avanzado); Francés (iniciante) \\ \hline
33. Sánchez Guevara, Jesús & 603340326 & jesus.\allowbreak{}sanchez\_g@\allowbreak{}ucr.\allowbreak{}ac.\allowbreak{}cr & 12 & Bach. Matemática;\newline Maestría con Esp. Lógica Mat.;\newline En Mat. Pura: Topología Algebraica y Teoría de Categorías & \textbullet~Extensión de una estructura de óperad definida sobre un sub-S-módulo (2025). Revista de Matemática: Teoría y Aplicaciones 32(2)\newline \textbullet~Introducción a la teoría de categorías, lógica y topos elementales para una mente curiosa (2024). Revista de Filosofía de la Universidad de Costa Rica 63(125)\newline \textbullet~On the relationship between expansion angle of Earth-directed CMEs and soft X-ray emission from their related flare (2023). Revista de Matemática: Teoría y Aplicaciones 30(2)\newline \textit{(4 en total; se listan las 3 principales; el detalle completo está en el CV)} & \textbullet~Investigación: Herramientas científicas para optimizar el estudio del clima espacial, código C0072, participación asociado (2020-2022) & Francés (conversación 100\%, lectura 100\%, escritura 100\%); Inglés (conversación 80\%, lectura 100\%, escritura 100\%); Italiano (conversación 20\%, lectura 70\%, escritura 50\%) \\ \hline
34. Sánchez Peña, Fabio & \textit{Pendiente} & fabio.\allowbreak{}sanchez@\allowbreak{}ucr.\allowbreak{}ac.\allowbreak{}cr & 16 & PhD in Computer Science;\newline Biological Statistics and Computational Biology & \textbullet~High-resolution climate-dengue modeling and mid-century projections under SSP5-8.5 in Costa Rica (2026). Science in One Health 5, 100157\newline \textbullet~Stochastic Dynamics in Epidemiological Models with Interconnected Populations (2026). Journal of Biological Systems 34\newline \textbullet~Managing populations after a disease outbreak: exploration of epidemiological consequences of managed host reintroduction following disease-driven host decline (2026). Philosophical Transactions of the Royal Society B 381(1945)\newline \textit{(27 en total; se listan las 3 principales; el detalle completo está en el CV)} & \textbullet~Investigación: El costo de una vida digna en Costa Rica: Ingreso vital y la construcción de metodologías para el cálculo de variaciones territoriales, UCREA, Universidad de Costa Rica (2023)\newline \textbullet~Investigación: Mathematical Models for Dengue Surveillance: Challenges and Opportunities, Seed Grant, University of California, Davis (2022)\newline \textbullet~Consultoría: PREDICDATA, firma de análisis de datos y consultoría, fundador y socio (2026-presente)\newline \textit{(5 en total; se listan las 3 principales; el detalle completo está en el CV)} & Inglés (fluido) \\ \hline
35. Segura Ugalde, Esteban & \textit{Pendiente} & esteban.\allowbreak{}seguraugalde@\allowbreak{}ucr.\allowbreak{}ac.\allowbreak{}cr & 12 & Bach. Matemática;\newline Maestría en Optimización Matemática;\newline Doctorado en Matemática & \textbullet~Construcción de diseños de bloques balanceados e incompletos (2026, aceptado). Matemática, Educación e Internet, TEC\newline \textbullet~Computation of Matrix p-th Root Using Moments (2024). Journal of Computational and Applied Mathematics\newline \textbullet~A Subspace Method for Time Series Anomaly Detection in Cyber-Physical Systems (2023). IFAC-PapersOnLine 55(41), 58-63\newline \textit{(5 en total; se listan las 3 principales; el detalle completo está en el CV)} & \textbullet~Consultoría: ACAVAMA, CIMPA, Universidad de Costa Rica (2021-presente) & Inglés (fluente, TOEIC 845/990); Francés (medio, DELF A2) \\ \hline
36. Solís Chacón, Maikol & \textit{Pendiente} & maikol.\allowbreak{}solis@\allowbreak{}ucr.\allowbreak{}ac.\allowbreak{}cr & 17 & Bach. Ciencias Actuariales;\newline Maestría Mat. Aplicada;\newline Master 2 Recherche;\newline Doctorat Statistique et Mathematiques Apliquées & \textbullet~Machine learning-driven COVID-19 early triage and large-scale testing strategies based on the 2021 Costa Rican Actualidades survey (2025). Mathematics in Medical and Life Sciences, 2(1), 2494001\newline \textbullet~Exploring the coffee fruit microbiome under different management practices and its correlation with beverage quality criteria in Costa Rica (2025). Current Microbiology, 83(2)\newline \textbullet~Relative weight analysis with residualization for detecting relevant non linear effects (2025). Communications in Statistics - Theory and Methods, 54(23), 7553–7568\newline \textit{(6 en total; se listan las 3 principales; el detalle completo está en el CV)} & \textbullet~Investigación: Machine learning-driven coffee quality modeling using its chemical, physical, and sensory attributes, investigador principal, CIMPA, N.º 821-C5240 (2025–2028)\newline \textbullet~Investigación: Clasificador de riesgo para individuos susceptibles a COVID-19, investigador principal, CIMPA, N.º 817-C2198 (2023–2025)\newline \textbullet~Investigación: Verificación de la sensibilidad in vitro de kits comerciales de RT-LAMP para la detección del virus SARS-CoV-2, investigador colaborador, INIFAR, N.º 817-C0530 (2022)\newline \textit{(4 en total; se listan las 3 principales; el detalle completo está en el CV)} & Inglés (fluido); Francés (fluido) \\ \hline
37. Tall, El Hadji Yaya & \textit{Pendiente} & el.\allowbreak{}tall@\allowbreak{}ucr.\allowbreak{}ac.\allowbreak{}cr & 2 & Bach. Matemática;\newline Lic. Matemática;\newline Maestría Matemática;\newline Doctorado en Matemática;\newline Diploma ICTP & \textbullet~Hölder regularity of the top Lyapunov exponent for random matrix products with low-dimensional equator (2026). Preprint (arXiv:2610.17153)\newline \textbullet~Log-Hölder continuity at zero Lyapunov gap for finite state Markov GL(2)-cocycles (2026). Preprint (arXiv:2608.22157)\newline \textbullet~Analyticity of Lyapunov Exponents for Mixed Markov Quasi-Periodic Cocycles (2026). Preprint (arXiv:2608.01569)\newline \textit{(4 en total; se listan las 3 principales; el detalle completo está en el CV)} & -- & Peul (lengua materna); Francés (nativo); Wolof (nativo); Inglés (fluido); Portugués (fluido); Italiano (intermedio) \\ \hline
38. Ugalde Gómez, William Javier & \textit{Pendiente} & william.\allowbreak{}ugalde@\allowbreak{}ucr.\allowbreak{}ac.\allowbreak{}cr & 33 & Bach. Matemática;\newline Lic. Matemática;\newline Master en Ciencias;\newline Diploma ICTP;\newline Doctor en Filosofía (PhD) & \textbullet~Symmetry of some noncommutative sphere algebras (2026). Preprint (arXiv:2602.06946), por aparecer en Journal of Algebra\newline \textbullet~Characterization of principal bundles: the noncommutative algebraic case (2024). Preprint (arXiv:2402.16852)\newline \textbullet~Characterization of principal bundles: the commutative case (2023). Preprint (arXiv:2310.03824) & \textbullet~Investigación (investigador principal, CIMPA, Vicerrectoría de Investigación UCR): Simetrías de campos cuánticos libres (enero 2021--diciembre 2023)\newline \textbullet~Investigación (investigador principal, CIMPA, Vicerrectoría de Investigación UCR): Aspectos de simetría en sistemas cuánticos (enero 2024--diciembre 2026) & -- \\ \hline
39. Ulloa Esquivel, Roberto & 4-0200-0107 & \textit{Pendiente} & 18 & Bachillerato en Matemática;\newline Doctor of Philosophy in Mathematics & \textit{Ninguna en 2022--2026 según el CV} & -- & Inglés (fluido; TOEFL 111/120, 2012); Alemán (certificación A2b, Goethe Center, 2011) \\ \hline
40. Vargas Ulate, Eny & 1-1440-0739 & \textit{Pendiente} & 2 & Bach. Ciencias Actuariales;\newline Lic. Ciencias Actuariales;\newline MBA con énfasis en Finanzas & \textit{Ninguna en 2022--2026 según el CV} & -- & Inglés (intermedio–avanzado); Portugués (básico) \\ \hline
41. Villalobos Alvarado, Jorge & \textit{Pendiente} & jorge.\allowbreak{}villalobosalvarado@\allowbreak{}ucr.\allowbreak{}ac.\allowbreak{}cr & 9 & Bach. Matemática;\newline Bach. Física;\newline Maestría en Matemática;\newline Doctorado en Matemática & \textbullet~Stable defect states in the continuous spectrum of bilayer graphene with magnetic field (2023). Physica D: Nonlinear Phenomena 455, 133891 & -- & -- \\ \hline
42. Villalobos Arias, Mario & 204210883 & mario.\allowbreak{}villalobos@\allowbreak{}ucr.\allowbreak{}ac.\allowbreak{}cr & 39 & Bachillerato en Matemática;\newline Doctorado en Matemática & \textbullet~Análisis estadístico y matemática aplicada en la construcción del índice de riesgo en la sostenibilidad de la prestación del servicio de agua y saneamiento, Costa Rica (2022). Revista Tecnología y Ciencias del agua, México & \textbullet~Investigación: Laboratorio de ingeniería en transferencia y arquitectura estratégica agropecuaria, Estación Experimental de Ganado Lechero Alfredo Volio, Universidad de Costa Rica (2023-2024)\newline \textbullet~Investigación: Red interdisciplinaria para el fortalecimiento de las capacidades en países centroamericanos para el desarrollo de una prueba diagnóstica para el tamizaje proactivo en focos de infección del virus SARS-CoV-2 empleando la tecnología RT-LAMP, INIFAR, Universidad de Costa Rica (2021-2022) & -- \\ \hline
43. Zamora Calero, Rafael & \textit{Pendiente} & rafael.\allowbreak{}zamora\_c@\allowbreak{}ucr.\allowbreak{}ac.\allowbreak{}cr & 11 & Bach. Matemática;\newline Maestría en Lógica Matemática;\newline Doctorado en Matemática & \textbullet~Strong Factorization Theorem for Smooth Vectors of Exponential Solvable Lie Group Representations (2026). Preprint (arXiv:2604.20998), enviado a Journal of Functional Analysis & \textbullet~Investigación (investigador principal, Fondo Semilla, Vicerrectoría de Investigación UCR): Nuevas perspectivas en el estudio de grupos dentro del universo de Shelah (2020--2023)\newline \textbullet~Investigación (colaborador, Fondo Semilla, Vicerrectoría de Investigación UCR): Avances en la Teoría de Factorización de Vectores Suaves y Analíticos de Representaciones de Grupos de Lie (2024) & -- \\ \hline
44. Zamora Luna, Oscar & \textit{Pendiente} & oscar.\allowbreak{}zamoraluna@\allowbreak{}ucr.\allowbreak{}ac.\allowbreak{}cr & 6 & Bachillerato en Matemática;\newline Doctor of Philosophy in Mathematics and its Applications & \textbullet~The Maximum Number of Pentagons in a Planar Graph (2025). Journal of Graph Theory 108(2), 229-256\newline \textbullet~Planar Turán Number of the theta6 (2024). Studia Scientiarum Mathematicarum Hungarica 61(2), 89-115\newline \textbullet~Packing independent cliques into planar graphs (2024). Preprint (arXiv:2408.03298)\newline \textit{(10 en total; se listan las 3 principales; el detalle completo está en el CV)} & -- & -- \\ \hline
45. Zúñiga Rojas, Ronald Alberto & \textit{Pendiente} & ronald.\allowbreak{}zunigarojas@\allowbreak{}ucr.\allowbreak{}ac.\allowbreak{}cr & 18 & Bachillerato en Matemática;\newline Doctor en Matemática & \textbullet~Introducción a la teoría de categorías, lógica y topos elementales para una mente curiosa (2024). Revista de Filosofía de la Universidad de Costa Rica 63(165), 71-85\newline \textbullet~Stratifications on the Nilpotent Cone of the Hitchin Map (2022). Revista Matemática Complutense 35, 311-321 & \textbullet~Investigación: Costarican-German Collaboration in Research: Character varieties and Higgs bundles (Colaboración en Investigación Costa Rica-Alemania: Variedades de caracter y fibrados de Higgs), proyecto 821-C4-650, iniciativa DFG-CONARE (vigente) & Portugués (alto, competencia bilingüe); Inglés (alto, competencia bilingüe); Italiano (alto, competencia bilingüe) \\ \hline
\end{longtable}
\endgroup
{\footnotesize\noindent\textit{Fuente:} currículos vitae aportados por el personal docente y Tabla de docentes del Departamento de Matemática (octubre de 2026). ``Pendiente'' indica que el dato no consta en el currículo aportado o que el currículo aún no se ha recibido.\par}
{\footnotesize\noindent\textsuperscript{a} Nombre, año y medio de publicación; se consideran 2022 a 2026. \textsuperscript{b} Tipo, nombre y vigencia; se consideran 2022 a 2026. Se omite la columna de años de aprendizaje en entornos virtuales porque la carrera se ofrece en modalidad presencial. Los años de experiencia docente universitaria se calcularon a 2026 a partir de los periodos que declara cada currículo.\par}
\end{landscape}


% %%%%%%%%%%%%%%%%%%%%%%%%%%%%%%%%%%%%%%%%%%%%%%%%%%
% \section{Otros aspectos}

% \subsection*{Plan de transición}
% El nuevo plan entra en vigencia en 2028. El plan 02 se mantiene activo hasta 2030 y se inactiva posteriormente; quienes no lo hayan concluido al año 2031 inclusive pasan de oficio al plan 2028, con las convalidaciones de la tabla correspondiente. La estrategia se define por cohorte (C5 y anteriores, C6, C7, y C8 en adelante, que ingresa directamente al nuevo plan), de acuerdo con la Circular VD-18-2015. La proyección de carga docente en transición es de entre 10 y 12 tiempos completos por año (2028--2031).

% \subsection*{Recursos y factibilidad}
% La Escuela cuenta con aulas y oficinas en el edificio de Física-Matemática y en el edificio del CIMPA, dos laboratorios de cómputo, una licencia de software por estudiante y acceso bajo demanda a recursos bibliográficos (SIBDI). El personal administrativo es el que ya atiende la carrera vigente. Las prácticas profesionales se ofrecen en modalidad académica o profesional, con apoyo de la Vicerrectoría de Vida Estudiantil para pasantías e intercambio estudiantil.

% \subsection*{Seguimiento y evaluación}
% El Departamento de Matemática estará a cargo de la implementación. Se propone una comisión de evaluación curricular que dé seguimiento desde el primer ciclo de vigencia (revisión anual de guías metodológicas, valoración de cursos al final de cada ciclo, análisis de la evaluación docente, consulta bianual sobre la carga académica, actualización anual de indicadores de admisión, permanencia, rezago y graduación, y autoevaluación). Se prevé un diagnóstico del perfil de egreso a más tardar cinco años después de la implementación completa, y una evaluación integral con asesoría del CEA seis años después de la implementación, con consulta a personas graduadas y empleadoras. 


%%%%%%%%%%%%%%%%%%%%%%%%%%%%%%%%%%%%%%%%%%%%%%%%%%
\newpage
\appendix
\renewcommand{\thesection}{Anexo \Alph{section}}
\section{Resumen de los currículos del personal docente}

Este anexo resume los currículos vitae aportados por las personas docentes del Departamento de Matemática incluidas en la tabla de profesorado de la propuesta. Las personas se ordenan alfabéticamente por apellido. El presente resumen se basa en currículos completos, que son documentos originales aportados por cada persona.

\begin{cvblock}{Acuña Larios, Jennifer}{Coordinadora del Área de Riesgos Financieros y No Financieros, SUGEF}
\cvsec{Formación académica}
\begin{cvlist}
  \item Maestría Académica en Matemática Pura, Universidad de Costa Rica, 2021.
  \item Bachillerato en Ciencias Actuariales, Universidad de Costa Rica, 2023.
  \item Bachillerato en Matemática, Universidad de Costa Rica, 2011.
\end{cvlist}
\cvsec{Experiencia profesional y docente}
\begin{cvlist}
  \item Coordinadora del Área de Riesgos Financieros y No Financieros, Superintendencia General de Entidades Financieras (SUGEF), 2025--actualidad. Coordina procesos de supervisión de modelos, metodologías y calidad de datos de las entidades del sistema financiero (razonabilidad metodológica, desempeño, robustez, gobernanza, integridad, consistencia y trazabilidad de la información) y equipos técnicos de supervisión. Antes, en el área de Estrategia y Modelación de Riesgos, desarrolló e implementó modelos, metodologías, indicadores y herramientas analíticas para riesgo de crédito, liquidez y mercado, y coordinó el ejercicio de pruebas de estrés Bottom-Up Stress Test (BUST).
  \item Docente y coordinadora de cátedra, Universidad de Costa Rica, 2010--2015 y 2020--2024. Cálculo con aplicaciones a ciencias económicas, matemática para computación e ingeniería, probabilidad y estadística, teoría matemática del interés, instrumentos de inversión, matemática actuarial y estadística actuarial.
  \item Docente, Instituto Tecnológico de Costa Rica, 2018. Estadística y matemática para Administración de Empresas e Ingenierías.
  \item Practicante, Banco Central de Costa Rica, 2022: seguimiento diario del Mercado Integrado de Liquidez (MIL), del mercado de moneda extranjera (MONEX) y de las transacciones bancarias, con Excel, Power BI y RStudio. Practicante, Dirección de Modelos Matemáticos, Banco Nacional de Costa Rica, 2016: modelos matemáticos con MATLAB.
\end{cvlist}
\cvsec{Áreas de especialidad e investigación}
Matemática y ciencias actuariales, modelación matemática, probabilidad y estadística, análisis de datos, gestión de riesgos financieros y desarrollo de metodologías e indicadores cuantitativos. Herramientas: R y MATLAB (avanzado); SQL, Python y Excel (intermedio); Power BI.
\cvsec{Producción académica}
El CV no lista publicaciones, libros ni preprints; su producción consta como estudios y proyectos de investigación en el CIMPA (véase la sección siguiente).
\cvsec{Proyectos, consultorías y acción social}
\begin{cvlist}
  \item Investigadora, Centro de Investigación en Matemática Pura y Aplicada (CIMPA), 2020--2022. Estudios actuariales para el Fondo de Becas de la Vicerrectoría de Administración (2020), el Colegio de Ciencias Económicas (2020), Funeraria Jardines del Recuerdo (2021), el Colegio de Enfermeros y Enfermeras de Costa Rica (2022, en sus dos etapas) y el Archivo Nacional (cálculo de cesantía, 2022); análisis de ciencia de datos para la Cooperativa de Productores de Leche Dos Pinos (2021).
  \item Investigadora y validadora, ACAVAMA, Escuela de Matemática, UCR, 2021--2023. Equipo de validación de modelos del Banco Nacional de Costa Rica: revisión, validación y recomendaciones sobre modelos del sector financiero.
  \item Acción social: taller sobre pensiones para la Embajada de Canadá; colaboración en la organización del XXIV SIMMAC; capacitación sobre tasas de interés para préstamos personales e hipotecarios en tiempos de crisis (2021).
\end{cvlist}
\cvsec{Formación complementaria, distinciones e idiomas}
\begin{cvlist}
  \item EMALCA, El Salvador (2011); cursos del IMPA, Brasil (2012); curso en CIMPA sobre análisis de Fourier dihedral (2012); VI Escuela de Matemática Pura y Aplicada, UNAN León (2019); XV Congreso Latinoamericano de Probabilidad y Estadística, Mérida (2019); taller Random Media and Random Structures, Lima (2020); capacitación en Python para finanzas matemáticas, CIMPA (2022, 20 horas).
  \item Inglés avanzado (IELTS B2, 2022); portugués avanzado (Centro de Estudios Brasileños, 2008).
\end{cvlist}
\end{cvblock}

\begin{cvblock}{Acuña Valverde, Luis}{Profesor, Universidad de Costa Rica}
\cvsec{Formación académica}
\begin{cvlist}
  \item Doctor of Philosophy in Mathematics, Purdue University, Indiana, Estados Unidos, estudios 2009--2015; tutor: Rodrigo Bañuelos. Tesis: ``Heat trace and heat content asymptotics for Schrödinger Operators of stable processes/fractional Laplacians''.
  \item Bachillerato en Matemática, Universidad de Costa Rica, estudios 2002--2007.
\end{cvlist}
\cvsec{Experiencia profesional y docente}
\begin{cvlist}
  \item Profesor, Universidad de Costa Rica, 2015--actualidad. Cursos de servicio: Cálculo en una variable III, Álgebra Lineal, Ecuaciones Ordinarias, Introducción a la probabilidad y Cálculo con Optimización. Cursos de carrera: Cálculo en una variable II (series y sucesiones), Cálculo en varias variables y Ecuaciones Diferenciales Ordinarias. Curso nivelatorio de Introducción a la probabilidad para la Maestría en Ciencias Actuariales.
  \item Asistente de los cursos de Cálculo multivariado y de Álgebra lineal y ecuaciones diferenciales, y asistente de investigación, Purdue University, 2009--2015.
  \item Profesor de cursos de servicio, Universidad de Costa Rica, 2005--2009: Cálculo en una variable I, II y III, Álgebra Lineal y Ecuaciones Ordinarias.
\end{cvlist}
\cvsec{Áreas de especialidad e investigación}
Probabilidad y análisis: procesos estables y de Lévy, laplacianos fraccionarios, operadores de Schrödinger fraccionarios, contenido de calor (heat content) y asintóticas de la traza del calor, con énfasis en dominios convexos, conjuntos de perímetro finito y movimientos brownianos subordinados.
\cvsec{Producción académica}
El CV lista 14 entradas, equivalentes a 13 trabajos distintos: 12 artículos y 1 preprint (la entrada 10, preprint de 2020, corresponde al artículo de J. Geom. Anal. 2023, entrada 12). Más recientes primero:
\begin{cvlist}
  \item ``On the one-dimensional weighted heat content over intervals'', Journal of Mathematical Analysis and Applications, vol. 566 (el CV consigna 2027).
  \item ``Estimates concerning the heat content for the Schrödinger operator related to a subordinate Brownian motion'', Journal of Mathematical Analysis and Applications, vol. 543, 2025.
  \item ``Spectral bounds for the heat flow related to subordinate Brownian motions over bounded regular domains'', preprint, 2022.
  \item ``On the Asymptotic Expansion of the Heat Content for Isotropic Stable Processes Over Convex Bodies'', Journal of Geometric Analysis, 33, 2023.
  \item ``On the heat content for the Poisson heat kernel over convex bodies'', Journal of Mathematical Analysis and Applications, 494, 2021.
  \item Heat content for the Poisson kernel over the unit ball, Bulletin of the London Mathematical Society, 2020; estimates for the fractional Schrödinger operator, Journal of Spectral Theory, 2020.
  \item Heat content sobre conjuntos de perímetro finito y para procesos estables, y descomposición de funcionales aditivos de procesos de Lévy: Journal of Mathematical Analysis and Applications (dos artículos, 2016), Journal of Geometric Analysis (2016) y Stochastic Processes and their Applications (2015).
  \item Trace asymptotics for fractional Schrödinger operators, Journal of Functional Analysis, 2014; con R. Bañuelos, Potential Analysis, 2014.
\end{cvlist}
\cvsec{Conferencias, seminarios e idiomas}
\begin{cvlist}
  \item Siete conferencias y seminarios como ponente: Probability Seminar, Texas A\&M (2021, por Zoom); CIMAT, México (2017); CLAPEM, UCR (2016); Purdue Probability Seminar (2013 y 2014); Third Ohio River Analysis Meeting (2012 y 2013).
  \item Inglés avanzado; francés e italiano intermedios.
\end{cvlist}
\end{cvblock}

\begin{cvblock}{Alfaro Murillo, Jorge}{Profesor, Escuela de Matemática, e investigador, CIMPA, Universidad de Costa Rica}
\cvsec{Formación académica}
\begin{cvlist}
  \item Doctorado (PhD) en Matemáticas, Purdue University, 12/2013. Tesis: ``An epidemic model structured by the time since last infection''; director: Zhilan Feng.
  \item Maestría (MSc) en Estadística, Purdue University, 12/2013.
  \item Bachillerato en Matemática, Universidad de Costa Rica, 02/2008; Bachillerato en Economía, Universidad de Costa Rica, 02/2006.
\end{cvlist}
\cvsec{Experiencia profesional y docente}
\begin{cvlist}
  \item Profesor, Escuela de Matemática (2025--actualidad), e investigador del CIMPA (2026--actualidad), Universidad de Costa Rica. Investigador principal de un proyecto sobre herramientas computacionales para estimar tasas de mutación e inferir selección en genómica del cáncer (2026--2028); coordinador de cátedra de Cálculo para Ciencias Económicas II y de Matemática para Ciencias Económicas I.
  \item Científico de Investigación Asociado, Yale University: Departamento de Bioestadística, 2020--2024 (epistasis entre mutaciones somáticas; efectos del tabaquismo y del alcohol en el cáncer); Centro para el Modelado y Análisis de Enfermedades Infecciosas, 2015--2019 (costo-efectividad de intervenciones contra zika, dengue, citomegalovirus, neumococo e influenza). Investigador posdoctoral, Yale, 2014.
  \item Asistente de investigación, Purdue University, 2008--2013, y Arizona State University, 2010.
\end{cvlist}
\cvsec{Áreas de especialidad e investigación}
Evolución somática del cáncer (tasas de mutación, selección y epistasis), inferencia bayesiana y genómica del cáncer; dinámica de transmisión de enfermedades infecciosas y costo-efectividad de intervenciones de salud pública; modelado ecológico.
\cvsec{Producción académica}
15 publicaciones numeradas (2012--2026), todas artículos de revista; dos figuran como publicación anticipada o en prensa (2026). Más recientes primero:
\begin{cvlist}
  \item Journal of Evolutionary Biology (publicación anticipada, 2026); Science Advances (en prensa, 2026); Expert Review of Vaccines, 25(1), 2026.
  \item ``Pairwise and higher-order epistatic effects among somatic cancer mutations across oncogenesis'', Mathematical Biosciences, 366, 2023.
  \item Multimodeling approach for influenza pandemics, Proceedings of the National Academy of Sciences, 120(28), 2023; Lung Cancer, 166, 2022.
  \item ``The case for replacing live oral polio vaccine with inactivated vaccine in the Americas'', The Lancet, 395, 2020.
  \item Proceedings of the National Academy of Sciences, 116, 2019 (vacunas universales contra la influenza); Journal of Differential Equations, 267, 2019.
  \item PLOS Neglected Tropical Diseases (zika, 2016 y 2017); Vaccine, 2016; The Lancet Infectious Diseases (Ébola, 2014); Journal of Biological Dynamics, 2013; Ecological Modelling, 2012.
\end{cvlist}
\cvsec{Proyectos, consultorías y software}
\begin{cvlist}
  \item Investigador principal: proyecto 821-C6-198, Vicerrectoría de Investigación, UCR (2026--2028); subvención Discovery, Engagement and Learning, Purdue (2012). Coordinador del proyecto ED-3480 de asesoría y capacitación en matemática del CIMPA (2026--actualidad).
  \item Consultorías: Merck \& Co. (2015--2018), Sanofi (2015--2016), Ministerio de Salud de Costa Rica (2017), Ministerio de Salud de Argentina (2020--2021), Ministerio de Economía de República Dominicana (2020).
  \item Software: paquetes de Python sigmutsel, sigmutselcovs y gdcfetch (2026); cancer-epistasis (2023).
\end{cvlist}
\cvsec{Docencia, servicio, distinciones e idiomas}
\begin{cvlist}
  \item Profesor en UCR (2005--2008 y 2025--actualidad) y en Purdue (2008--2013); mentoría de estudiantes de grado, maestría en salud pública y doctorado en Yale.
  \item Consejo Editorial (2026--actualidad) y Comité Científico Internacional (2020--2025) de la Revista de Matemática: Teoría y Aplicaciones; evaluador de propuestas de la Fundación ``la Caixa'' y UKRI; más de 20 conferencias invitadas. Beca de viaje SACNAS (2011--2013).
  \item Español nativo, inglés fluido, alemán avanzado.
\end{cvlist}
\end{cvblock}

\begin{cvblock}{Ávila Artavia, Josué David}{Profesor invitado, Universidad de Costa Rica}
\cvsec{Formación académica}
\begin{cvlist}
  \item Doctorado académico en Matemática Pura, University of Maryland, Estados Unidos, 2018--2024. El CV no indica título de tesis ni director.
  \item Bachillerato en Matemática Pura con honores (9.73/10), Universidad de Costa Rica, 2013--2017.
\end{cvlist}
\cvsec{Experiencia profesional y docente}
\begin{cvlist}
  \item Profesor invitado, Universidad de Costa Rica, agosto de 2024--actualidad.
  \item Asistente de docencia, University of Maryland, agosto de 2018--mayo de 2024. Cálculo II, Introducción al Álgebra Lineal, Estadística y Probabilidad Elementales y Criptografía; calificación de evaluaciones de cursos de nivel 400 (teoría de números, cálculo avanzado, criptografía); dos años de tutorías de cálculo elemental y avanzado en el Departamento de Matemáticas.
  \item Profesor temporal, Universidad de Costa Rica, 2016--agosto de 2018. Álgebra Lineal, Introducción a la Matemática para Ciencias de la Computación (lógica, teoría de conjuntos, funciones, álgebras de Boole), Cálculo para Ciencias Económicas II y Cálculo I.
  \item Asistente de docencia, Universidad de Costa Rica, 2014--2016. Principios de Matemática (redacción de demostraciones, lógica, conjuntos, funciones y relaciones) y Álgebra Lineal I y II para el Bachillerato en Matemática.
\end{cvlist}
\cvsec{Áreas de especialidad e investigación}
El CV no declara líneas de investigación; según sus publicaciones, teoría de números: invariantes de Iwasawa (módulo de Iwasawa e invariante lambda) de la Z2-extensión ciclotómica de cuerpos cuadráticos reales.
\cvsec{Producción académica}
Un artículo publicado y un preprint aceptado para publicación (2 trabajos, 2025--2026), este último en coautoría con Foivos Chnaras.
\begin{cvlist}
  \item ``On the Iwasawa lambda-invariant of the cyclotomic Z2-extension of a family of real quadratic fields in which 2 splits'', preprint aceptado en International Journal of Number Theory, arXiv:2605.09111 [math.NT], 2026.
  \item ``Iwasawa module of the cyclotomic Z2-extension of certain real quadratic fields'', The Ramanujan Journal, vol. 67, núm. 1, artículo 6, 2025.
\end{cvlist}
\cvsec{Docencia y formación de recursos humanos}
\begin{cvlist}
  \item Docencia universitaria desde 2014 en la Universidad de Costa Rica y en University of Maryland: cálculo, álgebra lineal, estadística elemental, criptografía, matemática para computación y fundamentos de demostración. El CV no consigna tesis dirigidas.
\end{cvlist}
\cvsec{Distinciones, becas e idiomas}
\begin{cvlist}
  \item Universidad de Costa Rica: mejor promedio de la Facultad de Ciencias y mejor promedio del Bachillerato en Matemática en 2015 (9.85/10) y en 2016 (9.69/10).
  \item University of Maryland: Seymour Goldberg Memorial Award (2021) y Hauptman Summer Research Award (2023).
  \item Español nativo; inglés avanzado.
\end{cvlist}
\end{cvblock}

\begin{cvblock}{Barboza Chinchilla, Luis Alberto}{Profesor Catedrático, Escuela de Matemática y CIMPA, Universidad de Costa Rica}
\cvsec{Formación académica}
\begin{cvlist}
  \item PhD Statistics, Purdue University, Estados Unidos, diciembre de 2012. Tesis: ``New Methods of Estimation of Long-Memory Models with an Application in Climatology''; asesores: Frederi G. Viens y Bo Li.
  \item MS Statistics with Specialization in Computational Finance, Purdue University, agosto de 2011.
  \item MSc Applied Mathematics, Universidad de Costa Rica, julio de 2007. Tesis: ``Interest Rate Swaps. Losses Estimation Model''; asesor: Jose A. Ramírez.
  \item Bachelor degree, Actuarial Science, Universidad de Costa Rica, diciembre de 2003.
\end{cvlist}
\cvsec{Experiencia profesional y docente}
\begin{cvlist}
  \item Profesor Catedrático, Escuela de Matemática, Universidad de Costa Rica, enero de 2013--actualidad. Docente de cursos de posgrado en el programa de Maestría Académica en Estadística, UCR, enero de 2014--actualidad.
  \item Director, Centro de Investigación en Matemática Pura y Aplicada (CIMPA), Universidad de Costa Rica, 2021--2024.
  \item Teaching Assistant, Department of Statistics, Purdue University, agosto de 2007--mayo de 2012.
  \item Profesor Instructor, Escuela de Matemática, Universidad de Costa Rica, agosto de 2003--junio de 2007.
\end{cvlist}
\cvsec{Áreas de especialidad e investigación}
Estadística ambiental; estimación bayesiana de modelos espacio-temporales y de sistemas dinámicos, con aplicaciones a modelos epidemiológicos; técnicas de aprendizaje estadístico aplicadas a modelado epidemiológico y climático. Herramientas: R, MATLAB, Python, C++ y Visual Basic.
\cvsec{Producción académica}
13 publicaciones seleccionadas (2014--2026), todas artículos de revista. Más recientes primero:
\begin{cvlist}
  \item ``Amortized Bayesian inference for spatio-temporal extremes: A copula factor model with autoregression'', Spatial Statistics, 76, 2026.
  \item ``Hurricane Impact Index: Capturing direct and indirect cyclone effects in Central America'', Earth Systems and Environment, 2026.
  \item ``Bayesian spatio-temporal model with INLA for dengue fever risk prediction in Costa Rica'', Environmental and Ecological Statistics, 2023.
  \item ``Spatio-temporal downscaling emulator for regional climate models'', Environmetrics, 34(7), 2023.
  \item ``Bayesian spatial modeling of childhood overweight and obesity prevalence in Costa Rica'', BMC Public Health, 23(1), 2023.
  \item ``Assessing dengue fever risk in Costa Rica by using climate variables and machine learning techniques'', PLoS Neglected Tropical Diseases, 17(1), 2023.
  \item ``Projecting the impact of Covid-19 variants and vaccination strategies in disease transmission using a multilayer network model in Costa Rica'', Scientific Reports, 12, 2022.
  \item ``The role of mobility and sanitary measures on the delay of community transmission of COVID-19 in Costa Rica'', Epidemiologia, 2(3), 2021.
  \item ``Efficient reconstructions of Common Era climate via integrated nested Laplace approximations'', Journal of Agricultural, Biological and Environmental Statistics, 24(3), 2019.
  \item ``Consistent multi-decadal variability in global temperature reconstructions and simulations over the Common Era'', Nature Geoscience, 12, 2019.
  \item ``Parameter estimates of the 2016--2017 Zika outbreak in Costa Rica: An Approximate Bayesian Computation approach'', Mathematical Biosciences and Engineering, 16(4), 2019.
  \item ``Parameter estimation of Gaussian stationary processes using the generalized method of moments'', Electronic Journal of Statistics, 11(1), 2017.
  \item ``Reconstructing past climate from natural proxies and estimated climate forcings using long memory models'', Annals of Applied Statistics, 8(4), 2014.
\end{cvlist}
\end{cvblock}

\begin{cvblock}{Barquero Sánchez, Adrián}{Profesor Catedrático, Escuela de Matemática, Universidad de Costa Rica}
\cvsec{Formación académica}
\begin{cvlist}
  \item Equiparación de grado: Doctorado Académico, Universidad de Costa Rica, Costa Rica, diciembre de 2018.
  \item Ph.D. in Mathematics, Texas A\&M University, Estados Unidos, agosto de 2016.
  \item Bachillerato en Matemática, Universidad de Costa Rica, Costa Rica, agosto de 2011.
\end{cvlist}
\cvsec{Experiencia profesional y docente}
\begin{cvlist}
  \item Escuela de Matemática, Universidad de Costa Rica: Profesor Interino (marzo--diciembre de 2010), Invitado (agosto de 2016--junio de 2019), Asociado (junio de 2019--febrero de 2025) y Catedrático (febrero de 2025--actualidad).
  \item Subdirector de la Escuela de Matemática, Universidad de Costa Rica, noviembre de 2019--noviembre de 2023.
  \item Graduate Teaching Assistant, Department of Mathematics, Texas A\&M University, enero de 2011--agosto de 2016.
  \item Coordinador del curso de Cálculo I (MA-1001/MA-1101), III ciclo de 2018 y I ciclo de 2019.
\end{cvlist}
\cvsec{Áreas de especialidad e investigación}
Teoría de números: fórmula de Chowla--Selberg, alturas de Faltings y conjetura de Colmez, cuerpos CM y curvas elípticas con multiplicación compleja (densidad y distribución), funciones de partición y de sobrepartición, series theta, grupos de Galois y productos entrelazados. Computación: Python, SageMath y \LaTeX{} (nivel intermedio).
\cvsec{Producción académica}
14 trabajos listados: 10 publicados, 3 en revisión (con arXiv) y 1 preprint. Más recientes primero:
\begin{cvlist}
  \item En revisión (arXiv): cálculo eficiente de coeficientes de Fourier de eta-cocientes; conteo de curvas elípticas sobre Q con altura ingenua acotada; distribución de ideales libres de potencias k en grupos de clases de rayos. Preprint: evaluaciones explícitas de alturas de Faltings de variedades abelianas CM.
  \item The density and distribution of CM elliptic curves over Q, Journal of Mathematical Analysis and Applications, 2025.
  \item On the embedding of Galois groups into wreath products, Communications in Algebra, 2024.
  \item Efficient computation of the overpartition function and applications, Journal of Mathematical Analysis and Applications, 2023; The distribution of G-Weyl CM fields and the Colmez conjecture, Research in the Mathematical Sciences, 2023.
  \item Theta series and number fields: theorems and experiments, The Ramanujan Journal, 2021; Stark units and special Gamma values, Research in Number Theory, 2021.
  \item On the Colmez conjecture for non-abelian CM fields, Research in the Mathematical Sciences, 2018; Faltings heights of CM elliptic curves and special Gamma values, Research in Number Theory, 2017.
  \item The Chowla--Selberg formula for abelian CM fields and Faltings heights, Compositio Mathematica, 2016; The asymptotic distribution of Andrews' smallest parts function, Archiv der Mathematik, 2015.
\end{cvlist}
\cvsec{Docencia y formación de recursos humanos}
\begin{cvlist}
  \item 36 cursos-ciclo registrados (29 en la Universidad de Costa Rica y 7 en Texas A\&M University): Cálculo I a III, Ecuaciones Diferenciales, Álgebra Lineal, Grupos y Anillos, Teoría de Galois, Variable Compleja, Teoría Analítica de Números, Teoría de Números Algebraicos, Teoría de Números y Tópicos en Teoría de Números: Curvas Elípticas (2020 y 2024).
  \item Mentor de la REU (Research Experience for Undergraduates) en Teoría de Números, Texas A\&M University, veranos de 2014, 2015 y 2016, con ocho estudiantes de bachillerato guiados en su primera experiencia de investigación.
  \item 23 charlas impartidas (coloquios, seminarios y congresos), entre ellas SIMMAC XXIV (2024), Bayou Arithmetic Research Days 3 (2023) y un coloquio en la Escuela de Matemática (octubre de 2026).
\end{cvlist}
\cvsec{Distinciones, actividades académicas e idiomas}
\begin{cvlist}
  \item Outstanding Teaching Assistant Award, Texas A\&M University, mayo de 2016.
  \item Árbitro de Proceedings of the American Mathematical Society, The Ramanujan Journal, Research in Number Theory, Acta Arithmetica, Mathematical Reviews, Houston Journal of Mathematics y Revista de Matemática: Teoría y Aplicaciones.
  \item Pasantía de investigación como Distinguished Visiting Professor, Bucknell University, abril de 2026; instructor de los tutoriales de SageMath en la escuela CIMPA (San José, julio--agosto de 2026); participación en escuelas y talleres en Estados Unidos, México, Argentina y Chile.
  \item Idiomas: español (lengua materna), inglés (avanzado), alemán (principiante--intermedio), francés y japonés (principiante).
\end{cvlist}
\end{cvblock}

\begin{cvblock}{Bermúdez Aguilar, Esteban}{Director del Departamento de Ciencias Actuariales, Universidad de Costa Rica}
\cvsec{Formación académica}
\begin{cvlist}
  \item Máster en Administración de Negocios con énfasis en Gerencia, Universidad Hispanoamericana, Costa Rica, 2012.
  \item Bachillerato en Ciencias Actuariales, Universidad de Costa Rica, Costa Rica, 2005.
  \item Formación complementaria: Stochastic Modeling of Mortality with Special Focus on Pandemic Risk (European Actuarial Academy, 2021); Estructura y riesgo de los sistemas de pensiones (Universidad de Costa Rica, 2022); Big Data y Data Science para la Administración Pública (Universidad de Costa Rica, 2018); Humanidades digitales (Universitat Autònoma de Barcelona/Coursera, 2023); Competencias Digitales para la Docencia (Universidad de Costa Rica, 2025); Gestión Lean y Ágil de Proyecto (Project Management Institute, 2020).
\end{cvlist}
\cvsec{Experiencia profesional y docente}
\begin{cvlist}
  \item Director del Departamento de Ciencias Actuariales, Universidad de Costa Rica, actualidad (el CV no indica la fecha de inicio). Gestión académica y participación en la formación de profesionales y en el desarrollo de la disciplina actuarial.
  \item Docente, Universidad de Costa Rica, 2005--2010 y 2022--actualidad. Cursos: Regímenes de Pensiones, Estadística Actuarial II, Modelos Lineales y de Sobrevivencia y el curso introductorio de la Maestría en Estadística (Probabilidad y Estadística), entre otros.
  \item Jefe del Departamento Actuarial, Junta de Pensiones y Jubilaciones del Magisterio Nacional (JUPEMA), 2014--marzo de 2023. Responsable de los estudios actuariales de dos regímenes de pensiones y de un equipo de 10 personas.
  \item Encargado de Servicios Actuariales, Ernst \& Young Centroamérica, 2011--2014. Estudios actuariales de beneficios a empleados, estimación y revisión de reservas de aseguradoras, valoración y deterioro de activos intangibles.
  \item Asesor actuarial, varias empresas, 2008--actualidad. Pasivos contingentes de funerarias, notas técnicas de aseguradoras y estudios actuariales de fondos de pensiones.
  \item Docente, Universidad Nacional (Costa Rica), 2009, 2010 y 2017: Investigación de Operaciones y Estructura Discreta.
\end{cvlist}
\cvsec{Áreas de especialidad e investigación}
Ciencias actuariales, pensiones y seguridad social, estadística y gestión de riesgos; análisis, modelación y valuación de sistemas de pensiones, con énfasis en sostenibilidad financiera y actuarial, suficiencia de beneficios, reformas y equidad intergeneracional.
\cvsec{Producción académica}
Dos artículos en revistas, más recientes primero:
\begin{cvlist}
  \item Metodología para establecer una escala salarial dinámica para el régimen de capitalización colectiva de los educadores de Costa Rica, Revista de Matemática, 25(2), 2018.
  \item Comparison of Nelson-Siegel and Svensson models in the optimization of yield curves in the Costa-Rican market, Revista Ciencia y Tecnología, 30(2), 2014.
\end{cvlist}
\cvsec{Cursos y seminarios impartidos}
\begin{cvlist}
  \item CISS-CIESS: Financiamiento en la seguridad social (2026); Lenguaje Actuarial en la Seguridad Social, con Regímenes Financieros en la Seguridad Social (2024); Tablas de mortalidad: construcción y uso en los seguros sociales (CIESS, 2019).
  \item Matemática Actuarial, Fundación de Investigación de la Universidad de Costa Rica (FUNDEVI), 2017; Cadenas de Markov y Simulación de Montecarlo, Tecnológico de Costa Rica, 2015; Técnicas para la determinación del Valor Razonable, Colegio de Contadores Públicos, 2013.
\end{cvlist}
\cvsec{Distinciones, arbitraje y afiliaciones}
\begin{cvlist}
  \item Revisor científico de Applied Stochastic Models in Business and Industry (Wiley), 2024; revisor por pares de tres trabajos del 1.er Congreso de Estudios Interamericanos de Seguridad Social (CISS--CIESS--RIUSS), 2026.
  \item Miembro del Colegio de Ciencias Económicas de Costa Rica desde 2006 y de la sección de Pensiones y Seguridad Social (PBSS) de la Asociación Internacional de Actuarios desde 2012.
\end{cvlist}
\end{cvblock}

\begin{cvblock}{Bolaños Guerrero, Raúl Esteban}{Profesor, Escuela de Matemática, Universidad de Costa Rica}
\cvsec{Formación académica}
\begin{cvlist}
  \item Doctor of Philosophy in Mathematics, University of Kansas, Estados Unidos de América, 16 de mayo de 2021.
  \item Master of Arts in Mathematics, University of Kansas, Estados Unidos de América, 14 de mayo de 2017.
  \item Bachillerato en Ciencias Actuariales, Universidad de Costa Rica, Costa Rica, 22 de abril de 2024.
  \item Bachillerato en Matemática Pura, Universidad de Costa Rica, Costa Rica, 22 de octubre de 2010.
  \item Cursos: Didáctica Universitaria (2022); Escuela de Verano ``Razonamiento Matemático y Formación de Profesores'', Universidad de Costa Rica, 25 horas (2022); XXIII Brazilian School of Probability, Universidade de São Paulo (2019); Competencias Digitales para la Docencia, 145 horas aprobadas, impartidas por Metics (septiembre de 2026).
\end{cvlist}
\cvsec{Experiencia profesional y docente}
\begin{cvlist}
  \item Profesor, Escuela de Matemática, Universidad de Costa Rica, desde el 16 de agosto de 2021 hasta la actualidad. Imparte cursos de Matemática. Nombramientos: invitado I (2021--2022) e invitado II (2022--2023), ambos como exbecario doctor; exbecario doctorado (2023--2026); instructor (desde agosto de 2026).
  \item Graduate Teaching Assistant, University of Kansas, agosto de 2015--diciembre de 2020. Cursos de Cálculo I, II y III.
  \item Profesor, Colegio Científico de San Pedro, febrero de 2013--julio de 2015. Matemática de décimo y undécimo año; salió al iniciar estudios de posgrado en el extranjero.
  \item Profesor interino, Universidad de Costa Rica, de agosto a diciembre de 2007 y de marzo de 2009 a julio de 2015, con nombramientos interinos sin título y luego como bachiller, en la Escuela de Matemática y en las sedes regionales del Pacífico (2007, 2011) y del Atlántico (2012).
\end{cvlist}
\cvsec{Áreas de especialidad e investigación}
Probabilidad y procesos estocásticos: integral de Skorohod, cálculo de Malliavin, teoremas límite y convergencia en distribución, ecuación de onda estocástica en dimensión 2, método de Stein--Malliavin y ecuación de calor estocástica.
\cvsec{Producción académica}
Dos artículos en revistas internacionales de matemática, ambos con David Nualart como coautor:
\begin{cvlist}
  \item Averaging 2D stochastic wave equation, Electronic Journal of Probability, vol. 26, 2021 (con Guangqu Zheng).
  \item Limit theorems for singular Skorohod integrals, Theory of Probability and Mathematical Statistics, vol. 102 (2020/2021), American Mathematical Society (con Denis Bell).
\end{cvlist}
\cvsec{Proyectos}
\begin{cvlist}
  \item C1257, Fluctuaciones gaussianas para la ecuación de calor estocástica con condición inicial aleatoria: investigador principal; estado terminado; del 11 al 31 de diciembre de 2023.
\end{cvlist}
\cvsec{Distinciones, becas e idiomas}
\begin{cvlist}
  \item Finalista, Florence Black Award for Excellence in Teaching, University of Kansas, 2021 (buen desempeño docente).
  \item Joan Kirkham Scholarship (2020), NSF Graduate Research Assistant Scholars (2019) y U.G. Mitchell Honor Scholarship (2018), University of Kansas, por excelente rendimiento académico.
  \item Incorporado al Colegio de Profesionales en Ciencias Económicas de Costa Rica, 12 de junio de 2026.
  \item Idiomas: inglés (conversación, lectura y escritura, 100\%); portugués (25\%, 50\% y 25\%, respectivamente).
\end{cvlist}
\end{cvblock}

\begin{cvblock}{Bustamante Medina, Ronald}{Profesor catedrático, Escuela de Matemática, e investigador del CIMPA, Universidad de Costa Rica}
\cvsec{Formación académica}
\begin{cvlist}
  \item Ph.D. en Lógica Matemática y Fundamentos de la Informática, Paris Diderot University, Francia, septiembre de 2001--octubre de 2005. Tesis: ``Théorie des modèles des corps différentiellement clos avec un automorphisme générique''; directora: Zoé Chatzidakis.
  \item M.Sc. en Lógica y Fundamentos de la Informática, Paris Diderot University, Francia, septiembre de 2001--octubre de 2005 (mismas fechas que el doctorado en el CV).
  \item Bachillerato en Matemática Pura, Universidad de Costa Rica, Costa Rica, julio de 1995--septiembre de 1999.
\end{cvlist}
\cvsec{Experiencia profesional y docente}
\begin{cvlist}
  \item Profesor catedrático, Escuela de Matemática, Universidad de Costa Rica, agosto de 1999--presente. Imparte cursos de matemática en grado y posgrado.
  \item Investigador, Centro de Investigación en Matemática Pura y Aplicada (CIMPA), Universidad de Costa Rica, febrero de 2005--presente. Varios proyectos de investigación en teoría de modelos de cuerpos.
  \item Profesor a tiempo parcial, Facultad de Filosofía, Universidad Nacional (Costa Rica), mayo--agosto de 2022. Curso de lógica en la Maestría en Bioética.
\end{cvlist}
\cvsec{Áreas de especialidad e investigación}
Lógica matemática, en particular la teoría de modelos de cuerpos y su interacción con la geometría algebraica y la teoría de números; actualmente, grupos definibles en teorías con operadores. En años recientes, Machine Learning aplicado a la clasificación de audio, el procesamiento de lenguaje natural y la ética, y a la composición algorítmica y la generación de sonido.
\cvsec{Producción académica}
Seis trabajos: cinco artículos y un preprint. Más recientes primero:
\begin{cvlist}
  \item Groups Definable in Partial Differential Fields with an Automorphism, Israel Journal of Mathematics, 2026 (con Zoé Chatzidakis y Samaria Montenegro).
  \item Zilber's Dichotomy for Differentially Closed Fields with an Automorphism, preprint, 2024.
  \item Definable Groups in DCFA, Revista de Matemática: Teoría y Aplicaciones, 26(2), 2019.
  \item Rank and Dimension in Difference-Differential Fields, Notre Dame Journal of Formal Logic, 52(4), 2011.
  \item Jet spaces in difference-differential fields, Revista de Matemática: Teoría y Aplicaciones, 17(1), 2010.
  \item Differentially Closed Fields with a Generic Automorphism, Revista de Matemática: Teoría y Aplicaciones, 14(1), 2007.
\end{cvlist}
\cvsec{Charlas, producción artística y becas}
\begin{cvlist}
  \item 16 charlas y seminarios registrados (2003--2021, además de una sin año), en Paris Diderot, Paris 6, Paris 7, Lyon 1, Barcelona, Urbana-Champaign, Florida, el CIMPA y SIMMAC XXII (2020, con un curso corto sobre análisis de audio para clasificación y agrupamiento).
  \item Compositor de música electrónica, con proyectos individuales, colaborativos y transdisciplinarios: 13 obras individuales entre 2012 y 2019 (EP y álbumes bajo los seudónimos Antisentido, Atsub y Visión Armónica del Caos, y piezas electroacústicas) y 12 colaboraciones (bandas sonoras de videomapping, cortometraje y videoarte, y un EP y un álbum con Saturno Devorando), de 2012 a 2025.
  \item Conciertos y performances en festivales, entre ellos Wassermusik Festival (Berlín, 2016), Festival Nrmal (México, 2015), Sonograna (San José, 2025) y festivales de música electroacústica de la Universidad de Costa Rica.
  \item Complemento de beca de la Universidad de Costa Rica y del Consejo Nacional de Ciencia y Tecnología para estudios en la Université Paris Diderot, 2000--2005.
\end{cvlist}
\end{cvblock}

\begin{cvblock}{Calvo Alpízar, Juan Gabriel}{Full professor, CIMPA y Escuela de Matemática, Universidad de Costa Rica}
\cvsec{Formación académica}
\begin{cvlist}
  \item Mathematics, Ph.D., Courant Institute, New York University, Estados Unidos, mayo de 2015. Tesis: Domain Decomposition Methods for Problems in H(curl).
  \item Mathematics, B.S., Universidad de Costa Rica, Costa Rica, febrero de 2010 (con honores).
  \item Electrical Engineering, B.S., Universidad de Costa Rica, Costa Rica, octubre de 2021.
  \item Business Administration, Diploma, Universidad Estatal a Distancia, Costa Rica, diciembre de 2020.
\end{cvlist}
\cvsec{Experiencia profesional y docente}
\begin{cvlist}
  \item Universidad de Costa Rica: Full professor (diciembre de 2018--presente), antes Invited Professor (agosto de 2015--diciembre de 2018) e Interim Professor (marzo de 2005--julio de 2010). Director del Departamento de Matemática y Ciencias Actuariales (enero de 2021--diciembre de 2024); subdirector del Posgrado en Matemática (abril de 2019--mayo de 2024); miembro del Consejo Editorial de la Revista de Matemática: Teoría y Aplicaciones (2018--2020).
  \item Director de la División de Gestión y Analítica de Datos, Banco Nacional de Costa Rica, agosto de 2019--julio de 2020.
  \item Postdoc, Institute of Mathematics, Czech Academy of Sciences, Praga, octubre de 2016--agosto de 2017; investigador invitado, University of Concepción, Chile, julio--agosto de 2016.
\end{cvlist}
\cvsec{Áreas de especialidad e investigación}
Análisis numérico: métodos de descomposición de dominio (BDDC, Schwarz con traslape) para problemas en H(curl) y H(grad), con subdominios irregulares y coeficientes de alto contraste; elementos virtuales; espacios gruesos y aprendizaje automático para precondicionadores. Modelación matemática de COVID-19, dengue y otras enfermedades infecciosas.
\cvsec{Producción académica}
31 entradas numeradas: 30 artículos (cuatro de ellos en actas de congresos de descomposición de dominio) y un preprint, más la tesis doctoral. Selección, más recientes primero:
\begin{cvlist}
  \item Stochastic dynamics in epidemiological models with interconnected populations, Journal of Biological Systems, 2026; Coarse spaces for virtual element methods on irregular 3D subdomain decompositions, actas de descomposición de dominio (en prensa).
  \item A Schwarz preconditioner for the smoothed lowest-order virtual element method, Computers and Mathematics with Applications, 2025; Robust domain decomposition methods for high-contrast multiscale problems, Journal of Computational Physics, 2024.
  \item Assessing dengue fever risk in Costa Rica by using climate variables and machine learning techniques, PLOS Neglected Tropical Diseases, 2023; Projecting the impact of Covid-19 variants and vaccination strategies, Scientific Reports, 2022.
  \item A BDDC algorithm with deluxe scaling for H(curl) for two dimensions with irregular subdomains, Mathematics of Computation, 2016.
\end{cvlist}
\cvsec{Proyectos y acción social}
\begin{cvlist}
  \item Investigador principal: C5054, Iterative substructuring methods for elliptic PDEs (desde 2025); B8290, Machine learning methods for the design of preconditioners (2021--2024); B8290, Preconditioners for the numerical solution of PDEs (2018--2021); B5A28, Numerical Methods for PDEs (2015--2018).
  \item Coordinador (2019--2023) y colaborador (desde 2024) del proyecto ED3493, Olimpiadas Universitarias de Matemática.
\end{cvlist}
\cvsec{Docencia y formación de recursos humanos}
\begin{cvlist}
  \item Tesis de maestría: Moisés Solano Córdoba (noviembre de 2024), precondicionador aditivo para métodos de Galerkin discontinuo; Fabián Mora Cordero (julio de 2019), plataforma computacional para EDP con FEM y BDDC; una más en curso (Josué Padilla Torres).
  \item Miembro de comités de 13 tesis de posgrado (11 de maestría y 2 de doctorado, cuatro en curso).
\end{cvlist}
\cvsec{Distinciones, becas e idiomas}
\begin{cvlist}
  \item Charla en el International Congress of Mathematicians 2026 (Filadelfia); conferencista plenario en SIMMAC XXI (2018); 17 charlas de congreso en total.
  \item MacCraken Fellowship (2010--2015) y beca de doctorado de la Universidad de Costa Rica (2010--2015); mejor estudiante de Ingeniería Eléctrica (2003).
  \item Idiomas: español nativo, inglés fluido.
\end{cvlist}
\end{cvblock}

\begin{cvblock}{Cambronero Villalobos, Santiago}{Profesor catedrático y subdirector (desde noviembre 2025), Escuela de Matemática, Universidad de Costa Rica}
\cvsec{Formación académica}
\begin{cvlist}
  \item Doctorado (Ph.D.) en Matemática, Courant Institute, New York University, Estados Unidos, 1996.
  \item Master de Ciencias (M.Sc.) en Matemática, Courant Institute, New York University, Estados Unidos, 1994.
  \item Licenciado en Matemática (1990) y Bachiller en Matemática (1988), Universidad de Costa Rica.
\end{cvlist}
\cvsec{Experiencia profesional y docente}
\begin{cvlist}
  \item Profesor catedrático, Escuela de Matemática, UCR, desde enero 2005 (medio tiempo 2011--2016; tiempo completo desde agosto 2016). Antes: profesor interino (1985--1992), invitado (febrero--noviembre 1996), instructor en propiedad (1996--1997) y asociado (1997--2004).
  \item Asistente y profesor, Courant Institute of Mathematical Sciences, New York University, 1992--1996.
  \item Banco Nacional de Costa Rica: analista de riesgo en la Dirección de Modelación Matemática (2007--2009), director a.i. de Modelación Matemática (agosto 2008--enero 2009), director de la Dirección de Modelos Matemáticos (enero 2011--julio 2016) y, desde agosto 2020, validador de modelos de riesgo.
\end{cvlist}
\cvsec{Gestión académica en la UCR}
\begin{cvlist}
  \item Director de la Escuela de Matemática (2003--2006); subdirector (2000--2002, 2017--2019 y desde noviembre 2025); vicedecano de la Facultad de Ciencias (julio--diciembre 2010).
  \item Director del Departamento de Enseñanza de la Matemática (1999--2002), de Matemática Aplicada (2018) y de Matemática Pura (2019--2020).
  \item Miembro de la Comisión de Posgrado (1996--mayo 2023; director a.i. en 1998 y 2009), de la Comisión de Credenciales y Reconocimientos (1996--2002), de la Comisión Compartida de Enseñanza de la Matemática (1999--2002) y de la Comisión Asesora (desde 1999).
  \item Coordinador de las cátedras de Cálculo I, Cálculo II y Ecuaciones Diferenciales (desde 2022).
\end{cvlist}
\cvsec{Áreas de especialidad e investigación}
Ecuaciones diferenciales, análisis, probabilidad y procesos estocásticos: ecuación de Hill con potencial de ruido blanco, operadores de Schrödinger, problemas de parada óptima, integración estocástica y EDP estocásticas; modelación matemática para la gestión de riesgo financiero.
\cvsec{Producción académica}
19 trabajos publicados (1991--2024; dos preprints de 2023 y un artículo de 2024), 13 obras didácticas y 10 modelos de riesgo implementados en el Banco Nacional.
\begin{cvlist}
  \item Cylindrical martingale-valued measures, stochastic integration and SPDE's, Stochastic Partial Differential Equations: Analysis and Computations, 2024; preprints de 2023 sobre el supremo de una familia de funciones de conjunto y la fórmula de Itô para procesos de Itô.
  \item The Riccati map in random Schrödinger and random matrix theory, MSRI Publications 55, 2007.
  \item On the shape of the density for the ground state eigenvalue of Hill's equation, Communications on Pure and Applied Mathematics 59(7), 2006.
  \item On the values of optimal stopping problems for generalized averages of discrete random variables, Brazilian Journal of Probability and Statistics 20, 2006.
  \item The ground state eigenvalue of Hill's equation with white noise potential, Communications on Pure and Applied Mathematics 52, 1999 (con H. P. McKean).
  \item Otros artículos en la Revista de Matemática: Teoría y Aplicaciones (1995--2007).
\end{cvlist}
\cvsec{Obras didácticas y profesionales}
\begin{cvlist}
  \item Apuntes de curso (1996--2025): Ecuaciones Diferenciales, Principios de Análisis I y II, Análisis I, Probabilidad, Introducción a la probabilidad y Cálculo en una variable I y II (última versión noviembre 2025); textos sobre construcción de conjuntos numéricos y análisis real (con A. Duarte).
  \item Implementación y revisión de modelos de score, rating, riesgo de crédito, mercado, liquidez, operativo y legal, y de asignación de capital, en el Banco Nacional de Costa Rica.
\end{cvlist}
\cvsec{Participación en eventos y capacitaciones}
Ponente en el SIMMAC, en los encuentros centroamericanos de investigadores en matemática, en el V Encuentro Franco-Centroamericano (1996) y en las Primeras Jornadas de Análisis Estocástico (UTEM, Chile, 2000); asistente a congresos de riesgo en Nueva York (2008, 2011, 2012) y al SIMMAC 2024; programa de especialista en derivados financieros (Tecnológico de Monterrey, 2007).
\end{cvblock}

\begin{cvblock}{Campos Fernández, José David}{Profesor e investigador, Escuela de Matemática y CIMPA, Universidad de Costa Rica}
\cvsec{Formación académica}
\begin{cvlist}
  \item Doctor en Matemática (distinción máxima), Pontificia Universidad Católica de Chile, Chile, 2012 (estudios 2008--2012, área de Probabilidad; defensa en diciembre de 2012). Tesis: Random Walks in static and Markovian time-dependent random environment; director: Dr. Alejandro Ramírez.
  \item Doctorado Académico en Ciencias, equiparación de grado y título (certificación), Universidad de Costa Rica, Costa Rica, 2014.
  \item Licenciatura en la Enseñanza de la Matemática (graduación de honor), Universidad de Costa Rica, Costa Rica, 2001. Tesis: Crecimiento de funciones no derivables; director: Dr. Carlos Torre Rodríguez.
\end{cvlist}
\cvsec{Experiencia profesional y docente}
\begin{cvlist}
  \item Profesor e investigador, Escuela de Matemática-CIMPA, Universidad de Costa Rica, 2023--actualidad; antes en la Escuela de Matemática-CIMM, 2013--2022.
  \item Profesor, Sede de Occidente, Universidad de Costa Rica, 1999--2007 (cursos de matemática); asistente en la misma sede, 1995--1998.
  \item Pasantías de investigación: Pontificia Universidad Católica de Chile (agosto--setiembre de 2015 y noviembre--diciembre de 2017, financiadas por el Núcleo Milenio Modelos Estocásticos de Sistemas Complejos y Desordenados) e Instituto Camille Jordan, Lyon, Francia (enero de 2013).
\end{cvlist}
\cvsec{Áreas de especialidad e investigación}
Caminatas aleatorias en medios aleatorios (comportamiento balístico, criterios de elipticidad, medida invariante del proceso ambiente, grandes desviaciones), procesos estocásticos y cálculo estocástico con respecto a procesos de tipo cilíndrico en espacios de Banach (medidas valuadas en martingalas cilíndricas, integración estocástica y EDP estocásticas).
\cvsec{Producción académica}
Ocho trabajos: cinco artículos (2013--2026) y tres preprints en arXiv (2024--2026). Los más recientes primero:
\begin{cvlist}
  \item The martingale representation theorem for cylindrical martingale valued measures, preprint en arXiv, julio de 2026.
  \item Itô's Formula for Itô processes defined with respect to a cylindrical-martingale valued measure, Stochastic Analysis and Applications, 2026.
  \item Markov property and path regularity for the solutions to SPDEs driven by cylindrical-martingale valued measures, preprint en arXiv, junio de 2026.
  \item Cylindrical Martingale-Valued Measures, Stochastic Integration and SPDEs, Stochastics and Partial Differential Equations: Analysis and Computations, 13, 887--955, diciembre de 2024.
  \item Riesz spaces of signed charges on semi-rings, preprint en arXiv, noviembre de 2024.
  \item Asymptotic expansion of the invariant measure for ballistic random walk in the low disorder regime, Annals of Probability, 45(6B), 4675--4699, 2017.
  \item Ellipticity criteria for ballistic behavior of random walks in random environment, Probability Theory and Related Fields, 160(1--2), 189--251, 2014.
  \item Level 1 quenched large deviation principle for random walks in dynamic random environment, Bulletin of the Institute of Mathematics, Academia Sinica (número especial), 8(1), 1--29, 2013.
\end{cvlist}
\cvsec{Docencia y formación de recursos humanos}
\begin{cvlist}
  \item Director de la tesis de maestría de Jennifer Acuña Larios, Universidad de Costa Rica (octubre de 2020): Condiciones de Balisticidad de las Caminatas Aleatorias sobre Medios Aleatorios.
  \item Charlas y simposios: expositor en los simposios SIMMAC XXII a XXV (2020, 2022, 2024 y 2026); charla en la Sede de Occidente (2024); workshop Random Media and Random Structures, IMCA (2020); seminarios en la Universidad de Costa Rica (2015, 2018); congreso latinoamericano de probabilidad y estadística matemática (2012, 2016); Somachi (2012); escuelas y congresos en Brasil y Estados Unidos (2010--2012).
\end{cvlist}
\cvsec{Distinciones, becas e idiomas}
\begin{cvlist}
  \item Premio de excelencia en tesis doctoral 2013, área Ciencias Exactas, Pontificia Universidad Católica de Chile; beca Conicyt de Chile para estudios de doctorado, 2008--2012.
  \item Graduación de honor (2001); reconocimiento de CONARE (1999); mejor promedio de las Sedes Regionales (1998).
  \item Idiomas: inglés, nivel intermedio-avanzado (examen en Lenguas Modernas, 2019). Herramientas: Wolfram Mathematica y Linux (intermedio), LaTeX (avanzado).
\end{cvlist}
\end{cvblock}

\begin{cvblock}{Campos Salas, Daniel}{Docente, Universidad de Costa Rica (oficina en CIMPA-OLCOMA)}
\cvsec{Formación académica}
\begin{cvlist}
  \item Doctorado en Matemática, University of Chicago, EE.UU., agosto de 2019. Tesis: Reconstruction of the magnetic field for a Schrödinger operator in a cylindrical setting; supervisor: Prof. Carlos Kenig.
  \item Maestría en Matemática, University of Chicago, EE.UU., junio de 2015.
  \item Bachillerato en Matemática (graduación de honor), Universidad de Costa Rica, Costa Rica, abril de 2013.
  \item Formación docente: College Fellow en University of Chicago (Basic Complex Variables, otoño de 2014; Math. Methods for Physical Sciences II, invierno de 2015; Basic Numerical Analysis, primavera de 2015) y curso FD0340 Didáctica Universitaria, Universidad de Costa Rica (I ciclo de 2020).
\end{cvlist}
\cvsec{Experiencia profesional y docente}
\begin{cvlist}
  \item Docente, Universidad de Costa Rica, 2013 (Cálculo I) y 2019--2026. Matemática Aplicada: Cálculo I y Cálculo para Ciencias Económicas II. Matemática Pura: Cálculo en Una Variable I y II, Álgebra Lineal I y II, Principios de Geometría y Variable Compleja. Educación Matemática: Álgebra Lineal, Teoría de Números, Geometría y Álgebra Lineal, Principios de Análisis I e Introducción a la Variable Compleja. Ciencias Actuariales: Introducción al Análisis Matemático en Una Variable (I-2026). Posgrado en Matemática: Complementos de Álgebra Lineal (I-2023).
  \item Instructor en el Department of Mathematics, University of Chicago: Calculus I--III (2015--2019) y Linear Algebra (verano de 2018).
  \item Miembro de la Comisión Organizadora de OLCOMA desde 2019: elaboración de problemas, revisión de pruebas, talleres para estudiantes de secundaria, capacitación de profesores y entrenamiento de los equipos nacionales.
\end{cvlist}
\cvsec{Áreas de especialidad e investigación}
Ecuaciones en derivadas parciales y problemas inversos para operadores de Schrödinger magnéticos; estimaciones de Carleman; análogos discretos de problemas de ecuaciones en derivadas parciales.
\cvsec{Producción académica}
Un artículo publicado y un preprint en arXiv, además de siete presentaciones (seis charlas y un póster):
\begin{cvlist}
  \item Reconstruction of the magnetic field for a Schrödinger operator in a cylindrical setting, preprint en arXiv (1908.01386).
  \item On the monotonicity of the correction term in Ramanujan's factorial approximation, The Mathematical Gazette, 97(539), 274--275 (con M. Villarino y J. Carvajal Rojas; el CV no indica el año).
  \item Charlas: Mathematical Congress of the Americas, Miami (2025); Ghent Methusalem Junior Analysis \& PDE Seminar, University of Ghent (2022); Coloquios Matemáticos, Universidad de Panamá (2020); Universidad de El Salvador (2020); SNAP, Northwestern University (póster, 2019); Calderón-Zygmund Seminar, University of Chicago (2019); Summer School on Unique Continuation and Inverse Problems, Kopp, Alemania (2018).
\end{cvlist}
\cvsec{Proyectos, consultorías y acción social}
\begin{cvlist}
  \item Proyecto ED-3493 Olimpiadas Universitarias de Matemática (ECM UCR): instructor de seminarios de resolución de problemas desde 2019 y coordinador general desde 2024 (seminarios de problemas, Club de Lectura y Concurso de Matemática Universitaria).
  \item Collegiate Scholars Program, University of Chicago: instructor de Introduction to Number Theory (junio--julio de 2017) para estudiantes de poblaciones subrepresentadas; instructor voluntario en Math Circles of Chicago (2016--2017).
  \item Servicio a la comunidad matemática: mentorías DRP y REU en University of Chicago (2014--2018); Cibercoloquio Latinoamericano de Matemáticas (2020--2023); sesión especial virtual en el Congreso Matemático de las Américas (julio de 2021); Competencia Iberoamericana Interuniversitaria de Matemáticas, Costa Rica (setiembre--octubre de 2025, 60 estudiantes internacionales); Escuela CIMPA Combinatorial Structures in Geometry, Graphs and Symmetry (julio de 2026, más de 70 participantes internacionales).
  \item Más de 30 problemas propuestos y seleccionados en competencias colegiales y universitarias (desde 2017).
\end{cvlist}
\end{cvblock}

\begin{cvblock}{Chaves Aguilar, Santiago}{Profesor invitado, Universidad de Costa Rica}
\cvsec{Formación académica}
\begin{cvlist}
  \item Ph.D. en Matemática, University of Chicago, 2018--2022. Área de especialización: Álgebra, Geometría Algebraica y Teoría de Representaciones. Tesis: On the Schmid-Vilonen conjecture for Irreducible Verma Modules and Discrete Series Representations.
  \item Msc. en Matemática, University of Chicago, 2016--2017.
  \item B.A. en Matemática, Universidad de Costa Rica, 2012--2015 (promedio 9,11/10).
  \item Cursos de Combinatoria (calificación A) y Probabilidad (calificación A+), Instituto Nacional de Matemática Pura e Aplicada (IMPA), Río de Janeiro, Brasil, 2013.
\end{cvlist}
\cvsec{Experiencia profesional y docente}
\begin{cvlist}
  \item Profesor, Universidad de Costa Rica, 2022--2026: Geometría Diferencial, Álgebra Lineal I y II, Geometría Algebraica, Cálculo en varias variables (para matemáticos y general), Ecuaciones Diferenciales, Variable Compleja, Cálculo 2 para ciencias actuariales, Funciones Derivables y Seminario en Teoría de Representaciones.
  \item Creación y conducción del seminario ``Teoría de Representaciones y Geometría Algebraica'' (2025--2026); participación en la conducción del seminario de álgebra y análisis de datos (2024); minicurso en Geometría Algebraica como instructor ad-honorem (2023); taller de Mecánica Clásica (2023, UCR/CIMPA).
  \item Minicursos: teoría geométrica de representaciones en la II Conference of Algebra and its Applications, Liberia (2025); Topología y Geometría Diferencial en las jornadas matemáticas de la UNED (2024).
  \item Instructor, University of Chicago, 2017--2022 (Cálculo en una y varias variables, MATH 152-153; MATH 19520 en 2019; Probabilidad y Combinatoria para el Collegiate Scholars Program en 2017); asistente de cursos de Ecuaciones Diferenciales Ordinarias, Teoría de Galois y Álgebra lineal abstracta, 2016--2017.
  \item Universidad de Costa Rica, 2015: profesor de Cálculo 2, Principios de Matemática y Cálculo 1 para matemáticos; instructor del Seminario en Topología Diferencial.
\end{cvlist}
\cvsec{Áreas de especialidad e investigación}
Álgebra, geometría algebraica y teoría (geométrica) de representaciones; conjetura de Schmid-Vilonen; factorización de vectores suaves y analíticos de representaciones de grupos de Lie. Proyectos de investigación en la Universidad de Costa Rica: La conjetura de Schmid-Vilonen: un estudio de formas invariantes en G-módulos irreducibles de dimensión infinita (2022--2024) y Avances en la teoría de factorización de vectores suaves y analíticos de representaciones de grupos de Lie, Proyecto Semilla (2024--2026).
\cvsec{Producción académica}
El CV no lista artículos ni preprints; registra un trabajo de REU y 13 charlas (2014--2026):
\begin{cvlist}
  \item A Survey on Representation Theory, trabajo del programa Research Experience for Undergraduates (REU), University of Chicago, 2015.
  \item Charlas recientes: Factorización fuerte para vectores suaves sobre álgebras de convolución, Seminario CIMPA (2026); Autovalores del Laplaciano en la esfera, Universidad de Costa Rica (2025); Sobre el quehacer matemático, charla motivacional para el día de pi (2024); Geometric Representation Theory, simposio conjunto CIMPA y CIMA (2023); Teoría Geométrica de Representaciones, CIMPA (2019).
  \item Otras charlas: Polinomios de Bernstein-Sato (2021); On the construction of the Nearby Cycles Functor (2019); Cohomología de las Fibras Cercanas (2018); Differential Operators on the Projective Space (2017); Sobre Integrales Elípticas y Geometría Algebraica (2016); The Laplacian on the Sphere and Representation Theory, REU (2015); Sobre el Método Probabilístico (2014).
\end{cvlist}
\cvsec{Distinciones, becas e idiomas}
\begin{cvlist}
  \item Lawrence and Josephine Prize for Excellence in Undergraduate Teaching, University of Chicago (2019); premio al mejor promedio de la carrera de Matemática, Universidad de Costa Rica (2014).
  \item Medalla de Oro, Olimpiada Iberoamericana para estudiantes Universitarios (OIMU, 2014); Medalla de Plata, Competencia Iberoamericana Interuniversitaria de Matemática (2013--2014); Medalla de Bronce, Olimpiada Iberoamericana de Matemática (2011) y Olimpiada Mundial de Matemática (IMO, 2010).
  \item Entrenador del equipo de Costa Rica para la IMO (2013--2015); tutor del equipo en la Olimpiada Iberoamericana (2014) y en la IMO (2024); jurado en PAGMO 2023 e Ibero 2024; miembro de la comisión de OLCOMA (2015 y 2022--2024); colaborador de círculos matemáticos del CIMPA (2023 y 2025); comité organizador del Directed Reading Program, University of Chicago (2019--2021).
  \item Idiomas: español (nativo), inglés y portugués.
\end{cvlist}
\end{cvblock}

\begin{cvblock}{Coto Mora, Edward}{Escuela de Matemática y CIMPA, Universidad de Costa Rica}
\cvsec{Formación académica}
\begin{cvlist}
  \item Doctorado en Matemáticas, Université Paris Cité, Francia, 2020--2024. Tesis: Conjeturas de Beilinson para GSp4 × GL2.
  \item Maestría en Matemáticas (ALGANT), doble titulación, 2018--2020: Université Paris-Saclay (M2, Mention Bien) y Universiteit Leiden (M1, Cum laude). Tesis de maestría ALGANT (agosto 2020): Computing finite Galois groups arising from automorphic forms.
  \item Licenciatura en Matemáticas, Universidad de Costa Rica, 2012--2015.
\end{cvlist}
\cvsec{Experiencia docente y otras experiencias}
\begin{cvlist}
  \item Docencia universitaria, Universidad de Costa Rica, 2016--2026. Cursos: SP3123 Criptografía; SP1304 Teoría de Galois; MA1005 Ecuaciones Diferenciales; MA1003 Cálculo III (multivariable); MA1002 Cálculo II; CA0151 Cálculo Diferencial e Integral I; MA0292 Álgebra Lineal para Ciencias de la Computación.
  \item Líder de equipo y miembro del jurado en la Olimpiada Internacional de Matemáticas (Rumanía, 2018) y en la Competencia Iberoamericana de Matemáticas para Estudiantes Universitarios, CIIM (Ecuador, 2017).
\end{cvlist}
\cvsec{Áreas de especialidad e investigación}
Teoría de números, álgebra, combinatoria, geometría algebraica, análisis, probabilidad y álgebra lineal; conjeturas de Beilinson y funciones L de Hasse--Weil, formas automorfas y grupos de Galois. Herramientas: Python, SageMath, NumPy y Pandas. Intereses adicionales: ajedrez, Go y diseño de problemas para competencias matemáticas.
\cvsec{Producción académica}
Un artículo publicado y tres preprints (2026), además de la tesis doctoral y la tesis de maestría.
\begin{cvlist}
  \item A weak form of Beilinson's conjecture for residues of Hasse--Weil L-functions attached to GSp4 × GL2, preprint, 2026.
  \item Classification of cyclic holomorphs with perfect order subsets, preprint, 2026.
  \item 169 and Sums of Positive Squares, Part II, preprint, 2026.
  \item 169 and Sums of Positive Squares, The American Mathematical Monthly 128(6), 2021 (con F. Viganò).
\end{cvlist}
\cvsec{Distinciones, becas e idiomas}
\begin{cvlist}
  \item FMJH Merit Scholarship, Fondation Mathématique Jacques Hadamard (2019--2020); Leiden Excellence Scholarship, Universiteit Leiden (2018--2019).
  \item Segundo premio en la International Mathematics Competition, IMC (Bulgaria, 2014 y 2015); medalla de oro en la CIIM (México, 2015) y en la Olimpiada Iberoamericana de Matemáticas para Universidades, OIMU (Costa Rica, 2015).
  \item Idiomas: español (nativo), inglés (C1) y francés (C1).
\end{cvlist}
\end{cvblock}

\begin{cvblock}{Fonseca Mora, Christian}{Profesor catedrático, Escuela de Matemática, Universidad de Costa Rica}
\cvsec{Formación académica}
\begin{cvlist}
  \item Doctorado en Matemática, Énfasis en Probabilidad y Estadística, The University of Sheffield, Inglaterra, diciembre de 2015. Tesis: Stochastic Analysis with Lévy Noise in the Dual of a Nuclear Space; supervisor: David B. Applebaum.
  \item Bachillerato en Matemática, Universidad de Costa Rica, abril de 2012; Bachillerato en Ciencias Actuariales, Universidad de Costa Rica, abril de 2019.
  \item Licenciatura en Economía (Summa Cum Laude), Universidad de Costa Rica, diciembre de 2010. Tesis: Valoración Económica del Volcán Arenal: Una Medición del Valor de su Belleza Escénica Mediante el Método de Valoración Contingente (2010); supervisor: Ronulfo Jiménez Rodríguez.
  \item Bachillerato en Economía, Universidad de Costa Rica, octubre de 2008.
\end{cvlist}
\cvsec{Experiencia profesional y docente}
\begin{cvlist}
  \item Profesor Catedrático, Universidad de Costa Rica, desde 05/2023; Profesor Asociado, 12/2019--05/2023.
  \item Profesor Adjunto, Universidad de Costa Rica, 08/2019--12/2023 (fechas tal como constan en el CV).
  \item Profesor Invitado, Universidad de Costa Rica, 01/2016--08/2019; Profesor Interino Licenciado, 03/2011--07/2012.
\end{cvlist}
\cvsec{Áreas de especialidad e investigación}
Análisis estocástico en el dual de espacios nucleares: procesos de Lévy y medidas infinitamente divisibles, semimartingalas, integración estocástica, ecuaciones de evolución y EDP estocásticas, convergencia débil y tensión de procesos; procesos y medidas cilíndricas valuadas en martingalas.
\cvsec{Producción académica}
15 artículos publicados (2018--2026), 5 preprints en arXiv y 2 tesis. Los más recientes primero:
\begin{cvlist}
  \item Itô's Formula for Itô processes defined with respect to a cylindrical-martingale valued measure, Stochastic Analysis and Applications, en prensa, 2026.
  \item Weak convergence for continuous stochastic processes in the dual of a nuclear space with applications to the convergence of SPDEs, Random Operators and Stochastic Equations, 33(4), 401--416, 2025.
  \item Cylindrical Martingale-Valued Measures, Stochastic Integration and SPDEs, Stoch PDE: Anal Comp, 13, 887--955, 2025.
  \item Convergence Uniform on Compacts in Probability with Applications to Stochastic Analysis in Duals of Nuclear Spaces, Stochastic Analysis and Applications, 42(4), 797--827, 2024.
  \item Stochastic Evolution Equations With Lévy Noise in Duals of Nuclear Spaces, Stoch PDE: Anal Comp, 12(1), 173--219, 2024.
  \item Almost Sure Uniform Convergence of Stochastic Processes in the Dual of a Nuclear Space, Journal of Theoretical Probability, 36(4), 2564--2589, 2023; Time regularity of stochastic convolutions and stochastic evolution equations in duals of nuclear spaces, Stochastic Analysis and Applications, 41(6), 1136--1154, 2023.
  \item Stochastic Integration With Respect to Cylindrical Semimartingales, Electronic Journal of Probability, 26, 2021; Stochastic Integration in Hilbert Spaces With Respect to Cylindrical Martingale-Valued Measures, ALEA, 18, 1267--1295, 2021.
  \item Semimartingales in the Dual of a Nuclear Space, Electronic Journal of Probability, 25, 2020; Tightness and Weak Convergence of Probabilities on the Skorokhod Space on the Dual of a Nuclear Space and Applications, Studia Mathematica, 254(2), 109--147, 2020; Lévy Processes and Infinitely Divisible Measures in the Dual of a Nuclear Space, Journal of Theoretical Probability, 33(2), 649--691, 2020; Regularization of Cylindrical Processes In Locally Convex Spaces, Infinite Dimensional Analysis, Quantum Probability and Related Topics, 23(4), 2020.
  \item Stochastic Integration and Stochastic PDEs Driven by Jumps on the Dual of a Nuclear Space, Stoch PDE: Anal Comp, 6(4), 618--689, 2018; Existence of Continuous and Càdlàg Versions for Cylindrical Processes in the Dual of a Nuclear Space, Journal of Theoretical Probability, 31(2), 867--894, 2018.
  \item Preprints en arXiv: Itô's formula for Lévy-Itô processes taking values in the dual of nuclear space; Markov property and path regularity for the solution to SPDEs driven by a cylindrical martingale-valued measure; The martingale representation theorem for cylindrical martingale valued measures; Vector-valued stochastic integration with respect to semimartingales in the dual of nuclear space; Riesz spaces of signed charges on semi-rings.
\end{cvlist}
\cvsec{Docencia y formación de recursos humanos}
\begin{cvlist}
  \item Dirección de una tesis de maestría: Anddy Alvarado Solano, Procesos de Lévy Cilíndricos: Regularización e Integración, Universidad de Costa Rica, finalizada en diciembre de 2019.
\end{cvlist}
\end{cvblock}

\begin{cvblock}{Guevara Villalobos, Álvaro}{Risk and Strategic Intelligence Director, Grupo Financiero CAFSA; profesor de Matemática y Ciencias Actuariales, Universidad de Costa Rica}
\cvsec{Formación académica}
\begin{cvlist}
  \item Ph.D. Mathematics, Louisiana State University, EE.\ UU., 2009.
  \item B.Sc. Mathematics, Universidad de Costa Rica, Costa Rica, 2002.
\end{cvlist}
\cvsec{Experiencia profesional y docente}
\begin{cvlist}
  \item Risk and Strategic Intelligence Director, Grupo Financiero CAFSA, San José, agosto de 2026--presente. Responsable de la gestión de riesgos y la inteligencia estratégica del grupo financiero supervisado; dirige a 8 profesionales encargados de identificar, monitorear y mitigar riesgos y de apoyar con análisis de datos las decisiones de la empresa.
  \item Director / Team Manager, Banco Nacional de Costa Rica, 2025--presente: supervisó un equipo analítico de 16 personas, con modelos de predicción de fuga de clientes (reducción de la fuga en 12 \%) y de segmentación (aumento de 15 \% en venta cruzada); representó al banco como expositor en foros locales e internacionales (Cámara de Bancos, Risk Evolve, Risk Live).
  \item Director / Team Manager, Banco Nacional de Costa Rica, 2016--2025: dirigió iniciativas de modelación cuantitativa de riesgo, con modelos de scoring crediticio (mejora de 18 \% en la exactitud de aprobación de préstamos), modelos de VaR y de precios para reportes internos, y cumplimiento de Basilea III, IFRS 9 y normativa local (reducción de 10 \% en el ciclo de validación de modelos).
  \item Professor of Mathematics and Actuarial Science, Universidad de Costa Rica, 2013--presente. Imparte optimización, teoría de juegos, análisis de regresión y riesgo actuarial.
\end{cvlist}
\cvsec{Áreas de especialidad e investigación}
Modelación de riesgo (crédito, mercado y operacional), cumplimiento de IFRS 9, Basilea III y regulación local, aprendizaje automático e inteligencia artificial, pruebas de estrés y optimización de portafolios. Herramientas: Python, R, SQL, MATLAB, PowerBI y Azure.
\cvsec{Producción académica}
Cuatro publicaciones destacadas (2011--2018), sin lista completa en el CV:
\begin{cvlist}
  \item Behavior Research Methods 50(1), 2018: batería de decisión basada en valor, con enfoque bayesiano adaptativo, para evaluar conducta impulsiva y riesgosa.
  \item NeuroImage 102(2), 2014: activación cerebral anticipatoria en el núcleo accumbens y el área tegmental ventral como predictor del esfuerzo.
  \item IEEE Southwest Symposium on Image Analysis and Interpretation (SSIAI), 2012: agrupamiento orientado por curvatura de campos de vectores de movimiento dispersos.
  \item IEEE Computer Society Conference on Computer Vision and Pattern Recognition Workshops (CVPRW), 2011: aprendizaje de correspondencias entre múltiples vistas a partir de coincidencias temporales.
\end{cvlist}
\cvsec{Docencia y formación de recursos humanos}
\begin{cvlist}
  \item Dirección de cuatro tesis de maestría en la Universidad de Costa Rica: calibración y análisis comparativo de modelos de probabilidad de incumplimiento bajo IFRS 9 (2026); score crediticio de pequeñas y medianas empresas para el Banco Nacional (2025); predicción de rotación de empleados en SITEL Costa Rica (2025); sistemas de votación en Costa Rica, con propuestas de matemática aplicada para elecciones legislativas (2025). Proyectos guiados con mejoras de precisión predictiva de hasta 25 \%.
\end{cvlist}
\cvsec{Desarrollo profesional, distinciones e idiomas}
\begin{cvlist}
  \item Expositor en Risk Live London (2025) y en la IFCS Conference (2024); panelista en Risk Evolve London (2025); organizador de una mesa redonda sobre riesgo de modelos (2024). Formación: Risk Americas (2023, 2026), Strategy \& Business Analytics (INCAE, 2025), Balance Sheet Management USA (2023), Enterprise Model Risk Management (2022), Microsoft Copilot Mentoring Program (2026).
  \item Idiomas: español (lengua materna), inglés C2, alemán B1 y francés (lectura).
\end{cvlist}
\end{cvblock}

\begin{cvblock}{Guier Acosta, Jorge}{Escuela de Matemática (Departamento de Matemática Pura), Universidad de Costa Rica}
\cvsec{Formación académica}
\begin{cvlist}
  \item Doctorat de l'Université de Paris 7, Spécialité: Logique et Fondements de l'Informatique, Francia, 5 de julio de 1999. Convalidado como Doctorado Académico en la Universidad de Costa Rica, 26 de julio de 2001.
  \item Diplôme d'Etudes Approfondies de Logique et Fondements de l'Informatique, Université Paris 7 -- Denis Diderot, Francia, septiembre de 1993.
  \item Bachillerato en Matemática, Universidad de Costa Rica, 31 de octubre de 1991.
\end{cvlist}
\cvsec{Cargos académico-administrativos y docencia}
\begin{cvlist}
  \item Subdirector de la Escuela de Matemática, UCR, noviembre 2013--diciembre 2017.
  \item Director del Departamento de Matemática Pura, Escuela de Matemática, UCR, enero 2007--diciembre 2010 y enero 2016--diciembre 2019.
  \item Coordinador de la comisión de docencia (enero 2016--diciembre 2021) y director de la comisión de reconocimientos (enero 2007--diciembre 2014), Escuela de Matemática, UCR.
  \item Profesor invitado, Escuela de Matemática UNAN León--UCR (julio de 2014 y agosto de 2016): minicursos sobre cuerpos ordenados y real cerrados, y sobre lógica proposicional.
  \item Seminario sobre su especialidad con charlas introductorias a la teoría de modelos, Escuela de Matemática UCR (verano de 2000); charlas frecuentes en el Coloquio de la Escuela sobre teoría de modelos de anillos real cerrados.
\end{cvlist}
\cvsec{Áreas de especialidad e investigación}
Lógica matemática y teoría de modelos; anillos real cerrados, f-anillos y eliminación de cuantificadores; productos booleanos de anillos de valuación; estructuras o-minimales y espectro real.
\cvsec{Producción académica}
Siete artículos (1999--2026), todos como autor único; dos de ellos en 2025--2026.
\begin{cvlist}
  \item Projectable reduced f-rings admitting elimination of quantifiers, Séminaires et Congrès (Société Mathématique de France) 35, 2026.
  \item Elimination of quantifiers of a theory of real closed rings, Annals of Pure and Applied Logic 176, 2025.
  \item Convex Lattice-Ordered Subrings of von Neumann Regular f-Rings, Revista Colombiana de Matemáticas 49(1), 2015.
  \item Anillos separables, de Baer y PP-anillos, Revista de Matemática: Teoría y Aplicaciones 13(2), 2006.
  \item Boolean products of real closed valuation rings and fields, Annals of Pure and Applied Logic 112, 2001.
  \item Produits booléens d'anneaux de valuations réels clos, Comptes Rendus de l'Académie des Sciences de Paris 330, 2000.
  \item Ultraproductos de f-anillos proyectables, Revista de Matemática: Teoría y Aplicaciones 6(2), 1999.
\end{cvlist}
\cvsec{Congresos y comités}
\begin{cvlist}
  \item Conferencista invitado, conferencia por el 40 aniversario del seminario Delon--Dickmann--Gondard, Banyuls-sur-Mer, Francia (2025); invitado como joven investigador por la ASL al Logic Colloquium 2000 y la European Logic Summer School 2000.
  \item Ponencias aceptadas en los SLALM XVII (Puebla, 2017), XX (Montevideo, 2024) y XXI (Bogotá, 2026); charlas en el Seminario de Estructuras Algebraicas Ordenadas de Paris 7 y Paris-Cité (1996, 1997, 2020 y 2022); conferencista en el XI SLALM (Mérida, 1998) y en la Summer School on Real Algebraic Geometry (ICTP, Trieste, 2005).
  \item Como estudiante de doctorado: charlas en el Grupo de Trabajo de Estudiantes de Doctorado en Lógica de Paris 7 (álgebras de Boole y secciones continuas de haces de L-estructuras sobre espacios booleanos) y en seminarios de Angers (1998); asistente a RESMOD School (Segovia, 1995), Order in Algebra and Logic VII y Colloque Franco-Qazaq de Théorie des Modèles (CIRM, Luminy, 1997) y Introductory Workshop in Model Theory of Fields (MSRI, Berkeley, 1998).
  \item Director del comité de organización local y miembro del comité científico del XII Simposio Latinoamericano de Lógica Matemática (UCR, 2004).
\end{cvlist}
\cvsec{Idiomas}
Español (lengua materna); francés (dominio avanzado, oral y escrito); inglés (dominio medio, oral y escrito).
\end{cvblock}

\begin{cvblock}{Gutiérrez Pavón, Jonathan}{Profesor, Escuela de Matemática, Universidad de Costa Rica}
\cvsec{Formación académica}
\begin{cvlist}
  \item Doctorado en ciencias en la especialidad de matemáticas, CINVESTAV, México, febrero de 2018. Tesis: Spectral decompositions for diffusions.
  \item Maestría en ciencias en la especialidad de matemáticas, CINVESTAV, México, febrero de 2015. Tesis: One dimensional diffusion processes and orthogonal polynomials.
  \item Licenciatura en la enseñanza de la matemática, Universidad de Costa Rica, diciembre de 2009. Tesis: Uso de metáforas en el discurso del aula de matemática: una investigación realizada en tres grupos de décimo año de tres colegios de la capital.
  \item Bachillerato en la enseñanza de la matemática, Universidad de Costa Rica, agosto de 2004.
\end{cvlist}
\cvsec{Experiencia profesional y docente}
\begin{cvlist}
  \item Profesor, Universidad de Costa Rica, desde 2004 hasta el presente.
  \item Coordinador durante dos años de la cátedra MA1210; coordinador del curso de matemática para economía y estadística 2 (MA313).
  \item Elaboración de exámenes para el proyecto ExMa.
\end{cvlist}
\cvsec{Áreas de especialidad e investigación}
Difusiones (difusión de Brox, difusión de Cox-Ingersoll-Ross generalizada), tiempos locales y de llegada, operadores aleatorios con ruido de semimartingala y procesos estocásticos, según sus tesis y publicaciones.
\cvsec{Producción académica}
8 artículos (2018--2025), todos en coautoría con C.G. Pacheco. Los más recientes primero:
\begin{cvlist}
  \item A formula for the density of local time of the Brox diffusion in a time-window, Random Operators and Stochastic Equations, 2025.
  \item A Schrodinger random operator with semimartingale potential, Random Operators and Stochastic Equations, 2023.
  \item The Killed Brox diffusion, Stochastic Models, 2022; Solving equations with semimartingale noise, Random Operators and Stochastic Equations, 2022.
  \item Quenched distributions for the maximum, minimum and local time of the Brox diffusion, Statistics and Probability Letters, 2021.
  \item Computations of hitting time densities for the generalized Cox-Ingersoll-Ross diffusion, Morfismos, 24, 15--34, 2020; A density for the local time of the Brox diffusion, Statistics and Probability Letters, 2020.
  \item Inverting Weak Random Operators, Random Operators and Stochastic Equations, 2018.
\end{cvlist}
\cvsec{Docencia y formación de recursos humanos}
\begin{cvlist}
  \item Cursos impartidos en la Escuela de Matemática: Cálculo 1, 2 y 3 (MA1001, MA1002, MA1003), Álgebra lineal (MA1004), Cálculo (MA1210), Introducción a la Probabilidad (MA1030), Funciones derivables (MA0012), Funciones en varias variables (MA0032), Principios de análisis 1 y 2 (MA0540, MA0551), Probabilidades 1 (MA0720), Cálculo diferencial e integral en varias variables (CA0351) y Procesos estocásticos (CA0821 y PF1347).
\end{cvlist}
\end{cvblock}

\begin{cvblock}{Hernández Alvarado, Alberto}{Director del CIMPA y profesor catedrático, Escuela de Matemática, Universidad de Costa Rica}
\cvsec{Formación académica}
\begin{cvlist}
  \item Ph.D. en matemática, Universidad de Porto, Portugal, 2016. Disertación: The quotient module, coring depth and factorisation algebras; supervisor: Lars Kadison.
  \item Bachillerato en matemática, Universidad de Costa Rica, 2010.
\end{cvlist}
\cvsec{Experiencia profesional y docente}
\begin{cvlist}
  \item Director del CIMPA y profesor catedrático de la Escuela de Matemática, Universidad de Costa Rica (designaciones actuales); subdirector del CIMPA, 2021--2023.
  \item Académico visitante: Departamento de Matemática de la Universidad de Évora, Portugal (2025) y de la Universidad de Gent, Bélgica (2023).
  \item Profesor e investigador invitado, Universidad de Costa Rica, 2017--2020.
  \item Miembro del grupo de álgebra, combinatoria y teoría de grupos de CMUP, Universidad de Porto, 2011--2016.
  \item Profesor asistente, Escuela de Matemática, Universidad de Costa Rica, 2010--2011.
  \item Coordinador del curso Ecuaciones Diferenciales Aplicadas a las Ciencias de la Salud, Escuela de Matemática, 2020--2022.
\end{cvlist}
\cvsec{Áreas de especialidad e investigación}
Álgebra: profundidad de subgrupos y de extensiones, álgebras de factorización, módulos de álgebras de Hopf y representaciones de grupos de Lie. También modelos epidémicos y geometría de datos (según publicaciones y proyectos); el CV no declara explícitamente sus áreas.
\cvsec{Producción académica}
8 publicaciones (2014--2025): 1 capítulo en libro de Springer y 7 artículos. Más recientes primero:
\begin{cvlist}
  \item UMAP projections and the survival of empty space: a geometric approach to high dimensional data, Springer (2025).
  \item Isomorphisms of relative holomorphs and matrix similarity, Bulletin of the Australian Mathematical Society (2024).
  \item Minimum depth of factorization algebra extensions, Revista Colombiana de Matemáticas 57(1) (2023).
  \item Geometric goodness of fit measure to detect patterns in data point sets, Journal of Computational Statistics (2022).
  \item A two-patch epidemic model with nonlinear relapse, Revista de Matemática: Teoría y Aplicaciones (2020).
  \item A Quantum Subgroup Depth, Acta Mathematica Hungarica 152(1) (2017).
  \item Subgroup Depth and Twisted Coefficients, Communications in Algebra 44(8) (2016); Algebraic Quotient Modules and Subgroup Depth, Abhandlungen aus dem Mathematischen Seminar der Universität Hamburg 84 (2014).
\end{cvlist}
\cvsec{Proyectos, consultorías y acción social}
\begin{cvlist}
  \item Proyectos de investigación: Advances in factorization theory for smooth and analytic vectors of Lie group representations (2024--2026); Indecomposable algebraic modules of a finite dimensional Hopf algebra (2021--2023); Finite dimensional extensions of finite depth (2017--2020).
  \item Coordinador del proyecto ACAVAMA, validación de modelos computacionales para finanzas (2021--2024), y del proyecto SACM, servicios de asesoría y capacitaciones en matemática del CIMPA (2021--2024).
\end{cvlist}
\cvsec{Gestión, servicio y organización de eventos}
\begin{cvlist}
  \item Organizador de la I Conferencia de Algebra y sus Aplicaciones (2023), de la II Conferencia (2025) y de Research Days III, conjunta con el CIMA de la Universidad de Évora (2025--2026).
  \item Miembro del Consejo Editorial de la Revista de Matemática: Teoría y Aplicaciones (2023); del Consejo Científico del CIMPA (2021--2023); del Consejo Asesor del CIMPA (2020--2023); de la Comisión de Posgrado en Matemática de la UCR (2020).
\end{cvlist}
\end{cvblock}

\begin{cvblock}{Jiménez López, David}{Profesor, Escuela de Matemática, y Director del CIMM, Universidad de Costa Rica}
\cvsec{Formación académica}
\begin{cvlist}
  \item Ph.D. en Matemáticas, Instituto Tecnológico de Georgia, EEUU, 2008. Tesis: Analysis of two problems in signal quantization for A/D conversion; asesor: Yang Wang.
  \item B.Sc. en Matemáticas, Universidad de Costa Rica, 2001.
\end{cvlist}
\cvsec{Experiencia profesional y docente}
\begin{cvlist}
  \item Profesor, Escuela de Matemática, Universidad de Costa Rica, desde 03/2013; Director del CIMM desde 12/2018.
  \item Investigador postdoctoral, Universidad de Houston, 01/2011--02/2013.
  \item Visiting Assistant Professor, Universidad de Texas A\&M, 09/2008--01/2011.
  \item Asistente graduado de enseñanza e investigación, Instituto Tecnológico de Georgia, 08/2003--07/2008.
  \item Instructor, Universidad de Costa Rica, 08/2001--07/2003.
\end{cvlist}
\cvsec{Áreas de especialidad e investigación}
Procesamiento de señales e imágenes digitales, imágenes médicas y neuronales, detección compresiva, análisis armónico aplicado, wavelets y teoría de frames, lingüística matemática y computacional, modelado de dinámica social.
\cvsec{Producción académica}
1 capítulo de libro, 7 artículos de investigación (2 marcados como enviados), 4 artículos de conferencia, 9 charlas de conferencia y 9 charlas de seminario (la última actualización del CV es de marzo de 2022). Más recientes primero:
\begin{cvlist}
  \item Numerical methods from bioinformatics applied to the diachronic study of the Chibchan languages (enviado, sin año).
  \item A Directional Representation for 3D Tubular Structures Resulting from Isotropic Well-Localized Atoms, Applied Computational Harmonic Analysis (2016).
  \item Automated Quantitative Analysis of Confocal Images of Neuronal Network Cultures (enviado, 2014); Improved Automatic Centerline Tracing for Dendritic Structures, Neuroinformatics 3(2) (sin año en el CV).
  \item Capítulo A Harmonic Analysis View on Neuroscience Imaging, en Excursions in Harmonic Analysis, Birkhäuser (2012).
  \item Conferencias: Constructive Theory of Functions, Sozopol (2013); SPIE Proc. 8858 (2013); ISBI (2013); SPIE Optics+Photonics (2007).
  \item Beta-alpha Encoders for Robust A/D Conversion, Acta Applicandae Mathematicae (2009); White Noise Hypothesis for Uniform Quantization Errors, SIAM Journal on Mathematical Analysis 38 (2007); Open Problems in Constructive Function Theory, Electronic Transactions on Numerical Analysis 25 (2006).
\end{cvlist}
\cvsec{Docencia y formación de recursos humanos}
\begin{cvlist}
  \item Universidad de Costa Rica: Ecuaciones Diferenciales, Álgebra Lineal, Cálculo I--III, Principios de Matemática, Probabilidad Aplicada, Análisis I y II, Análisis Armónico, Teoría de Números e Historia de la matemática.
  \item Universidad de Texas A\&M: Matemática para Negocios, Geometría Analítica, Matemática para Ingeniería (Cálculo I--III). Instituto Tecnológico de Georgia: Cálculo I--III, Álgebra Lineal, Matemática Discreta, Introducción al Cálculo.
  \item Charlas de conferencia (2005--2011): Edmonton, Canadá (2011); Nashville (2011); Baton Rouge (2011); Atlanta (2010); San Antonio (2010); San Diego (2007); Banff (2006).
  \item Charlas de seminario recientes, entre ellas ``Señales Digitales: Una perspectiva matemática'', III Escuela de Matemática Pura y Aplicada, Universidad Nacional Autónoma de Nicaragua (2015).
\end{cvlist}
\cvsec{Olimpiadas, distinciones e idiomas}
\begin{cvlist}
  \item Miembro de la Comisión Organizadora de las Olimpiadas Costarricenses de Matemática (OLCOMA) desde 2013; jefe de delegación costarricense a la Olimpiada Internacional de Matemática (IMO) 2014; miembro del comité selector del Banco de Problemas de la XVI OMCC (2014); juez de la Olimpiada Nacional de Robótica (2014).
  \item Coordinador de la competencia de matemática para secundaria del Instituto Tecnológico de Georgia (2006--2008); tutor voluntario de equipos costarricenses en olimpiadas iberoamericana y centroamericana (2001--2003).
  \item Festa Fellowship Award, Georgia Institute of Technology (2007). Idiomas fluidos: español e inglés.
\end{cvlist}
\end{cvblock}

\begin{cvblock}{Lacy Mora, Allan}{Instructor e investigador, Universidad de Costa Rica}
\cvsec{Formación académica}
\begin{cvlist}
  \item Licenciatura en Docencia, Universidad Estatal a Distancia, Costa Rica, 2021--2026. Tesis: Análisis de la percepción del profesorado sobre el uso de Sistemas de Álgebra Computacional en relación con los contenidos, la evaluación y los objetivos de aprendizaje del curso MA1004 Álgebra Lineal de la Universidad de Costa Rica, Sede Rodrigo Facio, durante el segundo semestre 2025.
  \item Ph.D. en Matemáticas, University of Georgia, Athens, Georgia, EEUU, 2009--2015. Tesis: On the index of genus one curves (Sobre el índice de curvas de género uno); profesor tutor: Pete L. Clark.
  \item B.S. en Matemáticas, Universidad de Costa Rica, 2002--2005.
\end{cvlist}
\cvsec{Experiencia profesional y docente}
\begin{cvlist}
  \item Instructor e investigador, Universidad de Costa Rica, desde 2009. Cursos: Teoría de Galois (II-2026), Introducción a la Topología (I-2025), Variable Compleja (I-2023), Seminario de Teoría de Números (II-2017), Grupos y anillos (I-2016).
  \item En la misma institución: Álgebra Lineal I y II, Álgebra Lineal para Computación y Álgebra Lineal para Ingeniería (2016--2026), esta última con coordinación de 2020 a 2023.
  \item Ecuaciones Diferenciales para Ingeniería (2015--2018; coordinador en 2017) y Ecuaciones Diferenciales Aplicadas (2015--2026); Cálculo para Ciencias Económicas III (III-2018) y Cálculo para Computación I (II-2019).
  \item Instructor y teaching assistant, University of Georgia, Athens, EEUU, 2011--2015: Precálculo, Geometría Analítica y Cálculo, Cálculo I para las Ciencias y la Ingeniería.
  \item Instructor, Universidad de Costa Rica, 2006--2009: Ecuaciones Diferenciales para Ingeniería y Aplicadas, Matemática para Computación I, Matemática para Ciencias Económicas I, Cálculo I y II.
\end{cvlist}
\cvsec{Áreas de especialidad e investigación}
Teoría de números, geometría aritmética, dinámica aritmética y docencia de las matemáticas.
\cvsec{Producción académica}
2 artículos de investigación en coautoría (2019 y 2024). Más recientes primero:
\begin{cvlist}
  \item Dynatomic Galois groups for a family of quadratic rational maps, International Journal of Number Theory 20(7) (2024), 1701--1724.
  \item There are genus one curves of every index over every infinite, finitely generated field, Journal für die reine und angewandte Mathematik (Crelle's Journal) 2019 (749), 65--86.
\end{cvlist}
\cvsec{Manejo de software}
Lenguajes de programación y herramientas: C++, Python y LaTeX.
\end{cvblock}

\begin{cvblock}{Loría Sorio, Jennifer de los Ángeles}{Profesora, Universidad de Costa Rica}
\cvsec{Formación académica}
\begin{cvlist}
  \item PH.D., Instituto de Matemática Pura e Aplicada (IMPA), Rio de Janeiro, Brasil, 2023.
  \item M. Sc., Instituto de Matemática Pura e Aplicada (IMPA), Rio de Janeiro, Brasil, 2017.
  \item B. Sc., Universidad de Costa Rica, San José, Costa Rica, 2014.
\end{cvlist}
\cvsec{Experiencia profesional y docente}
\begin{cvlist}
  \item Profesora, Universidad de Costa Rica, 2014 y 2023--actualidad.
  \item Asistente (álgebra y finanzas), IMPA, 2018--2021.
\end{cvlist}
\cvsec{Áreas de especialidad e investigación}
Matemáticas aplicadas: regularización de problemas inversos, análisis de datos, epidemiología y modelización computacional.
\cvsec{Producción académica}
5 artículos científicos (2021--2025), todos en coautoría. Más recientes primero:
\begin{cvlist}
  \item On the estimation of the time-dependent transmission rate in epidemiological models, Inverse Problems 41(6) (2025).
  \item Time-dependent vaccine efficacy estimation quantified by a mathematical model, Plos One 18(5) (2023).
  \item A specific prediction model for deep sternal wound infection after coronary artery bypass surgery, Scientific Reports 12(1) (2022).
  \item COVID-19 underreporting and its impact on vaccination strategies, BMC Infectious Diseases 21(1) (2021).
  \item The impact of COVID-19 vaccination delay: a data-driven modeling analysis for Chicago and New York City, Vaccine 39(41) (2021).
\end{cvlist}
\cvsec{Charlas y eventos}
\begin{cvlist}
  \item Expositora: ``Vaccine Efficacy Estimation Quantified by a Mathematical Model'', XXIV SIMMAC, Costa Rica (2024); ``Modelos matemáticos de enfermedades infecciosas'', EMALCA, Costa Rica (2024); ``Modelos matemáticos de enfermedades infecciosas: una introducción'', Colectivo Mujeres y + en matemáticas, Costa Rica (2021).
  \item Comité organizador del 1º Encontro de Mulheres Matemáticas do IMPA y del Seminário das mulheres do IMPA (ambos 2021).
  \item Participante: XXIV Simposio Internacional de Métodos Matemáticos Aplicados a las Ciencias, Liberia (2024); 33 y 32 Colóquio Brasileiro de Matemática, IMPA (2021, 2019); Encontro Brasileiro de Mulheres Matemáticas, IMPA (2019); International Congress of Mathematicians, Rio de Janeiro (2018); The Mathematics of Games in the Applied Sciences, UFF, Niterói (2018).
\end{cvlist}
\cvsec{Idiomas y habilidades técnicas}
\begin{cvlist}
  \item Español nativo; portugués: lectura y conversación avanzadas; inglés: comprensión lectora intermedia y conversación básica.
  \item Lenguajes de programación: Python y Matlab.
\end{cvlist}
\end{cvblock}

\begin{cvblock}{Mena Arias, Darío Alberto}{Profesor de la Escuela de Matemática e investigador del CIMPA, Universidad de Costa Rica; Director del Programa de Posgrado en Matemática}
\cvsec{Formación académica}
\begin{cvlist}
  \item Ph.D. en Matemática, Georgia Institute of Technology, Atlanta, Georgia, EEUU, 2018.
  \item M.Sc. en Matemática, Georgia Institute of Technology, Atlanta, Georgia, EEUU, 2015.
  \item B.S. en Matemática, Universidad de Costa Rica, 2013.
  \item B.S. en Enseñanza de la Matemática, Universidad de Costa Rica, 2010.
\end{cvlist}
\cvsec{Experiencia profesional y docente}
\begin{cvlist}
  \item Director, Programa de Posgrado en Matemática, Sistema de Estudios de Posgrado, Universidad de Costa Rica, 2022--actualidad.
  \item Profesor de la Escuela de Matemática e investigador adscrito al Centro de Investigación en Matemática Pura y Aplicada (CIMPA), Universidad de Costa Rica, 2018--actualidad.
  \item Instructor y asistente de investigación, Georgia Institute of Technology, Atlanta, 2013--2018.
  \item Profesor a tiempo completo para cursos de cálculo y miembro de la Comisión de Olimpiada Costarricense de Matemática, Universidad de Costa Rica, 2010--2013.
  \item Instructor en cursos de la carrera de Enseñanza de la Matemática, Universidad Estatal a Distancia, 2010.
\end{cvlist}
\cvsec{Áreas de especialidad e investigación}
Análisis armónico (integrales singulares, operadores maximales y operadores dispersos); probabilidad y análisis estocástico, incluidas medidas con valores martingala cilíndricas y EDP estocásticas.
\cvsec{Producción académica}
9 artículos publicados (2017--2025 con año consignado; uno sin año en el CV) y 5 preprints en arXiv (2025--2026). Más recientes primero:
\begin{cvlist}
  \item Preprints (2026): The martingale representation theorem for cylindrical martingale valued measures; Markov property and path regularity for the solutions to SPDEs driven by cylindrical-martingale valued measures; Szemerédi's theorem along Cantor sets of integers; Integer Cantor sets: arithmetic combinatorial properties. Preprint (2025): Weighted weak type estimates for non-integral square functions.
  \item Itô's formula for Itô processes defined with respect to a cylindrical-martingale valued measure, Stochastic Analysis and Applications (sin año en el CV).
  \item Cylindrical martingale-valued measures, stochastic integration and SPDEs, Stoch. Partial Differ. Equ. Anal. Comput. 13 (2025).
  \item A nonlinear relapse model with disaggregated contact rates: analysis of a forward-backward bifurcation, Infectious Disease Modelling 8 (2023).
  \item Sparse bounds for the discrete spherical maximal functions, Pure Appl. Anal. 2 (2020); Lacunary Discrete Spherical Maximal Functions, New York J. Math. 25 (2019).
  \item Characterization of a two-parameter matrix valued BMO by commutator with the Hilbert transform, Rocky Mountain J. Math. 48 (2018).
  \item Sparse Bounds for Bochner--Riesz Multipliers, J. Fourier Anal. Appl. (2017); Uniform sparse bounds for discrete quadratic phase Hilbert transforms, Anal. Math. Phys. (2017); The sparse T1 theorem, Houston J. Math. 43 (2017).
\end{cvlist}
\cvsec{Charlas, congresos y organización de eventos}
\begin{cvlist}
  \item Charlas: XXV SIMMAC, San José (2026); Mathematical Congress of the Americas, Miami (2025); seminarios de análisis de Georgia Tech (2015--2019, 2025); University of Minnesota (2021); University of Birmingham y University of Helsinki (2019); Hausdorff Center, Universität Bonn (2017); Ohio River Analysis Meeting (2017); AMS, Athens (2016).
  \item Organización: comité organizador de SIMMAC 2026 y 2024; coorganizador de la sesión Harmonic Analysis and Partial Differential Equations, Mathematical Congress of the Americas (2025); Harmonic Analysis and related topics, Centre de Recerca Matemàtica, Barcelona (2022).
  \item Asistencia a numerosos congresos y talleres (2015--2026), entre ellos MSRI, Berkeley (2017); AIM, San Jose (2017); IWOTA (2016); Universidad de Bordeaux (2015); Universidad Autónoma de Madrid (2018).
\end{cvlist}
\cvsec{Idiomas}
Español: lengua materna. Inglés: dominio del 100 \%.
\end{cvblock}

\begin{cvblock}{Méndez Hernández, Pedro}{Profesor Catedrático, Escuela de Matemática, Universidad de Costa Rica}
\cvsec{Formación académica}
\begin{cvlist}
  \item Ph.D. en Matemática, Purdue University, Estados Unidos, mayo de 2001.
  \item Licenciatura en Matemática, Universidad de Costa Rica, agosto de 1996.
  \item Bachillerato en Matemática, Universidad de Costa Rica, diciembre de 1994.
\end{cvlist}
\cvsec{Experiencia profesional y docente}
\begin{cvlist}
  \item Secretario Académico de Rectoría, Universidad de Costa Rica, 2022.
  \item Director del Programa de Posgrado en Matemática, UCR, 2015--2021 (y también 2007--2009); Director de la Escuela de Matemática, UCR, octubre de 2009--octubre de 2013.
  \item Profesor Catedrático (desde agosto de 2008), Profesor Asociado (agosto de 2006--julio de 2008), Profesor Invitado (agosto de 2004--julio de 2006) e Instructor (febrero de 1995--julio de 1996), Universidad de Costa Rica.
  \item Visiting Assistant Professor, Purdue University, agosto--diciembre de 2006; Wylie Assistant Professor, University of Utah, agosto de 2001--junio de 2004; Graduate Teaching Assistant, Purdue University, agosto de 1996--mayo de 2001.
  \item Council Member de la Bernoulli Society, 2013--2017.
\end{cvlist}
\cvsec{Áreas de especialidad e investigación}
Probabilidad y sus aplicaciones a la teoría espectral y a la teoría de potenciales. Los trabajos publicados tratan sobre procesos estables simétricos, movimiento browniano y desigualdades de tipo Brascamp--Lieb--Luttinger, ultracontractividad intrínseca, núcleos de calor y operadores de Schrödinger, y caminatas aleatorias dinámicas.
\cvsec{Producción académica}
Lista de 19 entradas: 18 artículos publicados entre 1992 y 2010 y un manuscrito (sin entradas posteriores a 2010). Representativas, más recientes primero:
\begin{cvlist}
  \item Journal of Functional Analysis 258 (2010): simetrización de procesos de Lévy y aplicaciones; Journal of Functional Analysis 257 (2009): pequeñas perturbaciones, ultracontractividad semi-intrínseca y representaciones integrales de soluciones de ecuaciones parabólicas.
  \item Stochastic Processes and their Applications 118 (2008): capacidades en el espacio de Wiener, funciones inferiores casi seguras y entropía de Kolmogorov.
  \item Transactions of the American Mathematical Society 359 (2007): propiedades de continuidad del proceso estable simétrico; Electronic Journal of Probability 12 (2007): tiempos de salida de procesos estables de dominios no acotados.
  \item Probability Theory and Related Fields 134 (2006): tiempos excepcionales e invariancia para caminatas aleatorias dinámicas; Potential Analysis 24 (2006), dos artículos; Rocky Mountain Journal of Mathematics 36 (2006).
  \item Annals of Probability 33 (2005): caminatas aleatorias gaussianas dinámicas.
  \item Indiana University Mathematics Journal 52 (2003); Illinois Journal of Mathematics 46 (2002); Duke Mathematical Journal 113 (2002): desigualdades de Brascamp--Lieb--Luttinger para dominios convexos de inradio finito.
  \item Journal of Functional Analysis 176 (2000) y Michigan Mathematical Journal 47 (2000), entre otros.
\end{cvlist}
Presentó 20 charlas y seminarios en el extranjero (Estados Unidos, Brasil, Japón, Alemania y México, entre 2001 y 2011), incluida la conferencia inaugural del Simposio Internacional de Métodos Matemáticos Aplicados a las Ciencias (UCR, 2006), y participó en 20 conferencias (1998--2014), entre ellas el Congreso Mundial de la Bernoulli Society (Guanajuato, 2000) y el Congreso Internacional de Matemáticos (Seúl, 2014).
\cvsec{Distinciones y becas}
\begin{cvlist}
  \item Premio TWAS-CONICIT 2005 (Third World Academy of Science), abril de 2005; CMI Liftoff Mathematician, Clay Mathematics Institute, mayo--agosto de 2001.
  \item Purdue Research Foundation Fellowship, Purdue University, 1998--2000; NSF Research Assistant, Purdue University (enero--agosto de 1998 y agosto--diciembre de 2000).
\end{cvlist}
\end{cvblock}

\begin{cvblock}{Montalto Cruz, Carlos}{Profesor Catedrático, Departamento de Matemáticas, Universidad de Costa Rica}
\cvsec{Formación académica}
\begin{cvlist}
  \item Ph.D. en Matemáticas, Purdue University, EEUU, agosto 2008--agosto 2024 (periodo tal como figura en el CV); asesor: Plamen Stefanov. Los artículos de su tesis versan sobre problemas inversos para ecuaciones elípticas e hiperbólicas.
  \item Bachillerato en Matemáticas, Universidad de Costa Rica, marzo 2003--marzo 2008.
\end{cvlist}
\cvsec{Experiencia profesional y docente}
\begin{cvlist}
  \item Profesor Catedrático, Universidad de Costa Rica, desde setiembre 2017. Desarrolla algoritmos de imagen fotoacústica en el Laboratorio de Investigación en Biomedicina; coordina el proyecto de Olimpiadas Costarricenses de Matemáticas (OLCOMA); imparte cursos de grado y posgrado; organizó el Seminario de Ecuaciones Diferenciales y Métodos Numéricos.
  \item Profesor Asistente (investigador postdoctoral con Gunther Uhlmann), University of Washington, setiembre 2014--setiembre 2017.
  \item Asistente de investigación y docencia, Purdue University, agosto 2008--agosto 2014.
  \item Asistente del Instituto de Investigaciones Psicológicas, UCR, junio--agosto 2011; profesor visitante, UCR, julio 2006--junio 2008; asistente de docencia, UCR, enero 2004--junio 2006; instructor de las Olimpiadas Nacionales de Matemáticas, Colegio Monterrey, 2006--2007.
\end{cvlist}
\cvsec{Áreas de especialidad e investigación}
Problemas inversos, modelación matemática e imagenología: imagen fotoacústica, imagen de la densidad de impedancia eléctrica, determinación estable de coeficientes a partir del operador Dirichlet-a-Neumann y algoritmos para mejorar la percepción visual en realidad aumentada.
\cvsec{Producción académica}
Siete artículos publicados (2013--2024), cinco preprints y 37 conferencias y presentaciones (2010--2026). Artículos, más recientes primero:
\begin{cvlist}
  \item A SharpView With Enhanced Out-of-Focus Text Legibility for Augmented Reality Systems, IEEE VR, 2024.
  \item Stability in conductivity imaging from partial measurements of one interior current, Inverse Problems and Imaging, 2017.
  \item Photoacoustic imaging taking into account thermodynamic attenuation, Inverse Problems, 2016.
  \item Multiwave imaging in an enclosure with variable sound speed, Inverse Problems, 2015.
  \item A Total Variation Approach for Customizing Imagery to Improve Visual Acuity, ACM Transactions on Graphics, 2015.
  \item Stable determination of a simple metric, a co-vector field and a potential from the hyperbolic Dirichlet-to-Neumann map, Comm. Partial Differential Equations, 2014.
  \item Stability of Coupled-Physics Inverse Problems with one Internal Measurement, Inverse Problems, 2013.
\end{cvlist}
\cvsec{Docencia y formación de recursos humanos}
\begin{cvlist}
  \item Cursos en la UCR (Cálculo Multivariable, Ecuaciones Diferenciales, Análisis Real, entre otros), en University of Washington (Microlocal Analysis, Linear Optimization, Linear Analysis, Discrete Mathematical Modeling) y en Purdue.
  \item Doctorado: Fernando Herrera (Colorado State University, en proceso) y Arefin Mohammed (Mississippi State University, coasesor Edward Swan II). Maestría en la UCR: Gerardo Acosta y Daniel Grajal (ambas en proceso).
\end{cvlist}
\cvsec{Distinciones, becas e idiomas}
AMS-Simons Travel Award (2016); Beca Complementaria para Estudios en el Extranjero de la UCR (2008); Beca de Honor de la UCR (2004); diploma de mejor estudiante de Matemática de la UCR (2003). Idiomas: español (nativo) e inglés (nivel profesional).
\end{cvblock}

\begin{cvblock}{Montenegro Guzmán, Samaria}{Directora del Departamento de Matemática y Profesora Asociada, Universidad de Costa Rica}
\cvsec{Formación académica}
\begin{cvlist}
  \item Doctorado en Lógica Matemática, Université Paris Diderot, París, Francia, 2011--2015. Tesis: ``Théorie des modèles des corps pseudo-réels clos et pseudo-p-adiquement clos''; directora: Zoé Chatzidakis.
  \item Maestría en Lógica Matemática y Fundamentos de la Informática, Université Paris Diderot, París, Francia, 2011.
\end{cvlist}
\cvsec{Experiencia profesional y gestión académica}
\begin{cvlist}
  \item Directora del Departamento de Matemática, Universidad de Costa Rica, 2025--presente; Profesora Asociada, UCR, 2017--presente.
  \item Postdoctorado, Universidad de los Andes, Bogotá, Colombia, hasta 2017 (el año de inicio aparece incompleto en el CV).
  \item Program Associate en el programa Arithmetic Geometry and Number Theory, MSRI, Berkeley, Estados Unidos, enero--mayo 2014.
\end{cvlist}
\cvsec{Áreas de especialidad e investigación}
Lógica matemática y teoría de modelos de cuerpos: cuerpos pseudo real cerrados, pseudo p-ádicamente cerrados y pseudo T-cerrados, propiedad NTP2, imaginarios, tipos invariantes y grupos definibles (incluso en cuerpos diferenciales parciales con automorfismo).
\cvsec{Producción académica}
Siete artículos (2017--2026), uno de ellos aceptado para publicación. Más recientes primero:
\begin{cvlist}
  \item Dividing lines among fields, en Women in Mathematics in Latin America, Trends in Mathematics (Birkhäuser/Springer), 2026.
  \item Pseudo T-closed fields, Model Theory 4(1), 2025.
  \item Groups definable in partial differential fields with an automorphism, Israel Journal of Mathematics (aceptado; consignado en 2023).
  \item Imaginaries, invariant types and pseudo-p-adically closed fields, Transactions of the American Mathematical Society 374, 2021.
  \item Stabilizers, NTP2 groups with f-generics and PRC fields, Journal of the Institute of Mathematics of Jussieu 19(3), 2020.
  \item Imaginaries in bounded pseudo real closed fields, Annals of Pure and Applied Logic 168(10), 2017.
  \item Pseudo real closed fields, pseudo p-adically closed fields and NTP2, Annals of Pure and Applied Logic 168(1), 2017.
\end{cvlist}
Charlas (2022--2026): A85 (Edimburgo, 2026), Model Theory at Notre Dame en honor de Zoé Chatzidakis (2025), XXI Brazilian Logic Conference (2025), XX Simposio Latinoamericano de Lógica Matemática (Montevideo, 2024), Logic Colloquium en UCLA (2024), CIRM-Luminy (2023), Université Paris Cité (2023), UNAM-Morelia (2022) y Université de Paris (2022).
\cvsec{Becas, fondos y proyectos}
\begin{cvlist}
  \item Pasantías de investigación en la École normale supérieure, París (enero--febrero 2023), y en el Institut de Mathématiques de Jussieu--Paris Rive Gauche (febrero--marzo 2022).
  \item Fondo UCREA 2022--2024 (metodología feminista para el análisis de género de la producción académica en la UCR); fondo semilla 2020--2022 sobre grupos dentro del universo de Shelah (investigadora asociada).
\end{cvlist}
\cvsec{Comités y otras actividades académicas}
\begin{cvlist}
  \item Miembro del comité ejecutivo y del consejo de la Association for Symbolic Logic (ASL), 2026--presente; de su comité para América Latina, 2022--presente; embajadora de la IMU para el comité de mujeres en matemáticas, 2021--presente.
  \item Miembro del comité científico del CIMPA de la UCR, 2020--2024; de los comités organizador y científico de simposios latinoamericanos de lógica (2018--2020), de la Escuela CIMPA de 2026 en Costa Rica, de EMALCA 2019 y de foros sobre género en STEM (2019--2021).
\end{cvlist}
\end{cvblock}

\begin{cvblock}{Mora Sáenz, Germán José}{Universidad de Costa Rica (2025--presente)}
\cvsec{Formación académica}
\begin{cvlist}
  \item Doctorado en Matemática, The Ohio State University, Columbus, Ohio, Estados Unidos, 2022--2025. Área: Mecánica de Fluidos.
  \item Maestría en Matemática Aplicada, The Ohio State University, Columbus, Ohio, Estados Unidos, 2019--2022.
  \item Bachillerato en Matemática, Universidad de Costa Rica, San Pedro, Costa Rica, 2013--2017, con graduación con honores.
\end{cvlist}
\cvsec{Experiencia profesional}
\begin{cvlist}
  \item Universidad de Costa Rica, 2025--presente.
  \item The Ohio State University, 2019--2025.
  \item Universidad de Costa Rica, 2017--2019.
\end{cvlist}
\cvsec{Idiomas}
Español e inglés.
\end{cvblock}

\begin{cvblock}{Ramírez González, José Alexander}{Profesor catedrático, Escuela de Matemática, Universidad de Costa Rica}
\cvsec{Formación académica}
\begin{cvlist}
  \item Ph.D., Courant Institute of Mathematical Sciences, New York University, Estados Unidos, 1996--2000; tutor: S.R.S. Varadhan.
  \item Bachiller en matemáticas (1991--1995) y bachiller en ingeniería eléctrica (1990--1995), Universidad de Costa Rica.
\end{cvlist}
\cvsec{Experiencia profesional y docente}
\begin{cvlist}
  \item Profesor, Escuela de Matemática, UCR, 1995--1996 y julio 2003--presente (actualmente catedrático). Director del CIMPA, 2018--2021; director del posgrado de matemáticas, 2005--2006; integrante de la Comisión de Doctorado en Ciencias, 2013--2025; director del departamento de matemática pura y ciencias actuariales (periodo consignado como «2011 a 2010» en el CV).
  \item Presidente y desarrollador, Quanto Soluciones Analíticas y Cuantitativas (software para análisis de riesgo), diciembre 2014--presente.
  \item Profesor, Escuela de Ingeniería Eléctrica, UCR, enero 2014--2016; Escuela de Matemática, Universidad Nacional, julio 2010--diciembre 2013.
  \item Director de modelación matemática, Dirección Corporativa de Riesgo, Banco Nacional de Costa Rica, enero 2004--julio 2010.
  \item Profesor visitante en Harvard University (setiembre 2008--enero 2009) y University of Colorado, Boulder (enero--mayo 2006); H.C. Wang Assistant Professor, Cornell University, 2000--2003.
\end{cvlist}
\cvsec{Áreas de especialidad e investigación}
Probabilidad, procesos estocásticos y ecuaciones en derivadas parciales, con relevancia en física matemática: núcleos del calor, gases reticulados, operadores diferenciales aleatorios y matrices aleatorias. Además, modelos de riesgo financiero (crédito, liquidez, mercado y operacional).
\cvsec{Producción académica}
Trece artículos (2001--2016) y 22 charlas principales (2000--2021). Representativos, más recientes primero:
\begin{cvlist}
  \item Spiking the random matrix hard edge, Probab. Theory Relat. Fields, 2016.
  \item Hard edge tail asymptotics, Electronic Communications in Probability, 2011.
  \item Bulk universality for Wigner hermitian matrices with subexponential decay, Mathematical Research Letters, 2010; Bulk Universality for Wigner Matrices, Comm. Pure Appl. Math., 2010.
  \item Beta ensembles, stochastic Airy spectrum, and a diffusion, Journal of the American Mathematical Society (aceptado, 2010).
  \item On the Shape of the Ground State Eigenvalue Density of a Random Hill's Equation, Comm. Pure Appl. Math., 2006; Superdiffusivity of Lattice Gas Models, J. Stat. Phys., 2005.
  \item Small Time Gaussian Behavior of Symmetric Diffusion Semigroups, Ann. Probab., 2003; Short Time Asymptotics in Dirichlet Spaces, Comm. Pure Appl. Math., 2001.
\end{cvlist}
\cvsec{Proyectos, consultorías y extensión}
Estudio Actuarial del Régimen de Invalidez, Vejez y Muerte de la CCSS (con Rodrigo Arias y Luis Barboza, 2016). Cursos y charlas para el sector financiero (2012--2026) sobre inteligencia artificial y riesgo de crédito, riesgo de liquidez, pruebas de estrés, lavado de dinero y gestión integral de riesgos, en Costa Rica, Panamá y República Dominicana.
\cvsec{Docencia y formación de recursos humanos}
Siete tesis dirigidas (una en proceso) en maestría en Matemática Aplicada (Luis Barboza, 2008; Maikol Solís, 2010) y licenciatura en Ciencias Actuariales (2012--2024, incluida una en proceso).
\cvsec{Distinciones y actividades}
Premio TWAS-Conicit (2012); Premio Wilhelm T. Magnus, NYU (2001); Investigador Liftoff, Clay Mathematics Institute (2000). Presidente de la Sociedad Latinoamericana de Probabilidad y Estadística Matemática (2026); presidente del comité organizador de congresos y escuelas latinoamericanas (2012--2016) y de jurados de olimpiadas matemáticas (2012, 2014).
\end{cvblock}

\begin{cvblock}{Rodríguez Rojas, Oldemar}{Profesor e investigador, Instituto Tecnológico de Monterrey (México); Profesor Catedrático de la Universidad de Costa Rica}
\cvsec{Formación académica}
\begin{cvlist}
  \item Postdoctorado en ``Data Mining and Applications'', Escuela de Estadística, Stanford University, California, EE.\ UU., 2009--2011.
  \item Doctor, U.F.R. Mathématique de la Décision, Université de Paris IX -- Dauphine, Francia, 2000, con mención ``Très Honorable''; título acreditado en EE.\ UU. como Ph.D. por WES. Tutor: Edwin Diday. Tesis: ``Classification et Modèles Linéaires en Analyse des Données Symboliques''.
  \item Diploma de Estudios a Profundidad en Informática (D.E.A.), Université de Paris IX -- Dauphine, 1998.
  \item Magister Scientae en Computación, Instituto Tecnológico de Costa Rica, 1994 (magna cum laude). Tesis: ``Desarrollo Orientado a Objetos: Una Aplicación al Análisis de Datos''.
  \item Licenciado en Matemática Pura (1991; tesis: ``Tópicos de Programación no Lineal'') y Bachiller en Matemática Pura (1988), Universidad de Costa Rica.
\end{cvlist}
\cvsec{Experiencia profesional y gestión académica}
\begin{cvlist}
  \item Profesor e investigador, Instituto Tecnológico de Monterrey, México, desde 2019. Profesor Catedrático de la Universidad de Costa Rica y de la Universidad Nacional.
  \item Vicerrector de Administración, UCR, enero--junio 2021. CIMPA: subdirector, febrero 2016--junio 2019, y director a.i., febrero 2016--2018.
  \item Consejo Universitario, UCR: miembro, octubre 2008--octubre 2011, y director, octubre 2009--octubre 2010.
  \item Decano de la Facultad de Ciencias, UCR, agosto 2004--octubre 2008; vicedecano, julio 2002--julio 2004; director del Programa de Maestría en Matemática, febrero 2001--enero 2003.
\end{cvlist}
\cvsec{Áreas de especialidad e investigación}
Ciencias de la Computación y Matemática Aplicada: análisis de datos simbólicos (componentes principales, regresión, escalamiento multidimensional y análisis de correspondencias para variables de tipo intervalo y multivaluadas), minería de datos, ciencia de datos y aprendizaje automático.
\cvsec{Producción académica}
41 entradas numeradas de publicaciones (artículos, actas, libros y tesis), 13 participaciones en eventos internacionales y 5 paquetes de R en CRAN. Representativas, más recientes primero:
\begin{cvlist}
  \item Regression Models for Symbolic Interval-Valued Variables, Entropy, 2021.
  \item Optimized dimensionality reduction methods for interval-valued variables and its application to face recognition, Entropy, 2019.
  \item Shrinkage linear regression for symbolic interval-valued variables, Revue MODULAD, 2018.
  \item Principal Curves and Surfaces to Interval Valued Variables, IBERAMIA 2016 (LNAI 10022); INTERSTATIS, Revista de Matemática: Teoría y Aplicaciones, 2013.
  \item I-Scal: Symbolic Multidimensional Scaling of Interval Dissimilarities, Computational Statistics \& Data Analysis; ponencias en ECML/PKDD 2004 e IFCS 2000 y 2004; libro de programación orientada a objetos en C++ (1997).
\end{cvlist}
Paquetes de R en CRAN, con descargas según el CV: RSDA (más de 30\,095), discoveR (22\,266), predictoR (15\,845), regressoR (13\,472) y traineR (11\,750).
\cvsec{Conferencias, distinciones e idiomas}
\begin{cvlist}
  \item Presentaciones en SDA 2012--2018, IFCS 2013 y SIMMAC 2017; cursos cortos sobre RSDA en la Université Paris-Dauphine (2015), Oporto (2018) y el Tecnológico de Monterrey (2019); conferencias de apertura en la Costa Rica Big Data School (2018) y el Encuentro de Ciencia de Datos del Grupo ICE (2018).
  \item Medalla de Oro al Inventor Destacado del Año 2004, OMPI, por ``Symbolic Personal Vector to detect frauds in credit cards''; Premio Nacional de Tecnología Clodomiro Picado Twight 2001. Idiomas: inglés y francés.
\end{cvlist}
\end{cvblock}

\begin{cvblock}{Rosales Ortega, José}{Profesor catedrático, Universidad de Costa Rica y Tecnológico de Costa Rica}
\cvsec{Formación académica}
\begin{cvlist}
  \item Doctorado en Matemáticas, CINVESTAV-IPN, México, 2006. Formación doctoral en acciones de grupos de Lie sobre variedades pseudo-riemannianas. Tesis: ``The Gromov's Centralizer Theorem for Semisimple Lie Group Actions''.
  \item Maestría en Matemática Pura, Universidad de Costa Rica, 1996.
  \item Bachillerato en Matemática Pura, Universidad de Costa Rica, 1989.
\end{cvlist}
\cvsec{Experiencia profesional y docente}
\begin{cvlist}
  \item Profesor catedrático, Universidad de Costa Rica (UCR) y Tecnológico de Costa Rica, con experiencia docente en ambas instituciones.
  \item Subdirector, Escuela de Matemática, UCR, noviembre 2023--noviembre 2025.
  \item Director de la Comisión de Trabajos Finales de Graduación e Investigación, Escuela de Matemática, UCR, 2022--2025.
  \item Director interino del Posgrado en Matemática, UCR, II semestre de 2012.
  \item Coordinador de Ecuaciones Diferenciales para Ingeniería, UCR, 2009--2016, y de Cálculo III, UCR, 2010--actualidad.
  \item Cursos representativos: Cálculo III; Cálculo y Geometría Analítica; Cálculo y Análisis 2; Ecuaciones Diferenciales para Ingeniería; Introducción a las Variables Complejas.
\end{cvlist}
\cvsec{Áreas de especialidad e investigación}
Geometría diferencial de curvas y superficies; acciones de grupos de Lie sobre variedades pseudo-riemannianas y estructuras geométricas; operadores de Toeplitz e integrales en espacios de Bergman y Fock; teoría de operadores y espacios de funciones analíticas.
\cvsec{Producción académica}
El CV presenta referencias seleccionadas: 7 artículos de investigación publicados (1998--2024), 1 capítulo aceptado, 2 publicaciones arbitradas adicionales (2017) y 1 texto de divulgación (2000). Representativas, de la más reciente a la más antigua:
\begin{cvlist}
  \item Capítulo aceptado sobre operadores de Toeplitz con símbolos invariantes en el espacio de Fock, en \emph{Tales and Trends in Operator Theory, Evolution Equations and Complex Variables: IWOTA 2024}, Birkhäuser/Springer (serie Operator Theory: Advances and Applications).
  \item ``Toeplitz operators with symmetric, alternating and anti-symmetric separately radial symbols on the unit ball'', Boletín de la Sociedad Matemática Mexicana, 2024 (Scopus y Web of Science).
  \item ``A geometric splitting theorem for actions of semisimple Lie groups'', Abhandlungen aus dem Mathematischen Seminar der Universität Hamburg, 2021.
  \item ``The Gromov's Centralizer Theorem for Semisimple Lie Group Actions'', Boletín de la Sociedad Matemática Mexicana, 2017.
  \item ``Foliation by G-orbits'', Boletim da Sociedade Paranaense de Matemática, 2015.
  \item ``G-estructuras de orden superior'', Revista de Matemática: Teoría y Aplicaciones, 2014 (SciELO, Redalyc).
  \item ``The Signature in Actions of Semisimple Lie Groups on Pseudo-Riemannian Manifolds'', Proyecciones, 2012.
  \item ``La factorización de una transformada de Fourier en el método de Wiener--Hopf'', Revista de Matemática: Teoría y Aplicaciones, 1998.
  \item Otras: ``Soluciones entre ecuaciones diferenciales de orden dos'', Revista MATUA, 2017 (TIndex); ``Numerabilidad y cardinalidad de conjuntos'', Revista Digital: Matemática, Educación e Internet, 2017; divulgación: ``O Princípio do Elemento Extremo'', Eureka!, 2000.
\end{cvlist}
\cvsec{Proyectos}
\begin{cvlist}
  \item ``Operadores Integrales en Espacios de Funciones Analíticas'', proyecto CIMPA-UCR, 2026--2029; investigador principal, con colaboración de Armando Sánchez-Nungaray y Carlos González-Flores.
  \item ``Interacciones del análisis y la geometría'', proyecto de investigación CIMPA-UCR, 2023--2024; investigador principal.
\end{cvlist}
\end{cvblock}

\begin{cvblock}{Sánchez Chavarría, Adriana}{Profesora Adjunta, Escuela de Matemática, Universidad de Costa Rica}
\cvsec{Formación académica}
\begin{cvlist}
  \item MSc. Estadística, Universidad de Costa Rica, 2024--actualidad (en curso).
  \item PhD. Matemáticas, especialidad Sistemas Dinámicos, Instituto de Matemática Pura e Aplicada (IMPA), 2014--2018. Orientador: Marcelo Viana. Tesis: ``Contribuciones al problema de continuidad de los Exponentes de Lyapunov''.
  \item MSc. Matemáticas, IMPA, 2012--2014. Orientador: Robert Morris.
  \item BSc. Ciencias Actuariales, Universidad de Costa Rica, 2004--2011; BSc. Matemáticas, Universidad de Costa Rica, 2004--2008.
\end{cvlist}
\cvsec{Experiencia profesional y docente}
\begin{cvlist}
  \item Profesora Adjunta, Escuela de Matemática, Universidad de Costa Rica (UCR), 2023--presente; profesora invitada exbecaria doctora, 2020--2023; profesora interina, 2009--2012.
  \item Posdoctorado, Instituto de Ciências Matemáticas e de Computação, Universidade de São Paulo, 2018--2020.
  \item Cursos impartidos en UCR (2020--2024): Ecuaciones diferenciales ordinarias (MA0455, varios periodos), Ecuaciones en derivadas parciales, Ecuaciones diferenciales II y Ecuaciones diferenciales, Ecuaciones diferenciales aplicadas, Análisis I, Introducción a los sistemas dinámicos y Herramientas para el cálculo actuarial.
\end{cvlist}
\cvsec{Áreas de especialidad e investigación}
Sistemas dinámicos, exponentes de Lyapunov, estados de equilibrio y teoría ergódica; aprendizaje automático aplicado al pronóstico de caudales extremos. Habilidades de programación: Python, R, Excel y Power BI.
\cvsec{Producción académica}
4 artículos o capítulos publicados (2018--2023), 1 preprint y la tesis doctoral; además, unas 16 charlas, pósteres y seminarios (2018--2023) en Brasil, México, Francia, Inglaterra, India, Italia y Chile.
\begin{cvlist}
  \item ``Analyticity of the Lyapunov exponents of random products of quasi-periodic cocycles'', Nonlinearity, 2023.
  \item ``The Continuity Problem of Lyapunov Exponents'', en International Workshop New Trends in Lyapunov Exponents, Springer, 2022.
  \item ``Equilibrium states for maps isotopic to Anosov'', Nonlinearity, 2021.
  \item ``The set of fiber-bunched cocycles with nonvanishing Lyapunov exponents over a partially hyperbolic map is open'', Mathematical Research Letters, 2018.
  \item ``Lyapunov exponents of probability distributions with non-compact support'', preprint, 2018 (arXiv).
  \item Charlas recientes: ``Caudales extremos y la recuperación de datos'', Data coffee, UCR, noviembre 2023; Americas XIII, febrero 2023; Encontro Mulheres Matemáticas do IMPA, mayo 2022.
\end{cvlist}
\cvsec{Proyectos de investigación}
\begin{cvlist}
  \item Investigadora principal: Proyecto de Educación Continua del CIMPA (2024--2026 y 2022--2023); ``Uso de modelos de aprendizaje automático para el pronóstico de caudales extremos (mínimos y máximos)'' (2023--2025); ``Estudio ergódico mediante herramientas analíticas'' (2022--2023); ``Estudio de estados de equilibrio mediante herramientas analíticas'' (2020--2022).
  \item Investigadora asociada: estudio de sistemas de producción en contexto de industria 4.0 con herramientas de aprendizaje profundo (2022--2025).
\end{cvlist}
\cvsec{Gestión académica e idiomas}
\begin{cvlist}
  \item Comité organizador de la 18th Conference of the International Federation of Classification Societies (2023--presente) y del XXIV (2022--2024) y XXIII (2021--2022) Simposio Internacional de Métodos Matemáticos Aplicados a las Ciencias.
  \item Consejo Editorial de la Revista de Matemática: Teoría y Aplicaciones, Comisión de la mujer en matemática y Comisión equipo de salud de la Escuela de Matemática (desde 2021); miembro de la Asociación Mujeres y + en Matemática (desde 2021).
  \item Idiomas: español (nativa), inglés y portugués (avanzado), francés (iniciante).
\end{cvlist}
\end{cvblock}

\begin{cvblock}{Sánchez Guevara, Jesús}{Profesor Asociado, Escuela de Matemática, Universidad de Costa Rica}
\cvsec{Formación académica}
\begin{cvlist}
  \item Doctorado Académico, Universidad de Costa Rica (registrado en 2022; el CV no indica año de finalización ni especialidad).
  \item Doctorado en Matemáticas, Université Paris 7, Francia, 2016. Especialidad: Topología Algebraica y Teoría de Categorías. Tesis: ``Estructuras E-infinito en L-álgebras''.
  \item Maestría con Especialidad Lógica Matemática y Fundamentos de la Informática, Université Paris 7, Francia, 2011.
  \item Licenciatura en Docencia, Universidad en Ciencias Administrativas San Marcos, Costa Rica, 2023. Tesis: ``Un abordaje de la definición de límite de funciones en la formación de docentes en matemáticas''.
  \item Bachillerato en Matemática Pura, Universidad de Costa Rica, 2007.
\end{cvlist}
\cvsec{Experiencia profesional y docente}
\begin{cvlist}
  \item Profesor Asociado, Escuela de Matemática, Universidad de Costa Rica (UCR), desde marzo de 2024; Asociado desde septiembre de 2023 y Adjunto desde diciembre de 2022.
  \item Profesor interino licenciado, UCR, febrero--diciembre de 2022; profesor exbecario doctor e invitado, Escuela de Matemática, UCR, diciembre de 2016--diciembre de 2021, y brevemente en el Sistema de Estudios de Posgrado (2019).
  \item Profesor interino bachiller, Escuela de Matemática, UCR, en varios periodos de 2007 a 2010 y en 2016; también en la sede de Alajuela, UCR (segundo semestre de 2008).
\end{cvlist}
\cvsec{Áreas de especialidad e investigación}
Topología algebraica y teoría de categorías (óperadas, L-álgebras, coálgebras E-infinito); lógica matemática y fundamentos de la informática; álgebra lineal; clima espacial (eyecciones de masa coronal).
\cvsec{Producción académica}
6 artículos (2015--2025) y 1 libro. El CV registra dos veces el artículo de 2024 en Revista de Filosofía UCR; aquí se cuenta una sola vez.
\begin{cvlist}
  \item ``Extensión de una estructura de óperad definida sobre un sub-s-módulo'', Revista de Matemática: Teoría y Aplicaciones, 2025.
  \item ``Introducción a la teoría de categorías, lógica y topos elementales para una mente curiosa'', Revista de Filosofía de la Universidad de Costa Rica, 2024.
  \item ``On the relationship between expansion angle of Earth-directed CMEs and soft X-ray emission from their related flare'', Revista de Matemática: Teoría y Aplicaciones, 2023.
  \item ``Operads libres sobre módulos diferenciales graduados'', Revista de Matemática: Teoría y Aplicaciones, 2022.
  \item ``A geometrical description for interplanetary propagation of Earth-directed CMEs based on radiative proxies'', Monthly Notices of the Royal Astronomical Society, 2020.
  \item Libro: ``Álgebra lineal fundamental (Teoría y ejercicios)'', Editorial UCR, 2020.
  \item ``E-infinity coalgebra structure on chain complexes with integer coefficients'', Revista Colombiana de Matemáticas, 2015.
\end{cvlist}
\cvsec{Proyectos de investigación}
\begin{cvlist}
  \item ``Herramientas científicas para optimizar el estudio del clima espacial'' (C0072), investigador asociado, 2020--2022 (terminado).
  \item ``Construcción de L-álgebras'' (B6A19), investigador principal, agosto--diciembre de 2018 (estado en el CV: cerrado por incumplimiento).
\end{cvlist}
\cvsec{Formación complementaria, incorporaciones e idiomas}
\begin{cvlist}
  \item Curso de Competencias Digitales para la Docencia (Didáctica Universitaria), 30 horas, UCR, 2023; curso de Didáctica Universitaria, 2019. Incorporado al Colegio de Licenciados y Profesores en Letras, Filosofía, Ciencias y Artes (2024).
  \item Idiomas: francés (conversación, lectura y escritura 100\%); inglés (conversación 80\%, lectura y escritura 100\%); italiano (conversación 20\%, lectura 70\%, escritura 50\%).
\end{cvlist}
\end{cvblock}

\begin{cvblock}{Sánchez Peña, Fabio}{Profesor, Escuela de Matemática, Universidad de Costa Rica (desde septiembre de 2025)}
\cvsec{Formación académica}
\begin{cvlist}
  \item Ph.D., Biological Statistics and Computational Biology, Cornell University, EE.\ UU., enero de 2007 (Sloan Scholar). Tesis: ``Theoretical Studies in Epidemiology and Social Dynamics''; director: Carlos Castillo-Chavez.
  \item B.S. in Computer Science, Minor in Mathematics, Universidad Metropolitana, Puerto Rico, junio de 2001.
\end{cvlist}
\cvsec{Experiencia profesional y docente}
\begin{cvlist}
  \item Universidad de Costa Rica (UCR): profesor visitante (2014--2017), invitado (2017--2018), adjunto (marzo--agosto de 2019), asociado (agosto de 2019--septiembre de 2025) y profesor (desde septiembre de 2025); director del Departamento de Matemática Aplicada, marzo de 2019--febrero de 2023. Fundador y socio de PREDICDATA (analítica de datos), 2026--presente.
  \item Arizona State University: instructor (2005--2006), profesor investigador asistente (2010--2012) y profesor adjunto (2007--2015). Manager, Credit Strategy, American Express, 2006--2009. Los Alamos National Laboratory: asistente de investigación (2003--2004) y coordinador del MTBI (2004--2005).
  \item Colegio Científico Costarricense, profesor (2016 y desde 2023); The European School, Heredia, profesor de matemática, ciencias e informática (julio de 2012--diciembre de 2014).
\end{cvlist}
\cvsec{Áreas de especialidad e investigación}
Biología matemática, modelos epidemiológicos, comportamiento humano adaptativo, salud pública, arbovirus (dengue, chikungunya, Zika) y dinámica social; modelos de recaída no lineal, redes multicapa y modelos climáticos de dengue.
\cvsec{Producción académica}
35 publicaciones (artículos y capítulos), 17 preprints y 22 informes técnicos. Representativas:
\begin{cvlist}
  \item ``High-resolution climate--dengue modeling and mid-century projections under SSP5-8.5 in Costa Rica'', Science in One Health, 2026; ``Managing populations after a disease outbreak'', Philosophical Transactions of the Royal Society B, 2026.
  \item ``Interplay between Foraging Choices and Population Growth Dynamics'', PLOS One, 2025; ``Source-sink dynamics in a two-patch SI epidemic model'', Bulletin of Mathematical Biology, 2024.
  \item ``Assessing dengue fever risk in Costa Rica by using climate variables and machine learning techniques'', PLOS Neglected Tropical Diseases, 2023; ``Bayesian spatio-temporal model with INLA for dengue fever risk prediction'', Environmental and Ecological Statistics, 2023.
  \item ``A multilayer network model of Covid-19'', Epidemics, 2022; ``Projecting the Impact of Covid-19 Variants and Vaccination Strategies'', Scientific Reports, 2022.
  \item ``Comparative analysis of dengue versus chikungunya outbreaks in Costa Rica'', Ricerche di Matematica, 2018; varios artículos en Revista de Matemática: Teoría y Aplicaciones (2019).
  \item Preprints de 2026 sobre dengue, el número reproductivo básico y conducta adaptativa (arXiv, Zenodo, SocArXiv, medRxiv).
\end{cvlist}
\cvsec{Docencia y formación de recursos humanos}
\begin{cvlist}
  \item Cursos en UCR desde 2013: Ecuaciones Diferenciales para Enseñanza, Ecuaciones Diferenciales Aplicadas y para Ingeniería, Análisis Numérico, Álgebra Lineal, Modelos Matemáticos y sus Aplicaciones, Probabilidad Aplicada, Investigación de Operaciones I, Cálculo para Ciencias Actuariales; en ASU, cursos de matemática aplicada y verano del MTBI (2002--2014).
  \item Dirigió 5 tesis de maestría concluidas en UCR (2020--2026) y tiene 5 en curso; asesoró 1 posdoctorado (2020).
\end{cvlist}
\cvsec{Proyectos, editorial y distinciones}
\begin{cvlist}
  \item Proyectos UCREA 2018 (modelos matemáticos contra \emph{Aedes aegypti}), 2020 (LaBsiq, Siquirres) y 2023 (costo de una vida digna); Seed Grant, UC Davis, 2022; beca posdoctoral de la Vicerrectoría de Investigación, 2018. Lideró el equipo EpiMEC de modelación de Covid-19 con las autoridades de salud.
  \item Editor jefe fundador de Epi-SCIENCE (desde 2021); editor académico de Scientific Reports (desde 2024), Frontiers in Epidemiology (desde 2023) y PLOS One (desde 2018); revisor de MathSciNet y de revistas de biología matemática.
  \item Unas 58 presentaciones orales (incluidas plenarias y conferencias magistrales) y 13 pósteres (1999--2025). Premios de presentación y póster (SIAM--SMB, AMS, SACNAS, MIT), 2000--2006. Miembro de AMS, SMB, SIAM y SACNAS. Español (nativo) e inglés (fluido).
\end{cvlist}
\end{cvblock}

\begin{cvblock}{Segura Ugalde, Esteban}{Director del Departamento de Matemática Aplicada (2023--2027) e investigador del CIMPA, Universidad de Costa Rica}
\cvsec{Formación académica}
\begin{cvlist}
  \item PhD. en Matemática, especialización en Álgebra Lineal Numérica / Análisis Numérico, XLIM -- DMI UMR 7252 CNRS, Université de Limoges, Francia, 2011--2015, mención Très Honorable. Tesis: ``Computation of Invariant Pairs and Matrix Solvents''; supervisores: Moulay Barkatou y Paola Boito.
  \item Master 2 ACSYON: Algorithmic, Symbolic Computation and Numerical Optimization, Université de Limoges, 2010--2011. Tesis: ``Polynomial Root-finding via Structured Matrices''; supervisora: Paola Boito.
  \item Estudiante de Master en Matemática, Universidad de Costa Rica, 2009--2010.
  \item Bachillerato en Matemática, Universidad de Costa Rica, 2002--2008.
\end{cvlist}
\cvsec{Experiencia profesional y docente}
\begin{cvlist}
  \item Profesor y coordinador de diversos cursos, Facultad de Matemática, Universidad de Costa Rica (UCR), 2009--2010 y 2015--presente; también en la Maestría Profesional de UCR, 2017 y 2019.
  \item Profesor y coordinador de diversos cursos, Facultad de Matemática, Universidad Nacional de Costa Rica, 2018--presente.
  \item Director del Departamento de Matemática Aplicada, UCR, 2023--2027. Investigador, CIMPA, UCR, 2015--presente.
  \item Consultor de proyectos del CIMPA: análisis de actualización de costos para el Colegio de Profesionales en Ciencias Económicas de Costa Rica (2017); ACAVAMA (2021--presente).
\end{cvlist}
\cvsec{Áreas de especialidad e investigación}
Cálculo simbólico-numérico, álgebra lineal numérica, matrices estructuradas y optimización numérica; en particular pares invariantes, solventes matriciales y búsqueda de raíces de polinomios.
\cvsec{Producción académica}
12 publicaciones (incluye la tesis doctoral y 2 libros), 1 artículo aceptado y 5 trabajos en preparación (4 artículos y 1 libro). Representativas, de la más reciente a la más antigua:
\begin{cvlist}
  \item ``Construcción de diseños de bloques balanceados e incompletos'', Matemática, Educación e Internet (TEC), 2026 (aceptado). En preparación (2026--2027): artículos sobre pares invariantes para el problema no lineal de valores propios, raíces de polinomios con matrices DPR1 y cuasi-separables, valores propios específicos dentro de un contorno y triangularización de polinomios matriciales; libro ``Análisis Numérico'', Editorial UCR, 2027.
  \item ``Computation of Matrix p-th Root Using Moments'', Journal of Computational and Applied Mathematics, 2024.
  \item ``A Subspace Method for Time Series Anomaly Detection in Cyber-Physical Systems'', IFAC-PapersOnLine, 2023; ``Fuzzy clustering by hyperbolic smoothing'', IFCS, 2023.
  \item Libro ``Álgebra Lineal: un enfoque práctico y conciso'', EUNA, 2022; libro ``Introducción a MATLAB'', Editorial UCR, 2020.
  \item ``Description and implementation of an algebraic multigrid preconditioner for H1-conforming finite element schemes'', UNICIENCIA, 2020.
  \item ``A Partial Differential Equation Model with Age-Structure and Nonlinear Recidivism'', Computers and Mathematics with Applications, 2019; ``Secuencias tipo Turyn'', Revista de Matemática: Teoría y Aplicaciones, 2019.
  \item ``A contour integral approach to the computation of invariant pairs'', Theoretical Computer Science, 2017.
  \item Tesis doctoral publicada (Université de Limoges, 2015); ``Some ideas for the computation of matrix solvents'', SNC 2014 (ACM); ``Polynomial Root-finding via structured matrices'', ACM Communications in Computer Algebra, 2012.
\end{cvlist}
\cvsec{Charlas y organización de eventos}
\begin{cvlist}
  \item Charlas en 11 conferencias internacionales (2012--2022), entre ellas ISSAC 2012, SNC 2014, ILAS 2014 y 2019 y ECAMI 2022; 5 en conferencias nacionales (SIMMAC 2016, 2018, 2024 y 2026; SANM2 2018) y 4 charlas internas o coloquios (2014--2018).
  \item Coorganizador de Structured Matrix Days (2014, 2015), del SIMMAC (2016, 2018, 2020 y 2022) y del ECAMI (2022).
\end{cvlist}
\cvsec{Comisiones editoriales, programación e idiomas}
\begin{cvlist}
  \item Comisión Editorial de la Revista de Matemática: Teoría y Aplicaciones (desde 2015); evaluador de artículos de esa revista (desde 2011) y de UNICIENCIA (desde 2019).
  \item Lenguajes y software: C, C++, Java, Python, Mathematica, MATLAB, Scilab y Maple.
  \item Idiomas: español (materna), inglés (TOEIC 845/990, fluente) y francés (DELF A2, medio).
\end{cvlist}
\end{cvblock}

\begin{cvblock}{Solís Chacón, Maikol}{Full Professor of Statistics, Universidad de Costa Rica (CIMPA)}
\cvsec{Formación académica}
\begin{cvlist}
  \item Ph.D. in Statistics, Université Paul Sabatier, Toulouse, Francia, 2010--2014. Tesis: ``Conditional covariance estimation for dimension reduction and sensitivity analysis''; directores: Jean-Michel Loubes y Clement Marteau.
  \item Master in Statistics (M2R), Université Paul Sabatier, Francia, 2009--2010. Tesis sobre estimación de direcciones efectivas de reducción de dimensión; director: Clement Marteau.
  \item Master in Applied Mathematics, Universidad de Costa Rica, 2007--2009. Tesis: ``Estudio de modelos de volatilidad estocástica en la valoración de opciones europeas''; director: José Alexander Ramírez González.
  \item Bachelor in Actuarial Sciences, Universidad de Costa Rica, 2002--2007.
\end{cvlist}
\cvsec{Experiencia profesional y docente}
\begin{cvlist}
  \item Full Professor, Depto. de Matemática, e investigador del CIMPA, Universidad de Costa Rica, agosto de 2014--actualidad. Estadística teórica y aplicada para ciencias actuariales: modelos de supervivencia, distribuciones de pérdida, análisis de datos de baja y alta dimensión y cálculo avanzado.
  \item Senior Data Analyst and Lead Actuary, Comisión Técnica Salarial, UCR, enero de 2021--diciembre de 2025: consultor actuarial y analista para el diseño de un nuevo sistema salarial bajo la Ley Marco de Empleo Público n.º 10159. Consultor del proyecto ACAVAMA para el Banco Nacional de Costa Rica (riesgo de crédito y de mercado), 2015--2021.
  \item Teaching assistant en el Institut de Mathématiques de Toulouse (2013); profesor en la Universidad de Costa Rica (2007--2009) y asistente docente (2003--2006). Analista matemático, Banco Nacional de Costa Rica (2007--2008).
\end{cvlist}
\cvsec{Áreas de especialidad e investigación}
Estadística no paramétrica, análisis de sensibilidad (índices de Sobol), análisis geométrico y topológico de datos, datos de alta dimensión, aprendizaje automático, modelos epidemiológicos y ciencias actuariales; desarrollo de paquetes de R y modelos de lenguaje.
\cvsec{Producción académica}
Ocho artículos arbitrados (2017--2025; cinco de ellos de 2023 en adelante), un capítulo (2025), tres paquetes de R y cuatro informes técnicos. Representativos:
\begin{cvlist}
  \item Communications in Statistics -- Theory and Methods, 2025 (análisis de peso relativo con residualización para efectos no lineales), 2020 y 2017; Communications in Statistics -- Simulation and Computation, 2021.
  \item Computational Statistics, 2023: medida geométrica de bondad de ajuste para detectar patrones en nubes de puntos; Mathematics in Medical and Life Sciences, 2025: triaje temprano de COVID-19 con aprendizaje automático; Revista de Biología Tropical, 2025: estrategias de pruebas masivas para COVID-19 en Costa Rica; Current Microbiology, 2025: microbioma del fruto de café.
  \item Capítulo en Data Science, Classification, and Artificial Intelligence for Modeling Decision Making (Springer), 2025.
  \item Paquetes de R: spatgeom (2023), ResidualRWA (2022) y sobolnp (2019).
\end{cvlist}
\cvsec{Proyectos de investigación}
\begin{cvlist}
  \item Investigador principal, CIMPA: modelación del café con aprendizaje automático según atributos químicos, físicos y sensoriales (2025--2028); clasificador de riesgo para individuos susceptibles a COVID-19 (2023--2025); modelos aditivos de carga alostática y análisis de sensibilidad con teoría topológica (2018--2020); modelos de tránsito en el área metropolitana (2015--2018).
  \item Investigador colaborador: verificación de kits RT-LAMP para SARS-CoV-2, INIFAR (2022).
\end{cvlist}
\cvsec{Docencia y formación de recursos humanos}
\begin{cvlist}
  \item Cinco tesis de maestría dirigidas (2023--2026) en Métodos Matemáticos y Aplicaciones, Ciencias Actuariales y Matemática, sobre bondad de ajuste geométrico, transición de salarios, inferencia causal en campañas publicitarias, aprendizaje profundo y datos de intervalo; una tesis de licenciatura en Ciencias Actuariales en curso (2026) sobre biodiversidad del suelo en cafetales.
  \item Charlas y cursos, entre ellos SIMMAC 2026, IFCS 2024 (San José), useR! 2019 (Toulouse) y LatinR 2018 (Buenos Aires).
\end{cvlist}
\cvsec{Becas e idiomas}
\begin{cvlist}
  \item Beca doctoral de la Universidad de Costa Rica (OAICE), 2010--2014, y Campus France (acuerdo Francia--Costa Rica), 2009--2014. Idiomas: español nativo, inglés y francés fluidos.
\end{cvlist}
\end{cvblock}

\begin{cvblock}{Tall, El Hadji Yaya}{Lecturer in Mathematics, Department of Mathematics, Universidad de Costa Rica (2024--presente)}
\cvsec{Formación académica}
\begin{cvlist}
  \item Ph.D. in Mathematics, IMPA, Rio de Janeiro, Brasil, 2018. Tesis: Moduli of continuity of Lyapunov exponents for random GL(2)-cocycles; asesor: Marcelo Viana.
  \item Postgraduate Diploma, ICTP, Trieste, Italia, 2014. Tesis: Hyperbolic dynamical systems and the Horshoes for the Hénon map; asesor: Stefano Luzzatto.
  \item M.Sc. in Mathematics, UCAD, Dakar, Senegal, 2012; B.Sc. in Mathematics, UCAD, Dakar, Senegal, 2010.
\end{cvlist}
\cvsec{Experiencia profesional y docente}
\begin{cvlist}
  \item Lecturer in Mathematics, Department of Mathematics, Universidad de Costa Rica, 2024--presente.
  \item Postdoctoral Researcher, Centro de Estudos em Analytics e Políticas de Segurança, FGV EAESP, Brasil, 2022--2024; supervisor de investigación: Eduardo Massad.
  \item Postdoctoral Researcher, Department of Mathematics and Statistics, Universidade de São Paulo (IME), Brasil, 2018--2022; supervisor de investigación: André de Carvalho.
\end{cvlist}
\cvsec{Áreas de especialidad e investigación}
Sistemas dinámicos y teoría ergódica. Sus trabajos tratan sobre exponentes de Lyapunov de cociclos aleatorios (continuidad, regularidad hölderiana y analiticidad), incluidos cociclos casi periódicos y cociclos markovianos.
\cvsec{Producción académica}
Dos artículos en revistas arbitradas y tres preprints en arXiv.
\begin{cvlist}
  \item Analyticity of the Lyapunov exponents of random products of quasi-periodic cocycles, Nonlinearity 36 (2023), 3467.
  \item Moduli of continuity for the Lyapunov exponents of random GL(2)-cocycles, Transactions of the American Mathematical Society 373 (2020), 1343--1383 (con M. Viana).
  \item Preprints: Hölder regularity of the top Lyapunov exponent for random matrix products with low-dimensional equator (arXiv:2610.17153); Log-Hölder continuity at zero Lyapunov gap for finite state Markov GL(2)-cocycles (arXiv:2608.22157); Analyticity of Lyapunov Exponents for Mixed Markov Quasi-Periodic Cocycles (arXiv:2608.01569).
\end{cvlist}
\cvsec{Docencia}
\begin{cvlist}
  \item En la Universidad de Costa Rica: Topología General, Introducción a la Teoría Ergódica y los cursos MA1004, MA1005, MA0293, MA0294, MA0321, MA0421, MA0551 y CA0151.
  \item Introduction to linear cocycles and Lyapunov exponents, IME, Universidade de São Paulo, 2020; Topics in Dynamical Systems, curso de verano en IMPA, 2018.
\end{cvlist}
\cvsec{Organización de eventos, idiomas y destrezas}
\begin{cvlist}
  \item Miembro del comité organizador y del comité científico de SIMMAC (2028); miembro del comité científico de EMALCA (2024).
  \item Idiomas: peul (lengua materna), francés y wolof (nativos), inglés, portugués y español (fluidos), italiano (intermedio).
  \item Programación: Python, MATLAB, R, Julia, C y Fortran; Mathematica; análisis numérico de bifurcaciones (AUTO-07p y métodos de continuación numérica); bases de datos SQL.
\end{cvlist}
\end{cvblock}

\begin{cvblock}{Ugalde Gómez, William Javier}{Catedrático, Escuela de Matemáticas, Universidad de Costa Rica}
\cvsec{Formación académica}
\begin{cvlist}
  \item Ph.D., Department of Mathematics, The University of Iowa, Estados Unidos, 1998--2003. Tesis: ``Conformal invariants and the Wodzicki residue''; director: Thomas P. Branson.
  \item Diploma in Mathematics, The Abdus Salam International Centre for Theoretical Physics, Trieste, Italia, 1997--1998. Tesis: ``Characterization of Banach Spaces by the duality map''; director: Charles E. Chidume.
  \item Magister Scientiae, Sistema de Estudios de Posgrado, Universidad de Costa Rica, 1994--1996. Tesis: ``Operadores de Dirac en fibrados de base esférica''; director: Joseph C. Várilly.
  \item Licenciatura, Escuela de Matemáticas, Universidad de Costa Rica, 1991--1992. Tesis: ``De la fórmula de sumación de Poisson a los procesos de muestreo y reconstrucción''; director: Vernor Arguedas. Bachiller, Universidad de Costa Rica, 1986--1991.
\end{cvlist}
\cvsec{Experiencia profesional y docente}
\begin{cvlist}
  \item Catedrático, Escuela de Matemáticas, UCR, diciembre 2008--presente; Profesor Asociado (2008) y Profesor Asociado Invitado (2006--2008), UCR. Catedrático Interino, Escuela de Matemáticas, Universidad Nacional, Heredia, 2011--presente.
  \item Research Assistant Professor, Purdue University, Estados Unidos, 2003--2006. Instructor en The University of Iowa (1999--2003), en Kirkwood College (enero--mayo 2001) y en la Escuela de Matemáticas, UCR (1992--1997); asistente en la UCR (1988--1991).
  \item Gestión académica: Director de la Escuela de Matemáticas, UCR, 2015--2021 (dos periodos); subdirector, 2009--2010; Programa de Posgrado en Matemáticas, Sistema de Estudios de Posgrado, UCR, 2011--2014 (el CV no indica el cargo).
\end{cvlist}
\cvsec{Áreas de especialidad e investigación}
Geometría no conmutativa, residuo de Wodzicki, invariantes conformes y operadores diferenciales conformemente invariantes, simetrías de sistemas cuánticos y álgebras de esferas no conmutativas; también enseñanza del análisis y del cálculo.
\cvsec{Producción académica}
Diez publicaciones (nueve artículos o capítulos y un libro) y cuatro preprints en arXiv.
\begin{cvlist}
  \item Symmetry of some noncommutative sphere algebras (con J. C. Várilly), por aparecer en Journal of Algebra, 2026 (arXiv:2602.06946); Characterization of principal bundles: the noncommutative algebraic case (arXiv:2402.16852, 2024) y the commutative case (arXiv:2310.03824, 2023).
  \item Fundamentos de Análisis Real, libro, UNED, San José, 2014; Una propuesta para la enseñanza del teorema de la función inversa, Matemática Educación e Internet 14 (2014).
  \item Differential forms and the noncommutative residue, Journal of Geometry and Physics 58 (2008); Conformally invariant differential operators and bilinear functionals in six dimensions, Revista de Matemáticas Teoría y Aplicaciones XV (2008).
  \item Some conformal invariants from the noncommutative residue for manifolds with boundary, SIGMA 3 (2007); A construction of GJMS operators using Wodzicki's residue, Communications in Mathematical Physics 261 (2006).
\end{cvlist}
\cvsec{Proyectos de investigación}
\begin{cvlist}
  \item Siete proyectos de la Vicerrectoría de Investigación, UCR: cinco en el CIMM (2007--2015; tres como investigador principal y dos como asociado, p. ej. ``Reformulaciones geométricas en el contexto no conmutativo'', 2012--2015) y dos en el CIMPA como principal: ``Simetrías de campos cuánticos libres'' (2021--2023) y ``Aspectos de simetría en sistemas cuánticos'' (2024--2026).
\end{cvlist}
\cvsec{Dirección de tesis y arbitraje}
\begin{cvlist}
  \item Cinco tesis dirigidas en enseñanza de las matemáticas: una de Licenciatura (UCR, 2012) y cuatro de Maestría (SEP, UCR, 2013--2015), sobre métodos numéricos, series, cálculo diferencial e integral, álgebra y espacios vectoriales.
  \item Árbitro de Journal of Geometry and Physics, Letters in Mathematical Physics y Journal of Mathematical Analysis and Applications.
\end{cvlist}
\end{cvblock}

\begin{cvblock}{Ulloa Esquivel, Roberto}{Profesor de Matemáticas, Universidad de Costa Rica (desde 2020)}
\cvsec{Formación académica}
\begin{cvlist}
  \item Doctor de Filosofía, énfasis en Matemáticas, Purdue University, West Lafayette, Indiana, Estados Unidos, agosto de 2020. Área de investigación: álgebra conmutativa, teoría de multiplicidad generalizada. Tesis: ``Epsilon Multiplicity of Modules with Noetherian Saturation Algebras''.
  \item Bachillerato en Matemáticas, Universidad de Costa Rica, 2012.
  \item Licenciatura en Ingeniería Química, Universidad de Costa Rica, 2011 (egresado, sin proyecto final de graduación).
  \item Bachillerato en Educación Media, Colegio Claretiano, Heredia, 2005.
\end{cvlist}
\cvsec{Experiencia profesional y docente}
\begin{cvlist}
  \item Profesor de Matemáticas, Universidad de Costa Rica, agosto 2020--actual. Imparte cursos de las carreras de Matemática, Ciencias Actuariales y Enseñanza de la Matemática, así como cursos de matemática aplicada.
  \item Profesor de Matemáticas, Purdue University, agosto 2013--junio 2020. Cursos de cálculo básico (límites, derivación, integración, series), cálculo multivariable y cursos enfocados en la enseñanza de las matemáticas, todos en inglés.
  \item Practicante, Teradyne, Heredia, Costa Rica, enero--junio 2013. Desarrolló modelos matemáticos para predecir variables y optimizar partes del proceso; presentó los resultados finales ante ejecutivos de la empresa.
  \item Asistente de docencia, Universidad de Costa Rica, marzo 2008--diciembre 2012. Cursos de Ingeniería Química y Matemáticas: Análisis de Procesos II, Química Analítica, Química General, Cálculo II, proyecto RAMA y Cálculo en una Variable II.
\end{cvlist}
\cvsec{Áreas de especialidad e investigación}
Álgebra conmutativa, teoría de enlaces (linkage theory), intersecciones residuales y comportamiento asintótico de ideales. Otras áreas de interés: ingeniería química, ciencias ambientales y ciencias políticas. Conocimientos informáticos: Sage, Mathematica, MATLAB, Magma y Singular (básico); LaTeX, C++ y Java (intermedio).
\cvsec{Producción académica}
Dos trabajos escritos.
\begin{cvlist}
  \item Epsilon Multiplicity of Modules with Noetherian Saturation Algebras, tesis doctoral, Purdue University, 2020.
  \item Mexican Transportation Systems: Centrality Analysis on Graph Theory, EMALCA Team, CIMAT, Guanajuato, México, 2012 (versión impresa en la revista del INEGI, México).
\end{cvlist}
\cvsec{Seminarios y actividades académicas}
\begin{cvlist}
  \item RTG Homological Conjectures in Commutative Algebra (University of Illinois, Chicago, 2017); Midwest Commutative Algebra and Algebraic Geometry Conference (University of Notre Dame, 2016); RTG in Local Cohomology (University of Illinois, Chicago, 2015).
  \item Seminarios en la UCR sobre homología persistente y teoría de Hurwitz (2012), geometría aritmética y EMALCA (2011) y teoría tauberiana (2010); Mexican Mathematicians in the World (CIMAT, 2012).
\end{cvlist}
\cvsec{Idiomas}
Español (lengua materna); inglés fluido (TOEFL 111/120, 2012); alemán (certificación A2b, Goethe Center, 2011).
\end{cvblock}

\begin{cvblock}{Vargas Ulate, Eny}{Directora Ejecutiva de Servicios Actuariales, Ernst \& Young}
\cvsec{Formación académica}
\begin{cvlist}
  \item Licenciatura en Ciencias Actuariales, Universidad de Costa Rica, Costa Rica, diciembre de 2024.
  \item MBA en Finanzas, Universidad de Costa Rica, Costa Rica, julio de 2019.
  \item Bachillerato en la carrera de Ciencias Actuariales, Universidad de Costa Rica, Costa Rica, agosto de 2014.
  \item Cursos hasta segundo año de la carrera de Estadística, Universidad de Costa Rica, 2010--2011.
\end{cvlist}
\cvsec{Experiencia profesional y docente}
\begin{cvlist}
  \item Ernst \& Young Centroamérica, Panamá y República Dominicana (áreas Valuation, Modeling and Economics, Strategy and Transactions), 2014--actualidad. Directora Ejecutiva de Servicios Actuariales desde julio de 2026.
  \item Valuaciones y auditorías actuariales de beneficios a empleados bajo normas de información financiera (NIC 19, NICSP 39, NIIF para Pymes Sección 28, NIC 26) y USGAAP (ASC 712 y 715), para Centroamérica, Panamá, República Dominicana, Jamaica, Trinidad y Tobago, Barbados, Santa Lucía y México.
  \item Valoraciones actuariales de fondos de mutualidad y planes de pensiones regulados; estimación y revisión de reservas técnicas y matemáticas según la legislación de los reguladores de seguros de Centroamérica, Panamá y República Dominicana; modelos de negocio para la gestión de riesgos de planes de pensiones y fondos de mutualidad.
  \item Análisis de brechas e implementación de la NIIF 17 en aseguradoras de Centroamérica, Panamá, República Dominicana y el Caribe; participación en asignación de precio de compra (PPA) y due diligence para compañías de seguros.
  \item Profesora de la carrera de Ciencias Actuariales, Universidad de Costa Rica, 2024--actualidad.
\end{cvlist}
\cvsec{Áreas de especialidad e investigación}
Valuación actuarial de beneficios a empleados, planes de pensiones y fondos de mutualidad, reservas técnicas y matemáticas, gestión de riesgos de planes de pensiones, implementación de la NIIF 17, PPA y due diligence en aseguradoras. El CV no registra publicaciones ni proyectos de investigación.
\cvsec{Liderazgo, inteligencia artificial y herramientas}
\begin{cvlist}
  \item Liderazgo estratégico y toma de decisiones, desarrollo y gestión de equipos de alto desempeño, comunicación ejecutiva, gestión de stakeholders y de relaciones con clientes, y liderazgo con enfoque inclusivo.
  \item Liderazgo de iniciativas de inteligencia artificial y transformación digital: identificación y priorización de casos de uso de IA en los servicios profesionales, automatización de procesos, integración de IA en análisis y toma de decisiones, y adopción responsable de IA en los equipos.
  \item Software y lenguajes: VBA para Excel, SQL, Matlab, R para estadística, LaTeX, Beamer y MS Office.
\end{cvlist}
\cvsec{Membresías e idiomas}
\begin{cvlist}
  \item Miembro activo del Colegio Profesional de Ciencias Económicas de Costa Rica, con reconocimiento del Bachillerato en Ciencias Actuariales y del MBA en Finanzas.
  \item Idiomas: inglés, nivel intermedio--avanzado; portugués, nivel básico.
\end{cvlist}
\end{cvblock}

\begin{cvblock}{Villalobos Alvarado, Jorge}{Docente e investigador, Escuela de Matemática y CIMPA, Universidad de Costa Rica}
\cvsec{Formación académica}
\begin{cvlist}
  \item Doctorado en Matemática (PhD), Universidad Estatal de Louisiana, Estados Unidos, mayo de 2024; reconocido y equiparado por la UCR al grado de Doctorado Académico en diciembre de 2025. Tesis: Reducibility of Schrödinger Operators on Multilayer Graphs. Promedio general de 4.0/4.0.
  \item Maestría en Matemática (MS), Universidad Estatal de Louisiana, mayo de 2022.
  \item Bachillerato en Física, Universidad de Costa Rica, febrero de 2018; Bachillerato en Matemática, Universidad de Costa Rica, abril de 2017.
\end{cvlist}
\cvsec{Experiencia profesional y docente}
\begin{cvlist}
  \item Docente e investigador, Universidad de Costa Rica, julio de 2024--actualidad. Docencia en la Escuela de Matemática e investigación en física matemática en el CIMPA. Cursos: Cálculo en una Variable I (MA-0250), Principios de Geometría (MA-0370), Cálculo en Varias Variables (MA-0450), Geometría Diferencial (MA-0870), Cálculo II (MA-1002) y Cálculo III (MA-1003).
  \item Asistente y docente, Universidad Estatal de Louisiana, agosto de 2018--mayo de 2024. Cursos de cálculo y precálculo; asistencia en cursos de posgrado en matemática; asistente de investigación de enero de 2020 a julio de 2021.
  \item Docente, Universidad de Costa Rica, agosto de 2017--julio de 2018: Cálculo II (MA-1002), Álgebra Lineal (MA-1004) y Cálculo I (MA-1210). Asistente en Principios de Matemática (MA-0150), agosto--diciembre de 2016.
\end{cvlist}
\cvsec{Áreas de especialidad e investigación}
Física matemática y teoría espectral; operadores de Schrödinger en grafos, con aplicaciones a redes cristalinas y a grafeno en campos magnéticos; análisis funcional, análisis de Fourier y ecuaciones diferenciales. Manejo de MATLAB y Mathematica.
\cvsec{Producción académica}
Un artículo publicado, un trabajo en curso, la tesis doctoral y seis presentaciones seleccionadas.
\begin{cvlist}
  \item Stable defect states in the continuous spectrum of bilayer graphene with magnetic field (con S. P. Shipman), Physica D: Nonlinear Phenomena 455 (2023), 133891.
  \item En curso: Magnetic regularity and defect states in multilayer honeycomb Hamiltonians (2026; título provisional). Tesis doctoral: Reducibility of Schrödinger Operators on Multilayer Graphs (Universidad Estatal de Louisiana, 2024).
  \item Póster en el SIAM Texas--Louisiana Sectional Meeting (University of Louisiana at Lafayette, noviembre de 2023); charla en el Spring Southeastern Sectional Meeting (Georgia Institute of Technology, marzo de 2023).
  \item Charlas en el Texas A\&M Weekly Math Seminar (noviembre de 2022); en el 5th Annual Meeting of the SIAM Texas--Louisiana Section (University of Houston, noviembre de 2022); en la Conference on Spectral Theory of Ergodic Quantum Systems (virtual, marzo de 2022); y en el 4th Annual Meeting de esa sección SIAM (University of Texas Rio Grande Valley, noviembre de 2021).
\end{cvlist}
\cvsec{Formación y participación en eventos}
\begin{cvlist}
  \item TXST Summer School in Mathematical Physics, Texas State University, julio de 2023; conferencia Algebraic Geometry in Spectral Theory, ICERM, Brown University, febrero de 2023; Sixth Great Lakes Mathematical Physics Meeting, Michigan State University, junio de 2022 (asistencia).
\end{cvlist}
\cvsec{Distinciones}
\begin{cvlist}
  \item Evaluación docente promedio en la Universidad de Costa Rica: 9.72 (II semestre de 2024) y 9.90 (I semestre de 2026).
  \item Graduate School Enrichment Award, Universidad Estatal de Louisiana, por excelencia académica.
  \item Mención honorífica en la Competencia Iberoamericana Interuniversitaria de Matemáticas (CIIM), Costa Rica, octubre de 2014; mejor resultado en el examen de admisión de la Universidad de Costa Rica en 2011.
\end{cvlist}
\end{cvblock}

\begin{cvblock}{Villalobos Arias, Mario}{Profesor Catedrático, Escuela de Matemática y CIMPA, Universidad de Costa Rica; profesor en la Escuela de Matemática, Instituto Tecnológico de Costa Rica}
\cvsec{Formación académica}
\begin{cvlist}
  \item Doctorado en Matemáticas, Centro de Investigación y de Estudios Avanzados del IPN (CINVESTAV), México, 2005. Tesis: ``Análisis de Heurísticas de Optimización para Problemas Multiobjectivo''; asesores: Onésimo Hernández-Lerma y Carlos A. Coello Coello.
  \item Maestría en Matemáticas, Universidad de Costa Rica, 1998. Tesis: ``Optimización estocástica para el análisis de proximidades''; asesor: Javier Trejos Zelaya.
  \item Licenciatura en Matemáticas, Universidad de Costa Rica, 1991. Tesis sobre elementos finitos aplicados a ecuaciones del tipo Navier--Stokes; asesor: Vernor Arguedas Troyo.
  \item Bachiller en Matemáticas, Universidad de Costa Rica, 1989.
\end{cvlist}
\cvsec{Experiencia profesional y docente}
\begin{cvlist}
  \item Profesor, Escuela de Matemática, Universidad de Costa Rica, 1987 a la fecha (Catedrático); profesor, Escuela de Matemática, Instituto Tecnológico de Costa Rica, 2010 a la fecha.
  \item Director del Centro de Investigación en Matemática Pura y Aplicada (CIMPA), 2012--2020; Director del Posgrado en Matemática, enero 2009--diciembre 2010; Vice-Decano de la Facultad de Ciencias Exactas, 2008--2010 (Decano a.i. octubre--diciembre 2008), Universidad de Costa Rica.
  \item Miembro de la Comisión de Negociación Salarial (2007 a la fecha) y de la Comisión de Negociación de la Convención Colectiva (2017--2018), por parte del Sindicato de Empleados, UCR.
\end{cvlist}
\cvsec{Áreas de especialidad e investigación}
Heurísticas de optimización (recocido simulado, búsqueda tabú, enjambres de partículas, algoritmos genéticos, sistema inmune artificial) y su convergencia, también multiobjetivo; clasificación, análisis de datos y escalamiento multidimensional; regresión lineal y no lineal; funciones especiales tipo integral de Fresnel; modelación de datos de COVID-19; índices de calidad de suelos y de riesgo.
\cvsec{Producción académica}
Lista de 44 entradas (1993--2022; artículos, capítulos, memorias, informes y preprints), 12 de ellas de 2016 en adelante y solo una de 2022. Representativas:
\begin{cvlist}
  \item Revista Tecnología y Ciencias del Agua (México), 2022: índice de riesgo en la sostenibilidad del servicio de agua y saneamiento en Costa Rica.
  \item Applied Mathematics and Computation 381, 2020: ceros de familias paramétricas de funciones tipo integral de Fresnel generalizada. Preprint en arXiv, 2020: regresión logística generalizada para pronosticar población infectada por Covid-19.
  \item Capítulo en Data Analysis and Rationality in a Complex World (Springer), 2020: agrupamiento de datos binarios con heurísticas de optimización combinatoria.
  \item Optimization 61(7), 2012; Mathematical Methods of Operations Research 64(2) y Soft Computing 10(11), 2006; FOGA 2005 (Lecture Notes in Computer Science 3469).
\end{cvlist}
\cvsec{Proyectos de investigación}
\begin{cvlist}
  \item Laboratorio de ingeniería en transferencia y arquitectura estratégica agropecuaria, Estación Experimental Alfredo Volio, UCR, 2023--2024; red para una prueba diagnóstica de SARS-CoV-2 con tecnología RT-LAMP, INIFAR, 2021--2022.
  \item CIMPA: funciones de Fresnel generalizadas (2018--2021); secuencias de Turyn (2016--2018); ceros de familias de funciones especiales (2015--2018).
  \item PIMAD-CIMPA: optimización global en regresión, estructuras de opinión pública y heurísticas multiobjetivo (2011--2013); clasificación y análisis de datos (2007--2010).
\end{cvlist}
\cvsec{Conferencias, cursos y formación}
\begin{cvlist}
  \item Alrededor de 48 conferencias y cursos en eventos internacionales (2008--2024) en México, Cuba, Bolivia, Chile, Japón, Alemania y otros países; p.\ ej., SLBM (Sonora, 2024) e IFCS (Tokio, 2017). Pasantía en la Universidad de Ciencias Aplicadas de Konstanz (2008).
\end{cvlist}
\cvsec{Distinciones}
\begin{cvlist}
  \item Premio Aportes a la creatividad y la excelencia, Florida Ice and Farm Co. (2008); Chikio Hayashi Award, Federación Internacional de Sociedades de Clasificación (2006).
\end{cvlist}
\end{cvblock}

\begin{cvblock}{Zamora Calero, Rafael}{Profesor Asociado, Escuela de Matemática, y Subdirector del CIMPA, Universidad de Costa Rica}
\cvsec{Formación académica}
\begin{cvlist}
  \item Doctorado en Matemáticas, Université Pierre et Marie Curie, 2011--2015. Tesis: Problèmes de séparation de relations analytiques; director: Dominique Lecomte.
  \item Máster (M2) en Lógica Matemática e Informática Fundamental, Université Paris Diderot, 2011.
  \item Licenciatura en Matemática, Universidad de Costa Rica, 2009.
\end{cvlist}
\cvsec{Experiencia profesional y docente}
\begin{cvlist}
  \item Profesor Asociado, Escuela de Matemática, Universidad de Costa Rica, 2022--presente; Profesor Adjunto, 2020--2022; profesor interino, instructor y adjunto, 2017--2022; profesor interino, 2008--2010.
  \item Subdirector del CIMPA, Universidad de Costa Rica, 2023--presente.
  \item Editor en jefe, Revista de Matemática: Teoría y Aplicaciones, UCR, 2022--presente.
  \item Académico invitado, Departamento de Matemática, Universidad de los Andes, Colombia, 2015--2017.
\end{cvlist}
\cvsec{Áreas de especialidad e investigación}
Según sus publicaciones y charlas: lógica matemática, teoría descriptiva de conjuntos y relaciones de equivalencia de Borel, teoría de modelos y grupos (grupos definiblemente promediables, grupos de isometrías de aleatorizaciones de Borel), y factorización de vectores suaves en representaciones de grupos de Lie.
\cvsec{Producción académica}
Cuatro artículos publicados (2015--2021) y dos preimpresiones (2021, 2026).
\begin{cvlist}
  \item Strong Factorization Theorem for Smooth Vectors of Exponential Solvable Lie Group Representations (con S. Chaves, A. Debrouwere, A. Hernández Alvarado y J. Vindas), preimpresión arXiv:2604.20998, 2026; enviado a Journal of Functional Analysis.
  \item A Fixed-Point Theorem for Definably Amenable Groups (con J. F. Carmona, K. Davila y A. Onshuus), Archive for Mathematical Logic 60 (2021), 413--424.
  \item On the Computational Content of the Theory of Borel Equivalence Relations (con N. Bazhenov, B. Monin y L. San Mauro), Oberwolfach Preprints OWP-2021-06, 2021.
  \item Isometry Groups of Borel Randomizations (con A. Berenstein), Notre Dame Journal of Formal Logic 61 (2020), 297--316.
  \item Injective Tests of Low Complexity in the Plane (con D. Lecomte), MLQ Mathematical Logic Quarterly 65 (2019), 134--169.
  \item Separation of Analytic Sets by Rectangles of Low Complexity, Topology and its Applications 182 (2015), 77--97.
\end{cvlist}
\cvsec{Proyectos y servicio académico}
\begin{cvlist}
  \item Investigador principal, ``Nuevas perspectivas en el estudio de grupos dentro del universo de Shelah'', Fondo Semilla, Vicerrectoría de Investigación, UCR, 2020--2023. Colaborador, ``Avances en la Teoría de Factorización de Vectores Suaves y Analíticos de Representaciones de Grupos de Lie'', Fondo Semilla, 2024.
  \item Organizador principal del XXVI SIMMAC (50 aniversario), CIMPA, UCR (2028); de EMALCA, Liberia (2024); y del XIX SLALM, San José (2021--2022). Organizador de EMALCA, Turrialba (2019). Comité científico del XX SLALM, Bogotá (2026).
\end{cvlist}
\cvsec{Charlas seleccionadas}
\begin{cvlist}
  \item ``Baire Methods for Expansions of Metric Structures'' (con A. Berenstein), Descriptive Set Theory in Turin, Italia, mayo de 2026; ``El Paraíso de Cantor'', Universidad Nacional Mayor de San Marcos, Lima, Perú (virtual), abril de 2026.
  \item ``Big Orbits of the Automorphism Group of a Randomization'' (con A. Berenstein), XX SLALM, Bogotá, julio de 2024; ``Dichotomy Theorems for Borel Relations'', Seminario de Lógica Computacional, TU Wien, Austria, febrero de 2023.
\end{cvlist}
\end{cvblock}

\begin{cvblock}{Zamora Luna, Oscar}{Assistant Professor, Department of Mathematics, Universidad de Costa Rica}
\cvsec{Formación académica}
\begin{cvlist}
  \item Ph.D. in Mathematics and its Applications, Central European University, 2016--2020. Tesis: Quantitative and Extremal Problems in Graphs and Hypergraphs; supervisor: Ervin Győri.
  \item B.Sc. in Mathematics, Universidad de Costa Rica, 2011--2014.
\end{cvlist}
\cvsec{Experiencia profesional y docente}
\begin{cvlist}
  \item Assistant Professor, Department of Mathematics, Universidad de Costa Rica, 2020--presente.
\end{cvlist}
\cvsec{Áreas de especialidad e investigación}
Combinatoria extremal, teoría de grafos, combinatoria probabilística, teoría de Ramsey y métodos algebraicos en combinatoria; en la práctica, problemas de Turán en grafos planares e hipergrafos (ciclos y caminos de Berge).
\cvsec{Producción académica}
Veinte artículos y dos preprints en arXiv (2019, 2024), la mayoría en coautoría con E. Győri.
\begin{cvlist}
  \item The Maximum Number of Pentagons in a Planar Graph, Journal of Graph Theory 108 (2025), 229--256.
  \item Planar Turán Number of the theta-6 graph (en el CV, letra griega theta con subíndice 6), Studia Scientiarum Mathematicarum Hungarica 61 (2024), 89--115.
  \item Turán numbers of Berge trees, Discrete Mathematics 346 (2023); The Minimum Number of 4-Cycles in a Maximal Planar Graph with Small Number of Vertices, Bulletin of the Iranian Mathematical Society 49 (2023).
  \item The Turn Number of the Triangular Pyramid of 3-Layers, Discrete Applied Mathematics 317 (2022); Inverse Turán numbers, Discrete Mathematics 345 (2022); The Maximum Number of Paths of Length Three in a Planar Graph, Journal of Graph Theory 101 (2022); The maximum number of induced C5's in a planar graph, Journal of Graph Theory 99 (2022); The Turán number of the square of a path, Discrete Applied Mathematics 307 (2022).
  \item Connected Hypergraphs without long Berge paths, European Journal of Combinatorics 96 (2021); Saturation problems in the Ramsey theory of graphs, posets and point sets, European Journal of Combinatorics 95 (2021); The Maximum Number of Paths of Length Four in a Planar Graph, Discrete Mathematics 344 (2021); A note on the Turn number of disjoint union of wheels, Discrete Mathematics 344 (2021).
  \item The Structure of Hypergraphs without long Berge cycles, Journal of Combinatorial Theory, Series B (2020); The Maximum Wiener Index of Maximal Planar Graphs, Journal of Combinatorial Optimization (2020).
  \item Avoiding long Berge cycles, the missing cases k = r + 1 and k = r + 2, Combinatorics, Probability and Computing (2019); Ramsey numbers of Berge-hypergraphs and related structures, Electronic Journal of Combinatorics 4 (2019); An Erdős-Gallai type theorem for vertex colored graphs, Graphs and Combinatorics 35 (2019).
  \item Preprints: Packing independent cliques into planar graphs (arXiv:2408.03298, 2024); Independent Chains in Acyclic Posets (arXiv:1912.03288, 2019).
\end{cvlist}
\cvsec{Presentaciones seleccionadas}
\begin{cvlist}
  \item CIMPA-CIMA Mathematics and Applications Workshop, enero de 2023; Instituto Rényi, Budapest, seminarios de teoría extremal de conjuntos y de teoría de grafos (2018--2019); Eurocomb, Bratislava, agosto de 2019; 27th British Combinatorial Conference, Birmingham, agosto de 2019; Karlsruhe Institute of Technology, julio de 2019.
  \item 26th Postgraduate Combinatorial Conference, Oxford, junio de 2019; 25th Postgraduate Combinatorial Conference, Londres, junio de 2018; 7th Polish Combinatorial Conference, Bedlewo, septiembre de 2018; International Colloquium on Graph Theory and Combinatorics, Lyon, julio de 2018; seminario de matemática, Universidad de Costa Rica, abril de 2018.
\end{cvlist}
\end{cvblock}

\begin{cvblock}{Zúñiga Rojas, Ronald Alberto}{Investigador y Profesor Asociado, Escuela de Matemática, Universidad de Costa Rica}
\cvsec{Formación académica}
\begin{cvlist}
  \item Ph.D. en Matemáticas (aprobado con distinción), Universidade de Coimbra y Universidade do Porto, Portugal, 2015. Tesis: ``Homotopy Groups of Moduli of Higgs Bundles''; supervisor: Peter B. Gothen (Centro de Matemática da Universidade do Porto).
  \item Doctor en Matemática, equiparado a Doctorado Académico (reconocimiento y validación del grado), Universidad de Costa Rica, 2017.
  \item Bachillerato en Matemática, Universidad de Costa Rica, 2009.
\end{cvlist}
\cvsec{Experiencia profesional y docente}
\begin{cvlist}
  \item Investigador y Profesor Asociado, Escuela de Matemática, UCR, enero 2019--presente. Coordinador del curso MA-1023 Cálculo y Optimización, para Economía y Estadística (ciclos I, II y III de 2017), y de la Comisión de Docencia de la Escuela (CoDo-EMat, 2021--2022).
  \item Investigador y profesor invitado, Escuela de Matemática, UCR, enero 2015--2019.
  \item Estudiante de Ph.D., Universidade de Coimbra y Universidade do Porto, septiembre 2010--septiembre 2014.
  \item Instructor (2009--2010) y asistente profesor de matemáticas (2003--2008), Escuela de Matemática, UCR.
\end{cvlist}
\cvsec{Áreas de especialidad e investigación}
Geometría y topología; también álgebra, geometría algebraica, geometría diferencial, topología algebraica y topología diferencial. Líneas: fibrados de Higgs, espacios móduli, estratificaciones del cono nilpotente de la aplicación de Hitchin, grupos de homotopía, variedades de caracteres y operadas modulares.
\cvsec{Producción académica}
Seis artículos en revistas internacionales con referees, un libro como autor, una comunicación en actas de encuentro con referees y cinco preprints en progreso.
\begin{cvlist}
  \item Introducción a la teoría de categorías, lógica y topos elementales para una mente curiosa, Revista de Filosofía UCR 63(165), 2024.
  \item Libro: Geometric Invariant Theory, holomorphic vector bundles and the Harder--Narasimhan filtration, Springer Briefs in Mathematics, 2021 (con A. Zamora).
  \item Stratifications on the Nilpotent Cone of the Hitchin Map, Revista Matemática Complutense 35, 2022.
  \item Variations of Hodge Structures of Rank Three k-Higgs Bundles and Moduli Spaces of Holomorphic Triples, Geometriae Dedicata 213, 2021.
  \item A Brief Survey of Higgs Bundles, Revista de Matemática: Teoría y Aplicaciones 26(2), 2019.
  \item Stabilization of the Homotopy Groups of the Moduli Space of k-Higgs Bundles, Revista Colombiana de Matemáticas 52(1), 2018.
  \item Stratifications on the Moduli Space of Higgs Bundles, Portugaliae Mathematica 74, 2017.
  \item Preprints en progreso: Character Varieties and Higgs Bundles; Representation Theory and Character Varieties; Differential Geometry and Hodge Star Operator; Modular Operads and the Moduli Space of k-Higgs Bundles; Homotopy of U(p,q)-Higgs Bundles.
\end{cvlist}
\cvsec{Proyectos y actividades académicas}
\begin{cvlist}
  \item Responsable del proyecto 821-C4-650 ``Costarican-German Collaboration in Research: Character varieties and Higgs bundles'' (iniciativa DFG-CONARE).
  \item Organizador del XXIII SIMMAC (2022); de la sesión Geometry and Topology de los SIMMAC 2020, 2022 y 2024; y del Workshop on Character Varieties and Higgs Bundles (Liberia, 2025). Organizó en 2016 el minicurso Teoría Geométrica de Invariantes en la UCR.
  \item Charlas, entre otras, en VBAC 2016 y 2018, V CLAM (2016), mca-2021 y mca-2025 (Miami), reunión RSME-UMA (Buenos Aires, 2017), coloquios de la Escuela de Matemática UCR y seminarios de IST Lisboa y CMUP; minicurso Geometría Euclídea, Inversiva y Proyectiva (UNAN-León, 2016) y taller de geometría en la Semana de la Educación Matemática (UNA, 2016); charla Lógica y Categorías en el I Coloquio Internacional de Lógica, Epistemología y Metodología (UCR, 2021); asistencia a escuelas y congresos en Europa y América (2012--2025).
\end{cvlist}
\cvsec{Idiomas}
Español (lengua nativa); portugués, inglés e italiano (nivel alto, competencia bilingüe).
\end{cvblock}

\end{document}
